% SPDX-FileCopyrightText: 2024-2026 Fabian Januszewski
% SPDX-License-Identifier: CC-BY-4.0

%======================================================================
% hida_hecke_algebra.tex --- [U]: Hida's universal ordinary Hecke algebra at tame level 1
%
%   NEW at v0.1.0-pre.1 (hecke-extensions slot 14, package WP-8; plan
%   docs/implementation/slot14/00-slot14-plan.md S2 "WP-8"; owner question Q5 at its default,
%   applied as a lead ruling in [X] v1.18.2: file name docs/latex/hida_hecke_algebra.tex,
%   citation key [U], NON-BINDING until the owner decides otherwise, scope tame level 1).
%   Contract discipline as in gl2_contract.tex ([C]) and hecke_extensions.tex ([E]): every
%   non-definitional assertion carries a status tag; [C] and [E] are cited by label, never
%   restated; the predictions (U1)--(U7) are numbered before any code that tests them
%   (rule 20); locators not checked against a text are marked VERIFY (rule A2, owner Q6 default).
%   Hida's papers are read in place from the author's local library (read-only) and are
%   not copied into this repository.
%
%   v0.1.0-pre.2 (slot-14 WP-9, repairs half): every finding of the Fable referee
%   docs/implementation/slot14/WP-9_referee.md applied (B1, B2, S1--S6, m1--m7) and the
%   testability section sec:u-test added; still NON-BINDING (Q5 default).
%
%   VERSION 0.1.0 (hecke-extensions slot 14, package WP-13/Q5, 2026-09-28: the OWNER-SIGNED
%   RELEASE of this document from v0.1.0-pre.2, owner decision 2026-09-28 "Go for publication
%   and release" ([S14] Q5 at the recommended option). [U] is now BINDING like [C] and [E]:
%   owner-released, cited by label, and code is gated on it wherever a future package reads a
%   statement of this document. Scope is tame level 1 only, as at pre.2; tame N >= 3 (CM \cap
%   GL(2), CM \cap Eisenstein) is out of scope until slot 15 (\S\ref{sec:u-scope}). No statement
%   changed from pre.2: the WP-9 referee's repairs (B1, B2, S1--S6, m1--m7) stand as applied, and
%   every VERIFY locator (rule A2, [S14] Q6 at its default) stays marked VERIFY, including
%   Theorem~\ref{thm:u-nocm}(b)'s dependence on Wiles \cite[Thm.~2.1.4]{wiles_ordinary} (case (b)
%   only; case (a) and the rest of the theorem do not need it). Report:
%   docs/implementation/slot14/WP-13_release_note.md \S6.
%
%   v0.1.1-pre.1 (slot-14 references edit, 2026-09-28, NOT signed): citations only. pr:u-U1's
%   Emerton / Ohta locators verified against the owner's library (Thm. 5; Thm. (1.5.5)), its "none
%   is in the project's library" corrected; prop:u-691 gains the Citro--Ghitza cross-check; bib
%   entries completed.  No statement, label or status tag moved; (U1) stays \CONJ.
%
%   v0.1.1-pre.2 (slot-14 WP-U1, 2026-09-28, NOT signed): (U1) proved from Emerton's Theorem 5 (owner decision:
%   write the matching out, re-tag if it holds).  New lem:u-emerton (the dictionary), a proof of pr:u-U1, new
%   rmk:u-ohta; pr:u-U1 re-tagged \CONJ -> \PROVED from \LIT (both halves, (U1-Lambda) and (U1-k), k >= 2);
%   rmk:u-classzero re-tagged \CONJ -> \PROVED from \LIT (Ohta (1.5.5), last clause); rmk:u-Jprincipal decided.
%   Standalone matching for the auditor: docs/implementation/slot14/WP-U1_emerton_matching.md.  A Fable audit follows.
%
%   v0.1.1-pre.3 (slot-14 WP-U1 audit repairs, 2026-09-28, NOT signed): the Fable audit
%   docs/implementation/slot14/WP-U1_fable_audit.md (verdict: the re-tag stands) applied, R1--R6: locators and wording
%   only (Hida AIF 38 (1988) Thm. 1.3, Thm. 2.2, (1.2)-(1.3); Hida ENS (1.11a), (1.15a); Washington 1st ed. Prop. 6.16,
%   Prop. 13.22, Thm. 7.10; Ohta's inputs to App. A.1).  No statement, label or status tag moved.
%
%   v0.1.1-pre.4 (slot-14 WP-C14, 2026-09-29, NOT signed): two \COMP additions, as the Fable referee
%   docs/implementation/slot14/wpc14/REFEREE_I3b_I7b_I8_I9.md S5 accepted them: rmk:u-zero points to the depth-2 cells
%   computed on cohomology ([EX] s14:depthlaw); a (U4)/(U7) data paragraph after sec:u-pred's list (sharpness at the new
%   ell at m = 1 only; g_E = 1 at the 11 computed irregular classes).  I-8's pr:u-U2 addition was rejected by the referee
%   and is not applied.  No statement, label or status tag moved.
%
%   v0.1.1-pre.5 (slot-14, owner rulings of 2026-09-29 on the contract-currency audit, NOT signed): Q-A [U] reconfirmed
%   binding.  ERRATUM (dated 2026-09-29) prop:u-loewy(4): 12 -> 9 non-reduced cells; (46,383), (50,677), (54,67) are reduced
%   (v(disc c_{k,2}) = 2, the double root lifts, two residual systems) -- WP-14c S6, refereed in
%   docs/implementation/slot14/wpc16/REFEREE_J8_WP-C16.md Part II (text II.3 applied); the (68,157) sentence kept and sharpened;
%   sec:u-test (U3) and pr:u-U2 / (U2) follow.  Q-B: L5 ((U4) at n = 4150, n = 4830).  [D1] P19 ((U7): theta^{12} Delta at 13;
%   281/353/745 on 287 rungs), P21 ((U4) exactly sharp: 7385 pairs at 7/11/13 -- W1's count, = I-8's 7412 minus its 27 pairs at
%   the new ell), P23 (NEW prop:u-HP, Harder-Pink's sub-tower, a second route to a special case of pr:u-U1; bib harder_pink).
%   J-8 [U] remark: NEW rmk:u-dihedral ((i) COMP; (ii) not CM via the level, not thm:u-nocm); Proposition R stays [EX].
%
%   v0.1.1-pre.6 (slot-14 inventory patch list, 2026-09-30, NOT signed): rmk:u-zero gains k = 914 ((912,59)); after
%   pr:u-U2 the referee's rewritten COMP sentence (320 WS rungs: every ordinary cuspidal block has g = 1; the ramified
%   rank-2 blocks of prop:u-loewy(4) are the next targets).  No statement, label or status tag moved.
%
%   VERSION 0.2.1 (hecke-extensions slot 17, the close, WP-18, 2026-10-02: the RELEASE from v0.2.1-pre.2 under the owner's ruling O3
%   of 2026-10-01 ("release [C] v1.17.0, [E] 1.17.0, [U] 0.2.1, [EX] 1.3.0 on F1 accept"), as 0.2.0 was released under O2.  The Fable referee
%   F1 (docs/implementation/slot17/WP-3_referee_F1.md) was ACCEPT WITH REPAIRS for [U] (0 blocking, 3 minor), all three applied at pre.2.
%   Over pre.2: the citation paragraph's pointers now read [C] v1.17.0 and [E] 1.17.0 (m-U3, as the paragraph said they must).  No
%   statement, label or status tag moved.
%
%   VERSION 0.2.1-pre.2 (hecke-extensions slot 17, WP-2 resumed, 2026-10-02, NOT signed): the slot-17 Fable referee F1's three
%   minors (docs/implementation/slot17/WP-3_referee_F1.md): m-U1 prop:u-691's l^2 argument moved out of the \LIT argument;
%   m-U2 the sec:u-scope COMP sentence's placement reworded; m-U3 the [C] pointer note (rmk:sattor; release pointers [C] v1.17.0).
%   No statement, label or status tag moved.
%
%   VERSION 0.2.1-pre.1 (hecke-extensions slot 17, WP-2, drafted 2026-10-01 PRE-START during the slot-16 close, NOT signed; plan
%   docs/implementation/slot17/00-slot17-plan.md S2 "WP-2"; release as 0.2.1 at the slot-17 close on the Fable referee F1 (WP-3) and the
%   owner's yes, ruling O3 of 2026-10-01).  Locators and one COMP sentence: (1) the two Washington VERIFYloc sites the slot-16 release left
%   (contract_release_report.md S2: old l. 550 Thm. 5.11 and old l. 698 Cor. 5.14) re-pointed to washington1st (Thm. 5.11 p. 57 verified
%   by WP-8; Cor. 5.14 p. 61 verified mod l by L16-3, the mod l^2 form by Thm. 5.12 p. 59's argument, written out); the prop:u-691 chain
%   sentence follows; (2) slot-16 WP-11's F1 sentence beside (U-1) in sec:u-scope (COMP, gated: the hidden blocks (220,17), (246,19)
%   are not free over the full Hecke algebra; non-ordinary, so no consequence for the universal algebra).  No statement, label or
%   status tag moved.  Report docs/implementation/slot17/WP-2_report.md.
%
%   VERSION 0.2.0 (hecke-extensions slot 16, the close, 2026-10-01: the OWNER-SIGNED RELEASE from v0.2.0-pre.2 under the owner's
%   ruling O2 of 2026-10-01 ("release [C] v1.16.0 and [U] 0.2.0 since WP-14 accepted and the repairs are confirmed; sign [U] 0.1.1").
%   The Fable referee docs/implementation/slot16/WP-14_referee.md was ACCEPT WITH REPAIRS (0 blocking); U-1..U-7 applied at pre.2 and
%   CONFIRMED by docs/implementation/slot16/WP-14_repair_confirmation.md.  Over pre.2, locators only (slot-16 WP-8,
%   docs/implementation/slot16/WP-8_citations.md S2): cor:u-irregular's Kummer congruence (Washington 1st ed. Cor. 5.14, p. 61) and weight
%   congruence (Swinnerton-Dyer Thm. 2(iv), p. 22; Serre Bourbaki 416 n. 1.2 Th. 2, p. 416-04), the same pair after lem:u-emerton, and
%   Washington Thm. 5.11, p. 57 in v0.1.1-pre.3's R5 entry; VERIFY removed at those sites.  [C] / [E] pointers -> v1.16.0 / v1.16.2.  No
%   statement, label or status tag moved.  Part 2 (sec:u-tameN) is binding with the rest.  Report
%   docs/implementation/slot16/contract_release_report.md.
%
%   VERSION 0.1.1 (signed 2026-10-01, owner ruling O2): the text of v0.1.1-pre.6, signed as 0.1.1 with no change of content; recorded at
%   the slot-16 close together with 0.2.0 (0.2.0 is built on the 0.1.1 text; its changes are listed in the changelog of part 2).
%
%   VERSION 0.2.0-pre.2 (hecke-extensions slot 16, WP-12 resumed, 2026-10-01, NOT signed): the repairs answering the Fable referee
%   docs/implementation/slot16/WP-14_referee.md (ACCEPT WITH REPAIRS; 0 blocking, 1 major U-1, 6 minor), every one applied: U-1 h_K = 1
%   in prop:u-cmeis(2) and (U10); U-2 prop:u-cmeis's set-up as three clauses; U-3 the lift of c in D_ell; U-4 Wiles p. 561 and Hida
%   Cor. 3.2 quotations (L16-3, locators.tsv); U-5 (U-1) re-tagged PROVED for theta^j-images, citing [EX] s15:hidden; U-6 Hida--Tilouine /
%   Mazur--Tilouine attribution and bib note; U-7 (U10)'s two Eisenstein systems, psi = chi~ in thm:u-cmN.  No label added, renamed or removed.
%
%   VERSION 0.2.0-pre.1 (hecke-extensions slot 16, WP-12, 2026-10-01, NOT signed; plan
%   docs/implementation/slot16/00-slot16-plan.md S2 "WP-12"): PART 2, tame level N -- NEW sec:u-tameN after sec:u-scope (no
%   section number moves): def:u-hN, thm:u-controlN (LIT, VERIFY), thm:u-cmN (PROVED from STD/LIT: CM needs ell split and
%   |d_K| | N; none at N in {1,2}), prop:u-cmexist, prop:u-cmeis (CM cap Eisenstein: a = 1 mod ell-1; the untwisted CM family
%   meets E_k(1,chi_K)), rmk:u-cmgl2 (CONJ/LIT), rmk:u-pred2 ((U8)-(U10) on [C] (NB6)'s cells).  sec:u-scope: out-of-scope
%   pointer, and proposal (U-1) of docs/implementation/slot15/WP-12_U_E_proposals.md as a CONJ sentence FOR THE REFEREE (owner
%   ruling O4, 2026-10-01).  Built on 0.1.1-pre.6, whose signature as 0.1.1 the owner approved on 2026-10-01 (ruling O2; the
%   signature is recorded at the slot close, not here).  Referee: slot-16 WP-14 (Fable).  No statement of v0.1.x changed, no
%   label renamed or removed, no tag moved.  Report docs/implementation/slot16/WP-12_report.md.
%
%   VERSION 0.3.0-pre.1 (hecke-extensions slot 19, WP-2, 2026-10-03, NOT signed; plan docs/implementation/slot19/00-slot19-plan.md
%   S2 "WP-2", headline 5; owner ruling O5 of 2026-10-03, "the Lambda-adic statement now"; referee slot-19 WP-3 (Fable F1)): PART 3, the
%   Lambda-adic algebra truncated -- NEW sec:u-lambda after part 2 (no section number moves): def:u-trunc, prop:u-truncbasic,
%   lem:u-wtprod, thm:u-interp, prop:u-lambdaprod, lem:u-twistvalues, lem:u-lambdaone, prop:u-eistrunc (COMP:
%   docs/implementation/slot19/wp2/lambda_check.py), rmk:u-lambdabasis, pr:u-lambda ((U11)-(U13)).  sec:u-status's and sec:u-scope's
%   out-of-scope sentence on the Lambda-adic algebra amended; sec:u-scope's code map points to part 3.  No statement of 0.2.1 changed,
%   no label renamed or removed, no tag moved.  Bibliography: hho.  Code: slot-19 WP-HL, gated on F1 (rule 22).
%
%   VERSION 0.3.0-pre.2 (hecke-extensions slot 19, WP-3 repairs, 2026-10-03, NOT signed): F1's six minors m-U1..m-U6 to part 3
%   (textual; docs/implementation/slot19/WP-3_repairs.md).  No label added, renamed or removed; no tag moved.
%
%   VERSION 0.3.0-pre.3 (hecke-extensions slot 19, WP-2 continued, 2026-10-04, NOT signed): the slot-19 research results folded in with
%   the tags of the Fable referee F4 (docs/implementation/slot19/R19-6_F4_referee_report.md, ACCEPT WITH REPAIRS; change table
%   docs/implementation/slot19/WP-2_pre3_changes.md): prop:u-eistrunc gains the depth-3 rungs 37580 / 914250 / 49928412 / 159113510,
%   PROVED from five exact Bernoulli numbers per class (F4 M1), and R19-2's COMP reading of (U12)'s checkable form; (U13) in pr:u-lambda
%   CONJ -> PROVED at (37,32), (103,24) (from the COMP d_Eis = 1 at one weight + Hida's control + Emerton; F4 m8); prop:u-691 a
%   cross-reference to R19-3's Z/691^2 field statement; rmk:u-lambdabasis, the code map of part 3 (F4 m9) and the changelog.  No label
%   added, renamed or removed.
%
%   VERSION 0.3.0-pre.4 (hecke-extensions slot 19, WP-2 continued, 2026-10-04, NOT signed; change table
%   docs/implementation/slot19/WP-2_pre5_changes.md): the code map of part 3 -> BUILT (slot-19 WP-HL, e3a4be35: arm (Q) primary by the
%   owner's ruling of 2026-10-03, HeckeHidaLambda{,Forms}, thm:u-interp's bound enforced, the tower route a cross-check at 37 / 103 only;
%   the validation is HeckeLambdaValidation, the planned name HeckeHidaLambdaValidation is superseded); the F4 m10 ratio corrected in
%   sign and count (F_1/G_32 = 2 mod 37 at the EIGHT weights 32..284; 35 is (a_2 - eps_2)/G_32, WP-HL); prop:u-eistrunc's depth-3 rungs
%   and alpha mod 37^3, 37^4, the depth-4 rung 1072544 read Bernoulli-free from the level-m Coates-Wiles measure (slot-19 R19-7, COMP).
%   No label added, renamed or removed; no tag moved.
%
%   VERSION 0.3.0 (hecke-extensions slot 19, the close, WP-18, 2026-10-04: the RELEASE from v0.3.0-pre.4 under the owner's ruling O3,
%   on the slot-19 referees' verdicts: F1 ACCEPT WITH REPAIRS for part 3 (repairs at pre.2), F4 ACCEPT WITH REPAIRS (the research
%   results; repairs at pre.3), F1-delta ACCEPT for [U] (its optional m-U1 not applied, docs/implementation/slot19/F1delta_repairs.md).
%   Part 3 is binding with the rest.  Release lines only: the citation paragraph's pointers read [C] v1.19.0 and [E] 1.19.0 (every label
%   [U] cites resolves there), part 3's status paragraph, the date line and the changelog of part 3.  No statement, label or status tag moved.
%
%   VERSION 0.4.0-pre.1 (hecke-extensions slot 20, WP-2, 2026-10-04, NOT signed; for the slot-20 Fable referee F1, WP-3): new
%   pr:u-rungs (U14) -- the depth-5 rung at (37,32) (55777784, alpha = 13400697 mod 37^5), the depth-4 rungs at (103,24) (60430740,
%   alpha = 3736634 mod 103^4) and (59,44) (8162384, alpha = 11649137 mod 59^4; lambda = 1 at 59 COMP), PROVED from exact Bernoulli
%   numbers (docs/implementation/slot20/wp2/rung_predict.py), stated before WP-RC20's layer windows; the 691 depth-3 rungs read on the
%   level-3 layer (slot-19 NF691, adc71e8d, COMP).  One label added, none renamed or removed, no tag moved.
%
%   VERSION 0.4.0-pre.2 (hecke-extensions slot 20, WP-3 repairs after the Fable referee F1, 2026-10-04, NOT signed; F1: [U] ACCEPT,
%   one optional note m-U1, applied): pr:u-rungs(2) names the power series whose divided differences are integral (G_a in Lambda).
%   No statement, label or status tag moved.
%
%   VERSION 0.4.0-pre.3 (hecke-extensions slot 20, the close delta for the Fable referee F1', 2026-10-04, NOT signed): pr:u-rungs gains
%   item (4), "Outcome" -- every prediction of (U14) MET by slot-20 WP-RC20's three layer windows (docs/implementation/slot20/RC20_levelm_report.md;
%   NumberFieldValidation, NUMBERFIELD_LAYER_CELLS, NUMBERFIELD_GATE checks=637 failing=0 each): the unique zero modulo ell^m and alpha equal (2)'s
%   values at (37,32,5), (103,24,4), (59,44,4).  Pointer [E] 1.20.0-pre.6.  No statement changed, no label added, renamed or removed; no tag moved
%   (the outcome is COMP beside the PROVED prediction).
%
%   VERSION 0.4.0 (hecke-extensions slot 20, the close, WP-18, 2026-10-05: the RELEASE from v0.4.0-pre.3 under the owner's ruling O6
%   of 2026-10-04 (releases after the Fable windows F1 and F1'), on the slot-20 referees: F1 (WP-3) ACCEPT for [U] (m-U1 applied at pre.2), F1' ACCEPT
%   (docs/implementation/slot20/F1prime_review.md; its optional m-U1, a wording of pr:u-rungs(4)'s cost figures, applied by WP-17 in pre.3's
%   text, docs/implementation/slot20/WP-17_repairs.md).  Release lines only: the citation paragraph's pointers read [C] v1.20.0 and [E] 1.20.0
%   (every label [U] cites resolves there), the date line and the changelog of part 3.  No statement, label or status tag moved.
%
%   VERSION 0.5.0-pre.1 (hecke-extensions slot 21, P8 text first, 2026-10-05, NOT signed; for the slot-21 Fable referee of part 4):
%   PART 4, Lambda-adic modular symbols truncated -- NEW sec:u-msym after part 3 (no section number moves): def:u-msymk,
%   lem:u-msymsigns, lem:u-latticeord (new: the ordinary part does not see the lattice Sym^n vs Hida's L_n), prop:u-msymfree,
%   def:u-lmsym, prop:u-lmsymcontrol, thm:u-lmsymfree, lem:u-msymmass, prop:u-msymnorm, pr:u-msym ((U15) PROVED, (U16), (U17) CONJ),
%   rmk:u-msymbasis.  Hida Invent. 85 Lemma 6.4, Cor. 4.5, Prop. 4.7, Thm. 6.6 read on the page.  Bibliography: greenberg_stevens
%   (not read).  No statement, label or status tag of 0.4.0 changed.  Implementation spec docs/implementation/slot21/P8_spec.md.
%
%   VERSION 0.5.0-pre.2 (hecke-extensions slot 21, F4 repairs to part 4, 2026-10-05, NOT signed): the slot-21 Fable referee of part 4
%   (docs/implementation/slot21/P8_F4_referee_report.md, ACCEPT WITH REPAIRS, 0 B / 3 M / 7 m) applied: M1 (BF) over all b mod ell-1
%   (lem:bdrytorus), M2 Hida's P_n = P_{n+2} read on p. 545 (VERIFYloc lifted), M3 (U17)(a) = Harder's index on the ordinary block
%   (PROVED from STD for n < ell-1), m1-m7 textual; a pointer sentence in rmk:u-gorenstein.  No label added, renamed or removed; no tag of
%   0.4.0 moved.  Spec amended: docs/implementation/slot21/P8_spec.md v0.2.
%   VERSION 0.5.0-pre.8 (hecke-extensions slot 21, 2026-10-06, NOT signed): the Fable referee F4 (docs/implementation/slot21/rc/
%   F4_referee_p5_report.md) of the third remark after lem:u-twistvalues: "lambda(G_b) = 0 iff ell does not divide B_b/b" PROVED
%   given mu(G_b) = 0 (Ferrero-Washington); the step lambda_p = sum_b lambda(G_b) behind HHO tagged STD (Iwasawa's theorem for the
%   minus part, Vandiver from HHO item (2)); HHO sec. 7 read by the referee (lambda_p, nu_p are class-group invariants). The seven
%   pairs re-derived by the referee's own code (tangent-number Bernoulli numbers, divided differences; rc/f4p5/f4p5_kl.py).
%   No statement, label, tag or number changed.  Companions [C] v1.21.0-pre.3, [E] 1.21.0-pre.8, [L37] 0.5.1-pre.6, [RC] 0.1.0-pre.11.
%   VERSION 0.5.0-pre.7 (hecke-extensions slot 21, 2026-10-06, NOT signed): the slot-21 P5 result --- a third unnumbered remark after
%   lem:u-twistvalues: lambda(G_b) = 1 at the irregular pairs used (LIT HHO, abstract item (3), read on the text layer for P5), mu(G_b) = 0
%   (Ferrero-Washington, STD recalled), re-derived from exact Bernoulli numbers by divided differences (COMP; [E] rmk:lviwasawa(6)).
%   No statement, label, tag or number changed.  Companions [C] v1.21.0-pre.3, [E] 1.21.0-pre.7, [L37] 0.5.1-pre.6, [RC] 0.1.0-pre.10.
%   VERSION 0.5.0-pre.6 (hecke-extensions slot 21, 2026-10-06, NOT signed): the measure-side Fable referee F4m's repairs to the two
%   remarks after lem:u-twistvalues (docs/implementation/slot21/rc/F4_referee_measure_report.md): (i) each G-factor is MINUS the branch
%   series of E_1 / iota_* E_{k-1} (E_1 = -x zeta_ell; the product unchanged), the residue's constant PROVED from the factorisation;
%   (ii) the wild extension written out as a proof from Thm. 7.10 / 5.11 (COMP -> PROVED, with COMP).  No statement, label or number
%   changed.  Companions [C] v1.21.0-pre.2, [E] 1.21.0-pre.4, [L37] 0.5.1-pre.5, [RC] 0.1.0-pre.8.
%   VERSION 0.5.0-pre.5 (hecke-extensions slot 21, 2026-10-06, NOT signed): the slot-21 P2 / P4 results --- two unnumbered remarks after
%   the proof of lem:u-twistvalues (the G-factors of the Eisenstein measure are the branch series of E_1 and iota_* E_{k-1}, [E] prop:lveismeas;
%   the wild extension G_b(zeta^t u^r - 1) = -B_{r, omega^{b-r} psi_t}/r, psi_t(u) = zeta^t, COMP, [E] rmk:lvwild) and one sentence in the
%   Evidence paragraph after pr:u-msym (P2 / P4 are single-weight: no evidence for (U16) / (U17)).  No statement, label or tag changed;
%   no number moves.  Companions [C] v1.21.0-pre.2, [E] 1.21.0-pre.3.
%   VERSION 0.5.0-pre.4 (hecke-extensions slot 21, 2026-10-05, NOT signed): the Evidence paragraph after pr:u-msym gains the (103,24)
%   env-arm gluing (j_k = L_k to k = 738, L = 2 at 228) and notes that F4's (691,12) is not in the committed _ALL arm.  Nothing else.
%   VERSION 0.5.0-pre.3 (hecke-extensions slot 21, 2026-10-05, NOT signed): status text only --- prop:u-691's sympy note names the
%   in-house reproduction (HeckeLambdaValidation HECKE_LAMBDA_ALL=1, HeckeLValueValidation (B)); pr:u-msym (U16) literature sentence
%   (Greenberg--Stevens read: a weaker, local, char-0, meromorphic form); an Evidence paragraph after pr:u-msym (P8, P9, P3; COMP); the
%   code map of part 4 BUILT (P8 96cdaaac, P9).  No statement, label or tag changed.  Companions [C] v1.21.0-pre.2, [E] 1.21.0-pre.1.
%
%   VERSION 0.8.1 (slot 26, 2026-11-11: SIGNED BY THE OWNER 2026-11-11 ("I sign the commit update. Fix the two in [U] please."); the text of 0.8.0, unchanged except two status sentences of part 7 (sec:u-status Part 7 sentence; the 'Status and scope' paragraph of part 7), 'pre-release, not signed' -> 'since v0.8.0 (signed by the owner 2026-11-11) released'; no statement, label or value moves)
%   VERSION 0.8.0 (slot 26, 2026-11-11: SIGNED BY THE OWNER 2026-11-11 ("I sign [U] 0.7.4-pre.2. I agree with the version number."); the text of 0.7.4-pre.2, unchanged: header, date, release record and changelog only -- no statement, label or value moves; numbered 0.8.0 rather than 0.7.4 per the referee's m11, accepted by the owner; the printed key of hida_aif2 became [Hid88b] (was [Hid88]) at pre.2, see the release record; the 0.7.4-pre.2 and earlier history below is kept)
%   VERSION 0.7.4-pre.2 (slot 26, Z6 fix round, 2026-11-11, pre-release, NOT signed): the U074pre1 referee's m1--m10 (m10 touches signed 0.7.3 text: the Part 6 release sentence of sec:u-status), n1 (Hida AJM 1988 Thm. 0.1, p. 329; VERIFY removed; bib hida_ajm88), n3, n4, n6; code audit F1 in pr:u-H2(a) (the subset agreement is not evidence); NEW rmk:u-H2data (COMP, the eleven rows; (36,181) case B, lambda = 2); code map Built.  No statement, label or number moves.  m11: the referee recommends 0.8.0 for the signed release; the owner chooses.  Companions [C] 1.27.0, [E] 1.25.3.
%   VERSION 0.7.4-pre.1 (slot 26, Z6 = H2, 2026-11-11, pre-release for a referee, NOT signed): NEW part 7, sec:u-H2, after part 6 (no statement number of 0.7.3 moves): def:u-DF, lem:u-DF, lem:u-DFread, pr:u-H2 ((U25)), code map; abstract count (U1)--(U25) and one clause; sec:u-status Part 7 sentence; release record and changelog of part 7.  Design docs/implementation/slot26/Z6_design.md v0.2.  Companions [C] 1.27.0, [E] 1.25.3.
%   VERSION 0.7.3 (slot 26, 2026-11-10: SIGNED BY THE OWNER 2026-11-10 ("I sign those title pages."); the text of 0.7.3-pre.1, unchanged: header, date and release record only -- no statement, label, number or tag moves; the title page names no sign-off)
%   VERSION 0.7.3-pre.1 (slot 26, 2026-11-10, pre-release for owner signature, NOT signed): title page only, no change of mathematical content (public release HACS 0.1.1): the \date line names only the version, the date and the revision history; the words 'signed by the owner 2026-11-09' leave the title page (the sign-off record stays in the release record of Appendix app:u-history). No statement, label, number or tag moves. (0.7.3-pre was the slot-26 plan's reservation for card Z6, which has not started: no branch, commit or scratch; Z6 now feeds 0.7.4-pre, as RELFIX did with 0.7.2-pre.1.)
%   VERSION 0.7.2 (slot 26, 2026-11-09: SIGNED BY THE OWNER 2026-11-09, effective on the referee's ACCEPT (RELFIX_signed_referee.md); the text of 0.7.2-pre.1, unchanged: header, date and release record only -- no statement, label, number or tag moves)
%   VERSION 0.7.2-pre.1 (slot 26, 2026-11-09, pre-release, NOT signed): text-only: private path removed, no change of mathematical content (release audit REL-AUDIT S1): seven scratch paths reduced to file names or to 'outside the tree'; two report file names that carried a cluster name replaced by 'the WP-C10 report'. No statement, label, number or tag moves.
%   VERSION 0.7.1 (slot 26, 2026-11-09: SIGNED BY THE OWNER 2026-11-09; the text of 0.7.1-pre.1, unchanged: header, date and release record only -- no statement, label, number or tag moves)
%   VERSION 0.7.1-pre.1 (slot 26, 2026-11-09, pre-release, NOT signed): no change of mathematical content: machine names and private paths removed for the public release (owner ruling 2026-11-09).
%
%   VERSION 0.7.0 (slot 25, 2026-11-07: SIGNED BY THE OWNER 2026-11-07; the text of 0.7.0-pre.4, unchanged: header, date and release record only -- no statement, label, number or tag moves)
%   VERSION 0.7.0-pre.4 (LaTeX polish stage 5a, 2026-11-07, NOT signed; NOT YET SEEN BY THE OWNER): front matter and layout only.
%   One-line \date; uniform author block (copied from cli_manual.tex); generic abstract over parts 1--6; new appendix app:u-history
%   "Revision history": the 0.7.0-pre.3 \date field (release record) and the changelog paragraphs of parts 1--6 moved there verbatim,
%   the old author line and abstract quoted there; sections of parts 2--6 retitled by subject (labels unchanged).  No statement,
%   label, number or status tag changes.  Report docs/hacs/LATEX_polish_s5a_report.md.
%
%   VERSION 0.7.0-pre.3 (hecke-extensions slot 25, H1 merge after its audit docs/implementation/slot25/H1_audit.md (PASS_WITH_MINORS), 2026-11-06,
%   NOT signed; BINDING CHANGES NOT YET SEEN BY THE OWNER): part 6 only, text only, no result changes.  pr:u-H1 (U23): the explicit anchor
%   f_{k*} = x^2 + 56702273x + 83587893 mod 547^3, v(disc) = 2, written before H1-R runs (audit F9, rule 22); new rmk:u-H12general (PROVED):
%   Steps 1-2 of prop:u-H12's proof at any g >= 2 (audit F5; the record's G8 word regular_deg_g, [C] sec:recordschema item 8); placed after
%   rmk:u-H3prime, so no statement number moves; the code map of part 6: Planned -> Built (H1, 006f60a9).  No label, number or tag of pre.2 moves.
%
%   VERSION 0.7.0-pre.2 (hecke-extensions slot 25, U6 fix round, 2026-11-06, NOT signed): referee docs/implementation/slot25/U6_referee.md
%   (ACCEPT_WITH_MINORS) F1-F13 applied, text only, no result changes: (U23) degree clause COMP, (U24) restated, prop:u-H12 Step 1 reordered and
%   tag completed, rmk:u-H0census block-dimension condition, lem:u-H0(3) translation, rmk:u-H12lit scope / pp. 282-283, notation, rmk:u-H3prime,
%   prop:u-loewy(3) tag; sec:u-status part-5 sentence fixed (released front matter, NOT YET SEEN BY THE OWNER).  Z0 on hecke-extensions:
%   \HA{prop:heckeZgen}, \HA{sec:heckeZ}; bib key bruns_herzog_cm -> [C]'s bruns_herzog.  Report docs/implementation/slot25/U6_report.md Sec. 5.
%
%   VERSION 0.7.0-pre.1 (hecke-extensions slot 25, package U6, 2026-11-06, NOT signed; card docs/implementation/slot24/heckez/06-work-packages.md
%   "U6", placement docs/implementation/slot25/00-slot25-plan.md v0.2 wave 3): NEW part 6, sec:u-eislambda, after part 5 (no section or
%   statement number of 0.6.0 moves): lem:u-H0, rmk:u-H0census, prop:u-H11, prop:u-H12, rmk:u-H12lit (owner-signed P-K2 of 2026-11-05),
%   pr:u-H1 ((U23) PROVED given data, (U24) CONJ), rmk:u-H3prime (COMP, slot-25 H3), code map, changelog of part 6; from the slot-24 HeckeZ
%   stage 05 v0.3.1 Sec. 1.2, 2.2-2.4.  rmk:u-classzero: owner-signed P-K1 (Ohta's adaptation of Harder-Pink / Kurihara).  prop:u-loewy(3): the
%   Bellaiche-Pollack citation (phenomenology catalogue Sec. 11 item 4; NOT YET SEEN BY THE OWNER).  Status paragraph: one sentence.  Bib keys
%   kurihara1993, calegari_eisdef, bellaiche_pollack, bruns_herzog_cm.  No statement, label, number or tag of 0.6.0 moved.  Companions [C] v1.24.0
%   ([C] 1.25.0-pre.3 sec:heckeZ cited by name only), [E] 1.24.0.  Report docs/implementation/slot25/U6_report.md.
%
%   VERSION 0.6.0 (hecke-extensions slot 22, CLOSE-T T3, 2026-10-12): the RELEASE CANDIDATE from 0.6.0-pre.3 -- NOT YET OWNER-SIGNED (owner
%   sign-off is CLOSE-T step T8).  Header, \date and one changelog paragraph only: no statement, label, number or tag moves.
%   Companions [C] v1.22.0, [E] 1.22.0.
%
%   VERSION 0.6.0-pre.3 (hecke-extensions slot 22, CLOSE-T T1 repairs, 2026-10-12, NOT signed): the F-T1 release referee's NIT-14 (the paragraph
%   'Wording repair at 0.6.0-pre.2' is retitled '(part 4)': it repairs part 4 but sits in part 5's changelog area) and NIT-15 (addendum (2) of
%   def:u-lmsym ends '(text v0.2 / [U] 0.6.0-pre.1, F-U5)', not 'v0.6.0-pre.2').  No statement, label, number or tag moves.
%   Companions [C] v1.22.0-pre.17, [E] 1.22.0-pre.3.
%
%   VERSION 0.6.0-pre.2 (hecke-extensions slot 22, DOC22, 2026-10-08, NOT signed): one wording repair in the 'Not in this part' sentence of part 4
%   (the paragraph before 'Notation of this part'): the L = 4 cells of [L37] were listed among the 'non-etale components'; they are non-ordinary
%   (U_37 nilpotent, [L37] prop:nonord; indeed U_37 in m_A^2, [L37] prop:Uposition), so they are not [U]-objects; the sentence now says so (docs/implementation/slot22/nonreduced_ordinary_survey.md,
%   'the census maximum').  No statement, label, number or status tag moves.
%
%   VERSION 0.6.0-pre.1 (hecke-extensions slot 22, P10-C, 2026-10-06, NOT signed; decisions taken as recommended): NEW part 5, sec:u-tower,
%   after part 4 (so no section or statement number of 0.5.0 moves): the text of docs/implementation/slot22/P10C_contract_text.tex.txt v0.2
%   part (B) verbatim, as refereed by the Fable referee F-U5 (P10C_F4_referee_report.md, ACCEPT WITH REPAIRS, 0 B / 3 M / 13 m, all applied;
%   the in-text "v0.6.0-pre.2" annotations name that refereed revision v0.2): def:u-towerlambda, lem:u-omegar, def:u-towerobj,
%   prop:u-towerorient (COMP at r = 1, gated by (U18)(o)), thm:u-towercontrol ((1), (5) PROVED given the orientation), prop:u-towereis,
%   prop:u-towerfree, cor:u-towertrunc, pr:u-tower ((U18)--(U21) PROVED, (U22) CONJ), rmk:u-towerbasis, rmk:u-towerkitagawa, code map (planned),
%   changelog of part 5.  def:u-lmsym gains the addendum of P10C_addenda.md (2) (labelling only); rmk:u-msymbasis points to part 5; the
%   status paragraph gains one sentence; macros \GaT, \Om, \IwT; bib key miyake (Thm. 4.6.17, locator unverified).  No statement, label or
%   status tag of 0.5.0 changed.  Companion [C] v1.22.0-pre.1 (sec:nebwildsec).
%
%   VERSION 0.5.0 (hecke-extensions slot 21, the close, 2026-10-06: the RELEASE from v0.5.0-pre.8 under the owner's decision D16
%   (release of the pre-release documents at the close; accepted as recommended, 2026-10-06).  Part 4 as refereed: F4 (P8_F4_referee_report.md,
%   repaired at pre.2), F4m (the two remarks after lem:u-twistvalues, repaired at pre.6), F4 of P5 (the third remark, repaired at pre.8).
%   Release lines only: the status paragraph's sentence on part 4, the citation pointers [C] v1.21.0 and [E] 1.21.0, the date line and the
%   changelog of part 4.  No statement, label, number or status tag moved; a release upgrades no tag.  Companions [C] v1.21.0,
%   [E] 1.21.0, [L37] 0.5.1, [RC] 0.1.0, [EX] 1.4.1, [LG] 0.1.5.
%
%   BUILD: ./build_hida_hecke_algebra.sh   (pdflatex, bibtex, pdflatex, pdflatex)
%======================================================================
\documentclass[11pt]{article}
% SPDX-FileCopyrightText: 2024-2026 Fabian Januszewski
% SPDX-License-Identifier: CC-BY-4.0

%======================================================================
% gl2_preamble.tex --- shared preamble and macro file
%
%   The SINGLE source of macros for the unified contract (gl2_contract.tex)
%   and for every companion document that cites it.  Documents \input this
%   file; they do NOT copy it.  (The pre-2026-09 corpus copied macros between
%   hecke_action.tex, hecke_extensions.tex and monodromy_v5.tex; that is what
%   let three incompatible meanings of \De, \tau and \iota coexist.)
%
%   Usage, after \documentclass and before \begin{document}:
%       % SPDX-FileCopyrightText: 2024-2026 Fabian Januszewski
% SPDX-License-Identifier: CC-BY-4.0

%======================================================================
% gl2_preamble.tex --- shared preamble and macro file
%
%   The SINGLE source of macros for the unified contract (gl2_contract.tex)
%   and for every companion document that cites it.  Documents \input this
%   file; they do NOT copy it.  (The pre-2026-09 corpus copied macros between
%   hecke_action.tex, hecke_extensions.tex and monodromy_v5.tex; that is what
%   let three incompatible meanings of \De, \tau and \iota coexist.)
%
%   Usage, after \documentclass and before \begin{document}:
%       % SPDX-FileCopyrightText: 2024-2026 Fabian Januszewski
% SPDX-License-Identifier: CC-BY-4.0

%======================================================================
% gl2_preamble.tex --- shared preamble and macro file
%
%   The SINGLE source of macros for the unified contract (gl2_contract.tex)
%   and for every companion document that cites it.  Documents \input this
%   file; they do NOT copy it.  (The pre-2026-09 corpus copied macros between
%   hecke_action.tex, hecke_extensions.tex and monodromy_v5.tex; that is what
%   let three incompatible meanings of \De, \tau and \iota coexist.)
%
%   Usage, after \documentclass and before \begin{document}:
%       \input{gl2_preamble}
%
%   Every macro below is normative: the notation chapter of the contract
%   (gl2_notation.tex) is the prose form of this file, and the two must agree.
%   A symbol that is not declared here is not part of the fixed notation.
%
%   RETIRED MACROS.  Symbols retired by the notation chapter are redefined
%   here to raise a LaTeX error rather than silently typeset, so that text
%   re-homed from the old documents cannot carry an old meaning in unnoticed.
%   See the block "Retired symbols" at the end.
%======================================================================

\usepackage{amsmath, amssymb, amsthm}
\usepackage{mathtools}
\usepackage{geometry}
\usepackage{booktabs}
\usepackage{array}
\usepackage{longtable}
\usepackage{enumitem}
\usepackage{tikz}
\usepackage{tikz-cd}
\usepackage[colorlinks=true,linkcolor=blue,citecolor=blue,urlcolor=blue]{hyperref}
\setlist{nosep}
\emergencystretch=3em

%----------------------------------------------------------------------
% Theorem environments.  One counter, sectioned, so every statement has a
% unique number of the form <section>.<n> that can be cited from code.
%----------------------------------------------------------------------
\newtheorem{theorem}{Theorem}[section]
\newtheorem{lemma}[theorem]{Lemma}
\newtheorem{proposition}[theorem]{Proposition}
\newtheorem{corollary}[theorem]{Corollary}
\newtheorem{claim}[theorem]{Claim}

\theoremstyle{definition}
\newtheorem{definition}[theorem]{Definition}
\newtheorem{algorithm}[theorem]{Algorithm}
\newtheorem{obligation}[theorem]{Obligation}
\newtheorem{certificate}[theorem]{Certificate}

\theoremstyle{remark}
\newtheorem{remark}[theorem]{Remark}
\newtheorem{assumption}[theorem]{Assumption}
\newtheorem{example}[theorem]{Example}
\newtheorem{prediction}[theorem]{Prediction}
\newtheorem{openitem}[theorem]{Open construction item}

%----------------------------------------------------------------------
% Normative boxes.
%----------------------------------------------------------------------
\newenvironment{contractbox}
  {\par\medskip\noindent\begin{quote}\textbf{CONTRACT.}\ }
  {\end{quote}\medskip}
\newenvironment{decisionbox}
  {\par\medskip\noindent\begin{quote}\textbf{OWNER DECISION.}\ }
  {\end{quote}\medskip}
\newenvironment{erratumbox}
  {\par\medskip\noindent\begin{quote}\textbf{ERRATUM.}\ }
  {\end{quote}\medskip}

%----------------------------------------------------------------------
% STATUS tags.  Every non-definitional assertion carries exactly one.
%----------------------------------------------------------------------
\newcommand{\PROVED}{\textsf{[proved here]}}
\newcommand{\STD}[1]{\textsf{[standard: #1]}}   % MUST carry a precise citation
\newcommand{\CONJ}{\textsf{[conjectural]}}
\newcommand{\COMP}{\textsf{[verified computationally]}}
\newcommand{\OPENSTAT}{\textsf{[open]}}

%----------------------------------------------------------------------
% Number systems, rings, groups
%----------------------------------------------------------------------
\newcommand{\Z}{\mathbb{Z}}
\newcommand{\Q}{\mathbb{Q}}
\newcommand{\Rl}{\mathbb{R}}              % the real numbers (NOT \R: see below)
\newcommand{\Cx}{\mathbb{C}}              % the complex numbers (NOT \C)
\newcommand{\F}{\mathbb{F}}
\newcommand{\Fl}{\F_\ell}
\newcommand{\Flf}{\F_{\ell^{f}}}          % residue field, residue degree f
\newcommand{\Hyp}{\mathcal{H}}            % the upper half plane
\newcommand{\GL}{\mathrm{GL}}
\newcommand{\SL}{\mathrm{SL}}
\newcommand{\Mat}{\mathrm{M}}             % matrix algebra: \Mat_2(\Z)
\newcommand{\Ga}{\Gamma}
\newcommand{\Gbdry}{\Ga_\infty}           % the boundary group <T> x <-I>
\newcommand{\De}{\Delta}                  % the monoid {g in M_2(Z) : det g > 0}
\newcommand{\Deprim}{\De^{\mathrm{prim}}}
\newcommand{\cring}{R}                    % the coefficient ring: F_ell, Z/ell^e, Z/m
\newcommand{\Rel}{\mathcal{R}}            % the relation set of the model (K,\Rel)

%----------------------------------------------------------------------
% Operators
%----------------------------------------------------------------------
\newcommand{\Hom}{\mathrm{Hom}}
\newcommand{\End}{\mathrm{End}}
\newcommand{\Ind}{\mathrm{Ind}}
\newcommand{\Sym}{\mathrm{Sym}}
\newcommand{\id}{\mathrm{id}}
\DeclareMathOperator{\res}{res}
\DeclareMathOperator{\cor}{cor}
\DeclareMathOperator{\conj}{conj}
\DeclareMathOperator{\adjg}{adj}
\DeclareMathOperator{\tr}{tr}
\DeclareMathOperator{\rank}{rank}
\DeclareMathOperator{\sgn}{sign}
\DeclareMathOperator{\fl}{fl}
\DeclareMathOperator{\im}{im}
\DeclareMathOperator{\coker}{coker}
\DeclareMathOperator{\soc}{soc}
\DeclareMathOperator{\rad}{rad}
\DeclareMathOperator{\Ext}{Ext}
\DeclareMathOperator{\Ann}{Ann}
\DeclareMathOperator{\MaxSpec}{MaxSpec}
\DeclareMathOperator{\Frob}{Frob}
\DeclareMathOperator{\supp}{supp}
\newcommand{\aug}{\varepsilon}            % augmentation C_0 -> Z

%----------------------------------------------------------------------
% The simplicial framework
%----------------------------------------------------------------------
\newcommand{\Ku}{K^{u}}                   % universal cover complex
\newcommand{\KR}{K/\Rel}                  % quotient complex
\newcommand{\prR}{\mathrm{pr}_{\Rel}}     % projection K -> K/\Rel
\newcommand{\pru}{\mathrm{pr}^{u}}        % projection K^u -> K/\Rel
\newcommand{\piOne}{\pi_1}
\newcommand{\LCS}{\mathbf{LCS}}
\newcommand{\Rep}{\mathbf{Rep}}
\newcommand{\LS}{\mathbf{LS}}
\newcommand{\cat}[1]{\mathbf{#1}}
\newcommand{\orient}[1]{o_{#1}}           % oriented q-simplices: \orient{q}(K)
% The distinguished lift of a simplex of K into the universal cover K^u, for the
% strict fundamental domain K \subseteq K^u of (H4)(iii).  The code spells this
% i(sigma); the document cannot, because i is sqrt(-1) (notation chapter, C15/C16).
\newcommand{\lift}[1]{\widetilde{#1}}
\newcommand{\Ceq}{C^{\bullet}_{\mathrm{eq}}}
\newcommand{\Cq}[1]{C^{#1}_{\mathrm{eq}}}
\newcommand{\Zq}[1]{Z^{#1}_{\mathrm{eq}}}
\newcommand{\Bq}[1]{B^{#1}_{\mathrm{eq}}}
\newcommand{\Heq}[1]{H^{#1}_{\mathrm{eq}}}

%----------------------------------------------------------------------
% Coefficients
%----------------------------------------------------------------------
\newcommand{\rhon}{\rho_n}                % the symmetric power representation
\newcommand{\rhodual}{\rho_n^{\vee}}      % the contragredient
\newcommand{\Vcoef}{V}                    % V := Sym^n(Z^2)
\newcommand{\VR}{V_{\cring}}              % V \otimes \cring
\newcommand{\Mn}{M_n}                     % Sym^n of F_ell^2 (synonym of \VR at R = F_ell)
\newcommand{\Dn}{D_n}                     % the divided-power / contragredient module
\newcommand{\cf}{f_{12}}                  % the normalised weight-12 cusp form (Ramanujan's Delta)

%----------------------------------------------------------------------
% The Hecke datum
%----------------------------------------------------------------------
\newcommand{\Th}{\Theta}                  % the fused cochain map Theta_g
\newcommand{\Thb}{\Theta^{\partial}}      % the boundary cochain map
\newcommand{\Xig}{\Xi_g}
\newcommand{\taub}{\tau_{\partial}}       % the boundary g-twisted chain datum
\newcommand{\Phig}{\Phi_g}                % the cone cochain map
\newcommand{\Amat}[1]{A^{(#1)}}           % cohomology-level Hecke matrix
\newcommand{\Acone}[1]{A_c^{(#1)}}
\newcommand{\Abdry}[1]{A_{\partial}^{(#1)}}
\newcommand{\Tdc}[1]{T_{\Ga #1 \Ga}}      % T_{Gamma g Gamma}
\newcommand{\diam}{T(p,p)}                % the diamond operator

%----------------------------------------------------------------------
% The pair, the cone, compact supports
%----------------------------------------------------------------------
\newcommand{\Xbulk}{X}                    % the Gamma-tessellation of the truncated half plane
\newcommand{\Xbdry}{\partial X}
\newcommand{\Kbdry}{K_{\partial}}         % the circle model 0-1-2-3
\newcommand{\rres}{r}                     % the restriction cochain map r = j^*
\newcommand{\Cone}[1]{C(#1)}              % the mapping cone
\newcommand{\Dcone}{D}                    % the cone differential
\newcommand{\coneinc}{\mathrm{inc}_c}     % C_eq(K_bdry) -> C(r)
\newcommand{\coneproj}{\mathrm{pr}_c}     % C(r) -> C_eq(K)[1]
\newcommand{\Hcpt}[1]{H^{#1}_{c}}
\newcommand{\etapair}{\eta}               % the pair homotopy
\newcommand{\hder}{h}                     % the derived cochain homotopy h = \tilde\eta
\newcommand{\Lam}{\Lambda}                % the unique C_0(bdry X) -> C_2(X) of eta-uniqueness

%----------------------------------------------------------------------
% The integral line (sec:integral, since [C] v1.8.0; notation
% sec:notation-integral)
%----------------------------------------------------------------------
\newcommand{\Iso}{\mathcal{I}}            % the isotropy-resolution model on X
\newcommand{\Isob}{\mathcal{I}^{\partial}} % its boundary subcomplex
\newcommand{\cmpI}{\mathrm{cmp}}          % the comparison map I_* -> C_*(X)
\newcommand{\CI}{C^{\bullet}_{\Iso}}      % Hom_{Z[Gamma]}(I_*, V_R)
\newcommand{\Hfr}[1]{H^{#1}_{\mathrm{fr}}} % H^q(Gamma;V) modulo torsion
\newcommand{\Zsix}{\widehat{\Z}_{6}}      % Z_2 x Z_3
\newcommand{\Qsix}{\widehat{\Q}_{6}}      % Q_2 x Q_3

%----------------------------------------------------------------------
% The Dickson line and the ordinary projector (sec:dicksontheta,
% sec:ordinary, since [C] v1.11.0-pre.1; notation sec:notation-dickson).
% Sans-serif L, D: the plain letters are taken (L the local system, D_n the
% contragredient module), rule A5.
%----------------------------------------------------------------------
\newcommand{\DkL}{\mathsf{L}}             % the first Dickson invariant x^ell y - x y^ell (the code's L)
\newcommand{\DkD}{\mathsf{D}}             % the second Dickson invariant, degree ell^2 - ell
\newcommand{\thL}{\theta_{\DkL}}          % multiplication by L on cochains and cohomology
\newcommand{\nplus}{n^{+}}               % n + ell + 1 (N is the Hecke level only)
\newcommand{\Qn}{Q_{n}}                   % the quotient Sym^{n+ell+1} / L Sym^n
\newcommand{\eord}{e_{\mathrm{ord}}}      % the ordinary projector of def:ordinary (e alone is lem:bulkperiodic's class)

%----------------------------------------------------------------------
% The involution
%----------------------------------------------------------------------
\newcommand{\inv}{\iota}                  % the involution diag(-1,1) / z -> -zbar
\newcommand{\invK}{\iota_K}               % its simplicial realisation on the 7-vertex model

%----------------------------------------------------------------------
% The Hecke algebra and its local pieces
%----------------------------------------------------------------------
\newcommand{\Talg}{\mathbb{T}}
\newcommand{\midl}{\mathfrak{m}}
\newcommand{\Htot}{H_{\mathrm{tot}}}

%----------------------------------------------------------------------
% Misc
%----------------------------------------------------------------------
\newcommand{\abs}[1]{\left| #1 \right|}
\newcommand{\smat}[4]{\begin{pmatrix} #1 & #2 \\ #3 & #4 \end{pmatrix}}
\newcommand{\code}[1]{\texttt{#1}}
\newcommand{\file}[1]{\texttt{#1}}

%----------------------------------------------------------------------
% Cross-document citation.  \HA{} survives for documents that still cite the
% OLD hecke_action.tex; inside the unified contract, use \ref{} instead.
%----------------------------------------------------------------------
\newcommand{\HA}[1]{[C,~\texttt{#1}]}
\newcommand{\vfive}[1]{[v5,~\texttt{#1}]}
\newcommand{\Espec}[1]{[E,~\texttt{#1}]}

%----------------------------------------------------------------------
% RETIRED SYMBOLS.  These macros existed in the pre-consolidation corpus with
% a meaning the notation chapter has retired.  Using one is an error, not a
% typographic choice.  (\K, \R, \C, \pR, \pu, \image were in the copied macro
% blocks of hecke_action.tex and monodromy_v5.tex.)
%----------------------------------------------------------------------
\newcommand{\gltwoRETIRED}[2]{%
  \PackageError{gl2_preamble}{Retired macro \string#1 -- #2}{%
    See the notation chapter (sec:notation) for the replacement.}}
\newcommand{\K}{\gltwoRETIRED{K}{ambiguous complex name; write K directly}}
\newcommand{\R}{\gltwoRETIRED{R}{R is the coefficient ring; write \string\Rl\space for the reals}}
\newcommand{\C}{\gltwoRETIRED{C}{write \string\Cx\space for the complex numbers}}
\newcommand{\pR}{\gltwoRETIRED{pR}{write \string\prR}}
\newcommand{\pu}{\gltwoRETIRED{pu}{write \string\pru}}
\newcommand{\image}{\gltwoRETIRED{image}{write \string\im}}
\newcommand{\M}{\gltwoRETIRED{M}{write \string\Mat\space for the matrix algebra}}

%
%   Every macro below is normative: the notation chapter of the contract
%   (gl2_notation.tex) is the prose form of this file, and the two must agree.
%   A symbol that is not declared here is not part of the fixed notation.
%
%   RETIRED MACROS.  Symbols retired by the notation chapter are redefined
%   here to raise a LaTeX error rather than silently typeset, so that text
%   re-homed from the old documents cannot carry an old meaning in unnoticed.
%   See the block "Retired symbols" at the end.
%======================================================================

\usepackage{amsmath, amssymb, amsthm}
\usepackage{mathtools}
\usepackage{geometry}
\usepackage{booktabs}
\usepackage{array}
\usepackage{longtable}
\usepackage{enumitem}
\usepackage{tikz}
\usepackage{tikz-cd}
\usepackage[colorlinks=true,linkcolor=blue,citecolor=blue,urlcolor=blue]{hyperref}
\setlist{nosep}
\emergencystretch=3em

%----------------------------------------------------------------------
% Theorem environments.  One counter, sectioned, so every statement has a
% unique number of the form <section>.<n> that can be cited from code.
%----------------------------------------------------------------------
\newtheorem{theorem}{Theorem}[section]
\newtheorem{lemma}[theorem]{Lemma}
\newtheorem{proposition}[theorem]{Proposition}
\newtheorem{corollary}[theorem]{Corollary}
\newtheorem{claim}[theorem]{Claim}

\theoremstyle{definition}
\newtheorem{definition}[theorem]{Definition}
\newtheorem{algorithm}[theorem]{Algorithm}
\newtheorem{obligation}[theorem]{Obligation}
\newtheorem{certificate}[theorem]{Certificate}

\theoremstyle{remark}
\newtheorem{remark}[theorem]{Remark}
\newtheorem{assumption}[theorem]{Assumption}
\newtheorem{example}[theorem]{Example}
\newtheorem{prediction}[theorem]{Prediction}
\newtheorem{openitem}[theorem]{Open construction item}

%----------------------------------------------------------------------
% Normative boxes.
%----------------------------------------------------------------------
\newenvironment{contractbox}
  {\par\medskip\noindent\begin{quote}\textbf{CONTRACT.}\ }
  {\end{quote}\medskip}
\newenvironment{decisionbox}
  {\par\medskip\noindent\begin{quote}\textbf{OWNER DECISION.}\ }
  {\end{quote}\medskip}
\newenvironment{erratumbox}
  {\par\medskip\noindent\begin{quote}\textbf{ERRATUM.}\ }
  {\end{quote}\medskip}

%----------------------------------------------------------------------
% STATUS tags.  Every non-definitional assertion carries exactly one.
%----------------------------------------------------------------------
\newcommand{\PROVED}{\textsf{[proved here]}}
\newcommand{\STD}[1]{\textsf{[standard: #1]}}   % MUST carry a precise citation
\newcommand{\CONJ}{\textsf{[conjectural]}}
\newcommand{\COMP}{\textsf{[verified computationally]}}
\newcommand{\OPENSTAT}{\textsf{[open]}}

%----------------------------------------------------------------------
% Number systems, rings, groups
%----------------------------------------------------------------------
\newcommand{\Z}{\mathbb{Z}}
\newcommand{\Q}{\mathbb{Q}}
\newcommand{\Rl}{\mathbb{R}}              % the real numbers (NOT \R: see below)
\newcommand{\Cx}{\mathbb{C}}              % the complex numbers (NOT \C)
\newcommand{\F}{\mathbb{F}}
\newcommand{\Fl}{\F_\ell}
\newcommand{\Flf}{\F_{\ell^{f}}}          % residue field, residue degree f
\newcommand{\Hyp}{\mathcal{H}}            % the upper half plane
\newcommand{\GL}{\mathrm{GL}}
\newcommand{\SL}{\mathrm{SL}}
\newcommand{\Mat}{\mathrm{M}}             % matrix algebra: \Mat_2(\Z)
\newcommand{\Ga}{\Gamma}
\newcommand{\Gbdry}{\Ga_\infty}           % the boundary group <T> x <-I>
\newcommand{\De}{\Delta}                  % the monoid {g in M_2(Z) : det g > 0}
\newcommand{\Deprim}{\De^{\mathrm{prim}}}
\newcommand{\cring}{R}                    % the coefficient ring: F_ell, Z/ell^e, Z/m
\newcommand{\Rel}{\mathcal{R}}            % the relation set of the model (K,\Rel)

%----------------------------------------------------------------------
% Operators
%----------------------------------------------------------------------
\newcommand{\Hom}{\mathrm{Hom}}
\newcommand{\End}{\mathrm{End}}
\newcommand{\Ind}{\mathrm{Ind}}
\newcommand{\Sym}{\mathrm{Sym}}
\newcommand{\id}{\mathrm{id}}
\DeclareMathOperator{\res}{res}
\DeclareMathOperator{\cor}{cor}
\DeclareMathOperator{\conj}{conj}
\DeclareMathOperator{\adjg}{adj}
\DeclareMathOperator{\tr}{tr}
\DeclareMathOperator{\rank}{rank}
\DeclareMathOperator{\sgn}{sign}
\DeclareMathOperator{\fl}{fl}
\DeclareMathOperator{\im}{im}
\DeclareMathOperator{\coker}{coker}
\DeclareMathOperator{\soc}{soc}
\DeclareMathOperator{\rad}{rad}
\DeclareMathOperator{\Ext}{Ext}
\DeclareMathOperator{\Ann}{Ann}
\DeclareMathOperator{\MaxSpec}{MaxSpec}
\DeclareMathOperator{\Frob}{Frob}
\DeclareMathOperator{\supp}{supp}
\newcommand{\aug}{\varepsilon}            % augmentation C_0 -> Z

%----------------------------------------------------------------------
% The simplicial framework
%----------------------------------------------------------------------
\newcommand{\Ku}{K^{u}}                   % universal cover complex
\newcommand{\KR}{K/\Rel}                  % quotient complex
\newcommand{\prR}{\mathrm{pr}_{\Rel}}     % projection K -> K/\Rel
\newcommand{\pru}{\mathrm{pr}^{u}}        % projection K^u -> K/\Rel
\newcommand{\piOne}{\pi_1}
\newcommand{\LCS}{\mathbf{LCS}}
\newcommand{\Rep}{\mathbf{Rep}}
\newcommand{\LS}{\mathbf{LS}}
\newcommand{\cat}[1]{\mathbf{#1}}
\newcommand{\orient}[1]{o_{#1}}           % oriented q-simplices: \orient{q}(K)
% The distinguished lift of a simplex of K into the universal cover K^u, for the
% strict fundamental domain K \subseteq K^u of (H4)(iii).  The code spells this
% i(sigma); the document cannot, because i is sqrt(-1) (notation chapter, C15/C16).
\newcommand{\lift}[1]{\widetilde{#1}}
\newcommand{\Ceq}{C^{\bullet}_{\mathrm{eq}}}
\newcommand{\Cq}[1]{C^{#1}_{\mathrm{eq}}}
\newcommand{\Zq}[1]{Z^{#1}_{\mathrm{eq}}}
\newcommand{\Bq}[1]{B^{#1}_{\mathrm{eq}}}
\newcommand{\Heq}[1]{H^{#1}_{\mathrm{eq}}}

%----------------------------------------------------------------------
% Coefficients
%----------------------------------------------------------------------
\newcommand{\rhon}{\rho_n}                % the symmetric power representation
\newcommand{\rhodual}{\rho_n^{\vee}}      % the contragredient
\newcommand{\Vcoef}{V}                    % V := Sym^n(Z^2)
\newcommand{\VR}{V_{\cring}}              % V \otimes \cring
\newcommand{\Mn}{M_n}                     % Sym^n of F_ell^2 (synonym of \VR at R = F_ell)
\newcommand{\Dn}{D_n}                     % the divided-power / contragredient module
\newcommand{\cf}{f_{12}}                  % the normalised weight-12 cusp form (Ramanujan's Delta)

%----------------------------------------------------------------------
% The Hecke datum
%----------------------------------------------------------------------
\newcommand{\Th}{\Theta}                  % the fused cochain map Theta_g
\newcommand{\Thb}{\Theta^{\partial}}      % the boundary cochain map
\newcommand{\Xig}{\Xi_g}
\newcommand{\taub}{\tau_{\partial}}       % the boundary g-twisted chain datum
\newcommand{\Phig}{\Phi_g}                % the cone cochain map
\newcommand{\Amat}[1]{A^{(#1)}}           % cohomology-level Hecke matrix
\newcommand{\Acone}[1]{A_c^{(#1)}}
\newcommand{\Abdry}[1]{A_{\partial}^{(#1)}}
\newcommand{\Tdc}[1]{T_{\Ga #1 \Ga}}      % T_{Gamma g Gamma}
\newcommand{\diam}{T(p,p)}                % the diamond operator

%----------------------------------------------------------------------
% The pair, the cone, compact supports
%----------------------------------------------------------------------
\newcommand{\Xbulk}{X}                    % the Gamma-tessellation of the truncated half plane
\newcommand{\Xbdry}{\partial X}
\newcommand{\Kbdry}{K_{\partial}}         % the circle model 0-1-2-3
\newcommand{\rres}{r}                     % the restriction cochain map r = j^*
\newcommand{\Cone}[1]{C(#1)}              % the mapping cone
\newcommand{\Dcone}{D}                    % the cone differential
\newcommand{\coneinc}{\mathrm{inc}_c}     % C_eq(K_bdry) -> C(r)
\newcommand{\coneproj}{\mathrm{pr}_c}     % C(r) -> C_eq(K)[1]
\newcommand{\Hcpt}[1]{H^{#1}_{c}}
\newcommand{\etapair}{\eta}               % the pair homotopy
\newcommand{\hder}{h}                     % the derived cochain homotopy h = \tilde\eta
\newcommand{\Lam}{\Lambda}                % the unique C_0(bdry X) -> C_2(X) of eta-uniqueness

%----------------------------------------------------------------------
% The integral line (sec:integral, since [C] v1.8.0; notation
% sec:notation-integral)
%----------------------------------------------------------------------
\newcommand{\Iso}{\mathcal{I}}            % the isotropy-resolution model on X
\newcommand{\Isob}{\mathcal{I}^{\partial}} % its boundary subcomplex
\newcommand{\cmpI}{\mathrm{cmp}}          % the comparison map I_* -> C_*(X)
\newcommand{\CI}{C^{\bullet}_{\Iso}}      % Hom_{Z[Gamma]}(I_*, V_R)
\newcommand{\Hfr}[1]{H^{#1}_{\mathrm{fr}}} % H^q(Gamma;V) modulo torsion
\newcommand{\Zsix}{\widehat{\Z}_{6}}      % Z_2 x Z_3
\newcommand{\Qsix}{\widehat{\Q}_{6}}      % Q_2 x Q_3

%----------------------------------------------------------------------
% The Dickson line and the ordinary projector (sec:dicksontheta,
% sec:ordinary, since [C] v1.11.0-pre.1; notation sec:notation-dickson).
% Sans-serif L, D: the plain letters are taken (L the local system, D_n the
% contragredient module), rule A5.
%----------------------------------------------------------------------
\newcommand{\DkL}{\mathsf{L}}             % the first Dickson invariant x^ell y - x y^ell (the code's L)
\newcommand{\DkD}{\mathsf{D}}             % the second Dickson invariant, degree ell^2 - ell
\newcommand{\thL}{\theta_{\DkL}}          % multiplication by L on cochains and cohomology
\newcommand{\nplus}{n^{+}}               % n + ell + 1 (N is the Hecke level only)
\newcommand{\Qn}{Q_{n}}                   % the quotient Sym^{n+ell+1} / L Sym^n
\newcommand{\eord}{e_{\mathrm{ord}}}      % the ordinary projector of def:ordinary (e alone is lem:bulkperiodic's class)

%----------------------------------------------------------------------
% The involution
%----------------------------------------------------------------------
\newcommand{\inv}{\iota}                  % the involution diag(-1,1) / z -> -zbar
\newcommand{\invK}{\iota_K}               % its simplicial realisation on the 7-vertex model

%----------------------------------------------------------------------
% The Hecke algebra and its local pieces
%----------------------------------------------------------------------
\newcommand{\Talg}{\mathbb{T}}
\newcommand{\midl}{\mathfrak{m}}
\newcommand{\Htot}{H_{\mathrm{tot}}}

%----------------------------------------------------------------------
% Misc
%----------------------------------------------------------------------
\newcommand{\abs}[1]{\left| #1 \right|}
\newcommand{\smat}[4]{\begin{pmatrix} #1 & #2 \\ #3 & #4 \end{pmatrix}}
\newcommand{\code}[1]{\texttt{#1}}
\newcommand{\file}[1]{\texttt{#1}}

%----------------------------------------------------------------------
% Cross-document citation.  \HA{} survives for documents that still cite the
% OLD hecke_action.tex; inside the unified contract, use \ref{} instead.
%----------------------------------------------------------------------
\newcommand{\HA}[1]{[C,~\texttt{#1}]}
\newcommand{\vfive}[1]{[v5,~\texttt{#1}]}
\newcommand{\Espec}[1]{[E,~\texttt{#1}]}

%----------------------------------------------------------------------
% RETIRED SYMBOLS.  These macros existed in the pre-consolidation corpus with
% a meaning the notation chapter has retired.  Using one is an error, not a
% typographic choice.  (\K, \R, \C, \pR, \pu, \image were in the copied macro
% blocks of hecke_action.tex and monodromy_v5.tex.)
%----------------------------------------------------------------------
\newcommand{\gltwoRETIRED}[2]{%
  \PackageError{gl2_preamble}{Retired macro \string#1 -- #2}{%
    See the notation chapter (sec:notation) for the replacement.}}
\newcommand{\K}{\gltwoRETIRED{K}{ambiguous complex name; write K directly}}
\newcommand{\R}{\gltwoRETIRED{R}{R is the coefficient ring; write \string\Rl\space for the reals}}
\newcommand{\C}{\gltwoRETIRED{C}{write \string\Cx\space for the complex numbers}}
\newcommand{\pR}{\gltwoRETIRED{pR}{write \string\prR}}
\newcommand{\pu}{\gltwoRETIRED{pu}{write \string\pru}}
\newcommand{\image}{\gltwoRETIRED{image}{write \string\im}}
\newcommand{\M}{\gltwoRETIRED{M}{write \string\Mat\space for the matrix algebra}}

%
%   Every macro below is normative: the notation chapter of the contract
%   (gl2_notation.tex) is the prose form of this file, and the two must agree.
%   A symbol that is not declared here is not part of the fixed notation.
%
%   RETIRED MACROS.  Symbols retired by the notation chapter are redefined
%   here to raise a LaTeX error rather than silently typeset, so that text
%   re-homed from the old documents cannot carry an old meaning in unnoticed.
%   See the block "Retired symbols" at the end.
%======================================================================

\usepackage{amsmath, amssymb, amsthm}
\usepackage{mathtools}
\usepackage{geometry}
\usepackage{booktabs}
\usepackage{array}
\usepackage{longtable}
\usepackage{enumitem}
\usepackage{tikz}
\usepackage{tikz-cd}
\usepackage[colorlinks=true,linkcolor=blue,citecolor=blue,urlcolor=blue]{hyperref}
\setlist{nosep}
\emergencystretch=3em

%----------------------------------------------------------------------
% Theorem environments.  One counter, sectioned, so every statement has a
% unique number of the form <section>.<n> that can be cited from code.
%----------------------------------------------------------------------
\newtheorem{theorem}{Theorem}[section]
\newtheorem{lemma}[theorem]{Lemma}
\newtheorem{proposition}[theorem]{Proposition}
\newtheorem{corollary}[theorem]{Corollary}
\newtheorem{claim}[theorem]{Claim}

\theoremstyle{definition}
\newtheorem{definition}[theorem]{Definition}
\newtheorem{algorithm}[theorem]{Algorithm}
\newtheorem{obligation}[theorem]{Obligation}
\newtheorem{certificate}[theorem]{Certificate}

\theoremstyle{remark}
\newtheorem{remark}[theorem]{Remark}
\newtheorem{assumption}[theorem]{Assumption}
\newtheorem{example}[theorem]{Example}
\newtheorem{prediction}[theorem]{Prediction}
\newtheorem{openitem}[theorem]{Open construction item}

%----------------------------------------------------------------------
% Normative boxes.
%----------------------------------------------------------------------
\newenvironment{contractbox}
  {\par\medskip\noindent\begin{quote}\textbf{CONTRACT.}\ }
  {\end{quote}\medskip}
\newenvironment{decisionbox}
  {\par\medskip\noindent\begin{quote}\textbf{OWNER DECISION.}\ }
  {\end{quote}\medskip}
\newenvironment{erratumbox}
  {\par\medskip\noindent\begin{quote}\textbf{ERRATUM.}\ }
  {\end{quote}\medskip}

%----------------------------------------------------------------------
% STATUS tags.  Every non-definitional assertion carries exactly one.
%----------------------------------------------------------------------
\newcommand{\PROVED}{\textsf{[proved here]}}
\newcommand{\STD}[1]{\textsf{[standard: #1]}}   % MUST carry a precise citation
\newcommand{\CONJ}{\textsf{[conjectural]}}
\newcommand{\COMP}{\textsf{[verified computationally]}}
\newcommand{\OPENSTAT}{\textsf{[open]}}

%----------------------------------------------------------------------
% Number systems, rings, groups
%----------------------------------------------------------------------
\newcommand{\Z}{\mathbb{Z}}
\newcommand{\Q}{\mathbb{Q}}
\newcommand{\Rl}{\mathbb{R}}              % the real numbers (NOT \R: see below)
\newcommand{\Cx}{\mathbb{C}}              % the complex numbers (NOT \C)
\newcommand{\F}{\mathbb{F}}
\newcommand{\Fl}{\F_\ell}
\newcommand{\Flf}{\F_{\ell^{f}}}          % residue field, residue degree f
\newcommand{\Hyp}{\mathcal{H}}            % the upper half plane
\newcommand{\GL}{\mathrm{GL}}
\newcommand{\SL}{\mathrm{SL}}
\newcommand{\Mat}{\mathrm{M}}             % matrix algebra: \Mat_2(\Z)
\newcommand{\Ga}{\Gamma}
\newcommand{\Gbdry}{\Ga_\infty}           % the boundary group <T> x <-I>
\newcommand{\De}{\Delta}                  % the monoid {g in M_2(Z) : det g > 0}
\newcommand{\Deprim}{\De^{\mathrm{prim}}}
\newcommand{\cring}{R}                    % the coefficient ring: F_ell, Z/ell^e, Z/m
\newcommand{\Rel}{\mathcal{R}}            % the relation set of the model (K,\Rel)

%----------------------------------------------------------------------
% Operators
%----------------------------------------------------------------------
\newcommand{\Hom}{\mathrm{Hom}}
\newcommand{\End}{\mathrm{End}}
\newcommand{\Ind}{\mathrm{Ind}}
\newcommand{\Sym}{\mathrm{Sym}}
\newcommand{\id}{\mathrm{id}}
\DeclareMathOperator{\res}{res}
\DeclareMathOperator{\cor}{cor}
\DeclareMathOperator{\conj}{conj}
\DeclareMathOperator{\adjg}{adj}
\DeclareMathOperator{\tr}{tr}
\DeclareMathOperator{\rank}{rank}
\DeclareMathOperator{\sgn}{sign}
\DeclareMathOperator{\fl}{fl}
\DeclareMathOperator{\im}{im}
\DeclareMathOperator{\coker}{coker}
\DeclareMathOperator{\soc}{soc}
\DeclareMathOperator{\rad}{rad}
\DeclareMathOperator{\Ext}{Ext}
\DeclareMathOperator{\Ann}{Ann}
\DeclareMathOperator{\MaxSpec}{MaxSpec}
\DeclareMathOperator{\Frob}{Frob}
\DeclareMathOperator{\supp}{supp}
\newcommand{\aug}{\varepsilon}            % augmentation C_0 -> Z

%----------------------------------------------------------------------
% The simplicial framework
%----------------------------------------------------------------------
\newcommand{\Ku}{K^{u}}                   % universal cover complex
\newcommand{\KR}{K/\Rel}                  % quotient complex
\newcommand{\prR}{\mathrm{pr}_{\Rel}}     % projection K -> K/\Rel
\newcommand{\pru}{\mathrm{pr}^{u}}        % projection K^u -> K/\Rel
\newcommand{\piOne}{\pi_1}
\newcommand{\LCS}{\mathbf{LCS}}
\newcommand{\Rep}{\mathbf{Rep}}
\newcommand{\LS}{\mathbf{LS}}
\newcommand{\cat}[1]{\mathbf{#1}}
\newcommand{\orient}[1]{o_{#1}}           % oriented q-simplices: \orient{q}(K)
% The distinguished lift of a simplex of K into the universal cover K^u, for the
% strict fundamental domain K \subseteq K^u of (H4)(iii).  The code spells this
% i(sigma); the document cannot, because i is sqrt(-1) (notation chapter, C15/C16).
\newcommand{\lift}[1]{\widetilde{#1}}
\newcommand{\Ceq}{C^{\bullet}_{\mathrm{eq}}}
\newcommand{\Cq}[1]{C^{#1}_{\mathrm{eq}}}
\newcommand{\Zq}[1]{Z^{#1}_{\mathrm{eq}}}
\newcommand{\Bq}[1]{B^{#1}_{\mathrm{eq}}}
\newcommand{\Heq}[1]{H^{#1}_{\mathrm{eq}}}

%----------------------------------------------------------------------
% Coefficients
%----------------------------------------------------------------------
\newcommand{\rhon}{\rho_n}                % the symmetric power representation
\newcommand{\rhodual}{\rho_n^{\vee}}      % the contragredient
\newcommand{\Vcoef}{V}                    % V := Sym^n(Z^2)
\newcommand{\VR}{V_{\cring}}              % V \otimes \cring
\newcommand{\Mn}{M_n}                     % Sym^n of F_ell^2 (synonym of \VR at R = F_ell)
\newcommand{\Dn}{D_n}                     % the divided-power / contragredient module
\newcommand{\cf}{f_{12}}                  % the normalised weight-12 cusp form (Ramanujan's Delta)

%----------------------------------------------------------------------
% The Hecke datum
%----------------------------------------------------------------------
\newcommand{\Th}{\Theta}                  % the fused cochain map Theta_g
\newcommand{\Thb}{\Theta^{\partial}}      % the boundary cochain map
\newcommand{\Xig}{\Xi_g}
\newcommand{\taub}{\tau_{\partial}}       % the boundary g-twisted chain datum
\newcommand{\Phig}{\Phi_g}                % the cone cochain map
\newcommand{\Amat}[1]{A^{(#1)}}           % cohomology-level Hecke matrix
\newcommand{\Acone}[1]{A_c^{(#1)}}
\newcommand{\Abdry}[1]{A_{\partial}^{(#1)}}
\newcommand{\Tdc}[1]{T_{\Ga #1 \Ga}}      % T_{Gamma g Gamma}
\newcommand{\diam}{T(p,p)}                % the diamond operator

%----------------------------------------------------------------------
% The pair, the cone, compact supports
%----------------------------------------------------------------------
\newcommand{\Xbulk}{X}                    % the Gamma-tessellation of the truncated half plane
\newcommand{\Xbdry}{\partial X}
\newcommand{\Kbdry}{K_{\partial}}         % the circle model 0-1-2-3
\newcommand{\rres}{r}                     % the restriction cochain map r = j^*
\newcommand{\Cone}[1]{C(#1)}              % the mapping cone
\newcommand{\Dcone}{D}                    % the cone differential
\newcommand{\coneinc}{\mathrm{inc}_c}     % C_eq(K_bdry) -> C(r)
\newcommand{\coneproj}{\mathrm{pr}_c}     % C(r) -> C_eq(K)[1]
\newcommand{\Hcpt}[1]{H^{#1}_{c}}
\newcommand{\etapair}{\eta}               % the pair homotopy
\newcommand{\hder}{h}                     % the derived cochain homotopy h = \tilde\eta
\newcommand{\Lam}{\Lambda}                % the unique C_0(bdry X) -> C_2(X) of eta-uniqueness

%----------------------------------------------------------------------
% The integral line (sec:integral, since [C] v1.8.0; notation
% sec:notation-integral)
%----------------------------------------------------------------------
\newcommand{\Iso}{\mathcal{I}}            % the isotropy-resolution model on X
\newcommand{\Isob}{\mathcal{I}^{\partial}} % its boundary subcomplex
\newcommand{\cmpI}{\mathrm{cmp}}          % the comparison map I_* -> C_*(X)
\newcommand{\CI}{C^{\bullet}_{\Iso}}      % Hom_{Z[Gamma]}(I_*, V_R)
\newcommand{\Hfr}[1]{H^{#1}_{\mathrm{fr}}} % H^q(Gamma;V) modulo torsion
\newcommand{\Zsix}{\widehat{\Z}_{6}}      % Z_2 x Z_3
\newcommand{\Qsix}{\widehat{\Q}_{6}}      % Q_2 x Q_3

%----------------------------------------------------------------------
% The Dickson line and the ordinary projector (sec:dicksontheta,
% sec:ordinary, since [C] v1.11.0-pre.1; notation sec:notation-dickson).
% Sans-serif L, D: the plain letters are taken (L the local system, D_n the
% contragredient module), rule A5.
%----------------------------------------------------------------------
\newcommand{\DkL}{\mathsf{L}}             % the first Dickson invariant x^ell y - x y^ell (the code's L)
\newcommand{\DkD}{\mathsf{D}}             % the second Dickson invariant, degree ell^2 - ell
\newcommand{\thL}{\theta_{\DkL}}          % multiplication by L on cochains and cohomology
\newcommand{\nplus}{n^{+}}               % n + ell + 1 (N is the Hecke level only)
\newcommand{\Qn}{Q_{n}}                   % the quotient Sym^{n+ell+1} / L Sym^n
\newcommand{\eord}{e_{\mathrm{ord}}}      % the ordinary projector of def:ordinary (e alone is lem:bulkperiodic's class)

%----------------------------------------------------------------------
% The involution
%----------------------------------------------------------------------
\newcommand{\inv}{\iota}                  % the involution diag(-1,1) / z -> -zbar
\newcommand{\invK}{\iota_K}               % its simplicial realisation on the 7-vertex model

%----------------------------------------------------------------------
% The Hecke algebra and its local pieces
%----------------------------------------------------------------------
\newcommand{\Talg}{\mathbb{T}}
\newcommand{\midl}{\mathfrak{m}}
\newcommand{\Htot}{H_{\mathrm{tot}}}

%----------------------------------------------------------------------
% Misc
%----------------------------------------------------------------------
\newcommand{\abs}[1]{\left| #1 \right|}
\newcommand{\smat}[4]{\begin{pmatrix} #1 & #2 \\ #3 & #4 \end{pmatrix}}
\newcommand{\code}[1]{\texttt{#1}}
\newcommand{\file}[1]{\texttt{#1}}

%----------------------------------------------------------------------
% Cross-document citation.  \HA{} survives for documents that still cite the
% OLD hecke_action.tex; inside the unified contract, use \ref{} instead.
%----------------------------------------------------------------------
\newcommand{\HA}[1]{[C,~\texttt{#1}]}
\newcommand{\vfive}[1]{[v5,~\texttt{#1}]}
\newcommand{\Espec}[1]{[E,~\texttt{#1}]}

%----------------------------------------------------------------------
% RETIRED SYMBOLS.  These macros existed in the pre-consolidation corpus with
% a meaning the notation chapter has retired.  Using one is an error, not a
% typographic choice.  (\K, \R, \C, \pR, \pu, \image were in the copied macro
% blocks of hecke_action.tex and monodromy_v5.tex.)
%----------------------------------------------------------------------
\newcommand{\gltwoRETIRED}[2]{%
  \PackageError{gl2_preamble}{Retired macro \string#1 -- #2}{%
    See the notation chapter (sec:notation) for the replacement.}}
\newcommand{\K}{\gltwoRETIRED{K}{ambiguous complex name; write K directly}}
\newcommand{\R}{\gltwoRETIRED{R}{R is the coefficient ring; write \string\Rl\space for the reals}}
\newcommand{\C}{\gltwoRETIRED{C}{write \string\Cx\space for the complex numbers}}
\newcommand{\pR}{\gltwoRETIRED{pR}{write \string\prR}}
\newcommand{\pu}{\gltwoRETIRED{pu}{write \string\pru}}
\newcommand{\image}{\gltwoRETIRED{image}{write \string\im}}
\newcommand{\M}{\gltwoRETIRED{M}{write \string\Mat\space for the matrix algebra}}

\geometry{a4paper, margin=1in}

% --- local macros of [U] (not part of the fixed notation of gl2_preamble) ---
\newcommand{\LIT}[1]{\textsf{[literature: #1]}}   % a cited theorem; locator VERIFIED or marked VERIFY
\newcommand{\VERIFYloc}{\textsf{VERIFY}}          % locator not checked against the text (rule A2)
\newcommand{\Iw}{\Lambda}                         % the Iwasawa algebra Z_ell[[X]]
\newcommand{\hcusp}{\mathbf{h}}                   % Hida's cuspidal ordinary algebra h^0(a)
\newcommand{\Tu}{\mathbf{T}}                      % the ordinary algebra with the Eisenstein component
\newcommand{\Pk}[1]{P_{#1}}                       % the weight-k prime (1+X)-u^k
\newcommand{\Eis}{\mathrm{E}}
\newcommand{\cusp}{\mathrm{S}}
\newcommand{\Cg}{\mathcal{C}}                     % congruence modules
\newcommand{\ee}{\epsilon}                        % the Lambda-adic Eisenstein eigenvalues
\newcommand{\Ek}{\mathcal{E}}                     % the set of ordinary eigenforms of weight k
\newcommand{\Qlb}{\overline{\Q}_\ell}
\newcommand{\Zl}{\Z_\ell}
\newcommand{\Ql}{\Q_\ell}
\newcommand{\LL}{\mathrm{LL}}                     % Loewy length
\newcommand{\Uspec}[1]{[U,~\texttt{#1}]}
% local macros of part 5 (v0.6.0-pre.1, slot-22 P10-C):
\newcommand{\GaT}[1]{\Ga^{\mathrm t}_{#1}}            % the tame-character group Gamma_{H_r}(ell^r) of part 5
\newcommand{\Om}[1]{\Omega_{#1}}                      % Hida's omega_{2,r}(X) = (1+X)^{ell^{r-1}} - u^{2 ell^{r-1}}
\newcommand{\IwT}[2]{\Iw^{(#2)}_{#1}}                 % Lambda/(ell^K, Omega_r)

\title{Hida's Universal Ordinary Hecke Algebra at Tame Level One\\[1ex]
\large Components, their intersections, no CM, and the saturation link to Loewy length}
\author{Fabian Januszewski\\ Universit\"at Paderborn\\ \texttt{fabian.januszewski@math.uni-paderborn.de}}
\date{Version 0.8.1 --- 2026-11-11 (revision history: Appendix~\ref{app:u-history})}

\begin{document}
\maketitle

\begin{abstract}
Let $\ell\ge5$. It fixes Hida's ordinary Hecke algebra $\Tu$ over $\Iw=\Zl[[X]]$ at tame level~$1$, split into the
even weight classes $a$ modulo $\ell-1$, and reads it at finite weight: quotients of finitely generated $\Iw$-modules modulo
$(P_k,\ell^m)$ depend only on a weight tower, so characteristic polynomials, mod-$\ell$ local algebras and truncated congruence
lengths are tower invariants, and one squarefree characteristic polynomial modulo $\ell$ proves a class étale over $\Iw$. At tame
level~$1$ it proves that no component is CM and that, in a class $a\not\equiv0$, the Eisenstein component meets the cuspidal locus
exactly at the irregular pairs, with (through Emerton's theorem) congruence length $L_k=v_\ell(B_k/k)$ and congruence module
$\Iw/(G_a)$; and it links saturation, Gorensteinness and Loewy length. It then determines where CM components can live, and where
they can meet other components, at tame level $N$ prime to $\ell$ with a tame character. Back at tame level~$1$, it treats truncations
of the $\Iw$-adic algebra and modular symbols, recovered from finitely many weights; the $\Iw$-action on the ordinary cohomology of
the tower $\Ga_1(\ell^{r})$, at one level and weight~$2$; and, for $a\not\equiv0$, the Eisenstein component over $\Iw$, whose
ring-theoretic type one weight decides, at cuspidal rank one and at $(\ell,a)=(547,486)$, the only pair of cuspidal rank two with
$\ell<10^{5}$; and the discriminant series of eleven ramified ordinary cuspidal families at tame level~$1$. Every non-definitional assertion carries a status tag; the predictions are numbered (U1)--(U25).
\end{abstract}

\tableofcontents

%----------------------------------------------------------------------
\section{Status, conventions, scope}\label{sec:u-status}
%----------------------------------------------------------------------

\paragraph{Status.} [U] is \textbf{binding} at v0.1.0 (owner question Q5 of the slot-14 plan, owner decision
2026-09-28 ``Go for publication and release''): like [C] and [E], it is owner-released, cited by label, and
code is gated on it wherever a future package reads one of its statements. \textbf{Scope}: tame level $1$
only; tame $N\ge3$ (CM $\cap$ $\GL(2)$, CM $\cap$ Eisenstein) is out of scope until slot~15
(\S\ref{sec:u-scope}). It was refereed by slot-14 WP-9
(Fable, \file{docs/implementation/slot14/WP-9\_referee.md}: accept with repairs); v0.1.0 applies every repair
of pre.2 unchanged (no statement moved at the release, \file{docs/implementation/slot14/WP-13\_release\_note.md}
\S6; the owner reconfirmed it as binding on 2026-09-29, audit question Q-A). \textbf{Since 2026-10-01} (owner ruling O2) v0.1.1 is signed (the text of
0.1.1-pre.6) and v0.2.0, with part~2, is the owner-signed release. [C] (\file{gl2\_contract.tex}, v1.20.0 since v0.4.0; v1.19.0 at v0.3.0; v1.17.0 at v0.2.1; v1.16.0 at v0.2.0; v1.14.1 through v0.1.1) is cited as \HA{\textit{label}},
[E] (\file{hecke\_extensions.tex}, 1.20.0 since v0.4.0; 1.19.0 at v0.3.0; 1.17.0 at v0.2.1; v1.16.2 at v0.2.0; v1.15.0 through v0.1.1) as \Espec{\textit{label}} (\HA{rmk:sattor}, cited here, is a label of
[C] from v1.17.0 on; the slot-17 referee F1's m-U3), the slot-14 Loewy note \file{docs/implementation/slot14/WP-M1\_loewy\_note.md}
as [M1], the slot-14 binding-document proposals \file{docs/implementation/slot14/WP-D1\_binding\_doc\_proposals.md} as [D1];
statements of [U] as \Uspec{\textit{label}}. \textbf{Part~2} (v0.2.0-pre.1, \S\ref{sec:u-tameN}) extends the scope to tame level $N$ with a tame character; it cites
[C] v1.15.0's \HA{sec:nebentypus} and \HA{NB6} (unchanged in [C] v1.16.0), and was a pre-release until the slot-16 referee (WP-14, ACCEPT WITH REPAIRS, repaired)
and the owner's signature of v0.2.0 (ruling O2, 2026-10-01); it is binding since v0.2.0. \textbf{Part~3} (v0.3.0-pre.1, \S\ref{sec:u-lambda}, the truncated $\Iw$-adic
algebra) was a pre-release through v0.3.0-pre.4 (refereed by the slot-19 F1, F4 and F1-delta); \textbf{since v0.3.0} (owner ruling O3) it is binding with the rest.
It cites [C] v1.21.0 and [E] 1.21.0 since v0.5.0 (v1.20.0 / 1.20.0 at v0.4.0; v1.19.0 / 1.19.0 at v0.3.0; every label it cites resolves there). \textbf{Part~4} (v0.5.0-pre.1, \S\ref{sec:u-msym}, the
truncated $\Iw$-adic modular symbols) was a pre-release, \textbf{not signed}, through v0.5.0-pre.8, written for the slot-21 Fable referee before the code of slot-21 P8 (refereed by F4, F4m and the F4 of P5, all repaired);
\textbf{since v0.5.0} it is released with the slot-21 close, 2026-10-06 (owner decision D16, accepted as recommended), with the rest; it cites [C] v1.21.0 and [E] 1.21.0 (v1.20.0 / 1.20.0 through v0.5.0-pre.8). \textbf{Part~5} (v0.6.0-pre.1, \S\ref{sec:u-tower}, the tower $\Ga_1(\ell^{r})$, wild characters
and the $\Om{}$-type truncations) was a pre-release through v0.6.0-pre.3, written before the code of slot-22 P10 and refereed by the slot-22 Fable referee F-U5 (ACCEPT WITH REPAIRS,
all applied in the inserted text); \textbf{since v0.6.0} (signed by the owner 2026-10-16) it is released with the rest; it cites [C] v1.22.0-pre.1 (\HA{sec:nebwildsec}). \textbf{Part~6} (v0.7.0-pre.1, \S\ref{sec:u-eislambda}, the Eisenstein component over $\Iw$ at irregular pairs) was a pre-release through v0.7.0-pre.4,
written for an adversarial referee before the code of H1 (refereed by the slot-25 U6 referee, ACCEPT WITH MINORS, repaired at v0.7.0-pre.2); \textbf{since v0.7.0} (signed by the owner 2026-11-07)
it is released with the rest; it cites [C] v1.24.0 by label and [C] v1.25.0 (\HA{sec:heckeZ}).
\textbf{Part~7} (v0.7.4-pre.1, \S\ref{sec:u-H2}, the discriminant series of the ramified ordinary families of Proposition~\ref{prop:u-loewy}(4)) was a pre-release through v0.7.4-pre.2,
written for a referee before the code of the slot-26 package Z6 (refereed at pre.1, ACCEPT WITH CHANGES, applied at v0.7.4-pre.2 together with the computed data, Remark~\ref{rmk:u-H2data});
\textbf{since v0.8.0} (signed by the owner 2026-11-11) it is released with the rest; parts~1--6 are as released in 0.7.3 (signed by the owner 2026-11-10). It cites no label of [C] beyond those of parts~1--6;
the companion is [C] 1.27.0.

\paragraph{Tags.} Every non-definitional assertion carries one tag. \PROVED: proved here, from the statements it
names --- when those include a \LIT{} or \STD{} statement the tag says so (``\PROVED{} from \dots''), and the proof is
complete relative to them. \STD{\dots}: a standard fact with a citation. \LIT{\dots}: a theorem of the literature,
with its locator; a locator not checked against the text is marked \VERIFYloc{} (rule~A2 of [C]; slot-14 Q6 at its
default), one checked against Hida's read-only papers says so. \CONJ: a prediction; a failure is a finding, not a
defect. \COMP: a computation, with its data source.

\paragraph{Conventions.} $\ell\ge5$ is prime, $u:=1+\ell$, $\Iw:=\Zl[[X]]$ (identified with $\Zl[[1+\ell\Zl]]$ by
$1+X\leftrightarrow u$), $\mathfrak m_\Iw=(\ell,X)$. Weights $k=n+2$ are even (odd $n$ contribute nothing at
$\ell\ge5$, \HA{def:hidatower}); a \emph{class} is $a\in2\Z/(\ell-1)\Z$ and $k$ lies in it when $k\equiv a$. For $k\ge2$
put $\Pk{k}:=(1+X)-u^{k}\in\Iw$, the weight-$k$ prime of \HA{thm:hidacontrol}(3); $\Iw/\Pk{k}=\Zl$ by $X\mapsto u^{k}-1$.
In [C]'s notation the tower of level~$m$ through $n_0$ is $\mathrm{Tw}_m(\ell;n_0)$ (\HA{def:hidatower}); for weights,
$k,k'$ lie in one tower of level $m$ iff $k\equiv k'\pmod{(\ell-1)\ell^{m-1}}$. $\omega$ is the Teichmüller character,
$\langle x\rangle=x\,\omega(x)^{-1}$. $v=v_\ell$ is the $\ell$-adic valuation, extended to $\Qlb$ with $v(\ell)=1$.
Eigenforms are normalised ($a_1=1$) and viewed in $\Qlb$ through a fixed embedding $\overline{\Q}\hookrightarrow\Qlb$.
Wherever $\Ek_k$, $d_k$ or $E_k$ appears, $k\ge3$ is assumed (it is, statement by statement): at $k=2$ Hida's $\hcusp_2$ also contains the
$\ell$-new forms of $\Ga_0(\ell)$ ($a_\ell=\pm1$, ordinary), which nothing here uses (\HA{def:hidatower} has $n\ge2$). \textbf{Warning:} $d_k$ here
(the number of ordinary level-one cuspidal eigenforms, Lemma~\ref{lem:u-levelone}) is \emph{not} [C]'s $d_k=\dim M_k$ of \HA{cl:dimH1};
[C]'s $h=2d_k-1$ is [U]'s $h^{1}(n)=2\dim S_k+1$.

\paragraph{Scope} (\S\ref{sec:u-scope}). Tame level $1$ only. The $\Iw$-adic algebra is not computed; only its
specialisations are (\S\ref{sec:u-finite}). \textbf{Since v0.3.0-pre.1} (part~3, \S\ref{sec:u-lambda}; a pre-release) its \emph{truncations} modulo
$(\ell^{K},X^{m})$ are in scope, reconstructed from finitely many specialisations with explicit precision; the whole algebra is not. CM $\cap$ $\GL(2)$ and CM $\cap$ Eisenstein live at
tame level $N\ge3$ and wait for composite $N$ (\HA{sec:levelN}; slot~15).

%----------------------------------------------------------------------
\section{The objects}\label{sec:u-setting}
%----------------------------------------------------------------------

\begin{definition}[The cuspidal ordinary Hida algebra]\label{def:u-hcusp}
For a class $a$, $\hcusp(a)$ is Hida's cuspidal ordinary Hecke algebra of tame level~$1$ in the class $a$ --- the
algebra $h^{0}(\ell,a)$ of \HA{thm:hidacontrol}(3) (Hida's $N=\ell$, $N_0=1$) --- a $\Iw$-algebra generated by the
images of $T(p)$, $p\ne\ell$, and $U=T(\ell)$ (their $\Zl$-span is dense and $\Iw$ acts through the diamond operators).
For $k\ge2$ in the class, $\hcusp_k:=h^{0}_k(\Ga_0(\ell))$ is the $\Zl$-algebra generated by the same operators on the
ordinary cusp forms of weight $k$ on $\Ga_0(\ell)$.
\end{definition}

\begin{theorem}[Hida: freeness and control]\label{thm:u-control}
\LIT{\cite[Thm.~3.1 and Cor.~3.2, p.~250]{hida_ens}, transcribed with verified locators as \HA{thm:hidacontrol}(3) and
\HA{thm:hidaweight}(1)} $\hcusp(a)$ is free of finite rank $d(a)$ over $\Iw$, and for every $k\ge2$ in the class the map
$\hcusp(a)\to\hcusp_k$ induces $\hcusp(a)/\Pk{k}\hcusp(a)\cong\hcusp_k$, with $T(p)\mapsto T_p$ and $U\mapsto U$. For $k\ge3$,
$d(a)=\rank_{\Zl}\hcusp_k$ is the number of $\ell$-ordinary normalised eigenforms of weight $k$ on $\Ga_0(\ell)$.
\end{theorem}

\begin{lemma}[The finite-weight algebra is a level-one algebra]\label{lem:u-levelone}
\PROVED{} from \HA{prop:stabilisation}(4), \HA{thm:hidaweight}(3) and \STD{\cite[Ch.~3]{shimura}: the $T_n$ on
$S_k(\SL_2(\Z))$ are self-adjoint for the Petersson product and commute; multiplicity one at level~$1$, an eigenform is
determined by its $T_n$-eigenvalues}. Let $k\ge3$ be even, $\Ek_k$ the set of normalised eigenforms $f\in S_k(\SL_2(\Z))$ with
$v(a_\ell(f))=0$, $d_k:=|\Ek_k|$, and for $f\in\Ek_k$ let $\alpha_f\in\Qlb$ be the unit root of $Y^{2}-a_\ell(f)Y+\ell^{k-1}$.
Then:
\begin{enumerate}
\item $\hcusp_k$ is isomorphic to the $\Zl$-subalgebra of $\prod_{f\in\Ek_k}\Qlb$ generated by the tuples $(a_p(f))_f$,
  $p\ne\ell$, and $(\alpha_f)_f$ (images of $T_p$ and $U$); it equals the $\Zl$-algebra $\Talg^{\mathrm{ord}}_k$ generated by all
  $T_n$ on $e\,S_k(\SL_2(\Z);\Zl)$, $e=\lim_jT_\ell^{\,j!}$, with $U\mapsto\alpha:=$ the unique root of $Y^{2}-T_\ell Y+\ell^{k-1}$ in
  $\Talg^{\mathrm{ord}}_k$ that is a unit.
\item $\hcusp_k$ is free over $\Zl$ of rank $d_k$, $\hcusp_k\otimes\Qlb\cong\Qlb^{\Ek_k}$, and for every $p\ne\ell$ the
  characteristic polynomial of multiplication by $T_p$ on $\hcusp_k$ is
  $c_{k,p}(X):=\prod_{f\in\Ek_k}\bigl(X-a_p(f)\bigr)\in\Zl[X]$, the characteristic polynomial of $T_p$ on $e\,S_k$.
\item $h^{1}_{\mathrm{ord}}(n)=2d_k+1$ and $\chi^{\mathrm{ord}}_{n,p}=\bigl(X-(1+p^{k-1})\bigr)\,c_{k,p}(X)^{2}$ for $n=k-2$ and
  $p\ne\ell$, with $h^{1}_{\mathrm{ord}}$, $\chi^{\mathrm{ord}}$ as in \HA{def:hidatower}.
\end{enumerate}
\end{lemma}

\begin{proof}
(1) By \HA{thm:hidaweight}(3) / \HA{prop:stabilisation}(4) the ordinary part at level $\ell$ in weight $k\ge3$ consists
of the ordinary $\ell$-stabilisations $f_\alpha(z)=f(z)-\ell^{k-1}\alpha_f^{-1}f(\ell z)$, $f\in\Ek_k$, one for each $f$, with
$a_p(f_\alpha)=a_p(f)$ for $p\ne\ell$ and $U f_\alpha=\alpha_f f_\alpha$; the $\ell$-new forms of weight $\ge3$ have
$a_\ell^{2}=\ell^{k-2}$ and are not ordinary. The $f_\alpha$ are a basis of eigenvectors, and their eigensystems are pairwise
distinct (multiplicity one), so the algebra generated by commuting operators diagonal in that basis is its image in
$\prod_f\Qlb$ (a torsion-free algebra acting faithfully on a module on which it is diagonalisable embeds into the product of its
eigencharacters). On the level-one side $\Talg^{\mathrm{ord}}_k$ is diagonal in the basis $\Ek_k$ with the same $a_p$; its
element $T_\ell$ is a unit ($e$ projects onto the $v(a_\ell)=0$ part), $Y^{2}-T_\ell Y+\ell^{k-1}\equiv Y(Y-T_\ell)$ modulo $\ell$
has the simple root $Y=T_\ell$ in the complete semilocal ring $\Talg^{\mathrm{ord}}_k$, so Hensel gives a unique unit root
$\alpha$, whose image in each factor is $\alpha_f$; and $T_\ell=\alpha+\ell^{k-1}\alpha^{-1}$. So both algebras are the
subalgebra of $\prod_f\Qlb$ generated by the same tuples. (2) $\hcusp_k$ is a finitely generated torsion-free $\Zl$-module, of
rank $\dim_{\Ql}(\hcusp_k\otimes\Ql)$; the latter is a subalgebra of $\prod_f\Qlb$ stable under $\mathrm{Gal}(\Qlb/\Ql)$, which
separates the $f$ (distinct eigensystems), so $\hcusp_k\otimes\Qlb\cong\Qlb^{\Ek_k}$; the characteristic polynomial of
multiplication by $T_p$ is then $\prod_f(X-a_p(f))$. (3) is \HA{thm:hidaweight}'s consequence ($h^{1}_{\mathrm{ord}}=2\rank
S^{\mathrm{ord}}+1$), with Eichler--Shimura (\HA{obl:oracle}): $\eord H(n)\otimes\Qlb$ is the Eisenstein line, with $T_p$-eigenvalue
$1+p^{k-1}$ (\HA{Hi5}), plus two copies of $e\,S_k\otimes\Qlb$.
\end{proof}

\begin{definition}[The $\Iw$-adic Eisenstein character]\label{def:u-eischar}
For a prime $p\ne\ell$ write $\langle p\rangle=u^{s(p)}$, $s(p)\in\Zl$, and put
\[
  \ee_p(X):=1+\omega(p)^{a-1}\langle p\rangle^{-1}(1+X)^{s(p)}\in\Iw,\qquad \ee_\ell(X):=1 .
\]
\end{definition}

\begin{lemma}\label{lem:u-eischar}
\PROVED{} For $k$ in the class $a$ and $p\ne\ell$: $\ee_p(u^{k}-1)=1+p^{k-1}$, the $T_p$-eigenvalue of the weight-$k$
Eisenstein series; and $1$ is the $U$-eigenvalue of its ordinary $\ell$-stabilisation.
\end{lemma}

\begin{proof}
$(1+X)^{s(p)}$ at $X=u^{k}-1$ is $u^{ks(p)}=\langle p\rangle^{k}$, so the value is
$1+\omega(p)^{a-1}\langle p\rangle^{k-1}=1+\omega(p)^{k-1}\langle p\rangle^{k-1}=1+p^{k-1}$, as $\omega(p)^{a-1}=\omega(p)^{k-1}$
($\omega$ has order dividing $\ell-1$). The Eisenstein series $E_k$ has $a_\ell=1+\ell^{k-1}$, roots $1$ and $\ell^{k-1}$ of
$Y^{2}-a_\ell Y+\ell^{k-1}$; the unit root is $1$ (\HA{rmk:hidacheckable}, \HA{Hi5}).
\end{proof}

\begin{definition}[The ordinary Hecke algebra of the class, with its Eisenstein component]\label{def:u-T}
$\Tu(a)$ is the $\Iw$-subalgebra of $\hcusp(a)\times\Iw$ generated by $(T(p),\ee_p)$, $p\ne\ell$, and $(U,1)$; $\Tu:=\prod_a\Tu(a)$.
For $k\ge2$ in the class, $\Tu_k(a)$ is the $\Zl$-subalgebra of $\hcusp_k\times\Zl$ generated by $(T_p,1+p^{k-1})$ and $(U,1)$.
$P_\Eis:=\ker(\Tu(a)\to\Iw)$ and $P_\cusp:=\ker(\Tu(a)\to\hcusp(a))$ (the two projections); $I_{\Eis,k}$, $I_{\cusp,k}$ likewise
for $\Tu_k(a)$.
\end{definition}

\begin{remark}[Relation to Hida's algebra of modular forms]\label{rmk:u-hidaM}
Hida's ordinary algebra of $\Iw$-adic \emph{modular} forms of tame level~$1$ has an Eisenstein component in every class
(\LIT{\cite[Thm.~7.2, p.~272]{hida_ens}, verified against the read-only PDF: for primitive $\chi,\psi$ of conductors prime to
$p$ there is a unique irreducible component isomorphic to the fraction field of $\Iw$ carrying the ordinary stabilisations of
$E_k(\chi\omega^{a-k},\psi)$, $k\ge3$}; here $\chi=\psi=1$). Its eigencharacter is $(\ee_p)_p$ by Lemma~\ref{lem:u-eischar}, and
$\Tu(a)$ is its image in $\hcusp(a)\times\Iw$. Nothing below uses the identification; $\Tu(a)$ is defined directly, and by
Lemma~\ref{lem:u-levelone} $\Tu_k(a)$ is the level-one algebra generated by the $T_n$ on $e\,M_k(\SL_2(\Z);\Zl)$
($E_k$ is ordinary, $a_\ell(E_k)=1+\ell^{k-1}$).
\end{remark}

%----------------------------------------------------------------------
\section{Components}\label{sec:u-comp}
%----------------------------------------------------------------------

\begin{lemma}[The Eisenstein eigencharacter is not cuspidal]\label{lem:u-deligne}
\PROVED{} from Lemma~\ref{lem:u-levelone} and \STD{Deligne's bound (the Ramanujan--Petersson conjecture, from the Weil conjectures): for a normalised
eigenform $f\in S_k(\SL_2(\Z))$ and every prime $p$, every complex absolute value of $a_p(f)$ is $\le2p^{(k-1)/2}$, \cite{deligne_formes} and
Deligne's proof of the Weil conjectures (\VERIFYloc)}. Let $k\ge3$ lie in the class $a$. For every $p\ne\ell$, $c_{k,p}(1+p^{k-1})\ne0$; hence
$(0,c_{k,p}(1+p^{k-1}))\in\Tu_k(a)$ is non-zero, the ideal $J_k=\{x:(0,x)\in\Tu_k(a)\}$ is non-zero, $L_k<\infty$ (Proposition~\ref{prop:u-eisspec}),
the Eisenstein eigencharacter of $\Tu_k(a)$ is distinct from every cuspidal one, and $\Tu_k(a)\subseteq\hcusp_k\times\Zl$ has finite index, so
$\rank_{\Zl}\Tu_k(a)=d_k+1$. Passing to $\Iw$: $J:=\{x\in\Iw:(0,x)\in\Tu(a)\}\ne0$.
\end{lemma}

\begin{proof}
$2p^{(k-1)/2}<1+p^{k-1}$ as $(p^{(k-1)/2}-1)^{2}>0$, so no $a_p(f)$, $f\in\Ek_k$, equals $1+p^{k-1}$ and $c_{k,p}(1+p^{k-1})=\prod_f(1+p^{k-1}-a_p(f))\ne0$
(Lemma~\ref{lem:u-levelone}(2)). The element $c_{k,p}(T_p)\in\Tu_k(a)$ is $(0,c_{k,p}(1+p^{k-1}))$ by Cayley--Hamilton on $\hcusp_k$ (the upper bound of
Theorem~\ref{thm:u-sandwich}). So $J_k\ni c_{k,p}(1+p^{k-1})\ne0$, $\ell^{L_k}\Zl=J_k$ gives $L_k<\infty$, and $\Tu_k(a)\supseteq I_{\cusp,k}=\{0\}\times J_k$ and $\Tu_k(a)/I_{\cusp,k}\cong\hcusp_k$, so the index of $\Tu_k(a)$ in $\hcusp_k\times\Zl$ is
$|\Zl/J_k|<\infty$. For $\Iw$: the element $c(T(p))\in\Tu(a)$ with $c$ the characteristic polynomial of multiplication by $T(p)$ on the free $\hcusp(a)$
(Theorem~\ref{thm:u-control}) is $(0,c(\ee_p))$, and $c(\ee_p)$ specialises at $\Pk{k}$ to $c_{k,p}(1+p^{k-1})\ne0$, so $c(\ee_p)\ne0$.
(Without Deligne, the weaker ``$a_p(f)=1+p^{k-1}$ for \emph{all} $p$ is impossible'' would not give $c_{k,p}(1+p^{k-1})\ne0$ at a single $p$.)
\end{proof}

\begin{proposition}[Components of $\Tu(a)$]\label{prop:u-components}
\PROVED{} from Theorem~\ref{thm:u-control}, Lemma~\ref{lem:u-levelone} and Lemma~\ref{lem:u-deligne}.
\begin{enumerate}
\item $\hcusp(a)$ and $\Tu(a)$ are reduced, finite and torsion-free over $\Iw$; $\Tu(a)\otimes_\Iw\mathrm{Frac}(\Iw)=
  \bigl(\hcusp(a)\otimes\mathrm{Frac}(\Iw)\bigr)\times\mathrm{Frac}(\Iw)$ is a finite product of fields.
\item The minimal primes of $\Tu(a)$ (its \emph{components}) are $P_\Eis$ --- the \textbf{Eisenstein component},
  $\Tu(a)/P_\Eis=\Iw$ --- and the preimages of the minimal primes of $\hcusp(a)$ (the \textbf{cuspidal components}). Each
  component $\Tu(a)/P$ is a domain, finite over $\Iw$, containing $\Iw$.
\item \textbf{(Disjointness at arithmetic weights.)} For $k\ge3$ in the class, the ring homomorphisms $\Tu(a)\to\Qlb$
  vanishing on $\Pk{k}$ are exactly the $d_k+1$ eigencharacters of $\Tu_k(a)$ (Lemma~\ref{lem:u-levelone} and $E_k$), and each lies
  on exactly one component.
\end{enumerate}
\end{proposition}

\begin{proof}
(1) $\hcusp(a)$ is free, so torsion-free. If $x\in\hcusp(a)$ is nilpotent and non-zero, write it in a $\Iw$-basis; a non-zero
coordinate $g\in\Iw$ lies in only finitely many $\Pk{k}\Iw$ (a non-zero power series over $\Zl$ has finitely many zeros in the disc
$v(X)>0$ by Weierstrass preparation, and the $u^{k}-1$ are pairwise distinct), so for some $k\ge3$ in the class the image of $x$ in
$\hcusp(a)/\Pk{k}\hcusp(a)=\hcusp_k$ is a non-zero nilpotent, which Lemma~\ref{lem:u-levelone}(2) excludes. $\Tu(a)$ is a
subalgebra of the reduced torsion-free $\hcusp(a)\times\Iw$. It contains $\{0\}\times J$ with $J\ne0$ (Lemma~\ref{lem:u-deligne}) and maps onto
$\hcusp(a)$, so $\Tu(a)\otimes\mathrm{Frac}(\Iw)\supseteq(0,1)\cdot\mathrm{Frac}(\Iw)$ and hence equals the full product
$\bigl(\hcusp(a)\otimes\mathrm{Frac}(\Iw)\bigr)\times\mathrm{Frac}(\Iw)$; a reduced finite algebra over a field is a product of fields.
(2) The minimal primes of a reduced subalgebra $\Tu\subseteq A\times B$ with both projections surjective ($A=\hcusp(a)$: it is
generated by the $T(p)$ and $U$; $B=\Iw$: it contains $1$) are the preimages of those of $A$ and of $B$: they correspond to the
factors of $\Tu\otimes\mathrm{Frac}(\Iw)=(A\otimes\mathrm{Frac})\times\mathrm{Frac}(\Iw)$ (a torsion-free $\Iw$-algebra's minimal primes
contain no non-zero element of $\Iw$ and are the contractions of those of its generic fibre). $\Tu/P$ is torsion-free and finite over
$\Iw$ and meets $\Iw$ in $0$. (3) A homomorphism $\lambda\colon\Tu(a)\to\Qlb$ killing $\Pk{k}$ factors through $\Tu(a)/\Pk{k}\Tu(a)$,
whose map onto $\Tu_k(a)$ (the image of the generators) has a torsion kernel --- both are finite over $\Zl$ of the same rank
$d(a)+1=d_k+1$ (Theorem~\ref{thm:u-control} and (1) for the source, Lemma~\ref{lem:u-deligne} for the target; the $\Zl$-rank of $M/\Pk{k}M$ is the generic $\Iw$-rank of a finite torsion-free $M$, as
$M_{(\Pk{k})}$ is free over the discrete valuation ring $\Iw_{(\Pk{k})}$) --- and $\lambda$ kills torsion; so $\lambda$ is an eigencharacter of $\Tu_k(a)$, and all
$d_k+1$ of them occur. For uniqueness, localise at the height-one prime $\Pk{k}$: $\Tu(a)_{(\Pk{k})}$ is finite and torsion-free, hence
free, over the discrete valuation ring $\Iw_{(\Pk{k})}$, and its special fibre is $\Tu_k(a)\otimes\Ql$ (the torsion kernel dies), a
reduced finite $\Ql$-algebra, hence étale. So $\Tu(a)_{(\Pk{k})}$ is étale over a regular ring, hence normal, and a normal
noetherian ring is a product of domains: each maximal ideal --- each Galois orbit of such $\lambda$ --- contains exactly one minimal
prime.
\end{proof}

\begin{definition}[Congruence modules]\label{def:u-congmod}
For disjoint sets $\mathcal S,\mathcal S'$ of components with intersections $P_{\mathcal S}$, $P_{\mathcal S'}$ put
$\Cg(\mathcal S,\mathcal S'):=\Tu(a)/(P_{\mathcal S}+P_{\mathcal S'})$ --- the ring of the scheme-theoretic intersection of the two
unions. The \textbf{Eisenstein--cuspidal congruence module} is $\Cg_\Eis(a):=\Tu(a)/(P_\Eis+P_\cusp)$, and at weight $k$,
$\Cg_{\Eis,k}(a):=\Tu_k(a)/(I_{\Eis,k}+I_{\cusp,k})$. Two components \emph{meet} iff their congruence module is non-zero, i.e.\ they
share a maximal ideal. For a component $P$ and $k\ge3$ in the class,
\[
  I_{P,k}:=\ker\Bigl(\Tu_k(a)\to\prod_{\lambda\in P(k)}\Qlb\Bigr),
\]
$P(k)$ the eigencharacters of $\Tu_k(a)$ lying on $P$ (Proposition~\ref{prop:u-components}(3): the weight-$k$ eigenforms of $P$) --- the \emph{kernel}
reading, not the image of $P$ in $\Tu_k(a)$ (the two differ by torsion in general). $I_{\Eis,k}$ and $I_{\cusp,k}$ are the cases $P=P_\Eis$ and
$P=P_\cusp$ (the union of the cuspidal components), since $\Tu_k(a)\subseteq\hcusp_k\times\Zl$ embeds into the product of its eigencharacters.
\end{definition}

\begin{lemma}[Fibre products]\label{lem:u-fibre}
\PROVED{} Let $\Tu\subseteq A\times B$ be a subring with both projections $\pi_A,\pi_B$ surjective, $I_A=\ker\pi_A$, $I_B=\ker\pi_B$,
$J_B:=\pi_B(I_A)=\{b:(0,b)\in\Tu\}$. Then $(A\times B)/\Tu\cong B/J_B\cong\Tu/(I_A+I_B)$, and $\Tu=A\times_{B/J_B}B$.
\end{lemma}

\begin{proof}
$\beta\colon A\to B/J_B$, $\pi_A(t)\mapsto\pi_B(t)+J_B$ is well defined ($\pi_A(t)=\pi_A(t')$ gives $t-t'\in I_A$) and a ring map;
$\psi(a,b):=b-\beta(a)$ is surjective onto $B/J_B$ with kernel $\Tu$: if $b-\beta(a)\in J_B$, pick $t$ with $\pi_A(t)=a$, then
$(0,b-\pi_B(t))\in\Tu$ and $(a,b)=t+(0,b-\pi_B(t))$. Finally $\pi_B$ induces $\Tu/I_B\cong B$, carrying $(I_A+I_B)/I_B$ to $J_B$.
\end{proof}

\begin{proposition}[The Eisenstein--cuspidal module specialises]\label{prop:u-eisspec}
\PROVED{} from Theorem~\ref{thm:u-control} and Lemma~\ref{lem:u-fibre}. $\Cg_\Eis(a)\cong\Iw/J$ for the ideal
$J=\{x\in\Iw:(0,x)\in\Tu(a)\}$, a cyclic $\Iw$-module; $\Cg_{\Eis,k}(a)\cong\Zl/\ell^{L_k}$ for an $L_k=L_k(a)\in\{0,1,\dots\}$ ($L_k<\infty$ at $k\ge3$ by Lemma~\ref{lem:u-deligne}),
and for every $k\ge2$ in the class
\[
  \Cg_{\Eis,k}(a)\;\cong\;\Cg_\Eis(a)/\Pk{k}\,\Cg_\Eis(a).
\]
In particular the Eisenstein component meets the cuspidal locus ($\Cg_\Eis(a)\ne0$) iff $L_k\ge1$ for one, equivalently every, $k$ in
the class.
\end{proposition}

\begin{proof}
Lemma~\ref{lem:u-fibre} for $\Tu(a)\subseteq\hcusp(a)\times\Iw$ gives $0\to\Tu(a)\to\hcusp(a)\times\Iw\to\Cg_\Eis(a)\to0$ with
$\Cg_\Eis(a)=\Iw/J$. Tensoring with $\Iw/\Pk{k}$ is right exact:
$\Tu(a)/\Pk{k}\to\hcusp_k\times\Zl\to\Cg_\Eis(a)/\Pk{k}\to0$ (Theorem~\ref{thm:u-control}), and the image of the first map is $\Tu_k(a)$.
So $\Cg_\Eis(a)/\Pk{k}\cong(\hcusp_k\times\Zl)/\Tu_k(a)$, which is $\Cg_{\Eis,k}(a)$ by Lemma~\ref{lem:u-fibre} at weight $k$ (both
projections of $\Tu_k(a)$ are onto). It is a cyclic quotient of $\Iw/\Pk{k}=\Zl$. Last sentence: Nakayama ($\Pk{k}\subseteq\mathfrak m_\Iw$).
\end{proof}

%----------------------------------------------------------------------
\section{The finite-level reading: what one cell determines, and what it only bounds}\label{sec:u-finite}
%----------------------------------------------------------------------

\begin{lemma}[The tower lemma]\label{lem:u-tower}
\PROVED{} For $k,k'\ge2$ in one class and $m\ge1$: $\Pk{k}\Iw+\ell^{m}\Iw=\Pk{k'}\Iw+\ell^{m}\Iw$ iff $k\equiv k'\pmod{(\ell-1)\ell^{m-1}}$,
i.e.\ iff $n=k-2$ and $n'=k'-2$ lie in one $\mathrm{Tw}_m$ (\HA{def:hidatower}). Consequently, for every finitely generated
$\Iw$-module $M$ (for instance $\hcusp(a)$, $\Tu(a)$, $\Cg_\Eis(a)$, any $\Tu(a)/\mathfrak a$), the module $M/(\Pk{k},\ell^{m})M$ --- with
its algebra structure and its named operators, when $M$ carries them --- depends only on the tower of level $m$ through $k$.
\end{lemma}

\begin{proof}
$\Pk{k}-\Pk{k'}=u^{k'}-u^{k}=u^{k}(u^{k'-k}-1)$ and $v(u^{x}-1)=1+v(x)$ for $u=1+\ell$, $\ell$ odd. If $\ell^{m-1}\mid k-k'$ the two
ideals contain each other. Conversely, if they are equal then $u^{k}-u^{k'}\in(\Pk{k},\ell^{m})\cap\Zl=\ell^{m}\Zl$ (the quotient by
$(\Pk{k},\ell^{m})$ is $\Z/\ell^{m}$ via $X\mapsto u^{k}-1$). Within one class $k\equiv k'\pmod{\ell-1}$ already.
\end{proof}

\begin{proposition}[What one cell determines]\label{prop:u-determines}
\PROVED{} from Theorem~\ref{thm:u-control}, Lemma~\ref{lem:u-levelone}, Proposition~\ref{prop:u-eisspec} and Lemma~\ref{lem:u-tower}.
Let $k,k'\ge3$ lie in one tower of level $m$.
\begin{enumerate}
\item $\hcusp_k/\ell^{m}\cong\hcusp_{k'}/\ell^{m}$ as $\Z/\ell^{m}$-algebras, $T_p\mapsto T_p$, $U\mapsto U$. Hence
  $c_{k,p}\equiv c_{k',p}\pmod{\ell^{m}}$ for every $p\ne\ell$ (Lemma~\ref{lem:u-levelone}(2): the characteristic polynomial of
  multiplication by $T_p$ on the free module $\hcusp_k/\ell^{m}$), and with Lemma~\ref{lem:u-levelone}(3) and $1+p^{k-1}\equiv1+p^{k'-1}
  \pmod{\ell^{m}}$: $\chi^{\mathrm{ord}}_{n,p}\equiv\chi^{\mathrm{ord}}_{n',p}\pmod{\ell^{m}}$ and $h^{1}_{\mathrm{ord}}(n)=h^{1}_{\mathrm{ord}}(n')$
  --- the statements of \HA{Hi3}, \HA{Hi4}, \HA{Hi7}, \HA{Hi9} of [C], derived here for $n\ge1$ along the chain Hida
  (forms) $\to$ stabilisation $\to$ Eichler--Shimura, which runs through characteristic polynomials only and so never meets the
  lattice question of \HA{rmk:stabconvention}(3). (\HA{Hi8}, the lattice datum, is \emph{not} derived: it is lattice-dependent.)
\item At $m=1$: the decomposition of $\hcusp_k/\ell$ into local factors, their residue fields, Hilbert functions and Loewy lengths
  depend only on the class.
\item $\min(L_k,m)=\min(L_{k'},m)$ (Proposition~\ref{prop:u-eisspec}: $\Cg_{\Eis,k}/\ell^{m}=\Cg_\Eis/(\Pk{k},\ell^{m})$ is
  $\Z/\ell^{\min(L_k,m)}$).
\end{enumerate}
\end{proposition}

\begin{proposition}[What one cell only bounds]\label{prop:u-bounds}
\PROVED{} from Proposition~\ref{prop:u-components}. Let $k\ge3$ lie in the class $a$.
\begin{enumerate}
\item \textbf{Components.} Distinct cuspidal components carry disjoint non-empty sets of weight-$k$ eigenforms, so
  $\#\{\text{cuspidal components}\}\le\#\{\mathrm{Gal}(\Qlb/\Ql)\text{-orbits on }\Ek_k\}$, and a component of $\Iw$-rank $r$ carries $r$
  eigenforms. If $c_{k,p}$ is squarefree for some $p$, the orbits are the $\Ql$-irreducible factors of $c_{k,p}$, whose number bounds the
  number of components from above. From below, the number of maximal ideals of $\hcusp(a)$ (= of $\hcusp_k/\ell$) is at least the number of
  distinct irreducible factors of $\bar c_{k,p}\in\Fl[X]$, and each maximal ideal lies on at least one component.
\item \textbf{Intersections.} For components $P\ne P'$ meeting at weight $k$ and $t\in\Tu_k(a)$, let $\chi_P(t)$, $\chi_{P'}(t)$ be the
  characteristic polynomials of $t$ on the $\Zl$-span of the eigenforms of $P$, $P'$ at weight $k$, and $I_{P,k}$, $I_{P',k}$ the kernels of
  Definition~\ref{def:u-congmod}. Then
  $\mathrm{length}_{\Zl}\,\Tu_k/(I_{P,k}+I_{P',k})\le v\bigl(\mathrm{Res}(\chi_P(t),\chi_{P'}(t))\bigr)$ for every $t$.
\end{enumerate}
\end{proposition}

\begin{proof}
(1) Proposition~\ref{prop:u-components}(3); $\Qlb$-points of a component above $\Pk{k}$ come in Galois orbits, and a component of rank $r$ has
$r$ of them counted with multiplicity, all simple there by the étaleness in the proof of (3). If $c_{k,p}$ is squarefree, $T_p$ separates
$\Ek_k$ and the orbits are the root sets of the irreducible factors. Distinct roots of $\bar c_{k,p}$ are distinct values of $T_p$ modulo
distinct maximal ideals. (2) Put $O':=\Tu_k/I_{P',k}$, a $\Zl$-order in the étale algebra $\prod_{f\in P'}\Qlb$ of rank $r'$. The element
$\chi_P(t)\in\Tu_k$ lies in $I_{P,k}$ (Cayley--Hamilton on $\Tu_k/I_{P,k}$), so its image in $O'$ lies in the image $J'$ of $I_{P,k}$, and
$\Tu_k/(I_{P,k}+I_{P',k})=O'/J'$ is a quotient of $O'/\chi_P(t)O'$, of length $v(\det(\chi_P(t)|O'))=v\bigl(\prod_{f\in P'}\chi_P(t)(f)\bigr)
=v(\mathrm{Res})$.
\end{proof}

\begin{theorem}[The étale criterion]\label{thm:u-etale}
\PROVED{} from Theorem~\ref{thm:u-control}, Lemma~\ref{lem:u-levelone}, Lemma~\ref{lem:u-tower}, Proposition~\ref{prop:u-components} and (for (3))
Theorem~\ref{thm:u-sandwich}. Suppose that for one $k\ge3$ in the
class $a$ and one prime $p\ne\ell$ the reduction $\bar c_{k,p}\in\Fl[X]$ is squarefree, with irreducible factors of degrees
$\delta_1,\dots,\delta_r$. Then:
\begin{enumerate}
\item $\hcusp(a)\cong\prod_{i=1}^{r}W(\F_{\ell^{\delta_i}})[[X]]$ is finite étale over $\Iw$; its components are these $r$ factors, pairwise
  disjoint, and it is generated over $\Iw$ by $T(p)$.
\item For every $k'\ge3$ in the class, $\hcusp_{k'}=\Zl[T_p]\cong\Zl[X]/(c_{k',p})$, étale over $\Zl$, and $\bar c_{k',p}=\bar c_{k,p}$.
\item If moreover $\bar c_{k,p}(1+p^{k-1})\ne0$ in $\Fl$, then $\Cg_\Eis(a)=0$: $\Tu(a)=\hcusp(a)\times\Iw$, and the Eisenstein component meets no
  cuspidal one. If instead $1+p^{k-1}$ is a root of $\bar c_{k,p}$ \textbf{and $\Cg_\Eis(a)\ne0$} (equivalently $L_k\ge1$; sufficient:
  $v(B_k/k)\ge1$ with $k\ge4$, $a\not\equiv0$, by Theorem~\ref{thm:u-sandwich}), it is a simple root and the Eisenstein component meets exactly one
  cuspidal component, the one of residue degree~$1$ carrying that root. (A root at one $p$ alone does not suffice: it says $a_p(f)\equiv1+p^{k-1}$ for
  one $f$, and nothing about the $a_q(f)$, $q\ne p$.)
\end{enumerate}
\end{theorem}

\begin{proof}
Let $d=d_k=\deg c_{k,p}$. The subalgebra $\Fl[\bar T_p]\subseteq\hcusp_k/\ell$ has dimension $d$: if $\mu(\bar T_p)=0$ for a proper divisor $\mu$ of
$\bar c_{k,p}$, lift $\mu$ to $\tilde\mu\in\Zl[X]$; then $\tilde\mu(T_p)\in\ell\hcusp_k$, so $v(\tilde\mu(a_p(f)))\ge1$ for all $f\in\Ek_k$; but a root
$\bar a_p(f)$ of the (squarefree) cofactor is not a root of $\mu$. As $\dim_{\Fl}\hcusp_k/\ell=d$ (Lemma~\ref{lem:u-levelone}(2)),
$\hcusp_k/\ell=\Fl[\bar T_p]\cong\Fl[X]/(\bar c_{k,p})$, a product of the fields $\F_{\ell^{\delta_i}}$. By Theorem~\ref{thm:u-control} this is
$\hcusp(a)/\mathfrak m_\Iw\hcusp(a)$. (1) By Nakayama $\hcusp(a)=\Iw[T(p)]$; $\Omega_{\hcusp(a)/\Iw}\otimes\hcusp(a)/\mathfrak m_\Iw=\Omega_{(\hcusp/\mathfrak
m_\Iw)/\Fl}=0$, so $\Omega=0$ by Nakayama, and $\hcusp(a)$ is free and unramified, i.e.\ finite étale, over the complete local $\Iw$; the idempotents of
the special fibre lift, and a local finite étale $\Iw$-algebra with residue field $\F_q$ is $W(\F_q)\otimes_{\Zl}\Iw=W(\F_q)[[X]]$, a domain.
(2) $\hcusp_{k'}=\hcusp(a)/\Pk{k'}$ is then étale over $\Zl$ with special fibre $\hcusp(a)/\mathfrak m_\Iw$; the same Nakayama argument gives
$\hcusp_{k'}=\Zl[T_p]$, and $\bar c_{k',p}=\bar c_{k,p}$ by Proposition~\ref{prop:u-determines}(1). (3) If $\bar e:=1+p^{k-1}\bmod\ell$ is not a root,
$\bar c_{k,p}(\bar e)\ne0$, so $L_k\le v(c_{k,p}(1+p^{k-1}))=0$ by Theorem~\ref{thm:u-sandwich} below (its upper bound needs only
Lemma~\ref{lem:u-levelone}), and $\Cg_\Eis=0$ by Proposition~\ref{prop:u-eisspec}. If it is a root it is simple (squarefreeness). Assume $\Cg_\Eis(a)\ne0$. Then $P_\Eis+P_\cusp\ne\Tu(a)$, and the maximal ideal of $\Tu(a)$
containing $P_\Eis$ --- $\mathfrak m_\Eis=(\mathfrak m_\Iw,\,T(q)-\ee_q,\,U-1)$, the only one, as $\Tu(a)/P_\Eis=\Iw$ is local --- contains
$P_\Eis+P_\cusp=P_\Eis+\bigcap_iP_i$ ($P_i$ the cuspidal minimal primes, Proposition~\ref{prop:u-components}(2)); being prime, it contains some $P_i$.
On $\Tu(a)/P_i=W(\F_{\ell^{\delta_i}})[[X]]$ (part (1)) the image of $\mathfrak m_\Eis$ has residue field $\Fl$ and $T(p)\equiv\bar e$ there, so $\delta_i=1$ and
$P_i$ is the factor of the simple root $\bar e$; it is unique, as distinct factors carry distinct roots of the squarefree $\bar c_{k,p}$.
(Without $\Cg_\Eis(a)\ne0$ this fails: then $J=\Iw$, $(y,0)\in\Tu(a)$ for every $y\in\hcusp(a)$ by Lemma~\ref{lem:u-fibre}, and the image of
$\mathfrak m_\Eis$ in $\hcusp(a)$ is the unit ideal.)
\end{proof}

%----------------------------------------------------------------------
\section{The Eisenstein component and its intersection with the cuspidal locus}\label{sec:u-eis}
%----------------------------------------------------------------------

\begin{theorem}[The sandwich]\label{thm:u-sandwich}
\PROVED{} from Lemma~\ref{lem:u-levelone} and \STD{the $q$-expansion formulas for $T_n$ on $M_k(\SL_2(\Z))$, \cite[Ch.~3]{shimura};
$E_4,E_6\in1+q\Z[[q]]$; von Staudt--Clausen, \cite[Thm.~5.10]{washington} (\VERIFYloc)}. Let $k\ge4$ lie in a class $a\not\equiv0\pmod{\ell-1}$. Then
\[
  v\!\left(\frac{B_k}{k}\right)\;\le\;L_k\;\le\;\min_{p\ne\ell}\ v\bigl(c_{k,p}(1+p^{k-1})\bigr).
\]
The upper bound holds for every $k\ge3$ in every class.
\end{theorem}

\begin{proof}
\emph{Upper bound.} $t:=c_{k,p}(T_p)\in\Tu_k(a)$ has cuspidal component $c_{k,p}(T_p|_{\hcusp_k})=0$ (Cayley--Hamilton for multiplication by
$T_p$ on $\hcusp_k$, Lemma~\ref{lem:u-levelone}(2)) and Eisenstein component $c_{k,p}(1+p^{k-1})$. So $(0,c_{k,p}(1+p^{k-1}))\in\Tu_k(a)$, and the
ideal $J_k=\ell^{L_k}\Zl$ of Lemma~\ref{lem:u-fibre} contains $c_{k,p}(1+p^{k-1})$.

\emph{Lower bound.} Put $c:=-B_k/2k$, so $E_k=c+\sum_{n\ge1}\sigma_{k-1}(n)q^{n}\in M_k(\SL_2(\Z);\Zl)$ (as $(\ell-1)\nmid k$, $B_k$ is $\ell$-integral), and
$\nu:=v(c)=v(B_k/k)$; assume $\nu\ge1$. With $k=4\alpha+6\beta$, $g:=E_k-cE_4^{\alpha}E_6^{\beta}\in S_k(\SL_2(\Z);\Zl)$ and $g\equiv E_k$ coefficientwise
modulo $\ell^{\nu}$. The $T_n$ preserve $M_k(\Zl)$ and $T_nE_k=\sigma_{k-1}(n)E_k$, so the image $\bar g$ of $g$ in $S_k(\Z/\ell^{\nu})$ satisfies
$T_n\bar g=\sigma_{k-1}(n)\bar g$; as $a_1(\bar g)=1$, $\Z/\ell^{\nu}\bar g$ is free of rank one and stable, and the cuspidal level-one Hecke algebra
$\Talg_k\otimes\Zl$ (generated by all $T_n$ on $S_k(\Zl)$) acts on it through a ring map $\varphi\colon\Talg_k\otimes\Zl\to\Z/\ell^{\nu}$,
$\varphi(T_n)=\sigma_{k-1}(n)$. Since $\varphi(T_\ell)=1+\ell^{k-1}$ is a unit, $\varphi(e)=\lim\varphi(T_\ell)^{j!}=1$, so $\varphi$ factors through
$e(\Talg_k\otimes\Zl)=\hcusp_k$ (Lemma~\ref{lem:u-levelone}(1)); $\varphi(U)$ is a unit root of $(Y-1)(Y-\ell^{k-1})$ in $\Z/\ell^{\nu}$, so
$\varphi(U)=1$. The two ring maps $\Tu_k(a)\to\Z/\ell^{\nu}$, $t\mapsto\varphi(t_\cusp)$ and $t\mapsto t_\Eis\bmod\ell^{\nu}$, agree on the generators
$(T_p,1+p^{k-1})$, $(U,1)$, hence everywhere; for $(0,x)\in\Tu_k(a)$ this gives $x\equiv0\pmod{\ell^{\nu}}$, i.e.\ $J_k\subseteq\ell^{\nu}\Zl$.
\end{proof}

\begin{corollary}[The Eisenstein locus is the irregular locus]\label{cor:u-irregular}
\PROVED{} from Theorem~\ref{thm:u-sandwich}, Proposition~\ref{prop:u-eisspec} and \STD{Kummer's congruence $B_k/k\equiv B_{k'}/k'\pmod\ell$ for
$k\equiv k'\not\equiv0\pmod{\ell-1}$, \cite[Cor.~5.14, p.~61]{washington1st} (verified, L16-3); the weight congruence of mod-$\ell$ modular forms, over $\overline{\F}_\ell$: for $\ell\ge5$ the algebra of mod-$\ell$ modular forms of level~$1$
is $\overline{\F}_\ell[E_4,E_6]/(E_{\ell-1}-1)$, graded modulo $\ell-1$, so if $f\in M_k(\overline{\F}_\ell)$, $f'\in M_{k'}(\overline{\F}_\ell)$ have the same
non-zero $q$-expansion then $k\equiv k'\pmod{\ell-1}$, \cite[Thm.~2(iv), p.~22]{swinnerton_dyer}, \cite[n$^\circ$~1.2, Th.~2, p.~416-04]{serre_bourbaki416} (verified, slot-16 L16-3 / WP-8)}. For a class $a\not\equiv0\pmod{\ell-1}$ the Eisenstein component meets the cuspidal locus iff $v(B_k/k)\ge1$ for one (equivalently
every) $k\ge4$ in the class, i.e.\ iff $(\ell,a_0)$ is an irregular pair ($\ell\mid B_{a_0}$), $a_0\in\{2,\dots,\ell-3\}$ the representative of $a$.
\end{corollary}

\begin{proof}
If $v(B_k/k)\ge1$ then $L_k\ge1$ by Theorem~\ref{thm:u-sandwich}. Conversely, if $\Cg_\Eis(a)\ne0$ then $L_k\ge1$ (Proposition~\ref{prop:u-eisspec}), so a
maximal ideal of $\Tu_k(a)$ contains $I_{\Eis,k}+I_{\cusp,k}$: some $f\in\Ek_k$ and a prime $\lambda$ of $\Qlb$ above $\ell$ with $a_p(f)\equiv1+p^{k-1}$ for
all $p\ne\ell$ and $\alpha_f\equiv1$, hence $a_\ell(f)=\alpha_f+\ell^{k-1}\alpha_f^{-1}\equiv1+\ell^{k-1}$, modulo $\lambda$. Both $(a_n(f))$ and
$(\sigma_{k-1}(n))$ are determined by their values at primes through the same recursions, so $a_n(f)\equiv\sigma_{k-1}(n)$ for all $n\ge1$ and
$E_k-f\equiv c\pmod\lambda$ as $q$-series. If $v(c)=0$, then reducing modulo $\lambda$, as forms over the residue field $\F_\lambda\subseteq\overline{\F}_\ell$ (the coefficients of
$f$ lie in the ring of integers $\mathcal O_\lambda$ of a finite extension of $\Ql$), the constant $\bar c\ne0$ would be a mod-$\lambda$ modular form of weight
$k$ with the $q$-expansion of the weight-$0$ form $\bar c$, forcing $k\equiv0\pmod{\ell-1}$. So $v(B_k/k)=v(c)\ge1$ at this $k$, hence at every $k$ of the class by Kummer.
\end{proof}

\begin{lemma}[Emerton's objects in the notation of {[U]}]\label{lem:u-emerton}
\PROVED{} from \LIT{\cite{emerton_eis}, library copy (an $8$-page preprint), read in place: \S1, pp.~2--3 (the set-up), \S3, pp.~5--6
(ordinary parts, Theorem~5), p.~7 (the shift against Mazur--Wiles); the identifications with Hida's objects are Emerton's own, p.~3 and p.~5,
checked against \cite[(1.12), (1.14)--(1.19), \S3 pp.~249--250]{hida_ens} (read-only PDF), and closed by \cite[(1.2), p.~9; (1.3), p.~10; Thm.~1.3,
p.~17; Thm.~2.2, p.~20]{hida_aif2} (verified, read-only PDF; item (2))} and \STD{Kubota--Leopoldt interpolation,
\cite[Thm.~5.11, p.~57]{washington1st} (verified on the page images, first edition, slot-16 WP-8; re-pointed from the second edition's key at
v0.2.1-pre.1)}.
Let $a$ be an even class, $a\not\equiv0\pmod{\ell-1}$, and write Emerton's $p$ as $\ell$.\footnote{Emerton's ``elementary description'' of $V$
\cite[p.~2]{emerton_eis} (the $\ell$-adic completion of $M\cap\Zl[[q]]$, $M$ the graded ring of level-one forms over $\Ql$) is \emph{not used} in [U], and is
inaccurate as literally stated (the integral $q$-series $\sum_{n\ge1}\sigma_{\ell-2}(n)q^{n}$ of $G_{\ell-1}-a_0(G_{\ell-1})\in M\cap\Zl[[q]]$ does not reduce to a
mod-$\ell$ modular form of level one; slot-14 WP-U1 audit, \S3). The objects used are Katz's geometric $V$ and $V_{\mathrm{par}}$, identified with Hida's
spaces in item~(2); the correct elementary description is the completion of $\bigoplus_k M_k(\Zl)$.}
\begin{enumerate}
\item \textbf{(Components and $\Iw$.)} Emerton's $\Iw=\Zl[[\Zl^\times]]=\prod_{i\bmod\ell-1}\Iw_{(i)}$, with $\F_\ell^\times\cong\mu_{\ell-1}$ acting on
  $\Iw_{(i)}$ through $\omega^{i}$, and $\langle\gamma\rangle$ acting on $M_k$ as $\gamma^{k}$ (p.~3). So $\mu_{\ell-1}$ acts on $M_k$ through
  $\zeta\mapsto\zeta^{k}$ and weight $k$ lies in Emerton's component $i$ iff $k\equiv i$: \textbf{his $i$ is our class $a$}, with no shift. The ring
  isomorphism $\iota_a\colon\Iw_{(a)}\to\Iw$, $\langle\gamma\rangle\mapsto(1+X)^{s(\gamma)}$ for $\gamma=u^{s(\gamma)}\in1+\ell\Zl$, sends
  $\langle x\rangle$, $x\in\Zl^\times$, to $\omega(x)^{a}(1+X)^{s(x)}$ with $\langle x\rangle_\Gamma:=x\omega(x)^{-1}=u^{s(x)}$, and carries Emerton's
  weight-$k$ specialisation ($\langle\gamma\rangle\mapsto\gamma^{k}$) to $X\mapsto u^{k}-1$, i.e.\ to our $\Pk{k}$. Hida's $\Iw$-action is the same
  ($f|z=z_p^{k}f$ on weight $k$, \cite[(1.12)]{hida_ens}; $P_k=(1+X)-u^{k}$, \cite[p.~250]{hida_ens}).
\item \textbf{(The cuspidal algebra.)} Emerton's $\mathbf T_{\mathrm{par}}$ is the quotient of his $\mathbf T$ (the closed $\Iw$-algebra generated by
  the $T_q$, $q\ne\ell$, and $U_\ell$ on Katz's generalised $\ell$-adic modular functions of tame level~$1$) acting faithfully on the cuspidal subspace;
  he identifies it with Hida's $h$ (p.~3), and his superscript ord with Hida's superscript $0$ (p.~5). The identification is checkable at three
  locators: (i) Hida's space of $\ell$-adic modular forms of tame level~$1$ embeds in Katz's $V$, \cite[(1.11a), p.~240]{hida_ens}, and is all of it by
  Katz's density theorem (\VERIFYloc); Hida's algebra does not depend on the power of $\ell$ in the level, \cite[(1.15a), p.~243]{hida_ens}; (ii) Emerton's
  $V_{\mathrm{par}}=\{f:a_0(f|\langle x\rangle)=0\ \forall x\}$ is Hida's $S(1;\Zl)$: \cite[Thm.~2.2, p.~20]{hida_aif2}, the exact sequence
  $0\to S(N;\mathcal O)\to M(N;\mathcal O)\xrightarrow{\delta}C(S\times\Z_\ell^\times;\mathcal O)\to0$ with $\delta(f)(s,z)=a(0,f|z)$, $S$ a point at $N=1$, and
  \cite[(1.2), p.~9; (1.3), p.~10]{hida_aif2} for the independence of the $\ell$-power level (Katz); (iii) the parabolic duality behind Emerton's
  Proposition~4 (Gouv\^ea's Cor.~III.1.3, not in the library) is \cite[Thm.~1.3, p.~17]{hida_aif2}: $h(N;\mathcal O)\cong\mathrm{Hom}(S(N;\mathcal O),\mathcal O)$
  and $S(N;\mathcal O)\cong\mathrm{Hom}(h(N;\mathcal O),\mathcal O)$. Hence, through $\iota_a$,
  \[
    \mathbf T^{\mathrm{ord}}_{\mathrm{par},(a)}\;=\;\hcusp(a)\qquad(T_q\mapsto T(q),\ U_\ell\mapsto U)
  \]
  as $\Iw$-algebras (Definition~\ref{def:u-hcusp}: $\hcusp(a)=h^{0}(\ell,a)$, the component on which $\mu_{\ell-1}$ acts through $\zeta\mapsto\zeta^{a}$,
  \cite[p.~249]{hida_ens}; Hida's algebra does not depend on the power of $\ell$ in the level, \cite[(1.19)]{hida_ens}). ``Parabolic'' is Emerton's
  word for cuspidal; the ordinary algebra with the Eisenstein component is his $\mathbf T^{\mathrm{ord}}_{(a)}$, whose image in $\hcusp(a)\times\Iw$
  is our $\Tu(a)$ (Remark~\ref{rmk:u-hidaM}), and which is not needed below.
\item \textbf{(The Eisenstein ideal.)} Emerton's generator $T_q-1-\langle q\rangle q^{-1}$ ($q\ne\ell$) goes under $\iota_a$ to
  $T(q)-\bigl(1+\omega(q)^{a}(1+X)^{s(q)}\,\omega(q)^{-1}\langle q\rangle_\Gamma^{-1}\bigr)=T(q)-\ee_q$ (Definition~\ref{def:u-eischar}), and $U_\ell-1$ to
  $U-1$. So his $I^{0}_{\Eis,(a)}$, the image in $\mathbf T^{\mathrm{ord}}_{\mathrm{par},(a)}$ of the ideal generated by these elements (p.~6), is
  \[
    I^{0}(a):=\bigl(T(q)-\ee_q\ (q\ne\ell),\ U-1\bigr)\subseteq\hcusp(a).
  \]
\item \textbf{(The $L$-function.)} Emerton's $L_\ell(\omega^{a},1-s)\in\Iw_{(a)}$ is the element whose weight-$k$ specialisation is the Kubota--Leopoldt value
  $L_\ell(1-k,\omega^{a})=-(1-\ell^{k-1})B_k/k$ for $k\equiv a$, $k\ge1$ (the character is written first and $s$ is the weight). Its image under $\iota_a$ is
  therefore the $G_a$ of Prediction~\ref{pr:u-U1} (which lies in $\Iw$ for $a\not\equiv0$: \LIT{\cite[{Thm.~7.10, p.~123, with the power series $f(T,\theta)\in\mathcal O_\theta[[T]]$
  of p.~122 for $\theta=\omega^{a}\ne1$}]{washington1st}, verified, first edition}): two elements of $\Iw$ with the same values at the infinitely many points $u^{k}-1$ are equal (Weierstrass
  preparation). Emerton's ``the image of $A_{0,(a)}$ \dots\ is equal to $L_\ell(\omega^{a},1-s)$'' (p.~6) holds up to the unit $\tfrac12$: the ordinary
  Eisenstein family normalised by $a_1=1$ has constant term $\tfrac12L_\ell(1-k,\omega^{a})$ at weight $k$; only the ideal is used.
\end{enumerate}
\end{lemma}

\begin{proof}
(1) and (3) are the computations shown: $\langle q\rangle=\omega(q)\langle q\rangle_\Gamma$ in $\Zl^\times=\mu_{\ell-1}\times(1+\ell\Zl)$, $\mu_{\ell-1}$ acts on
$\Iw_{(a)}$ through $\omega^{a}$, so $\iota_a(\langle q\rangle q^{-1})=\omega(q)^{a}(1+X)^{s(q)}\omega(q)^{-1}\langle q\rangle_\Gamma^{-1}$. (2) is the cited
identification. (4): both sides specialise at $X=u^{k}-1$, $k\equiv a$, to $L_\ell(1-k,\omega^{a})$, which for $k\equiv a$ is $-(1-\ell^{k-1})B_{k,\omega^{a-k}}/k=
-(1-\ell^{k-1})B_k/k$; a non-zero element of $\Iw$ has finitely many zeros in $v(X)>0$. The constant term of $E_k-\ell^{k-1}E_k(\ell z)$ is
$(1-\ell^{k-1})(-B_k/2k)$.
\end{proof}

\paragraph{What the tag of Prediction~(U1) rests on.} Two steps of \cite{emerton_eis} are printed as sketches and are load-bearing; both are
closed by standard published results (slot-14 WP-U1 audit, \S3). \emph{Corollary~2} (p.~4, $V_{(a)}\hat\otimes R\cong\mathrm{Hom}_{\mathrm{cont}}(\mathbf T_{(a)},R)$,
``by the methods of'' Gouv\^ea \S III.1.2): its content is the finite-level duality $M_k(\Zl)\cong\mathrm{Hom}(\mathbf T_k,\Zl)$, $f\mapsto(T\mapsto a_1(f|T))$, in
weights $k\not\equiv0\pmod{\ell-1}$ --- injective as there are no constants in weight $k>0$, onto by the rank count $\mathrm{rank}\,\mathbf T_k=\dim M_k$ and
because $a_0(f)\in\Zl$ once $a_n(f)\in\Zl$ for $n\ge1$, by the weight congruence (\cite[Thm.~2(iv)]{swinnerton_dyer}, \cite[Th.~2]{serre_bourbaki416}, verified; the argument of
Corollary~\ref{cor:u-irregular}) --- together with Katz's density theorem and $q$-expansion principle (\VERIFYloc), which give
$\mathbf T_{(a)}/\ell^{m}\cong\varprojlim_k\mathbf T_k/\ell^{m}$, Pontryagin dual to $V_{(a)}/\ell^{m}$, hence the statement for every $\ell$-adically complete $R$,
$R=\Iw_{(a)}$ included. \emph{Lemma~3}'s ``easily seen'' description of $V_{\mathrm{par}}$ is \cite[Thm.~2.2]{hida_aif2} (item~(2)(ii)). So the tag rests on
Emerton's Theorem~5 and these standard results; none of them is a Gorenstein property, a freeness hypothesis on a module over the Hecke algebra, or a main
conjecture.

\begin{prediction}[(U1) The congruence length is the Bernoulli valuation]\label{pr:u-U1}
\PROVED{} from \LIT{Emerton \cite[Thm.~5, p.~6]{emerton_eis} (verified)} through Lemma~\ref{lem:u-emerton}, Lemma~\ref{lem:u-fibre} and
Proposition~\ref{prop:u-eisspec}; \textbf{no Gorenstein, freeness, $\mu$- or $\lambda$-invariant hypothesis}, and no main conjecture
(\CONJ{} through v0.1.1-pre.1; corroborated by \LIT{Ohta \cite[Thm.~(1.5.5)(II), p.~238]{ohta_cong} (verified)}, whose proof in this case uses the
Mazur--Wiles main conjecture, Remark~\ref{rmk:u-ohta}). Let $\ell\ge5$, let $a$ be an even class with $a\not\equiv0\pmod{\ell-1}$, and let $G_a\in\Iw$ be the
Iwasawa power series of the Kubota--Leopoldt function, $G_a(u^{k}-1)=L_\ell(1-k,\omega^{a})=-(1-\ell^{k-1})B_k/k$ for $k\equiv a$. Then:
\begin{enumerate}
\item[(U1-$\Iw$)] $J=G_a\Iw$ for the ideal $J$ of Proposition~\ref{prop:u-eisspec}; equivalently $\Cg_\Eis(a)\cong\Iw/(G_a)$, as rings and as $\Iw$-modules, and
  $\hcusp(a)/I^{0}(a)\cong\Iw/(G_a)$ with $I^{0}(a)$ of Lemma~\ref{lem:u-emerton}(3). In particular $J$ is principal and $\Tu(a)$ is $\Iw$-free of rank $d(a)+1$
  (Remark~\ref{rmk:u-Jprincipal}).
\item[(U1-$k$)] For \textbf{every} $k\ge2$ in the class, $\Cg_{\Eis,k}(a)\cong\Zl/\ell^{v(B_k/k)}$, i.e.\ $L_k=v(B_k/k)$; the lower bound of
  Theorem~\ref{thm:u-sandwich} is an equality. No condition on $k$ beyond $k\equiv a\pmod{\ell-1}$ (at $k=2$, i.e.\ $a\equiv2$: $L_2=v(B_2/2)=0$).
\end{enumerate}
\emph{Checkable form} (now a check of the implementation): at a computed cell, $\min_{p\in\mathcal P}v(c_{k,p}(1+p^{k-1}))\ge v(B_k/k)=L_k$ for every computed set
$\mathcal P$ (Theorem~\ref{thm:u-sandwich}), with equality once $\mathcal P$ is large enough; a value \emph{below} $v(B_k/k)$ is a defect, a value above it at every
computed $p$ only says that $\mathcal P$ is too small. \emph{Not covered}: the class $a\equiv0$ (Emerton's Proposition~1 needs $i\not\equiv0$, and $G_0$ has a
pole; Remark~\ref{rmk:u-classzero}), $\ell\in\{2,3\}$, tame level $N>1$.
\end{prediction}

\begin{proof}
\emph{Step 1: $P_\Eis$ is generated by the Eisenstein differences, and $\pi_\cusp(P_\Eis)=I^{0}(a)$.} Let $I\subseteq\Tu(a)$ be the ideal generated by
$g_q:=(T(q),\ee_q)-\ee_q\cdot1$, $q\ne\ell$, and $g_\ell:=(U,1)-1$. Each generator of the $\Iw$-algebra $\Tu(a)$ (Definition~\ref{def:u-T}) is congruent modulo $I$ to an
element of $\Iw\cdot1$, so every monomial is, and $\Tu(a)=\Iw\cdot1+I$. The composite $\Iw\to\Tu(a)/I\to\Tu(a)/P_\Eis\cong\Iw$ (the second map exists as
$I\subseteq P_\Eis$, the isomorphism is $\pi_\Eis$) is the identity, and the first map is onto; so $\Tu(a)/I\to\Tu(a)/P_\Eis$ is an isomorphism and $I=P_\Eis$. As
$\pi_\cusp\colon\Tu(a)\to\hcusp(a)$ is onto, $\pi_\cusp(P_\Eis)$ is the ideal generated by the $\pi_\cusp(g_q)=T(q)-\ee_q$ and $\pi_\cusp(g_\ell)=U-1$, i.e.\ $I^{0}(a)$.

\emph{Step 2: two readings of one kernel.} Let $\phi\colon\Iw\to\Cg_\Eis(a)=\Tu(a)/(P_\Eis+P_\cusp)$ be the structure map ($\Iw$ embeds diagonally in $\Tu(a)$). Through
$\Tu(a)/P_\Eis\cong\Iw$ ($\pi_\Eis$, the identity on $\Iw$), $(P_\Eis+P_\cusp)/P_\Eis$ is $\pi_\Eis(P_\cusp)=J$, so $\ker\phi=J$. Through $\Tu(a)/P_\cusp\cong\hcusp(a)$
($\pi_\cusp$, which is $\Iw$-linear), $(P_\Eis+P_\cusp)/P_\cusp$ is $\pi_\cusp(P_\Eis)=I^{0}(a)$ (Step~1), so $\ker\phi=\{x\in\Iw:x\cdot1_{\hcusp(a)}\in I^{0}(a)\}$, the
pull-back of $I^{0}(a)$ along the structure map $\Iw\to\hcusp(a)$. (This is Lemma~\ref{lem:u-fibre} read from both sides; $\hcusp(a)=0$ is allowed, and then
$J=\Iw$.)

\emph{Step 3: Emerton.} By Lemma~\ref{lem:u-emerton}(1)--(3) the pull-back in Step~2 is, through $\iota_a$, the pull-back of $I^{0}_{\Eis,(a)}$ along
$\Iw_{(a)}\to\mathbf T^{\mathrm{ord}}_{\mathrm{par},(a)}$, which \cite[Thm.~5]{emerton_eis} computes as the principal ideal generated by $L_\ell(\omega^{a},1-s)$ (hypotheses: $\ell\ge5$,
$a\not\equiv0\pmod{\ell-1}$). By Lemma~\ref{lem:u-emerton}(4) that ideal is $G_a\Iw$. So $J=G_a\Iw$, $\Cg_\Eis(a)\cong\Iw/J=\Iw/(G_a)$ as rings (by $\phi$, which is onto), and
$\hcusp(a)/I^{0}(a)\cong\Cg_\Eis(a)$ by Step~2. Freeness: $0\to\{0\}\times J\to\Tu(a)\to\hcusp(a)\to0$ is exact ($P_\cusp=\{0\}\times J$) and splits, $\hcusp(a)$ being
$\Iw$-free (Theorem~\ref{thm:u-control}); $J=G_a\Iw\cong\Iw$ as $G_a\ne0$ ($J\ne0$, Lemma~\ref{lem:u-deligne}).

\emph{Step 4: weight $k$.} Proposition~\ref{prop:u-eisspec} (valid for every $k\ge2$ in the class) gives $\Cg_{\Eis,k}(a)\cong\Cg_\Eis(a)/\Pk{k}\cong\Iw/(G_a,\Pk{k})\cong
\Zl/G_a(u^{k}-1)\Zl$, and $v(G_a(u^{k}-1))=v\bigl((1-\ell^{k-1})B_k/k\bigr)=v(B_k/k)$, as $1-\ell^{k-1}$ is a unit for $k\ge2$ ($B_k/k$ is $\ell$-integral, $(\ell-1)\nmid k$).
\end{proof}

\begin{remark}[Ohta's theorem, and why no Gorenstein hypothesis enters]\label{rmk:u-ohta}
\LIT{\cite[Thm.~(1.5.5), p.~238; (1.1.4), (1.1.6), p.~229--230; (3.2.2)--(3.2.3), p.~253--254; (1.4.5), p.~235; Appendix~A.1, p.~265]{ohta_cong}, verified against
the library copy.} At tame level $N=1$ Ohta's pair is $(\theta,\psi)=(\omega^{a-2},\mathbf1)$: its weight-$k$ specialisations have nebentypus $\theta\omega^{2-k}=\omega^{a-k}$,
trivial for $k\equiv a$, and $T(q)$-eigenvalue $\theta(q)q\langle q\rangle_\Gamma^{k-2}+1=1+q^{k-1}$. Hypothesis (1.5.2) holds for every even $a$ ($\ell\nmid\varphi(1)$, $uv\in\{1,\ell\}$ ---
$\psi=\mathbf1$ has conductor $v=1$, and $\theta=\omega^{a-2}$ has conductor $u=1$ for $a\equiv2$, so $uv=1=N$, and $u=\ell$ otherwise, so $uv=\ell=N\ell$ ---
and $(\theta,\psi)$ is not exceptional, $\theta|_{(\Z/\ell)^\times}=\omega^{a-2}\ne\omega^{-1}$). Case~(II) of (1.5.5) is the sequence $0\to I(\theta,\psi)_{\mathfrak M}\to e\mathbf
H_{\mathfrak M}\to\Iw\to0$ of the full ordinary algebra, and its congruence module is, by (3.2.2)--(3.2.3), $(\mathbf H_1\oplus\mathbf H_2)/\mathbf H=e h_{\mathfrak M}/I(\theta,\psi)_{\mathfrak M}$ ---
the ring $\Tu/(P_\Eis+P_\cusp)$ of Definition~\ref{def:u-congmod}, i.e.\ Lemma~\ref{lem:u-fibre} in Ohta's language. The theorem gives $\Iw/(A)$ with
$A=G(T,\xi\omega^{2})=G(T,\omega^{a})$ ($N=1$: no Euler factor in (1.5.4)); Ohta's variable is weight-$2$ normalised, $1+T$ acting as $u^{k-2}$ at weight $k$
((1.2.5), (1.4.5)), so $G(T,\omega^{a})=G_a(u^{2}(1+T)-1)$ --- the shift $\langle\gamma\rangle\mapsto\gamma^{-2}\langle\gamma\rangle$ of \cite[p.~7]{emerton_eis}. This is (U1-$\Iw$)
again, for $a\not\equiv0$, \emph{provided} Ohta's $e\mathbf H(1;\Zl)_{\mathfrak M}$ is identified with $\Tu(a)_{\mathfrak m_\Eis}$ (Remark~\ref{rmk:u-hidaM}: it acts faithfully on
the ordinary $\Iw$-adic forms, which span $e S\otimes\mathrm{Frac}(\Iw)\oplus\mathrm{Frac}(\Iw)E$ by (I), so it embeds in $\hcusp(a)\times\Iw$ with image generated by the Hecke pairs;
the step is recorded, not audited). \textbf{Two routes, one statement}: in this case ($\psi=\mathbf1$) Ohta's proof of (II) runs through the appendix, i.e.\ Harder--Pink /
Kurihara plus the Mazur--Wiles main conjecture (p.~254, p.~263); his elementary route (the duality (3.2.7)) needs $\psi\ne\mathbf1$, and he notes (p.~256) that a direct proof of
the duality (3.3.1) ($em(N;\Iw)_{\mathfrak M}=eM(N;\Iw)_{\mathfrak M}$, p.~256) would give (II) directly --- at $N=1$ that is Emerton's Corollary~2 \emph{at $R=\Iw_{(a)}$}: a
$\Iw$-linear continuous functional on $\mathbf T_{(a)}$ is an element of $V_{(a)}\hat\otimes\Iw_{(a)}$, i.e.\ a $\Iw$-adic form with integral constant term. It is
\emph{not} Proposition~1 applied weight by weight: integrality of a $\Iw$-adic constant term at every weight does not give integrality in $\Iw$ ($(X^{\ell}-X)/\ell$
is $\Zl$-valued on $\ell\Zl$). \emph{Why Gorenstein never enters}: [U]'s $\Cg_\Eis$ is the congruence module of the Hecke \emph{algebra}, whose Eisenstein factor is $\Iw$ itself;
Gorensteinness (or freeness of a module over the Hecke algebra) is what relates it to congruence modules of \emph{modules} --- Ohta's (III), (IV), Harder's
denominator on $H^{1}_{!}$ (compare Harder--Pink's Satz~6--7 and \S5, [EX] \texttt{s14:hp}) --- and none of those is used here.
\end{remark}

\begin{remark}[The Iwasawa zero and the ladders]\label{rmk:u-zero}
If $G_a=(\text{unit})\cdot(X-\xi)$ ($\lambda$-invariant $1$, $\mu=0$), (U1) --- a theorem since v0.1.1-pre.2 --- reads $L_k=v(u^{k}-1-\xi)=1+v(k-s_a)$ with $u^{s_a}=1+\xi$, $s_a\in\Zl$:
$L_k=1$ on the class except on the sub-tower $k\equiv s_a\pmod{\ell}$ (where $L_k\ge2$), $\ge3$ on $k\equiv s_a\pmod{\ell^{2}}$, and so on --- the
congruence deepens exactly along the $\ell$-adic approach of $k$ to the zero $s_a$. Independently of (U1), Proposition~\ref{prop:u-determines}(3) makes
$\min(L_k,m)$ constant on every $\mathrm{Tw}_m$. \COMP{} Computed on cohomology at $k=284$ ($\ell=37$), $k=228$ ($\ell=103$) and, since
v0.1.1-pre.6, $k=914$ ($\ell=59$, the cell $(912,59)$), where $L_k=2$: [EX] \texttt{s14:depthlaw} (slot-14 WP-C14 I-3b).
\end{remark}

\begin{proposition}[The two irregular classes at $\ell=691$]\label{prop:u-691}
\COMP{} (Bernoulli numbers: \file{docs/implementation/slot14/wp8/bern.py}, exact $B_k$ by \texttt{sympy} (slot~14, before the owner's no-CAS rule of 2026-10-03;
every value used here since reproduced in-house, v0.5.0-pre.3: $B_k/k\bmod691^{8}$ at $k=12+690t$ and $200+690t$, $t=0,\dots,5$, by tangent numbers in \code{HeckeLambdaValidation}'s
read-only arm \code{HECKE\_LAMBDA\_ALL=1}, reproducing slot-19 R19-2's $691$ data field by field ($\tilde G$, $\alpha$, the rungs $342252=12+690\cdot496$, $342440=200+690\cdot496$), the
tangent route itself gated against WP-NF's exact \code{BigInt} Bernoulli numbers for $k\le200$ in the default arm; and $v_{691}(B_k/k)=2$ exactly at both rungs and $=1$ at
rung${}\pm690$ in \code{HeckeLValueValidation}'s arm~(B), by Voronoi and the Kummer--Mahler route; no value differs). \emph{Primary route} for the class residues
(slot-14 WP-9, \file{bern\_check.py}, outside the tree): the three exact values $B_{k_0+690t}$, $t=0,1,2$, fix the
linear model of $(1-\ell^{k-1})B_k/k\bmod\ell^{2}$ on each class, the third point checking linearity (slopes $80,80$ on the class of $12$, $99,99$ on the class
of $200$), and its zero is at $t=496$ on both classes; Kummer modulo $\ell^{2}$, the one \VERIFYloc{} this chain had through v0.2.0, is re-pointed at v0.2.1-pre.1 to verified
first-edition locators (below). \emph{Cross-check}: the scans by
Voronoi's congruence \cite[Prop.~9.5]{washington} (\VERIFYloc), \file{voronoi\_scan.py} with $a=2$ (WP-8) and \file{voronoi\_a3.py} with $a=3$ (WP-9),
agreeing at every test point, and against the exact $B_{12}/12$, $B_{702}/702$ modulo $691^{2}$).
$691$ is irregular exactly for the classes of $k\equiv12$ and $k\equiv200\pmod{690}$ \cite{washington} (\VERIFYloc; cross-check, verified against the
library copy: \cite[Prop.~14 and the appendix table, row $p=691$, {`E [12, 200]'}]{citro_ghitza}, the weights $k\le p+1$ of the Eisenstein--cuspidal congruences mod~$691$; recomputed here: $v(B_k/k)=1$
at $k=12,200,702,890,1392,1580$, and $0$ at $k=102,202,302,402$). On the class of $12$, $v_{691}(B_k/k)\ge2$ iff $k\equiv342252\pmod{690\cdot691}$; on the class of
$200$, iff $k\equiv342440\pmod{690\cdot691}$; everywhere else on the two classes $v_{691}(B_k/k)=1$. (By Kummer's congruences
$(1-\ell^{k-1})B_k/k\bmod\ell^{2}$ depends only on $k\bmod(\ell-1)\ell$ for $k\not\equiv0\pmod{\ell-1}$, so the $691$ residues scanned cover each class.
The inputs are \LIT{\cite[Thm.~5.11, p.~57; Thm.~5.12, p.~59; Cor.~5.14, p.~61]{washington1st}, first edition, page images (slot-16 L16-3 and WP-8): Cor.~5.14
is the Kummer congruence, verified as a congruence modulo~$\ell$; whether it states the $\ell^{a+1}$ form itself was not recorded by the reader}, and the modulus
$\ell^{2}$ follows from Thm.~5.12 by the argument of its proof (v0.2.1-pre.2: moved out of the \LIT{} argument, the slot-17 referee F1's m-U1): for
$k\equiv k'\pmod{(\ell-1)\ell}$ put $\chi=\omega^{k}=\omega^{k'}\ne1$; Thm.~5.11 gives $(1-\ell^{k-1})B_k/k=-L_\ell(1-k,\chi)$, and Thm.~5.12,
$L_\ell(s,\chi)=\sum_ia_i(s-1)^{i}$ with $\ell\mid a_i$ for $i\ge1$, gives $L_\ell(1-k,\chi)-L_\ell(1-k',\chi)=\sum_{i\ge1}a_i\bigl((-k)^{i}-(-k')^{i}\bigr)\equiv0
\pmod{\ell^{2}}$, as $\ell\mid k-k'$.)
\emph{A field-theoretic reading of the two rungs} (v0.3.0-pre.3; a cross-reference, nothing here rests on it): the rungs $342252$ and $342440$ are the unique lifts
$\kappa$ of $\omega^{12}$, resp.\ $\omega^{200}$ ($\kappa\equiv12$, resp.\ $200\pmod{690}$) for which the cyclotomic unit $\eta_\kappa\in\Q(\zeta_{691^{2}})$ is a local
$691^{2}$-th power at $(1-\zeta_{691^{2}})$ --- \COMP{} by an equivariant group-ring computation that uses no Bernoulli number (slot-19 R19-3,
\file{docs/implementation/slot19/r193/eqlayer2.py}; its method reproduced by the slot-19 referee F4 with a second implementation at $37$ and $59$) --- so
$X^{691^{2}}-\eta_{342252}$ and $X^{691^{2}}-\eta_{342440}$ define the $\omega^{-11}$- and $\omega^{-199}$-parts of the $691$-Hilbert class field of $\Q(\zeta_{691^{2}})$
(\PROVED{} there given the computation and \LIT{} inputs: [L37], \texttt{sec:fieldsribet}).
\end{proposition}

\begin{theorem}[$L_k$ at the computed irregular cells]\label{thm:u-L691}
\PROVED{} from Theorem~\ref{thm:u-sandwich}, Lemma~\ref{lem:u-levelone}(3), Proposition~\ref{prop:u-determines}(3) and the \COMP{} data named.
At $\ell=691$: $L_{12}=L_{200}=L_{702}=1$; and $L_k=1$ for every $k\equiv12\pmod{690\cdot691}$ and every $k\equiv200\pmod{690\cdot691}$.
\emph{Data}: $c_{k,2}(1+2^{k-1})$ has valuation exactly~$1$ at $k=12$ ($\tau(2)-(1+2^{11})=-2073=-3\cdot691$), at $k=200$ (the cell $(198,691)$
over $\Z/691^{2}$, this package's run, report \S3) and at $k=702$ (the cell $(700,691)$ over $\Z/691^{2}$, \file{docs/matrix-engine/results/hida\_towers/},
slot-13 WP-10), $c_{k,2}$ read from $\chi^{\mathrm{ord}}_{n,2}=(X-(1+2^{n+1}))c_{k,2}^{2}$ modulo $691^{2}$ (the monic square root modulo $\ell^{m}$ is unique:
monic polynomials are non-zero-divisors); $v(B_k/k)=1$ at the three weights (Proposition~\ref{prop:u-691}).
\emph{Independent cross-check} (\COMP, slot-14 WP-9, \file{forms\_side.py}, outside the tree, not from the project's
cohomology data): $T_2$ on $S_k(\SL_2(\Z))$ from $q$-expansions (Miller basis, exact over $\Z$) gives $v_{691}(c_{200,2}(1+2^{199}))=1$ (all $16$ forms
$691$-ordinary) and $v_{691}(\det(1+2^{701}-T_2\mid S_{702}))=1$, an upper bound for the ordinary factor, so $L_{200}=1$ and $L_{702}\le1$ from
$q$-expansions alone.
\end{theorem}

\begin{proof}
The sandwich at each weight reads $1\le L_k\le1$. For the towers: $\min(L_k,2)=\min(L_{12},2)=1$ on $\mathrm{Tw}_2$ through $12$, and likewise through $200$.
\end{proof}

\begin{proposition}[Harder--Pink's sub-tower]\label{prop:u-HP}
\PROVED{} from \LIT{Harder--Pink \cite[Satz~5, p.~92; \S4.2, pp.~92--93]{harder_pink} (verified, slot-14 WP-Q2)} with the repair in the
proof, Lemma~\ref{lem:u-levelone}(1), Theorem~\ref{thm:u-sandwich}, Proposition~\ref{prop:u-determines}(3), Kummer's congruences and the analytic
class number formula \cite[Thm.~4.17]{washington} (\VERIFYloc); Harder--Pink's own inputs, their Satz~3 (iv), (v), are \VERIFYloc. Let
$\ell\ge5$, let $a\not\equiv0\pmod{\ell-1}$ be an even class, and put $\delta_a:=v(B_{1,\omega^{a-1}})$. Then $L_k=v(B_k/k)=\delta_a$ for every
$k\ge4$ in the class with $k\equiv1\pmod{\ell^{\delta_a}}$. At $\ell=691$ ($\delta_a=1$ on both irregular classes, Proposition~\ref{prop:u-691}):
$L_k=1$ for $k\equiv7602$ and for $k\equiv137510\pmod{690\cdot691}$.
\emph{Since v0.1.1-pre.2 a special case of Prediction~\ref{pr:u-U1}} (\PROVED{} from Emerton's Theorem~5); recorded as a second route,
independent of Emerton and Ohta. Its weights are disjoint from those of Theorem~\ref{thm:u-L691}.
\end{proposition}

\begin{proof}[Proof (slot-14 WP-Q2 \S3)]
Harder--Pink's $\mathcal T_\Eis/I\cong\Zl/\ell^{\delta'}$ is $\Cg_{\Eis,k}(a)$, so $\delta'=L_k$: their algebra acts faithfully on the torsion-free
$H^{1}_{!}$, which after $\otimes\Q$ is two copies of $S_k$, so its image is the level-one algebra; localised at $\mathfrak m_\Eis$ it is ordinary,
$\hcusp_k$ by Lemma~\ref{lem:u-levelone}(1); and their Eisenstein ideal contains $T_\ell-1-\ell^{k-1}=(\alpha-1)(1-\ell^{k-1}\alpha^{-1})$ with the second factor a
unit (their Hilfssatz~3.1.5). Their \S4.2 concludes $\delta'=\delta:=v(B_k/k)$ at $k=1+j\ell^{N}$, $j=a-1$, $N\gg0$, under the hypothesis
$\ell^{\delta'}\mid k-1$, which the paper does not show. \emph{Repair}: run their \S3 with $\mathcal T_\Eis/(I,\ell^{N})=\Z/\ell^{\min(\delta',N)}$; Satz~5's Fitting
argument gives an unramified extension of degree $\ge\ell^{\min(L_k,N)}$ in the $\omega^{-j}$-eigenspace of the (finite) class group of $\Q(\zeta_\ell)$, so
$L_k<N$ for $N$ large. Over all odd $j$ at once: $\sum_jL_{k_j}\le\sum_jv\#A(\omega^{-j})=v(h^{-})=\sum_jv(B_{1,\omega^{j}})$ (the class number formula),
$v(B_{1,\omega^{j}})=v(B_{k_j}/k_j)$ for $N$ large (the $\ell$-adic limit $\zeta(-j\ell^{N})\to-B_{1,\omega^{j}}$), and $\sum_jv(B_{k_j}/k_j)\le\sum_jL_{k_j}$ (Theorem~\ref{thm:u-sandwich}); so every inequality is an
equality, and $L_{k^*}=\delta_a$ at $k^*=1+(a-1)\ell^{N}$. With $m=\delta_a+1$, Proposition~\ref{prop:u-determines}(3) gives $\min(L_k,m)=\min(L_{k^*},m)=\delta_a$ for
every $k$ of the class with $k\equiv k^*\pmod{\ell^{m-1}}$, i.e.\ $k\equiv1\pmod{\ell^{\delta_a}}$, and Kummer's congruence modulo $\ell^{m}$ gives $v(B_k/k)=\delta_a$
there. At $691$: $7602\equiv12\pmod{690}$, $137510\equiv200\pmod{690}$, both $\equiv1\pmod{691}$.
\end{proof}

\begin{remark}[Class $a\equiv0$]\label{rmk:u-classzero}
For $k\equiv0\pmod{\ell-1}$ the Eisenstein series is $\equiv1$ modulo $\ell$ after normalisation, Corollary~\ref{cor:u-irregular}'s argument does not
apply, and Emerton's theorem excludes the class ($i\not\equiv0$; $G_0$ has a pole). \LIT{\cite[Thm.~(1.5.5), last clause, p.~238, proved in Appendix~A.1,
p.~265]{ohta_cong} (verified): for $(\theta,\psi)=(\omega^{-2},\mathbf1)$ --- at $N=1$ the class $a\equiv0$ (Remark~\ref{rmk:u-ohta}) --- the Eisenstein ideal of
the cuspidal ordinary algebra is the unit ideal, by Ohta's adaptation \cite[App.~A.1]{ohta_cong} of the construction of Harder--Pink and
Kurihara \cite[\S3]{kurihara1993} (whose own range $0<k<\ell-1$ excludes this class) and the triviality of the $\omega$-part of the unramified Iwasawa
module of $\Q(\zeta_{\ell^\infty})$ (``well known'' there)}. The proof is by contradiction and \emph{does not use the main conjecture}, which enters Ohta's
Appendix~A only for $(\theta,\psi)\ne(\omega^{-2},\mathbf1)$ (Thm.~(A.1.13), p.~265): a proper $I^{*}$ would give an unramified abelian pro-$\ell$ extension $L$ of
$\Q(\zeta_{\ell^\infty})$ on which $\mathrm{Gal}(\Q(\zeta_\ell)/\Q)$ acts through $\omega$, with $\mathrm{Gal}(L/\Q(\zeta_{\ell^\infty}))\otimes r\twoheadrightarrow B/I^{*}B$ for a faithful
$h^{*}$-module $B$ ((A.1.5), (A.1.10), (A.1.11), pp.~264--265). Its inputs are Ohta's (3.4.2) (p.~257: geometric Frobenius acts as $T^{*}(p)$ on the inertia
invariants) and \emph{Ohta's freeness theorem} --- $X^{+}=X^{I_\ell}$ free of rank one over $h^{*}$, $X/X^{+}\otimes\mathcal L$ free of rank one (p.~258, citing his
[16]~(2.3.6) and [15]~(5.1.2)) --- a theorem of Ohta's about an \'etale cohomology module, not a hypothesis of [U]. The ``well known'' vanishing is
\LIT{\cite[Prop.~6.16, p.~102]{washington1st} ($A_0=A_1=0$: the $\omega$-part of the $\ell$-class group of $\Q(\zeta_\ell)$ vanishes) and \cite[Prop.~13.22,
p.~284]{washington1st} ($A_n\cong X/((1+T)^{\ell^{n}}-1)X$ for a $\Z_\ell$-extension with exactly one ramified prime, totally ramified), verified, first edition}:
$e_1X/Te_1X=e_1A_0=0$, so $e_1X=0$ by Nakayama ($X$ the unramified Iwasawa module of $\Q(\zeta_{\ell^\infty})$), and $L$, a quotient of $e_1X$, is trivial.
Ohta's ideal matches ours: $I^{*}(\omega^{-2},\mathbf1)$ is generated by $T(l)-(\theta(l)l\iota(l)+\psi(l))$, $T(\ell)-\psi(\ell)$ and $T(q,q)-(\theta\psi)(q)\iota(q)$
(p.~237, Definition~(1.5.1) and the generators after it, read at primes); on the $\theta\psi=\omega^{-2}$-component the last vanishes and the others are $T(q)-\ee_q$, $U-1$, so
$I^{*}(\omega^{-2},\mathbf1)\cdot(\text{that component})=I^{0}(0)\otimes r$, and $I^{0}(0)=\hcusp(0)$ follows over $\Zl$ from the statement over $r$ by faithfully flat
descent ($r$ is free over $\Zl$). With Steps~1--2 of the proof of Prediction~\ref{pr:u-U1} (which do not use
$a\not\equiv0$): $\hcusp(0)/I^{0}(0)=0$, so $J=\Iw$ and $\Cg_\Eis(0)=0$ for every $\ell\ge5$ --- \PROVED{} from that \LIT{} statement, through the identification of
Ohta's cohomological cuspidal algebra $e^{*}h^{*}$ with $\hcusp(0)$ (his end of \S1.2, via Hida; eigenvalues $1+q^{k-1}$ as in Remark~\ref{rmk:u-ohta}).
(\CONJ{} through v0.1.1-pre.1, with the heuristic: the residual representation would be an extension between $\mathbf 1$ and $\omega^{-1}$, whose
$\omega^{\pm1}$-parts of the class group of $\Q(\mu_\ell)$ vanish, \cite[Ch.~6]{washington} (\VERIFYloc).) Independently \PROVED{} at $\ell\in\{5,7,11,13\}$ by
Theorem~\ref{thm:u-ell13} below.
\end{remark}

%----------------------------------------------------------------------
\section{No CM components at tame level one}\label{sec:u-cm}
%----------------------------------------------------------------------

\begin{definition}[CM component]\label{def:u-cm}
A component $P$ of $\Tu(a)$ has \emph{CM by} an imaginary quadratic field $K$ if the image of $T(p)$ in $\Tu(a)/P$ is $0$ for every prime $p\ne\ell$ inert
in $K$.
\end{definition}

\begin{theorem}[No CM at tame level one]\label{thm:u-nocm}
\PROVED{} from Proposition~\ref{prop:u-components}(3), Lemma~\ref{lem:u-levelone}, and \STD{Deligne: for a normalised eigenform $f$ of weight $k\ge2$ and
level $1$ there is a continuous semisimple $\rho_f\colon G_\Q\to\GL_2(\Qlb)$ unramified outside $\ell$ with $\det(1-\rho_f(\Frob_p)Y)=1-a_pY+p^{k-1}Y^{2}$,
$\det\rho_f=\chi_{\mathrm{cyc}}^{k-1}$ \cite{deligne_formes}; as stated in \cite[(2.1a,b), p.~556]{hida_invent} (verified against the read-only PDF)},
\STD{Chebotarev density and Brauer--Nesbitt} and \LIT{Wiles \cite[Thm.~2.1.4, p.~561]{wiles_ordinary} (verified, page image, slot-16 L16-3, \file{docs/implementation/slot16/locators.tsv}; v0.2.0-pre.2): if moreover $f$ is $\ell$-ordinary then
$\rho_f|_{G_{\Ql}}$ is reducible with an unramified quotient character}. No component of $\Tu(a)$ has CM, for any $\ell\ge5$ and any class $a$.
\end{theorem}

\begin{proof}
The Eisenstein component has $\ee_p(u^{k}-1)=1+p^{k-1}\ne0$. Let $P$ be cuspidal with CM by $K$. Take $k\ge3$ in the class and a point of $P$ above $\Pk{k}$
(lying over, $\Iw\subseteq\Tu/P$ integral); by Proposition~\ref{prop:u-components}(3) and Lemma~\ref{lem:u-levelone} it is the eigencharacter of some $f\in\Ek_k$,
so $a_p(f)=0$ for every $p\ne\ell$ inert in $K$. Let $\chi_K$ be the quadratic character of $K$. For every prime $p\nmid\ell d_K$,
$\chi_K(p)a_p(f)=a_p(f)$ (split: $\chi_K(p)=1$; inert: $a_p=0$), so $\mathrm{tr}\,\rho_f(\Frob_p)=\mathrm{tr}\,(\rho_f\otimes\chi_K)(\Frob_p)$ on a density-one set, and by
Chebotarev and Brauer--Nesbitt $\rho_f\cong\rho_f\otimes\chi_K$ (both semisimple).
\emph{(a) Some prime $q\ne\ell$ divides $d_K$.} Then $\rho_f$ is unramified at $q$, so the inertia group $I_q$ acts trivially on $\rho_f$ and through the non-trivial
$\chi_K|_{I_q}$ on $\rho_f\otimes\chi_K$; isomorphic representations have isomorphic restrictions to $I_q$. Contradiction.
\emph{(b) Otherwise} $\ell$ is the only prime ramified in $K$, so $K=\Q(\sqrt{-\ell})$, $\ell\equiv3\pmod4$, and $\chi_K|_{I_\ell}\ne1$. By Wiles,
$\rho_f|_{G_{\Ql}}$ has semisimplification $\varepsilon_1\oplus\varepsilon_2$ with $\varepsilon_2$ unramified; and $\{\varepsilon_1,\varepsilon_2\}=
\{\varepsilon_1\chi_K,\varepsilon_2\chi_K\}$ on $G_{\Ql}$. $\varepsilon_2=\varepsilon_2\chi_K$ fails on $I_\ell$, so $\varepsilon_2=\varepsilon_1\chi_K$, whence
$\varepsilon_1|_{I_\ell}=\chi_K|_{I_\ell}$ and $\det\rho_f|_{I_\ell}=\chi_K|_{I_\ell}$ has order $2$ --- but it is $\chi_{\mathrm{cyc}}^{k-1}|_{I_\ell}$, of infinite
order for $k\ge2$. Contradiction.
\end{proof}

\begin{remark}[The standard route, and where CM lives]\label{rmk:u-cmroute}
The plan's route \STD{a newform with CM by $K$ has level divisible by $|d_K|$, \cite[\S3]{ribet_nebentypus} (\VERIFYloc); an ordinary CM family by $K$ needs
$\ell$ split in $K$, \cite[(7.1), p.~267]{hida_ens} (verified: Hida's CM components are constructed under that hypothesis)} gives the same conclusion: $\ell$
unramified in $K$, so $|d_K|\ge3$ divides the tame level. Hida's CM components \cite[Thm.~7.1, pp.~268--269]{hida_ens} (verified) have conductor
$N_0=d\,N(\mathfrak c)$, divisible by $d=|d_K|$. The proof above avoids the level formula and case (b) needs ordinarity only at $\ell$; with the crystallinity
of $\rho_f|_{G_{\Ql}}$ at level one (Faltings, Scholl; \VERIFYloc) it would exclude CM for every level-one eigenform. CM $\cap$ $\GL(2)$ and CM $\cap$ Eisenstein
therefore start at tame level $N\ge3$ (\S\ref{sec:u-scope}).
\end{remark}

\begin{remark}[Dihedral residual systems at tame level one]\label{rmk:u-dihedral}
(v0.1.1-pre.5; slot-14 WP-C16, \file{docs/implementation/slot14/wpc16/J-7\_report.md} \S2, \S4 and its addendum, refereed in
\file{wpc16/REFEREE\_J8\_WP-C16.md} \S H.) Let $\ell\equiv3\pmod4$, $K=\Q(\sqrt{-\ell})$, $h=h(-\ell)$; the ordinary dihedral blocks lie in the class
$k\equiv(\ell+1)/2\pmod{\ell-1}$. \textbf{(i)} \COMP{} At every weight of that class run, the ordinary dihedral part of $H^{1}(\Ga;\Mn)$ has
$\Fl$-dimension $h-1$, whether the Hecke eigenvalues are $\Fl$-rational or not ($\ell=47$, $103$, $127$, $151$, $167$, $191$, $359$, $367$; one
block of dimension $2r_d$ per Galois orbit of non-trivial class-group characters $\chi$ up to $\chi\leftrightarrow\chi^{-1}$, so that
$\sum 2r_d\cdot\#\{\text{orbits}\}=h-1$ is arithmetic, and $m=1$ on these blocks is data, $20$ / $20$). The $\Iw$-rank is not claimed.
\textbf{(ii)} The non-ordinary dihedral ideals $\mathfrak m$ (the weights $k_\alpha$, $\alpha\ge1$) are non-reduced: the full local Hecke algebra of $S_k$ at
$\mathfrak m$, modulo $\ell$, is $\F_q[\varepsilon]/(\varepsilon^{2})$, $\F_q$ the residue field (\COMP{}) --- and their first-order direction is anticyclotomic (\COMP{} on the
$56$ measured ideals). They are \emph{not} CM in characteristic~$0$, but not by Theorem~\ref{thm:u-nocm}, which concerns the ordinary
components: by the level, \STD{a CM newform has level divisible by $|d_K|$} (Remark~\ref{rmk:u-cmroute}'s locator, \VERIFYloc), and the
level is~$1$. (That there is no level-one dihedral $\bar\rho$ at all when $\ell\equiv1\pmod4$ --- WP-C16's Proposition~R --- is recorded in
[EX], not here: [U] is about the ordinary $\Iw$-adic algebra.)
\end{remark}

%----------------------------------------------------------------------
\section{Cuspidal \texorpdfstring{$\cap$}{cap} cuspidal, and the computed component structure}\label{sec:u-cc}
%----------------------------------------------------------------------

\begin{theorem}[The ordinary algebra at tame level one for $\ell\le13$]\label{thm:u-ell13}
\PROVED{} from Theorem~\ref{thm:u-control}, Lemma~\ref{lem:u-levelone}, Theorem~\ref{thm:u-etale}, with the \COMP{} ordinary ranks and polynomials of
\file{docs/matrix-engine/results/hida\_towers/hida\_towers.tsv} (published 2026-09-27; every class of each $\ell$ occurs there with $n\ge2$). For
$\ell\in\{5,7\}$, $\Tu(a)=\Iw$ for every class (no ordinary cusp form: $h^{1}_{\mathrm{ord}}=1$). For $\ell=11$, $\Tu(a)=\Iw$ except in the class $k\equiv2
\pmod{10}$, where $\Tu(a)=\Iw\times\Iw$ (the $\Delta$-family and the Eisenstein component, disjoint); for $\ell=13$, likewise except in the class $k\equiv0
\pmod{12}$. In particular every component is $\cong\Iw$, no two components meet, and every ordinary local Hecke algebra of weight $k\ge3$ at $\ell\le13$ is
$\Zl$ (Loewy length $1$).
\end{theorem}

\begin{proof}
Lemma~\ref{lem:u-levelone}(3) gives $d_k=(h^{1}_{\mathrm{ord}}-1)/2$: $0$ everywhere at $5,7$; at $11$, $1$ exactly for $n\equiv0\pmod{10}$, at $13$, $1$
exactly for $n\equiv10\pmod{12}$ (the published ranks, \HA{rmk:hidacheckable} \HA{Hi3}'s values). With $d(a)=d_k$ (Theorem~\ref{thm:u-control}), $\hcusp(a)=0$ or
$\hcusp(a)$ free of rank~$1$, i.e.\ $\Iw$. For the Eisenstein intersection in the rank-one classes the catalogue gives $v(c_{k,2}(1+2^{k-1}))=0$ on all $12$
($\ell=11$) resp.\ $16$ ($\ell=13$) cells at $m=1$ --- including the class $a\equiv0$ at $13$ --- so $L_k=0$ (Theorem~\ref{thm:u-sandwich}) and
$\Cg_\Eis(a)=0$ (Proposition~\ref{prop:u-eisspec}); $\Tu(a)=\hcusp(a)\times\Iw$ by Lemma~\ref{lem:u-fibre}.
\end{proof}

\begin{theorem}[Six classes at $\ell=691$, one at $\ell=2411$]\label{thm:u-691classes}
\PROVED{} from Theorem~\ref{thm:u-etale} with the \COMP{} polynomials of the catalogue and of this package's run of $(198,691)$ (report \S3).
Writing $\Tu(k_0)$ for the class of $k_0$ and $\delta$-lists for the residue degrees of the cuspidal components:
\begin{center}\small
\begin{tabular}{@{}l l l p{0.23\textwidth} p{0.28\textwidth}@{}}
\toprule
class of $k$ & cell used & $d(a)$ & cuspidal components (residue degrees) & Eisenstein \\
\midrule
$12$  & $(10,691)$, $(700,691)$ & $1$  & $[1]$ ($\Delta$-family) & meets it, $\Cg_\Eis\ne0$, $L_{12}=L_{702}=1$ \\
$200$ & $(198,691)$           & $16$ & $[1,1,2,2,10]$ & meets one degree-$1$ component, $L_{200}=1$ \\
$102$ & $(100,691)$           & $8$  & $[1,7]$ & disjoint \\
$202$ & $(200,691)$           & $16$ & $[1,15]$ & disjoint \\
$302$ & $(300,691)$           & $24$ & $[5,19]$ & disjoint \\
$402$ & $(400,691)$           & $33$ & $[6,6,21]$ & disjoint \\
\midrule
$12$ at $\ell=2411$ & $(10,2411)$ & $0$ & none & $\Tu=\Iw$ \\
\bottomrule
\end{tabular}
\end{center}
In every class listed, $\hcusp(a)$ is finite étale over $\Iw$, its components are pairwise disjoint, and the only intersections of components are the Eisenstein
ones of the two irregular classes; there $\Tu(a)_{\mathfrak m_\Eis}\cong\Iw\times_{\Iw/J}\Iw$ with $J$ the ideal of
Proposition~\ref{prop:u-eisspec}, $\Iw/(J,\Pk{k})=\Z/\ell^{L_k}$ for every $k$ in the class; hence
$\Tu_k(a)_{\mathfrak m_\Eis}\cong\Zl\times_{\Z/\ell^{L_k}}\Zl$. (That $J=(G_a)$ is principal comes from (U1), not from the data: Remark~\ref{rmk:u-Jprincipal}.)
\end{theorem}

\begin{proof}
$\bar c_{k,2}$ (and $\bar c_{k,3}$) are squarefree at the cells named, with the listed factor degrees (report \S2, \S3); Theorem~\ref{thm:u-etale}(1),(3).
The Eisenstein value $1+2^{k-1}$ is not a root of $\bar c_{k,2}$ at $k=102,202,302,402$ and is a simple root at $k=12,200$, and $\Cg_\Eis\ne0$ there by
Theorem~\ref{thm:u-sandwich} with $v(B_k/k)=1$ (Proposition~\ref{prop:u-691}), so Theorem~\ref{thm:u-etale}(3) applies. In the irregular classes the cuspidal
component through $\mathfrak m_\Eis$ has residue degree $1$ and is étale, so it is $\Iw$; Lemma~\ref{lem:u-fibre} with $A=B=\Iw$ gives
$\Tu(a)_{\mathfrak m_\Eis}=\Iw\times_{\Iw/J}\Iw$ ($J\ne0$ by Lemma~\ref{lem:u-deligne}), and Proposition~\ref{prop:u-eisspec} gives
$\Iw/(J,\Pk{k})=\Cg_{\Eis,k}=\Z/\ell^{L_k}$; tensoring the fibre-product sequence with $\Iw/\Pk{k}$ (right exact) and comparing with
Lemma~\ref{lem:u-fibre} at weight $k$ gives the finite-level statement.
\end{proof}

\begin{remark}[Is $J$ principal?]\label{rmk:u-Jprincipal}
Every quotient of $\Iw$ is cyclic, but a non-zero ideal of the two-dimensional regular local ring $\Iw$ need not be principal ($\mathfrak m_\Iw$ is not), and the
data cannot decide it: $J=\mathfrak m_\Iw$ also gives $L_k=1$ for every $k$, exactly what is measured on the level-$2$ towers through $12$ and $200$. As a
$\Iw$-module $\Tu(a)_{\mathfrak m_\Eis}\cong\Iw\oplus J$ (via $(x,y)\mapsto(y,x-y)$), so \textbf{$J$ is principal iff $\Tu(a)_{\mathfrak m_\Eis}$ is
$\Iw$-free}. That would follow from Hida's freeness of the ordinary algebra of \emph{modular} forms (\HA{thm:hidacontrol}(3), \cite[Thm.~3.1]{hida_ens},
verified) \emph{once} $\Tu(a)$ is identified with Hida's algebra (Remark~\ref{rmk:u-hidaM}); that identification needs the reducedness of Hida's algebra and that its
only non-cuspidal component is the one of \cite[Thm.~7.2]{hida_ens}, and is not carried out here. \textbf{Decided since v0.1.1-pre.2}: for
$a\not\equiv0$, $J=(G_a)$ is principal by (U1-$\Iw$) (Prediction~\ref{pr:u-U1}, \PROVED{} from Emerton's Theorem~5), so $\Tu(a)_{\mathfrak m_\Eis}\cong\Iw\oplus\Iw$ is
$\Iw$-free in the situation above (and $\Tu(a)$ is $\Iw$-free in every class $a\not\equiv0$); for $a\equiv0$, $J=\Iw$ (Remark~\ref{rmk:u-classzero}). The data's ambiguity ($J=\mathfrak m_\Iw$ versus $J=(G_a)$) is resolved by the literature,
not by the data. The finite-level statements (Theorem~\ref{thm:u-691classes}, Proposition~\ref{prop:u-eisblock}) do not depend on it.
\end{remark}

\begin{prediction}[(U2) Cuspidal congruence lengths along towers]\label{pr:u-U2}
For two cuspidal components $P\ne P'$ of one class that meet, let $L_k(P,P'):=\mathrm{length}\,\Tu_k/(I_{P,k}+I_{P',k})$. \PROVED{} (as
Proposition~\ref{prop:u-eisspec}, with Lemma~\ref{lem:u-fibre} for $\Tu/(P\cap P')\subseteq\Tu/P\times\Tu/P'$) when $\Tu/P$ and $\Tu/P'$ are $\Iw$-free (e.g.\ of
$\Iw$-rank~$1$) --- then $(\Tu/P)/\Pk{k}$ is torsion-free of rank $r=\rank_\Iw\Tu/P$, hence equal to $\Tu_k/I_{P,k}$ (the kernel reading of
Definition~\ref{def:u-congmod}) by the étaleness at $\Pk{k}$ shown in the proof of Proposition~\ref{prop:u-components}(3), and likewise for $P'$: $L_k(P,P')=\mathrm{length}\,\Cg(P,P')/\Pk{k}$ and $\min(L_k(P,P'),m)$ is constant on $\mathrm{Tw}_m$. \CONJ{} in general: the same constancy
without freeness. Upper bounds at one cell: Proposition~\ref{prop:u-bounds}(2). \textbf{No cuspidal $\cap$ cuspidal instance is known at tame level~$1$}: by
Theorems~\ref{thm:u-ell13}, \ref{thm:u-691classes} no two cuspidal components meet in any computed class, and on [M1] \S5.5's forms-side list
($k\le298$, $\ell\le23170$) every ordinary non-reduced block is a \emph{ramified component}, not an intersection (\COMP, slot-14 WP-9 \S3.4;
Proposition~\ref{prop:u-loewy}(4): nine cells and $(68,157)$; three further rows of the screened list are reduced).
\COMP{} (v0.1.1-pre.6; slot-14 WP-C14 I-8, in the wording of its referee \file{docs/implementation/slot14/wpc14/REFEREE\_I3b\_I7b\_I8\_I9.md}
\S5): in the $320$ weight-shift rungs of WP-C14 ($\ell\in\{7,11,13,23,31,37,59,67,101,103,131,149,157,691\}$, the listed cells of [EX]
\texttt{s14:wshida}), every ordinary cuspidal block has $g=1$; the rank-$2$ ramified blocks of
Proposition~\ref{prop:u-loewy}(4) ($(30,67)$, $(26,131)$, \dots) are not among these cells and are the next targets of the intra-block reading.
\end{prediction}

%----------------------------------------------------------------------
\section{Saturation, Gorensteinness and Loewy length}\label{sec:u-sat}
%----------------------------------------------------------------------

Let $\mathfrak m$ be a maximal ideal of $\Tu_k(a)$, $\Tu_{k,\mathfrak m}$ its localisation, and $\Lambda_{\mathfrak m}$ a $\Tu_{k}$-stable $\Zl$-lattice on which
$\Tu_{k,\mathfrak m}$ acts faithfully --- for instance $\eord H(n)_{\mathfrak m}$ on the contract's lattice $\Sym^{n}$ or on Hida's lattice $L_n$
(\HA{def:hidalattice}, \HA{rmk:stabconvention}(3)). As in [M1] Prop.~1, $\Tu^{\mathrm{sat}}_{\mathfrak m}:=(\Tu_{k,\mathfrak m}\otimes\Q)\cap\End(\Lambda_{\mathfrak m})$,
$\delta(\mathfrak m):=\mathrm{length}(\Tu^{\mathrm{sat}}_{\mathfrak m}/\Tu_{k,\mathfrak m})$, and the image $A$ of $\Tu_{k,\mathfrak m}$ in
$\End(\Lambda_{\mathfrak m}/\ell)$ is $\Tu_{k,\mathfrak m}/\ell$ iff $\delta(\mathfrak m)=0$ ([M1] Prop.~1(2); \PROVED{} there).

\begin{proposition}[Maximal orders are saturated]\label{prop:u-maxorder}
\PROVED{} If $\Tu_{k,\mathfrak m}$ is a maximal $\Zl$-order in $\Tu_{k,\mathfrak m}\otimes\Q$ --- e.g.\ étale over $\Zl$ --- then $\delta(\mathfrak m)=0$ on
\emph{every} faithful stable lattice, and $A=\Tu_{k,\mathfrak m}/\ell$. This holds at every ordinary block of every class in
Theorems~\ref{thm:u-ell13}, \ref{thm:u-691classes} except the Eisenstein blocks of the two irregular classes at $691$.
\end{proposition}

\begin{proof}
$\Tu^{\mathrm{sat}}_{\mathfrak m}$ is an order (a finitely generated $\Zl$-algebra, being a submodule of $\End(\Lambda_{\mathfrak m})$) containing the maximal one.
A block that is a local factor of an étale algebra ($\hcusp_k$ by Theorem~\ref{thm:u-etale}(2), or the factor $\Zl$ of $\hcusp_k\times\Zl$) is étale.
\end{proof}

\begin{proposition}[The Eisenstein blocks]\label{prop:u-eisblock}
\PROVED{} Let $\Tu_{k,\mathfrak m}=\Zl\times_{\Z/\ell^{L}}\Zl$ with $L\ge1$ (Theorem~\ref{thm:u-691classes}). Then $\Tu_{k,\mathfrak m}=\Zl[t]$ with $t=(\ell^{L},0)$,
$t^{2}=\ell^{L}t$, so $\Tu_{k,\mathfrak m}/\ell\cong\Fl[t]/(t^{2})$ (Loewy length $2$, independent of $L$); $0\le\delta(\mathfrak m)\le L$, with $\delta=L$ iff
$\Lambda_{\mathfrak m}$ is the direct sum of its Eisenstein and cuspidal parts, and $\delta=0$ iff $\Lambda_{\mathfrak m}/\ell$ is a faithful
$\Fl[t]/(t^{2})$-module. \COMP{} ([M1] \S5.1, the census cell $(10,691)$ on $\Sym^{10}$: $A_1=A_3/691$ of Hilbert function $(1,1)$, $\varepsilon_3=0$):
$\delta=0$ at $(10,691)$.
\end{proposition}

\begin{proof}
$\Tu_{k,\mathfrak m}=\{(x,y):x\equiv y\ (\ell^{L})\}=\Zl(1,1)\oplus\Zl t$; $\Tu^{\mathrm{sat}}\subseteq\Zl\times\Zl$, of index $\ell^{L}$ over $\Tu$, and it is all of
$\Zl\times\Zl$ iff the idempotent $(1,0)$ preserves $\Lambda_{\mathfrak m}$. The last equivalence is [M1] Prop.~1(2).
\end{proof}

\begin{remark}[Gorenstein, multiplicity one, and the lattice]\label{rmk:u-gorenstein}
\LIT{Mazur--Tilouine \cite{mazur_tilouine}, Wiles \cite[Ch.~2]{wiles_flt} (\VERIFYloc): for a non-Eisenstein ordinary $\mathfrak m$ with $\ell$-distinguished residual
representation, the $\Iw$-adic $\Tu_{\mathfrak m}$ is Gorenstein and the ordinary cohomology on Hida's lattice is free over it} --- at tame level $1$ the residual
representation of an ordinary form of even weight $k$ is automatically $\ell$-distinguished ($\bar\varepsilon_1/\bar\varepsilon_2|_{I_\ell}=\omega^{k-1}\ne1$, $k-1$
odd, $\ell-1$ even). Freeness on Hida's lattice gives $\delta=0$ there ([M1] Prop.~1(2), sufficient condition); on $\Sym^{n}$ with $n\ge\ell$ the lattice differs
(\HA{rmk:stabconvention}(3)), and nothing is transferred. In every computed class the ordinary non-Eisenstein blocks are étale, so Proposition~\ref{prop:u-maxorder}
gives $\delta=0$ on both lattices without the literature. \emph{Since v0.5.0-pre.2} (slot-21 F4): on \emph{ordinary} parts the two lattices give isomorphic Hecke modules,
$e_1H^{1}(\Ga;L_n(\Zl))\cong e_1H^{1}(\Ga;\Sym^{n}\Z^{2}\otimes\Zl)$ (Lemma~\ref{lem:u-latticeord}(3), part~4), so freeness on Hida's lattice \emph{is} transferred to the
contract's lattice on ordinary parts, and ``nothing is transferred'' above is superseded for them; the sentence is kept as the record of v0.1--0.4.
\end{remark}

\begin{prediction}[(U3) Ordinary blocks are saturated]\label{pr:u-U3}
\CONJ{} On every ordinary block of $\eord H^{1}(\Ga;\Sym^{n}\Z^{2})\otimes\Zl$, $\ell\ge5$, $\delta(\mathfrak m)=0$; equivalently the census-type algebra $A$ of the block
equals $\Tu_{k,\mathfrak m}/\ell$. Proved for the étale blocks (Proposition~\ref{prop:u-maxorder}) and \COMP{} at $(10,691)$ (Proposition~\ref{prop:u-eisblock});
open for non-étale ordinary blocks, of which the computed ones are the Eisenstein blocks of Theorem~\ref{thm:u-691classes}. Testable by slot-14 WP-6's route
($\varepsilon_K=0$ at the stable $K$) or by \HA{Hi8}(c)'s comparison $D_\ast$.
\end{prediction}

\begin{proposition}[Loewy length and components]\label{prop:u-loewy}
\PROVED{} from Proposition~\ref{prop:u-determines}(2). Let $\mathfrak m$ be an ordinary maximal ideal of $\hcusp_k$, $k\ge3$.
\begin{enumerate}
\item $\LL(\hcusp_{k,\mathfrak m}/\ell)\le\dim_{\Fl}\hcusp_{k,\mathfrak m}/\ell=\rank_{\Zl}\hcusp_{k,\mathfrak m}$, the number of eigenforms in the block; and it is a
  class invariant.
\item $\LL=1$ iff $\hcusp(a)_{\mathfrak m}$ is étale over $\Iw$ iff a single component $P$ passes through $\mathfrak m$ and $(\Tu/P)_{\mathfrak m}$ is étale over $\Iw$.
  So $\LL\ge2$ means at least one of: \textbf{(i)} an intersection of components at $\mathfrak m$; \textbf{(ii)} a single component whose normalisation is
  ramified over $\Iw$ at $\mathfrak m$; \textbf{(iii)} a single non-normal component (normalisation étale over $\Iw$, the component a proper suborder of it).
  If all components through $\mathfrak m$ are $\cong\Iw$ (rank one), $\LL$ is at most their number.
\item Consequently an ordinary block with $\LL\ge3$ needs at least three eigenforms in one residual system and, if its components are rank one, at least three
  components through $\mathfrak m$. At $\ell\le13$ none exists (Theorem~\ref{thm:u-ell13}); this is the forms-side reading of [M1] Prop.~3 (non-reduced blocks at
  $\ell\le13$ are non-ordinary) and agrees with [M1] \S5.5 (no ordinary level-one block with $\LL\ge3$ for $\ell<300$, \COMP{} there).
  (\emph{Added at 0.7.0-pre.1}, the phenomenology proposal of \file{docs/phenomenology/phenomena\_catalogue.md} \S11 item~4; citation only, no statement moves.)
  The ordinary level-one Eisenstein block $\Tu_{k,\mathfrak m}$ at $(n,\ell)=(484,547)$, $k=486$, has $\Tu_{k,\mathfrak m}/\ell\cong\Fl[y]/(y^{3})$, Loewy length~$3$, and
  $\delta=0$ on $\Sym^{484}$ (part~6, Proposition~\ref{prop:u-H12}(3), (6), \PROVED{} given the data (D1)--(D4) there, not from Proposition~\ref{prop:u-determines}(2));
  its cuspidal rank $2$, and that $(547,486)$ is the only irregular pair with $\ell<10^{5}$ of cuspidal rank $\ge2$, are \LIT{\cite[\S1.1, arXiv v1, p.~2]{bellaiche_pollack} (page
  image verified)}, the pair first in \LIT{\cite[p.~68]{calegari_eisdef} (page image verified)}; this is a block \emph{with} the Eisenstein component, so it is not an
  instance of the cuspidal statement above.
\item \COMP{} (slot-14 WP-9 \S3.4 on [M1] \S5.5's screened list \file{fl\_search.txt} --- the rows with $\ell\mid\operatorname{disc}c_{k,2}$
  ---, corrected by slot-14 WP-14c \S6 and its referee \file{docs/implementation/slot14/wpc16/REFEREE\_J8\_WP-C16.md} Part~II; the scripts
  \file{u2\_candidates.py} and, for the non-Eisenstein check,
  \file{s6\_eis\_check.py}, both outside the tree)
  \textbf{The first ordinary $\LL=2$ cuspidal blocks at tame level~$1$ are all of kind (ii).} Of the $12$ rows of that list with
  $\ell\le1000$, $k\le60$, nine are non-reduced --- the cells $(n,\ell)=(30,67)$, $(26,131)$, $(26,139)$, $(32,479)$, $(36,181)$,
  $(36,349)$, $(38,73)$, $(44,227)$, $(52,71)$ --- and each of these carries one $2$-dimensional block with a double root of
  $\bar c_{k,2}$, on which $T_\ell$ is invertible modulo $\ell$ (ordinary), which is not Eisenstein (the double root is not
  $1+2^{k-1}\bmod\ell$), and at which $v_\ell(\operatorname{disc}c_{k,2})=1$ with no root of $c_{k,2}$ modulo $\ell^{2}$ lifting the
  double root: a ramified quadratic order, a single component, $\hcusp_{k,\mathfrak m}/\ell=\Fl[t]/(t^{2})$. Example $(30,67)$:
  $c_{32,2}=X^{2}-39960X-2235350016$, $v_{67}(\operatorname{disc})=1$, $\det T_{67}\equiv56\pmod{67}$. The other three rows,
  $(46,383)$, $(50,677)$, $(54,67)$, are \emph{reduced}: there $v_\ell(\operatorname{disc}c_{k,2})=2$, the double root of
  $\bar c_{k,2}$ lifts modulo $\ell^{2}$, and $T_3$ separates two distinct residual systems on the $T_2$-generalised eigenspace
  (a conjugate pair over $\F_{\ell^{2}}$ at $383$ and $67$, two $\Fl$-rational systems at $677$); the double root is a coincidence
  of $T_2$-eigenvalues, not a non-reduced block, and no statement of this proposition applies to them. (The earlier text listed
  all $12$ as non-reduced: [M1]'s list is screened by $\ell\mid\operatorname{disc}c_{k,2}$, and the three reduced rows were
  carried over without the Loewy-length filter.) Among the $215$ rows with $\ell\le23170$ the only one with
  $v_\ell(\operatorname{disc}c_{k,2})\ge2$ \emph{and} a non-reduced algebra $\hcusp_k/\ell$ is $(68,157)$: two double roots
  ($111$, $96$ modulo $157$), both blocks ordinary, both ramified (no lift modulo $157^{2}$); every other row with
  $v_\ell(\operatorname{disc}c_{k,2})\ge2$ has reduced $\hcusp_k/\ell$. So that list has \textbf{no} cuspidal $\cap$ cuspidal
  instance up to $k\le298$, $\ell\le23170$. At the nine cells with $v_\ell(\operatorname{disc})=1$ (and at the two blocks of
  $(68,157)$, checked directly) the block's $\Zl[T_2]$ is the maximal order of a ramified quadratic extension (its discriminant
  has odd valuation~$1$), so $\hcusp_{k,\mathfrak m}$ is that maximal order and $\delta(\mathfrak m)=0$ on every lattice
  (Proposition~\ref{prop:u-maxorder}); slot-14 WP-14c \S6 confirms $\delta=0$ on both shipped routes at all eleven blocks. These
  are cheap tower cells for (U4), (U6) and (U3).
\end{enumerate}
\end{proposition}

\begin{proof}
(1) The chain $\hcusp_{k,\mathfrak m}/\ell\supsetneq\mathfrak m\supsetneq\mathfrak m^{2}\supsetneq\cdots$
strictly decreases in $\Fl$-dimension until it reaches $0$ (Nakayama), so its length is at most $\dim_{\Fl}$; class invariance is Proposition~\ref{prop:u-determines}(2).
(2) $\LL=1$ iff the special fibre $\hcusp(a)_{\mathfrak m}/\mathfrak m_\Iw$ is a field, iff (as in the proof of Theorem~\ref{thm:u-etale}) $\hcusp(a)_{\mathfrak m}$ is étale over
$\Iw$, which is a domain, so one component, étale. If one component passes through $\mathfrak m$ and is not étale there, either its normalisation is ramified at
$\mathfrak m$, or the normalisation is étale and the component is a proper suborder of it; these and (i) exhaust the cases. For rank-one components $\hcusp(a)_{\mathfrak m}\subseteq\Iw^{r}$ has rank $r$, so (1) bounds $\LL$ by $r$.
(3) follows.
\end{proof}

%----------------------------------------------------------------------
\section{Predictions for the weight ladders}\label{sec:u-pred}
%----------------------------------------------------------------------

A \emph{ladder} of level $m$ through $n_0$ is $n_j=n_0+j(\ell-1)\ell^{m-1}$, $j\ge0$ --- the cells of $\mathrm{Tw}_m(\ell;n_0)$ in order (\HA{def:hidatower}); a sibling
package proposes cluster runs along such ladders. The predictions below are numbered before any run (rule~20). A failure of a \PROVED{} one is a defect of the
implementation or of the cited chain (Hida's control, stabilisation, Eichler--Shimura) and stops the run's analysis; a failure of a \CONJ{} one is a finding.

\begin{enumerate}[label=(U\arabic*)]
\item \PROVED{} from \LIT{Emerton, Thm.~5} (Prediction~\ref{pr:u-U1}; \CONJ{} through v0.1.1-pre.1) $L_k=v(B_k/k)$ at every rung of a class $a\not\equiv0$. At $\ell=691$: on the ladder $n_j=10+690j$, $L=1$ for every $j\not\equiv496
  \pmod{691}$ and $L\ge2$ at $j\equiv496$ ($n=342250$); on $n_j=198+690j$, likewise with the exceptional rungs $j\equiv496$ ($n=342438$)
  (Proposition~\ref{prop:u-691}). Checkable at a rung computed over $\Z/691^{2}$ as $\min_{p}v(c_{k,p}(1+p^{k-1}))=1$ --- now a check of the implementation: a value $0$ is a defect, a
  rung with all computed $p$ giving $\ge2$ means too few $p$. \emph{Already a theorem} (Theorem~\ref{thm:u-L691}) on the level-$2$ ladders $n_j=10+476790j$, $198+476790j$.
\item \CONJ{} / \PROVED{} under freeness (Prediction~\ref{pr:u-U2}) cuspidal congruence lengths: $\min(L_k(P,P'),m)$ constant along a level-$m$ ladder. No cuspidal $\cap$ cuspidal instance is known at tame level~$1$: on
  [M1] \S5.5's list ($k\le298$, $\ell\le23170$) every ordinary non-reduced block is a ramified component (slot-14 WP-9 \S3.4,
  Proposition~\ref{prop:u-loewy}(4): nine cells and $(68,157)$; three further rows of the screened list are reduced).
\item \CONJ{} (Prediction~\ref{pr:u-U3}) $\delta=0$ on ordinary blocks; at $691$ class $12$ it says the census algebra of the Eisenstein block of every rung $n=10+690j$ is
  $\Fl[t]/(t^{2})$ (Hilbert function $(1,1)$), as at $(10,691)$.
\item \PROVED{} (Proposition~\ref{prop:u-determines}, Theorem~\ref{thm:u-etale}) \textbf{which components survive}: all of them, with the same residue degrees, at every rung
  of a level-$1$ ladder --- $h^{1}_{\mathrm{ord}}$ is constant and $\bar c_{k,p}$ is the same polynomial in $\Fl[X]$ at every rung, for every $p\ne\ell$. Concretely at
  $\ell=691$, with $\bar c_{k,2}$ factor degrees: ladder through $n=10$: $h^{1}_{\mathrm{ord}}=3$, $[1]$ (predicted by \HA{Hi11}; confirmed at $n=700$ by \file{hida\_towers/} and at $n=1390,2080,2770$ by the slot-14 cluster runs,
  \file{docs/implementation/slot14/WP-C3\_cluster\_run1\_report.md} \S1--2 (\texttt{hi11\_1390\_691}, \texttt{hi11\_2080\_691}) and
  \file{WP-C4\_cluster\_run1\_part2\_report.md} \S3 (\texttt{hi11\_2770\_691}), recorded in [C] as [D1]'s data remarks P2, P8;
  and at $n=4150$ by the September second-cluster run, the WP-C10 report, \S1 (\texttt{hi11\_4150\_691}), \HA{Hi11}'s
  v1.14.1 sentence; \COMP, data signed off by the owner 2026-09-29); through
  $198$: $33$, $[1,1,2,2,10]$; through $100$: $17$, $[1,7]$; through $200$: $33$, $[1,15]$; through $300$: $49$, $[5,19]$; through $400$: $67$, $[6,6,21]$. At
  $\ell=2411$ through $10$: $1$, no cuspidal factor (predicted by \HA{Hi12}; confirmed at $n=2420$, \file{WP-C3\_cluster\_run1\_report.md} \S1--2, \texttt{hi12\_2420\_2411}; [D1] P1;
  and at $n=4830$, \texttt{hi12\_4830\_2411}, the WP-C10 report, \S1, \HA{Hi12}'s v1.14.1 sentence). On level-$m$ ladders the polynomials agree modulo $\ell^{m}$.
\item \PROVED{} (Theorem~\ref{thm:u-etale}(3), Lemma~\ref{lem:u-tower}) \textbf{Eisenstein disjointness is a class invariant}: $v(c_{k,2}(1+2^{k-1}))=0$ at every rung of
  the ladders through $100,200,300,400$ at $691$ and of every rank-one class at $\ell=11,13$; and $\ge1$ at every rung through $10$ and $198$ at $691$.
\item \PROVED{} (Proposition~\ref{prop:u-determines}(2), Proposition~\ref{prop:u-eisblock}, Theorem~\ref{thm:u-691classes}) \textbf{which algebra}: $\Tu_{k,\mathfrak m}/\ell$ ---
  the algebra slot-14 WP-6's $\Z/\ell^{K}$ route reads as $A_K/\ell$ at the stable $K$, \emph{not} the census image $A$ on $H^{1}(\Fl)$ (the two agree iff
  $\delta(\mathfrak m)=0$, which is (U3), \CONJ). It is constant along a level-$1$ ladder on every non-Eisenstein block (there $\Tu_{k,\mathfrak m}=\hcusp_{k,\mathfrak m'}$
  and control at $m=1$, Proposition~\ref{prop:u-determines}(2)) and on the Eisenstein blocks of the classes through $10$ and $198$ at $691$ (Proposition~\ref{prop:u-eisblock},
  $L_k\ge1$ on the class, the cuspidal partner étale of residue degree~$1$): at $691$ it is a field on every block of the classes through $100,200,300,400$ and on
  the non-Eisenstein blocks through $198$; $\Fl[t]/(t^{2})$ on the Eisenstein blocks of the classes through $10$ and $198$. Constancy of $\Tu_{k,\mathfrak m}/\ell$ at a maximal
  ideal through which the Eisenstein component and a non-étale cuspidal component both pass is \emph{not} claimed ($\Tu(a)/\Pk{k}\to\Tu_k(a)$ may have a torsion
  kernel, e.g.\ $\mathrm{Tor}_1(\Iw/J,\Iw/\Pk{k})\ne0$ for $J=\mathfrak m_\Iw$). For the census $A$ see (U3).
\item \PROVED{} (Theorem~\ref{thm:u-ell13}, \HA{cor:ordprecision}(1), Proposition~\ref{prop:u-maxorder}) \textbf{the ordinary blocks of the ladders at
  $\ell\in\{5,7,11,13\}$}, at every rung $n\ge2$ of every shift that preserves the class ($s\equiv0\pmod{\ell-1}$): the ordinary part of
  $H^{1}(\Ga;\Mn)$ is the Eisenstein line alone, except at $\ell=11$, $n\equiv0\pmod{10}$ and at $\ell=13$, $n\equiv10\pmod{12}$, where it is the Eisenstein
  line plus the $2$-dimensional block of $\Delta$'s residual system ($(a_2,a_3,a_5)\equiv(9,10,1)$ at $11$, $a_{11}\equiv1$), whose algebra $A$ is $\Fl$
  (Loewy length~$1$) and which is disjoint from the Eisenstein system. Every other block met on these ladders (residual systems cut by every $T_p$, $p\le p_{\max}$,
  \textbf{and by $U_\ell$}, as the census does) --- every block with Loewy length $\ge2$ in particular --- is non-ordinary ($U_\ell$ nilpotent on it). (A block cut by
  the $T_p$, $p\ne\ell$, alone can mix the ordinary $\Delta$-part with non-ordinary classes of the same prime-to-$\ell$ system, e.g.\ the mod-$11$ form
  $\theta^{10}\Delta$, filtration $132$, $n=130$ on the ladder, $a_{11}=0$; at $\ell=13$ the mod-$13$ form $\theta^{12}\Delta$, filtration $156$,
  $n=154$, $a_{13}=0$, enters the $\Delta$ class the same way.)
\end{enumerate}
\paragraph{(U4) and (U7) at every $\ell$ with weight-shift data.} \COMP{} (slot-14 WP-C14 I-8,
\file{docs/implementation/slot14/wpc14/hida\_types.py}, \file{I-8\_report.md}; refereed; [EX] \texttt{s14:wshida}, \texttt{s14:depthlaw}) (U4) holds on
$7412$ ordinary pairs ($\ell=7$, $11$, $13$, $23$, $31$, $37$, $103$, $131$), all sharp; sharpness at the new $\ell$ ($23$, $31$, $37$, $103$, $131$) is at $m=1$
only, since no two computed weights there differ by $(\ell-1)\ell$. The ordinary rank is constant on $37$ classes and equals
$2\dim M^{\mathrm{ord}}_{k_c}(\Fl)-1$, $k_c\ge4$ the least weight of the class. It exceeds the regular value by exactly $2$ at the $11$ computed irregular
classes (Corollary~\ref{cor:u-irregular}): $g_E=1$ there, read at one weight of each class ($D=3$ on its Eisenstein block) and constant on the class by
Proposition~\ref{prop:u-determines}(2); whether $g_E=1$ at every irregular class is not claimed. On (U7)'s pattern at the other $\ell$: every ordinary
non-Eisenstein block met ($96$, all of rank one) and every Eisenstein block with $g=0$ ($297$) has $A=\Fl$, and the $17$ Eisenstein blocks with $g=1$ have $A=\Fl[t]/(t^{2})$ (Proposition~\ref{prop:u-eisblock}).

\paragraph{(U4) and (U7) on the weight-shift run at $\ell\in\{7,11,13\}$.} \COMP{} (v0.1.1-pre.5; second-cluster jobs \code{3465617} and \code{3467350},
all $287$ rungs, $n\le1220$, data signed off by the owner 2026-09-29; \file{docs/implementation/slot14/wpl1/w1\_287/} (\file{c9\_stdout.txt},
\file{hida\_pairs.tsv}); [EX] \texttt{s14:wsverdicts}, \texttt{s14:sept}, \texttt{s14:wshida}; [D1] P19, P21 with the counts re-derived on the
$287$ rungs) \textbf{(U7)}: the ordinary rank equals the prediction on $281$ / $281$ rungs (the other $6$ are controls), the ordinary blocks
are exactly those named on $353$ / $353$, and the $745$ / $745$ blocks of Loewy length $\ge2$ are non-ordinary. \textbf{(U4), sharp}: of the
$7412$ pairs above, the $7385$ at $\ell\in\{7,11,13\}$ --- $4540$ at $7$ (Eisenstein), $1122+946$ at $11$ (Eisenstein $+$ $\Delta$), $399+378$
at $13$, with $m=1$, $2$, $3$ on $6575$, $778$, $32$ --- are every pair of rungs of one class with a common ordinary system, and each is
congruent modulo $\ell^{m}$ and \emph{not} modulo $\ell^{m+1}$ (the least level over $p\le23$ and $\ell$): the level is exactly $m$, the extreme
pair being $(10,11)\to(1220,11)$, exactly modulo $11^{3}$ on both its Eisenstein and its $\Delta$ block. (The other $27$ pairs are
I-8's, at the five new $\ell$, $m=1$.) On the Eisenstein line this is lifting the exponent; on the $\Delta$-family it says that $a_p(k)$ has a
unit derivative in $k$ at some $p\le23$.

The \textbf{non-ordinary} part is \emph{not} governed by $\Tu$: nothing above predicts the census (non-ordinary) blocks or C-F9 (\Espec{pr:latticesplit}) along a ladder.

%----------------------------------------------------------------------
\section{Testability of (U1)--(U3) with the shipped routes}\label{sec:u-test}
%----------------------------------------------------------------------

The shipped tower route (\code{HeckeHidaHeavyCellProbe}, \HA{rmk:hidacheckable2}) and slot-14 WP-6's $\Z/\ell^{K}$ Loewy route; from the referee's
\S4 (\file{docs/implementation/slot14/WP-9\_referee.md}).

\begin{itemize}
\item \textbf{(U1): a theorem since v0.1.1-pre.2} (Prediction~\ref{pr:u-U1}); the cells below, computed before the re-tag (confirmed, [EX]
  \texttt{s14:zt}), are now checks of the implementation. \textbf{Testable at small irregular primes.} The $L_k=1$ rungs are theorems at $m=2$ (Theorem~\ref{thm:u-L691}); further rungs of the
  $691$ classes over $\Z/691^{2}$ only repeat that. A genuine test of the Bernoulli-zero prediction $L_k\ge2$ needs a rung with $v(B_k/k)=2$, out of reach at
  $691$ ($n=342250$) but cheap at smaller irregular pairs. \COMP{} (exact Bernoulli numbers: \file{u1\_feasible.py},
  each the first such rung of its class, linearity checked on a third point; the two cells re-checked in
  \file{u1\_cells\_check.py}; both outside the tree):
  $v_{103}(B_{228}/228)=2$ and $v_{37}(B_{284}/284)=2$ (also $v_{59}(B_{914}/914)=2$, $v_{271}(B_{1434}/1434)=2$). \textbf{Test cells}: $(n,\ell)=(226,103)$ and
  $(282,37)$ \textbf{over $\Z/\ell^{3}$} ($103^{3}=1092727$ and $37^{3}=50653$ are $<2^{32}$, the narrow band; the tower route at $m=3$), reading
  $\min_pv(c_{k,p}(1+p^{k-1}))$ from $\chi^{\mathrm{ord}}_{n,p}$ as in Theorem~\ref{thm:u-L691}. The value $2$ confirms (U1) as the theorem $L_k=2$ at that cell
  (Theorem~\ref{thm:u-sandwich}); a value $\ge3$ at every computed $p$ means too few $p$, a value $\le1$ is a defect. Recommended for WP-14c / the ladder proposal.
\item \textbf{(U2): not testable yet.} It needs two cuspidal components sharing a maximal ideal, and the natural candidate list has none
  (Proposition~\ref{prop:u-loewy}(4)). What the tower route \emph{can} test on the ramified cells $(30,67)$, $(26,131)$, \dots, $(68,157)$ (all tiny) is (U4)/(U6) on
  a ramified ordinary block --- $\bar c_{k,2}$ with a double root and $\hcusp_{k,\mathfrak m}/\ell=\Fl[t]/(t^{2})$ at every rung --- and, at $m=2$, that the double
  root stays unlifted. $L_k(P,P')$ itself is readable neither from characteristic polynomials (Proposition~\ref{prop:u-bounds}(2) only bounds it) nor from WP-6's route
  as shipped (its idempotents split residual systems, not components inside one system); it needs the $\ell$-adic factorisation of $c_{k,p}$ to precision $K$.
\item \textbf{(U3): testable now} with WP-6's \code{HECKE\_RING\_LOEWY\_CELLS}: on the Eisenstein blocks $(198,691)$ ($K=2$, $\approx3$\,GB) and $(700,691)$ ($K=2$,
  behind the $U_\ell$ lever), prediction $\varepsilon_K=0$, $A_1=A_K/691$ of Hilbert function $(1,1)$ as at $(10,691)$; on the ramified ordinary blocks of
  Proposition~\ref{prop:u-loewy}(4) (e.g.\ $(30,67)$, $K\le4$), $\varepsilon_K=0$, $A_1=A_K/\ell$ of Hilbert function $(1,1)$ --- there \PROVED, not only
  predicted, at the nine cells with $v_\ell(\operatorname{disc})=1$ and the two blocks of $(68,157)$ (maximal orders,
  Proposition~\ref{prop:u-maxorder}; held on both routes at all eleven, slot-14 WP-14c \S6), so a failure there is a defect.
\end{itemize}

%----------------------------------------------------------------------
\section{Scope, and the code map}\label{sec:u-scope}
%----------------------------------------------------------------------

\paragraph{Out of scope.} Tame level $N\ge3$ (CM $\cap$ $\GL(2)$, CM $\cap$ Eisenstein, the Eisenstein ideal at level $N$), after composite $N$ (\HA{sec:levelN}; slot~15)
--- \textbf{since v0.2.0-pre.1 treated in part~2, \S\ref{sec:u-tameN}}, except the Eisenstein ideal at level $N$;
$\ell\in\{2,3\}$; the $\Iw$-adic algebra as a computed object --- only its specialisations are computed, and every $\Iw$-adic statement here is read from them through
Theorem~\ref{thm:u-control} (\textbf{amended at v0.3.0-pre.1}: its truncations $\Tu(a)\otimes\Iw/(\ell^{K},X^{m})$ are in scope, part~3, \S\ref{sec:u-lambda}; the whole
$\Iw$-adic algebra, every $X$-adic order at once, stays out of scope); the class $a\equiv0$ in general (Remark~\ref{rmk:u-classzero}); the non-ordinary part.
\textbf{(U-1), a scope note} (v0.2.0-pre.1, owner ruling O4 of 2026-10-01; re-tagged at v0.2.0-pre.2 by the WP-14 referee's U-5): the hidden congruences of the
Frobenius-layer kind --- the ordinary block of an eigenform $f$ meeting its Frobenius partner $V(f)$ after a $\theta$-twist, [EX] \texttt{s15:hidden} --- are
\emph{non-ordinary}, so they have no consequence for $\Tu$ and leave the control statements of this document untouched. \PROVED{} for every block in the image of
$\theta^{j}$, $j\ge1$ ($U_\ell\theta=0$, \HA{prop:thetatwist}: such a block has $a_\ell=0$); that the Frobenius-layer hidden blocks are such images is [EX]
\texttt{s15:hidden} (\textsf{CMPD}). No statement of [U] changes with it.
\COMP{} (v0.2.1-pre.1; slot-16 WP-11, \file{docs/implementation/slot16/WP-11\_report.md}~\S2.1, \code{HeckeNebentypusValidation} arm (M2), gated) At the
two first Frobenius-layer hidden blocks, $(n,\ell)=(220,17)$ and $(246,19)$ (the system of $\theta\Delta$, $\dim=4$, $a_\ell=0$), every $T_p$, $p\le97$, acts on
$H^{1}(\Ga;\Sym^{n}\Fl^{2})_{\mathfrak m}$ as the scalar $p\,\tau(p)\bmod\ell$, while $\Talg_{\mathfrak m}/\ell\cong\Fl[x]/x^{2}$ (rank $2$ by \HA{lem:heckerankone}; for
$\Z_\ell[T_2,T_3,T_5,T_7,T_\ell]$ certified by \HA{cor:stablecert}): so the lattice $\Lambda_{\mathfrak m}$ needs $4$ generators and is \emph{not} free over
$\Talg_{\mathfrak m}$, the full Hecke algebra (a generation certificate by the $T_p$, $p\le\dim S_k$), and the saturation defect is $\mathsf A_1=\Fl\subsetneq\Talg_{\mathfrak m}/\ell$
(\HA{rmk:sattor}: $\dim I=1$). By the scope note above these blocks are non-ordinary, so this concerns neither $\Tu$ nor any statement of [U]; it is
recorded here, beside (U-1), and not at Proposition~\ref{prop:u-loewy}, because the blocks are non-ordinary (v0.2.1-pre.2 wording, the slot-17 referee F1's
m-U2); it bears on that proposition only by contrast: there the freeness of the ordinary blocks is what links saturation and Loewy length, and the
Frobenius-layer blocks show that the link does not extend to them.

\paragraph{Code map.} \textbf{Not built.} The readings of \S\ref{sec:u-finite}--\S\ref{sec:u-cc} were made by out-of-tree scripts over the published catalogue
(\file{docs/implementation/slot14/wp8/}); a component-reading validation is slot-14 WP-14c (stretch) or slot~15. The data came from the shipped
\code{HeckeHidaHeavyCellProbe} (\HA{rmk:hidacheckable2}) and the published \file{hida\_towers/} catalogue. \textbf{Since v0.3.0-pre.1}: the truncated $\Iw$-adic arm is planned as slot-19 WP-HL, gated on the referee of part~3 (its code map,
\S\ref{sec:u-lambda}); \textbf{built} at v0.3.0-pre.4.

%----------------------------------------------------------------------
\section{Tame level \texorpdfstring{$N$}{N} --- CM components, CM \texorpdfstring{$\cap$}{cap} GL(2), CM \texorpdfstring{$\cap$}{cap} Eisenstein (part~2)}\label{sec:u-tameN}
%----------------------------------------------------------------------
% NEW at v0.2.0-pre.1 (hecke-extensions slot 16, WP-12; plan docs/implementation/slot16/00-slot16-plan.md S2 "WP-12", [U] part 2).
% Placed after sec:u-scope so that no section number of v0.1.x moves.  Contract only; nothing is built.  The finite-level
% cohomological side is [C] sec:nebentypus (released v1.15.0; code slot-16 WP-10), and the sanity cells are [C] (NB6).

\paragraph{Status and scope of this part.} Written as a pre-release (v0.2.0-pre.1), refereed by slot-16 WP-14 (Fable, ACCEPT WITH REPAIRS, repaired at
pre.2 and confirmed), \textbf{released and owner-signed at v0.2.0} (ruling O2 of 2026-10-01). It extends \S\ref{sec:u-cm} from tame level $1$ to tame level $N\ge1$, $\ell\nmid N$, $\ell\ge5$, with a
tame character. Its three questions are the ones \S\ref{sec:u-scope} deferred: \emph{which} components have CM (Theorem~\ref{thm:u-cmN}),
\emph{whether} a CM component meets a non-CM cuspidal one (CM $\cap$ $\GL(2)$, Remark~\ref{rmk:u-cmgl2}) and \emph{whether} it meets an Eisenstein one
(CM $\cap$ Eisenstein, Proposition~\ref{prop:u-cmeis}). Every statement here is about eigensystems and residual representations; the
$\Iw$-adic algebras are cited, not computed, as in Part~1. The finite-level reading is the ordinary part of the character model of
\HA{def:nebmodel} --- \HA{sec:nebentypus}'s $H^{1}(\Ga_0(N);\VR(\chi))$ over $\F_q$ or $W_K(\F_q)$ --- and its sanity cells are \HA{NB6}.

\paragraph{Conventions of this part.} $N\ge1$ with $\ell\nmid N$; $\chi\colon G_N=(\Z/N)^{\times}\to\F_q^{\times}$ a character (\HA{def:nebdata}(4); $\ell\nmid\operatorname{ord}\chi$),
$\mathcal O:=W(\F_q)$, $\tilde\chi:=\mathfrak t\circ\chi$ its Teichm\"uller lift (\HA{lem:teichmuller}), $\Iw_{\mathcal O}:=\mathcal O[[X]]$. $K$ is an imaginary quadratic
field of discriminant $d_K$, $\chi_K$ its quadratic character, $c$ complex conjugation, $w_K:=\abs{\mathcal O_K^{\times}}$. A fixed embedding
$\overline\Q\hookrightarrow\Qlb$ determines, when $\ell$ splits in $K$, the prime $\mathfrak l\mid\ell$ of $K$ it induces, and $\bar{\mathfrak l}$ its conjugate.
$\omega_{\mathfrak l}$ is the Teichm\"uller character of $I_{\mathfrak l}$. Symbols new to this part: $\hcusp(N,\chi;a)$, $\hcusp_k(N,\chi)$, $\theta_\lambda$, $\varepsilon_\lambda$,
$\rho_\lambda$; they mean nothing else.

\begin{definition}[The ordinary algebra at tame level $N$ with character]\label{def:u-hN}
For a class $a\in\Z/(\ell-1)$ with $\chi(-1)=(-1)^{a}$, $\hcusp(N,\chi;a)$ is Hida's cuspidal ordinary Hecke algebra of tame level $N$, tame character
$\tilde\chi$ and class $a$: the $\Iw_{\mathcal O}$-algebra generated by $T(p)$, $p\nmid N\ell$, and $U(q)$, $q\mid N\ell$, on the ordinary $\Iw$-adic cusp forms of
that tame level and character. For $k\ge2$ with $k\equiv a\pmod{\ell-1}$, $\hcusp_k(N,\chi)$ is the $\mathcal O$-algebra generated by the same operators on
the ordinary part $e\,S_k(\Ga_0(N\ell),\tilde\chi;\mathcal O)$, $e=\lim U(\ell)^{j!}$. At $N=1$ it is Definition~\ref{def:u-hcusp}. \emph{Components} and
\emph{meeting} are as in Proposition~\ref{prop:u-components}(2) and Definition~\ref{def:u-congmod}: two components meet iff they share a maximal ideal.
A component has \emph{CM by $K$} if the image of $T(p)$ vanishes for every $p\nmid N\ell$ inert in $K$ (Definition~\ref{def:u-cm}).
\end{definition}

\begin{theorem}[Hida: freeness and control at tame level $N$]\label{thm:u-controlN}
\LIT{\cite[Thm.~3.1 and Cor.~3.2, p.~250]{hida_ens}, verified (page image) by slot-16 L16-3, \file{docs/implementation/slot16/locators.tsv}: Cor.~3.2 reads
``$\mathfrak h_0(N,a;\mathcal O_K)/P_k\mathfrak h_0(N,a;\mathcal O_K)\cong\mathfrak h^{0}_k(\Ga_1(N_0p),a;\mathcal O_K)$ if $k\ge2$'', i.e.\ the source is at tame level $N_0$ on
$\Ga_1(N_0p)$. The character-$\chi$ algebra of Definition~\ref{def:u-hN} is the direct factor cut out by the idempotent of $\chi$ on the prime-to-$\ell$ part of
$(\Z/N)^{\times}$ ($\ell\nmid\operatorname{ord}\chi$), so freeness and control pass to it; \VERIFYloc{} only for the identification of Hida's $\mathfrak h_0(N,a)$ with the
$\Ga_1(N_0p)$-algebra of all tame characters} $\hcusp(N,\chi;a)$ is free of finite rank over $\Iw_{\mathcal O}$, and for every $k\ge2$ in the class,
$\hcusp(N,\chi;a)/\Pk{k}\cong\hcusp_k(N,\chi)$ with $T(p)\mapsto T_p$, $U(q)\mapsto U_q$.
\end{theorem}

\begin{theorem}[Where CM can live]\label{thm:u-cmN}
\PROVED{} from \STD{Deligne and Deligne--Serre/Ribet: for a normalised cuspidal eigenform $f$ of weight $k\ge2$, level $N\ell^{r}$ and character $\psi$
there is a continuous irreducible $\rho_f\colon G_\Q\to\GL_2(\Qlb)$, unramified outside $N\ell$, with $\mathrm{tr}\,\rho_f(\Frob_p)=a_p$ and
$\det\rho_f=\psi\chi_{\mathrm{cyc}}^{k-1}$ (\cite{deligne_formes}, \cite[\S2]{ribet_nebentypus}; \VERIFYloc)}, \STD{Chebotarev density, Brauer--Nesbitt, and Clifford
theory: an irreducible $\rho\cong\rho\otimes\chi_K$ is induced from $G_K$} and \LIT{Wiles \cite[Thm.~2.1.4, p.~561]{wiles_ordinary} (verified, page image, slot-16 L16-3, \file{docs/implementation/slot16/locators.tsv}): an $\ell$-ordinary $\rho_f$ is
reducible on $G_{\Ql}$ with an unramified quotient}. Let a component of $\hcusp(N,\chi;a)$ have CM by $K$. Then:
\begin{enumerate}
\item $\ell$ \textbf{splits} in $K$;
\item every prime dividing $d_K$ divides $N$;
\item \STD{a CM newform by $K$ has level $\abs{d_K}\,\mathrm{N}(\mathfrak c)$, $\mathfrak c$ the conductor of its Hecke character, \cite[\S3]{ribet_nebentypus} (\VERIFYloc)} hence
  $\abs{d_K}\mid N$; in particular $N\ge3$, and \textbf{tame levels $1$ and $2$ carry no CM component} (at $N=1$ this is Theorem~\ref{thm:u-nocm}).
\end{enumerate}
\end{theorem}

\begin{proof}
Take $k\ge3$ in the class and a point of the component above $\Pk{k}$ (Theorem~\ref{thm:u-controlN}): it is the eigensystem of an $\ell$-ordinary eigenform $f$ of weight $k$,
level $N\ell^{r}$, character $\psi=\tilde\chi$ (as $k\equiv a$; Definition~\ref{def:u-hN}'s $e\,S_k(\Ga_0(N\ell),\tilde\chi)$), with $a_p(f)=0$ for $p\nmid N\ell$ inert in $K$. As in the proof of Theorem~\ref{thm:u-nocm},
$\rho_f\cong\rho_f\otimes\chi_K$; $\rho_f$ is irreducible, so $\rho_f\cong\mathrm{Ind}_{G_K}^{G_\Q}\lambda$ for a character $\lambda$ of $G_K$. On $I_\ell$,
$\det\rho_f=\omega^{a-k}\chi_{\mathrm{cyc}}^{k-1}$ ($\tilde\chi$ is unramified at $\ell$, as $\ell\nmid N$), of infinite order for $k\ge2$; and by Wiles,
$\rho_f|_{G_{\Ql}}$ has a sub $\varepsilon_1$ and an unramified quotient $\varepsilon_2$.
\emph{(2)} A prime $q\ne\ell$ dividing $d_K$ but not $N$: $\rho_f$ is unramified at $q$ while $\chi_K|_{I_q}\ne1$ --- the contradiction of
Theorem~\ref{thm:u-nocm}(a).
\emph{(1), $\ell$ ramified in $K$:} $\{\varepsilon_1,\varepsilon_2\}=\{\varepsilon_1\chi_K,\varepsilon_2\chi_K\}$ on $G_{\Ql}$ and $\chi_K|_{I_\ell}\ne1$ force
$\varepsilon_1|_{I_\ell}=\chi_K|_{I_\ell}$ (Theorem~\ref{thm:u-nocm}(b)), so $\det\rho_f|_{I_\ell}$ has order $2$, against its infinite order.
\emph{(1), $\ell$ inert in $K$:} $K_\ell/\Ql$ is unramified, so $I_\ell\subseteq G_{K}$ (after the embedding), $\rho_f|_{G_{K_\ell}}=\lambda\oplus\lambda^{c}$, and
$\{\varepsilon_1,\varepsilon_2\}|_{I_\ell}=\{\lambda,\lambda^{c}\}|_{I_\ell}$; $\varepsilon_2|_{I_\ell}=1$ makes $\lambda|_{I_\ell}$ or $\lambda^{c}|_{I_\ell}$ trivial, hence both (choose the lift $\tilde c$ of $c$ in $D_\ell$ --- possible, $D_\ell\to\mathrm{Gal}(K_\ell/\Ql)=\mathrm{Gal}(K/\Q)$ being onto as $\ell$ is inert ---, which
normalises $I_\ell$; WP-14 referee U-3), so $\rho_f|_{I_\ell}$ is unipotent and $\det\rho_f|_{I_\ell}=1$ --- a contradiction. So $\ell$ splits.
\emph{(3)} By (1) $\ell\nmid d_K$; the newform underlying $f$ has level $\abs{d_K}\mathrm{N}(\mathfrak c)\mid N\ell^{r}$, so $\abs{d_K}\mid N$; $\abs{d_K}\ge3$.
\end{proof}

\begin{proposition}[CM components exist where Theorem~\ref{thm:u-cmN} allows]\label{prop:u-cmexist}
\LIT{\cite[Thm.~7.1, pp.~268--269]{hida_ens}, verified for the construction under $\ell$ split (Remark~\ref{rmk:u-cmroute}); the exact tame character and class
of the family in the source's notation \VERIFYloc} If $\ell$ splits in $K$ and $\mathfrak c$ is an ideal of $\mathcal O_K$ prime to $\ell$, there is a CM component by $K$ of tame
level $\abs{d_K}\mathrm{N}(\mathfrak c)$. \PROVED{} (classical, from \STD{Hecke's theta series: for a Hecke character $\lambda$ of $K$ of conductor $\mathfrak c$ and
$\infty$-type $\alpha\mapsto\alpha^{k-1}$, $\theta_\lambda=\sum_{\mathfrak a}\lambda(\mathfrak a)q^{\mathrm N\mathfrak a}\in S_k(\Ga_0(\abs{d_K}\mathrm N\mathfrak c),\chi_K\varepsilon_\lambda)$,
$\varepsilon_\lambda$ the restriction of $\lambda$'s finite part to $\Z$; \cite[\S3]{ribet_nebentypus} (\VERIFYloc)}) At a weight $k$, $\theta_\lambda$ is $\ell$-ordinary iff
$\ell$ splits in $K$: for $\ell=\pi\bar\pi$, $a_\ell=\lambda(\pi)+\lambda(\bar\pi)$ is a unit at $\mathfrak l$ (exactly one summand is), and for $\ell$ inert $a_\ell=0$.
\end{proposition}

\begin{proposition}[CM $\cap$ Eisenstein]\label{prop:u-cmeis}
\PROVED{} from Theorem~\ref{thm:u-cmN}'s inputs and Proposition~\ref{prop:u-cmexist}. Let $P$ be a CM component by $K$ ($\ell=\mathfrak l\bar{\mathfrak l}$ split), whose weight-$k$
points are $\theta_\lambda$ with $\lambda$ of $\infty$-type $k-1$, $\rho_{\theta_\lambda}=\mathrm{Ind}_{G_K}^{G_\Q}\rho_\lambda$, where $\bar\rho_\lambda$ is unramified outside $\mathfrak c\ell$, $\bar\rho_\lambda|_{I_{\mathfrak l}}=\omega^{k-1}$ and
$\bar\rho_\lambda|_{I_{\bar{\mathfrak l}}}=1$ (the $\infty$-type through the chosen embedding; $\mathfrak c$ prime to $\ell$; WP-14 referee U-2).
\begin{enumerate}
\item If $P$ meets an Eisenstein component (of the ordinary algebra with its Eisenstein components, Definition~\ref{def:u-T}'s analogue), then
  $\mathrm{Ind}\,\bar\rho_\lambda$ is reducible, which holds iff $\bar\rho_\lambda=\bar\rho_\lambda^{c}$; and that forces \textbf{$a\equiv1\pmod{\ell-1}$} (odd weights); then $\bar\rho_\lambda$ is unramified at $\ell$, and $\bar\varepsilon:=\bar\rho_\lambda$ is a $c$-invariant
  character of $G_K$ unramified outside $\mathfrak c$. Then $\bar\rho_{P}^{\mathrm{ss}}\cong\mu\oplus\mu\chi_K$ for a character $\mu$ of $G_\Q$ with $\mu|_{G_K}=\bar\varepsilon$.
\item \textbf{The untwisted family meets Eisenstein.} Assume $h_K=1$ (WP-14 referee U-1: for $h_K>1$ the character below is defined on principal ideals only, its
  $h_K$ extensions differ by class-group characters, and the congruence is the one of (1) with $E_k(\mu,\mu\chi_K)$, $\mu$ the character of (1), $\mu=1$ when
  $h_K=1$). Let $\mathfrak c=1$, $\lambda=\lambda_k$ the character $(\alpha)\mapsto\alpha^{k-1}$ (defined iff $w_K\mid k-1$), so
  $\theta_{\lambda_k}\in S_k(\Ga_0(\abs{d_K}),\chi_K)$. If $k\equiv1\pmod{\ell-1}$ then, modulo $\mathfrak l$, $a_p(\theta_{\lambda_k})\equiv1+\chi_K(p)\,p^{k-1}$ for every
  $p\nmid d_K\ell$: the eigensystem of $\theta_{\lambda_k}$ is congruent to that of the Eisenstein series $E_k(\mathbf 1,\chi_K)$ of the same level and character.
\end{enumerate}
\end{proposition}

\begin{proof}
(1) Components meeting share a maximal ideal, hence a residual eigensystem; an Eisenstein one is $p\mapsto\psi_1(p)+\psi_2(p)p^{k-1}$, the trace of a reducible
representation, so by Brauer--Nesbitt $\mathrm{Ind}\,\bar\rho_\lambda$ is reducible; by Mackey/Clifford an induced representation of a character of an index-$2$
subgroup is reducible iff the character is $c$-invariant. $\rho_\lambda$ is unramified at $\bar{\mathfrak l}$ and equals $\chi_{\mathrm{cyc}}^{k-1}$ times an unramified character on
$I_{\mathfrak l}$ (the $\infty$-type through the chosen embedding; $\mathfrak c$ prime to $\ell$), so $\bar\rho_\lambda|_{I_{\mathfrak l}}=\omega_{\mathfrak l}^{k-1}$ and
$\bar\rho^{c}_\lambda|_{I_{\mathfrak l}}=\bar\rho_\lambda|_{I_{\bar{\mathfrak l}}}\circ c=1$. Equality forces $\omega^{k-1}=1$, i.e.\ $(\ell-1)\mid(k-1)$; then $\bar\rho_\lambda$ is unramified at $\ell$,
$\bar\varepsilon:=\bar\rho_\lambda$, and $\bar\varepsilon=\bar\varepsilon^{c}$. A $c$-invariant character of $G_K$ extends to $G_\Q$ in two ways, $\mu$ and $\mu\chi_K$, and $\mathrm{Ind}\,\bar\varepsilon=\mu\oplus\mu\chi_K$.
(2) For $p=\pi\bar\pi$ split ($\mathfrak p=(\pi)$ principal, as $h_K=1$), $a_p=\pi^{k-1}+\bar\pi^{k-1}$ and $\pi,\bar\pi$ are units at $\mathfrak l$ with residue field $\F_\ell$, so both powers are $\equiv1$ when
$(\ell-1)\mid k-1$; $1+\chi_K(p)p^{k-1}\equiv2$. For $p$ inert, $a_p=0$ and $1+\chi_K(p)p^{k-1}=1-p^{k-1}\equiv0$.
\end{proof}

\begin{remark}[CM $\cap$ $\GL(2)$]\label{rmk:u-cmgl2}
\CONJ{} with \LIT{\cite{hida_tilouine}, \cite{mazur_tilouine}; \VERIFYloc} (1) \PROVED{} A CM component $P$ meets a non-CM cuspidal component only at a maximal ideal whose
residual $\bar\rho$ is $\mathrm{Ind}_{G_K}\bar\rho_\lambda$ --- dihedral if $\bar\rho_\lambda\ne\bar\rho_\lambda^{c}$, Eisenstein (Proposition~\ref{prop:u-cmeis}) otherwise; a
non-CM component through a dihedral maximal ideal is the phenomenon Remark~\ref{rmk:u-dihedral}(ii) met at tame level~$1$ (there with no CM component at all). (2) The
congruence module of $P$ against the rest of the algebra is governed by an anticyclotomic Katz $\ell$-adic $L$-function of $K$ --- one divisibility by Hida--Tilouine,
the converse the anticyclotomic main conjecture (Mazur--Tilouine), both \VERIFYloc{} (WP-14 referee U-6) --- so CM $\cap$ $\GL(2)$ at $P$ is predicted exactly where $\ell$ divides that $L$-value; no numerical form of this is stated here.
\COMP{} reading for the code: at a cell of the character model on $P$'s weight ladder, the ordinary generalised eigenspace at $P$'s residual system has $\F_q$-dimension
$2\cdot\#\{$CM eigenforms of $P$ at that weight$\}$ plus the non-CM contribution; a positive non-CM contribution is a CM $\cap$ $\GL(2)$ instance.
\end{remark}

\begin{remark}[(U8)--(U10): the predictions of part 2, on \HA{NB6}'s cells]\label{rmk:u-pred2}
Committed at v0.2.0-pre.1, before slot-16 WP-10's and WP-11's runs (rule~20); a failing \PROVED{} item is a defect, a failing \CONJ{} item a finding. Cells are
$(N,\chi,k,\ell)$ of the character model, $\ell\ge5$, $\ell\nmid N$.
\begin{enumerate}[label=(U\arabic*),start=8]
\item\label{U8} \PROVED{} (Proposition~\ref{prop:u-cmexist}, with \HA{NB6}'s classical input) \textbf{Ordinarity of the CM cells.} At $(4,\chi_{-4},5)$, $(3,\chi_{-3},7)$,
  $(7,\chi_{-7},3)$ --- the forms $\theta_{\lambda_k}$, $K=\Q(i),\Q(\sqrt{-3}),\Q(\sqrt{-7})$ --- the CM eigensystem lies in the ordinary part of $H^{1}$ iff $\ell$ splits in
  $K$ ($\ell\equiv1\pmod4$; $\ell\equiv1\pmod3$; $\ell\equiv1,2,4\pmod7$); for inert $\ell$ its $T_\ell$-eigenvalue is $0$.
\item\label{U9} \PROVED{} (Theorem~\ref{thm:u-cmN}) \textbf{No CM at tame level $1$ or $2$, and what the code must not read as CM.} No ordinary
  eigensystem in characteristic~$0$ at $N\in\{1,2\}$ has CM. Modulo $\ell$ an ordinary block can still have $T_p$ nilpotent at every $p$ inert in some $K$: at
  $N=1$ these are the residually dihedral blocks of Remark~\ref{rmk:u-dihedral}(i), $K=\Q(\sqrt{-\ell})$. A block with that property is \emph{residually
  dihedral}, not CM; CM is a statement about the component (Definition~\ref{def:u-hN}), read on characteristic-$0$ or $\Z/\ell^{K}$ data at large $K$.
\item\label{U10} \PROVED{} (Proposition~\ref{prop:u-cmeis}(2); $h_K=1$ at \HA{NB6}'s three fields $\Q(i)$, $\Q(\sqrt{-3})$, $\Q(\sqrt{-7})$) \textbf{CM $\cap$ Eisenstein.} At $(4,\chi_{-4},5,5)$ and $(3,\chi_{-3},7,7)$ --- weight $\equiv1\pmod{\ell-1}$,
  $\ell$ split --- and at $(7,\chi_{-7},11,11)$ (on the CM ladder of $(7,\chi_{-7},3)$, $w_K=2$), the CM eigensystem of $H^{1}(\Ga_0(N);V_{\F_\ell}(\chi))$ coincides modulo
  $\ell$ with the Eisenstein system $p\mapsto1+\chi(p)p^{k-1}$: the joint generalised eigenspace of $(T_p)_{p\nmid N\ell}$ at that system contains the Eisenstein
  classes and the two CM classes; both Eisenstein systems $1+\chi(p)p^{k-1}$ and $\chi(p)+p^{k-1}$ coincide modulo $\ell$ there ($p^{k-1}\equiv1$ for $(\ell-1)\mid(k-1)$),
  so ``the Eisenstein classes'' means both and the eigenspace has dimension $\ge4$ (WP-14 referee U-7). \CONJ{} The local algebra of the $T_p$ ($p\nmid N\ell$) on that eigenspace is non-reduced (Loewy length $\ge2$), as at the
  Eisenstein blocks of tame level~$1$ with $g=1$ (Proposition~\ref{prop:u-eisblock}).
\end{enumerate}
\end{remark}

\paragraph{Code map of part 2.} \textbf{Not built.} The readings are WP-10's (\code{HeckeNebentypusValidation}, the \HA{NB6} arm) and WP-11's runs; no
validation of [U] is scheduled in slot~16.

%----------------------------------------------------------------------
\section{The \texorpdfstring{$\Iw$}{Lambda}-adic algebra, truncated (part~3)}\label{sec:u-lambda}
%----------------------------------------------------------------------
% NEW at v0.3.0-pre.1 (hecke-extensions slot 19, WP-2; plan docs/implementation/slot19/00-slot19-plan.md S2 "WP-2", headline 5;
% owner ruling O5 of 2026-10-03: the Lambda-adic statement now).  Placed after part 2 so that no section number of v0.2.x moves.
% Contract only; the code arm is slot-19 WP-HL (upper core), built only if the referee F1 accepts this text (rule 22).

\paragraph{Status and scope of this part.} Written as a pre-release (v0.3.0-pre.1) for the slot-19 Fable referee F1 (WP-3), revised at v0.3.0-pre.2 on its ACCEPT WITH REPAIRS; \textbf{binding
since v0.3.0} (the owner's release at the slot-19 close, ruling O3). Part~1 computes only specialisations of the $\Iw$-adic algebras and reads every $\Iw$-adic statement off
them through Theorem~\ref{thm:u-control} (\S\ref{sec:u-scope}). This part makes a \emph{finite} piece of the $\Iw$-adic algebra itself a computed
object: its truncations modulo $(\ell^{K},X^{m})$ --- the ``$\mathbf T^{\mathrm{ord}}\otimes\Iw/(\ell^{k},X^{m})$'' of the slot-19 plan, written
with $K$ for the $\ell$-power because $k$ is a weight here. It states (i) what the truncation is and that finitely many weights determine it,
with explicit precision (\PROVED, elementary, from Part~1); (ii) which control statement is used and which literature it rests on;
(iii) what is \PROVED{} at the slot-19 targets $\ell=691$, $(282,37)$, $(226,103)$, and what is \COMP{} or \CONJ. Tame level~$1$, $\ell\ge5$, as in Part~1.

\begin{definition}[Truncations]\label{def:u-trunc}
For $K,m\ge1$ put $\Iw_{K,m}:=\Iw/(\ell^{K},X^{m})=(\Z/\ell^{K})[X]/(X^{m})$, and for $x\in\ell\Zl$, $\Iw_{K,m}(x):=\Iw/(\ell^{K},(X-x)^{m})$. For a
class $a$ and a finitely generated $\Iw$-module $M$ put $M_{K,m}:=M\otimes_\Iw\Iw_{K,m}$ and $M_{K,m}(x):=M\otimes_\Iw\Iw_{K,m}(x)$; in particular
\[
  \Tu_{K,m}(a):=\Tu(a)/(\ell^{K},X^{m})\Tu(a),\qquad \hcusp_{K,m}(a),\qquad \Cg_{\Eis;K,m}(a):=\Cg_\Eis(a)/(\ell^{K},X^{m}),
\]
finite $\Z/\ell^{K}$-algebras (Definitions~\ref{def:u-hcusp}, \ref{def:u-T}, \ref{def:u-congmod}), with the images of $T(p)$, $U$ as named operators.
For an integer $k$ write $x_k:=u^{k}-1$, so that $\Pk{k}=X-x_k$ and $v(x_k)=1+v(k)$ (the proof of Lemma~\ref{lem:u-tower}). The truncation
\emph{centred at the weight $k$} is $M_{K,m}(x_k)$. For the named algebras the argument is written after a semicolon:
$\Tu_{K,m}(a;x):=\Tu(a)_{K,m}(x)$, $\hcusp_{K,m}(a;x):=\hcusp(a)_{K,m}(x)$, $\Cg_{\Eis;K,m}(a;x):=\Cg_\Eis(a)_{K,m}(x)$ (v0.3.0-pre.2).
\end{definition}

\begin{proposition}[First properties]\label{prop:u-truncbasic}
\PROVED{} from Theorem~\ref{thm:u-control}, Proposition~\ref{prop:u-components}, Lemma~\ref{lem:u-deligne} and Prediction~\ref{pr:u-U1} (itself \PROVED{} from \LIT{Emerton}).
Let $a$ be an even class.
\begin{enumerate}
\item $\hcusp_{K,m}(a)$ is free over $\Iw_{K,m}$ of rank $d(a)$, hence free over $\Z/\ell^{K}$ of rank $m\,d(a)$. If $a\not\equiv0$, $\Tu_{K,m}(a)$ is free over
  $\Iw_{K,m}$ of rank $d(a)+1$, and $\Cg_{\Eis;K,m}(a)\cong\Iw_{K,m}/(G_a)$.
\item \textbf{(Control for $\Tu$.)} If $a\not\equiv0$, then $\Tu(a)/\Pk{k}\Tu(a)\cong\Tu_k(a)$ for every $k\ge3$ in the class, $T(p)\mapsto T_p$, $U\mapsto U$.
\item \textbf{(At $m=1$ the truncation is a finite-level object.)} $\Tu_{K,1}(a;x_k)\cong\Tu_k(a)/\ell^{K}$ ($a\not\equiv0$) and $\hcusp_{K,1}(a;x_k)\cong\hcusp_k/\ell^{K}$
  for $k\ge3$ in the class; and $\Tu_{K,1}(a)=\Tu_{K,1}(a;x_k)$ for any such $k$ with $\ell^{K-1}\mid k$ (one exists by the Chinese remainder theorem).
\item \textbf{(Centres.)} If $v(x-x')\ge K$, then $(\ell^{K},(X-x)^{m})=(\ell^{K},(X-x')^{m})$. Hence $M_{K,m}(x_k)$ depends only on $k$ modulo
  $(\ell-1)\ell^{K-1}$, i.e.\ only on the tower of level $K$ through $k$ (\HA{def:hidatower}); at $m=1$ this is Lemma~\ref{lem:u-tower}.
\end{enumerate}
\end{proposition}

\begin{proof}
(1) $\hcusp(a)$ is $\Iw$-free of rank $d(a)$ (Theorem~\ref{thm:u-control}) and, for $a\not\equiv0$, $\Tu(a)$ is $\Iw$-free of rank $d(a)+1$ (Prediction~\ref{pr:u-U1},
(U1-$\Iw$)); base change preserves freeness and rank. $\Cg_\Eis(a)\cong\Iw/(G_a)$ is (U1-$\Iw$), and $\otimes\,\Iw_{K,m}$ is right exact.
(2) $\Tu(a)/\Pk{k}$ is $\Zl$-free of rank $d(a)+1$ by (1) and maps onto $\Tu_k(a)$ (the generators go to the generators); $\Tu_k(a)$ is $\Zl$-free of rank
$d_k+1=d(a)+1$ (Lemma~\ref{lem:u-deligne}, Theorem~\ref{thm:u-control}). A surjection between free $\Zl$-modules of the same finite rank is an isomorphism.
(3) $(\ell^{K},X-x_k)$ is the ideal $(\ell^{K},\Pk{k})$, and the quotients are $\Tu(a)/\Pk{k}$ and $\hcusp(a)/\Pk{k}$ reduced modulo $\ell^{K}$: (2) and
Theorem~\ref{thm:u-control}. $(\ell^{K},X)=(\ell^{K},\Pk{k})$ iff $v(x_k)\ge K$ iff $\ell^{K-1}\mid k$.
(4) $(X-x)^{m}=\sum_{j}\binom mj(X-x')^{m-j}(x'-x)^{j}$, and every term with $j\ge1$ lies in $\ell^{K}\Iw$; symmetrically. $v(x_k-x_{k'})=1+v(k-k')$, and
$k\equiv k'\pmod{\ell-1}$ inside a class.
\end{proof}

\begin{lemma}[A product of weight primes]\label{lem:u-wtprod}
\PROVED{} Let $k_1,\dots,k_N$ be integers, $x_j:=x_{k_j}$. If the sum of the $N-m+1$ smallest of the $v(x_j)$ is $\ge K$ --- in particular if
$N\ge K+m-1$ --- then $\prod_{j=1}^{N}(X-x_j)\in(\ell^{K},X^{m})$.
\end{lemma}

\begin{proof}
The coefficient of $X^{t}$, $t<m$, in $\prod_j(X-x_j)$ is $\pm e_{N-t}(x_1,\dots,x_N)$, a sum of products of $N-t\ge N-m+1$ distinct $x_j$; each $v(x_j)\ge1$.
\end{proof}

\begin{theorem}[Interpolation from finitely many weights]\label{thm:u-interp}
\PROVED{} (elementary; Weierstrass division in $\Iw$). Let $g\in\Iw$, let $k_1,\dots,k_N$ be pairwise distinct integers with $N\ge K+m-1$, $x_j:=x_{k_j}$, and let
$r\in\Q_\ell[X]$ be the interpolation polynomial of degree $<N$ with $r(x_j)=g(x_j)$.
\begin{enumerate}
\item $r\in\Zl[X]$, and $g\equiv r\pmod{(\ell^{K},X^{m})}$: the truncation of $g$ is determined by the $N$ values $g(x_{k_j})$.
\item \textbf{(Precision.)} Suppose $k_j=k_0+(j-1)(\ell-1)$, $1\le j\le N\le\ell$ (consecutive weights of one class), and the values are known only modulo
  $\ell^{K'}$: $\tilde g_j\equiv g(x_j)\pmod{\ell^{K'}}$. Let $\tilde r$ be computed from the $\tilde g_j$ by Newton's divided differences (the $\tilde D_i$ formed in $\Q_\ell$: the $\tilde g_j$ lifted to $\Z$ before
  dividing by $\ell^{i}$; the error analysis below covers it). If $K'\ge K+m-1$,
  then $\tilde r\equiv g\pmod{(\ell^{K},X^{m})}$.
\item \textbf{(A further weight is a check.)} In the situation of~(2), for every integer $k'$ in the class, $\tilde r(x_{k'})\equiv g(x_{k'})\pmod{\ell^{\min(K',N)}}$.
\end{enumerate}
\end{theorem}

\begin{proof}
Write $r=\sum_{i=0}^{N-1}D_i\,B_i$ in Newton form, $B_i:=\prod_{j\le i}(X-x_j)$, $D_i:=g[x_1,\dots,x_{i+1}]$ the divided difference. (1) For $g=X^{n}$,
$D_i=h_{n-i}(x_1,\dots,x_{i+1})$ (complete homogeneous), integral; for $g=\sum c_nX^{n}$ the series $\sum c_nh_{n-i}(x)$ converges in $\Zl$ ($v(h_{n-i}(x))\ge
n-i$), so $D_i\in\Zl$ and $r\in\Zl[X]$. $g-r\in\Iw$ vanishes at the distinct points $x_j\in\ell\Zl$; division by the distinguished polynomial $X-x_1$ gives
$g-r=(X-x_1)h_1$ with $h_1\in\Iw$, and $h_1(x_j)=0$ for $j\ne1$ ($\Zl$ is a domain); by induction $g-r\in\prod_j(X-x_j)\,\Iw$, which lies in $(\ell^{K},X^{m})$ by
Lemma~\ref{lem:u-wtprod}. (2) Here $v(x_a-x_b)=1+v((a-b)(\ell-1))=1$ for $a\ne b$ ($|a-b|<\ell$), so $D_i=\sum_{a\le i+1}g(x_a)/\prod_{b\le i+1,b\ne a}(x_a-x_b)$
has denominators of valuation exactly $i$, and an error of valuation $\ge K'$ in the values changes $D_i$ by an error of valuation $\ge K'-i$. Modulo $X^{m}$
the coefficients of $B_i$ are $\pm e_{i-t}(x_1,\dots,x_i)$, $t<m$, of valuation $\ge\max(0,i-m+1)$. So the error of $D_iB_i$ modulo $X^{m}$ has valuation
$\ge K'-i+\max(0,i-m+1)\ge K'-(m-1)\ge K$; with~(1). (3) $g(x_{k'})-r(x_{k'})=q(x_{k'})\prod_j(x_{k'}-x_j)$ with $q\in\Iw$ (the proof of~(1)), of valuation $\ge N$;
and the error of $D_iB_i(x_{k'})$ has valuation $\ge(K'-i)+i=K'$, as $v(x_{k'}-x_j)\ge1$.
\end{proof}

\noindent For general distinct $k_j$ the same argument gives~(2) with $K'\ge K+\max_{i}\max_{a\le i+1}\sum_{b\le i+1,\,b\ne a}v(x_{k_a}-x_{k_b})$ (cruder; the
$X^{m}$-gain is not used). Consecutive weights of one class are the economical choice, and the one WP-HL implements.

\begin{proposition}[The truncated algebra from the weights]\label{prop:u-lambdaprod}
\PROVED{} from Proposition~\ref{prop:u-truncbasic} and Lemma~\ref{lem:u-wtprod}. Let $a\not\equiv0$, let $k_1,\dots,k_N\ge3$ be pairwise distinct weights in
the class, $N\ge K+m-1$, and $\delta:=\max_a\sum_{b\ne a}v(x_{k_a}-x_{k_b})$ ($=N-1$ for consecutive weights, $N\le\ell$).
\begin{enumerate}
\item $\Tu(a)/\prod_j\Pk{k_j}\to B:=\prod_j\Tu_{k_j}(a)$ is injective; its image $A$ is the $\Iw$-subalgebra of $B$ generated by the tuples $(T_p)_j$, $p\ne\ell$, and
  $(U)_j$, with $X$ acting as $(x_{k_j})_j$; and $\ell^{\delta}B\subseteq A$.
\item $\Tu_{K,m}(a)\cong A/(\ell^{K},X^{m})A$.
\item \textbf{(Precision.)} If the $\Tu_{k_j}(a)$ are given modulo $\ell^{K'}$ with $K'\ge K+\delta$ (e.g.\ as the algebras of the Hecke matrices on a faithful
  $\Zl$-free module such as $e\,M_{k_j}(\SL_2(\Z);\Zl)$, Remark~\ref{rmk:u-hidaM}), the $\Iw$-subalgebra $\bar A'$ of $B/\ell^{K'}B$ generated by the images of the tuples $(T_p)_j$, $(U)_j$ \emph{and the weight tuple $(x_{k_j})_j$}, through which $X$
  acts (generated by the Hecke tuples alone over $\Z/\ell^{K'}$ it can be smaller; v0.3.0-pre.2, F1 m-U3),
  satisfies $\bar A'/\ell^{K'-\delta}\bar A'\cong A/\ell^{K'-\delta}A$, so it determines $\Tu_{K,m}(a)$.
\item \textbf{(The monogenic case.)} If a direct factor $\Tu(a)_{\mathfrak m}$ (a maximal ideal $\mathfrak m$) is generated over $\Iw$ by one element $t$ --- by
  Nakayama, iff $\Tu(a)_{\mathfrak m}/\mathfrak m_\Iw$ is generated by $\bar t$ over $\Fl$ --- then $\Tu(a)_{\mathfrak m}\cong\Iw[Y]/(F)$, $F$ the characteristic polynomial
  of multiplication by $t$ on the free $\Iw$-module $\Tu(a)_{\mathfrak m}$, $F(x_k,Y)$ is the characteristic polynomial of $t$ on $\Tu_k(a)_{\mathfrak m}$ (the
  $\mathfrak m$-factor of $\Tu_k(a)$) for $k\ge3$ in the class, and $\Tu_{K,m}(a)_{\mathfrak m}\cong\Iw_{K,m}[Y]/(F)$. The coefficients of $F$ modulo $(\ell^{K},X^{m})$ are
  determined by Theorem~\ref{thm:u-interp} from the characteristic polynomials at $N$ consecutive weights modulo $\ell^{K+m-1}$. The same holds for
  $\hcusp(a)_{\mathfrak m}$, every $a$.
\end{enumerate}
\end{proposition}

\begin{proof}
(1) $\Tu(a)$ is $\Iw$-free (Proposition~\ref{prop:u-truncbasic}(1)), so $\Tu(a)/\prod\Pk{k_j}\cong(\Iw/\prod\Pk{k_j})^{d(a)+1}$ and $\Tu(a)/\Pk{k_j}\cong\Tu_{k_j}(a)$
(Proposition~\ref{prop:u-truncbasic}(2)). The $\Pk{k_j}$ are pairwise non-associate primes of the factorial ring $\Iw$, so $\bigcap_j\Pk{k_j}\Iw=\prod_j\Pk{k_j}\Iw$, and
$\Iw/\prod\Pk{k_j}\to\prod_j\Iw/\Pk{k_j}=\Zl^{N}$ is injective; by Weierstrass division its source is $\Zl[X]_{<N}$, and the Lagrange polynomials
$\prod_{b\ne a}(X-x_b)/\prod_{b\ne a}(x_a-x_b)$ show that $\ell^{\delta}$ kills the cokernel. The image of $\Tu(a)$ is the image of the $\Iw$-algebra generated by
$T(p)$, $U$ (Definition~\ref{def:u-T}). (2) $\prod_j\Pk{k_j}\in(\ell^{K},X^{m})$ (Lemma~\ref{lem:u-wtprod}). (3) If $\ell^{K'}b\in A$ with $b\in B$, then
$\ell^{K'}b=\ell^{K'-\delta}(\ell^{\delta}b)\in\ell^{K'-\delta}A$; so $\ell^{K'}A\subseteq A\cap\ell^{K'}B\subseteq\ell^{K'-\delta}A$, and $\bar A'=A/(A\cap\ell^{K'}B)$.
(4) $\Tu(a)$ is finite over the complete local ring $\Iw$, hence the product of its localisations at maximal ideals; each factor is a direct summand of a
free $\Iw$-module, hence free. Nakayama gives $\Tu(a)_{\mathfrak m}=\Iw[t]$, and a monogenic algebra free of rank $r$ over $\Iw$ is $\Iw[Y]/(F)$ with $F$ the
characteristic polynomial of $t$ (Cayley--Hamilton and the rank count). Specialisation commutes with the characteristic polynomial, and
$\Tu(a)_{\mathfrak m}/\Pk{k}$ is the $\mathfrak m$-factor of $\Tu_k(a)$ by~(2) of Proposition~\ref{prop:u-truncbasic}. For $\hcusp(a)$: Theorem~\ref{thm:u-control}.
\end{proof}

\begin{lemma}[$G_a$ at every integer, and the twisted Bernoulli numbers]\label{lem:u-twistvalues}
\PROVED{} from \LIT{\cite[Thm.~5.11, p.~57; Thm.~7.10, p.~123]{washington1st} (both verified for [U], see Lemma~\ref{lem:u-emerton}(4); the
variable of Thm.~7.10 read on the page by the slot-19 referee F1: ``Let $\chi=\theta\psi$ be an \emph{even} Dirichlet character ($\theta$ of the first kind, $\psi$ of the
second kind), $\zeta_\psi=\psi(1+q_0)^{-1}$. Then $L_p(s,\chi)=f(\zeta_\psi(1+q_0)^{s}-1,\theta)$'', $f(T,\theta)\in\mathcal O_\theta[[T]]$ for $\theta\ne1$ (p.~122);
and Thm.~5.11, p.~57: $L_p(1-n,\chi)=-(1-\chi\omega^{-n}(p)p^{n-1})B_{n,\chi\omega^{-n}}/n$, $n\ge1$)}. Let $a\not\equiv0\pmod{\ell-1}$ be even. For \emph{every} integer
$k\ge1$,
\[
  G_a(x_k)=L_\ell(1-k,\omega^{a})=-\bigl(1-\omega^{a-k}(\ell)\,\ell^{k-1}\bigr)\frac{B_{k,\omega^{a-k}}}{k},
\]
which is $-(1-\ell^{k-1})B_k/k$ for $k\equiv a$ (Prediction~\ref{pr:u-U1}) and $-B_{k,\omega^{a-k}}/k$ for $k\not\equiv a\pmod{\ell-1}$ ($\omega^{a-k}$ has conductor $\ell$).
So the generalised Bernoulli numbers $B_{k,\omega^{j}}$, $j\not\equiv0$, are values of the same $G_{k+j}$ at points of other classes.
\end{lemma}

\begin{proof}
Thm.~7.10 is stated for \emph{even} characters; $\chi=\omega^{a}$ with $a$ even is even. With $\theta=\omega^{a}\ne1$, $\psi=1$ (so $\zeta_\psi=1$) and $1+q_0=1+\ell=u$, it
reads $L_\ell(s,\omega^{a})=f(u^{s}-1)$ for $s\in\Zl$, with $f\in\Zl[[T]]$ --- the $u^{s}$ convention (verified on the page, slot-19 F1; at v0.3.0-pre.1 marked
\VERIFYloc). Put $G(X):=f\bigl(u(1+X)^{-1}-1\bigr)\in\Iw$ ($u(1+X)^{-1}-1=(\ell-X)(1+X)^{-1}\in\mathfrak m_\Iw$). Then $G(x_k)=f(u^{1-k}-1)=L_\ell(1-k,\omega^{a})$ for
every integer $k$, and the value is Thm.~5.11 with $\chi=\omega^{a}$, $n=k$: $-(1-\chi\omega^{-k}(\ell)\ell^{k-1})B_{k,\chi\omega^{-k}}/k$. $G$ and $G_a$ agree at the
infinitely many $x_k$, $k\equiv a$, so $G=G_a$ (Weierstrass preparation, as in Lemma~\ref{lem:u-emerton}(4)).
\end{proof}

\noindent\emph{Two remarks} (v0.5.0-pre.5, slot~21; unnumbered, so that no number moves; no change of statement). (i) \emph{The Eisenstein measure} (slot-21 P2, \Espec{prop:lveismeas}).
The two factors $G_{a+1}(u(1+T)-1)$ and $G_{k-1-a}(u^{k-1}(1+T)^{-1}-1)$ of [E]'s Eisenstein congruence of measures are, each with the sign $-1$, the branch series of the Bernoulli
distribution $E_1$ and of $\iota_\ast E_{k-1}$, $\iota(x)=x^{-1}$ ($E_1=-x\,\zeta_\ell$ on $\Zl^{\times}$, $\zeta_\ell$ the pseudo-measure with $\int\omega^{b}(x)\langle x\rangle^{s}d\zeta_\ell=G_b(u^{s}-1)$;
0.5.0-pre.5 omitted the signs --- F4m, 0.5.0-pre.6): the measure $\mu^{\mathrm{Eis}}=-E_1\star\iota_\ast E_{k-1}$ on $\Zl^{\times}$ has branch series exactly their product (with a minus sign) in $\Iw\otimes\Q$ ([E]:
\PROVED{} given Washington Thm.~12.2, p.~239, read on the page by F4m; \COMP{} against the two group-ring routes of $G_b$). On the two pole classes one factor is $G_0$, and the polar term has residue
$-\rho_\infty\,G_k(x_k)$, $\rho_\infty=(1-\ell)\log_\ell(u)/\ell$ --- the value of this lemma at $x_k$ times $\zeta_\ell$'s residue (the constant \PROVED{} from the factorisation and $G_0=\rho_\infty/T+(\Iw)$, \Espec{prop:lveismeas}(4) in [E] 1.21.0-pre.4; \COMP{} as a calibration). (ii) \emph{Wild points} (slot-21 P4,
\Espec{rmk:lvwild}(5)). For $\psi_t$ the character of $(\Z/\ell^{m})^{\times}$, $m\ge2$, with $\psi_t(u)=\zeta^{t}$, $\zeta$ of order $\ell^{m-1}$, $\ell\nmid t$, and $G_b$ read as the
group-ring element with $1+X\leftrightarrow[u]$,
\[
  G_b(\zeta^{t}u^{r}-1)=-\frac{B_{r,\omega^{b-r}\psi_t}}{r}\qquad(b\not\equiv0\ \text{even},\ r\ge1),
\]
with no Euler factor ($\ell$ divides the conductor): \PROVED{} from Theorem~7.10 and Theorem~5.11 as quoted above (F4m, 0.5.0-pre.6; pre.5 called it ``a reading''): $G_b(X)=f(u(1+X)^{-1}-1,\omega^{b})$
by the proof of this lemma, and at $1+X=\zeta^{t}u^{r}$ the argument is $\zeta^{-t}u^{1-r}-1=\zeta_\psi(1+q_0)^{s}-1$ with $s=1-r$ and $\zeta_\psi=\psi_t(1+q_0)^{-1}=\zeta^{-t}$ ($\chi=\omega^{b}\psi_t$ is
even, $\psi_t$ of the second kind), so $G_b(\zeta^{t}u^{r}-1)=L_\ell(1-r,\omega^{b}\psi_t)$, and Theorem~5.11 gives $-(1-\omega^{b-r}\psi_t(\ell)\ell^{r-1})B_{r,\omega^{b-r}\psi_t}/r$ with
$\omega^{b-r}\psi_t(\ell)=0$. \COMP{} (slot-21 P4, both group-ring routes, $r\le10$, $t\in\{1,2\}$, at $(\ell,m)=(5,2)$, $(5,3)$, $(7,2)$, $(11,2)$, $(37,2)$, against the
definition of $B_{r,\chi}$; and independently by F4m at $(5,2)$, $(7,2)$: $48/48$, the convention $\zeta^{-t}$ $0/48$); the convention $\zeta^{-t}$ fails at every non-zero value. On the pole class $b=0$, $v_\pi(B_{r,\omega^{-r}\psi_t}/r)=-1$, $\pi=1-\zeta$ (\COMP).

\noindent\emph{A third remark} (v0.5.0-pre.7, slot~21 P5, \Espec{rmk:lviwasawa}(6); unnumbered; no change of statement; refereed by F4 at v0.5.0-pre.8). \emph{The Iwasawa invariants of $G_b$ at the pairs used.}
$\lambda(G_b)=0$ iff $\ell\nmid B_b/b$ (\PROVED{} \emph{given} $\mu(G_b)=0$: the constant term is $\equiv G_b(x_b)=-(1-\ell^{b-1})B_b/b\pmod\ell$, by this lemma; a unit constant term gives $(\mu,\lambda)=(0,0)$
unconditionally, a divisible one gives $\lambda\ge1$ only once $\mu=0$ is known --- F4, pre.8); $\lambda(G_b)=1$ at the irregular pairs used in [U] and [E] ---
$(37,32)$, $(59,44)$, $(67,58)$, $(103,24)$, $(131,22)$, $(691,12)$, $(691,200)$ --- \LIT{\cite{hho}, abstract item~(3), p.~1: ``the Iwasawa invariants satisfy $\lambda_p=\nu_p=i_p$'', read on the text layer
for slot-21 P5 and again by F4, who also read \S7 of the manuscript ($\lambda_p$, $\nu_p$ are the Iwasawa invariants of the $p$-class groups of $\Q(\zeta_{p^{n}})$; the test is four incongruences per irregular index)}; the step
$\lambda_p=\sum_{b\text{ even}}\lambda(G_b)$ is \STD{Iwasawa's theorem for the minus part through the analytic class number formula, with the plus part trivial by Kummer--Vandiver, verified in \cite{hho} item~(2); recalled, locator unverified},
and with $\lambda(G_b)\ge1$ at each irregular index (Ferrero--Washington again) this forces $\lambda(G_b)=1$; $\mu(G_b)=0$ for every even $b$ (\STD{Ferrero--Washington, recalled;
not read for [U]}). Both re-derived (\COMP, slot-21 P5, \texttt{HeckeLValueIwasawaInvariantsValidation} arm I0; and independently by the referee F4 with own code, tangent-number Bernoulli numbers and divided
differences at $s\le5$, \file{docs/implementation/slot21/rc/f4p5/f4p5\_kl.py}: the seven pairs $(0,1)$, every other even class of the $16$ primes $(0,0)$, $\mu=0$ throughout): from exact Bernoulli numbers alone, by the divided differences of this lemma's values
$G_b(x_r)$ at $r=b+s(\ell-1)$, $s=0,\dots,3$ ($v(x_r-x_{r'})=1$; the reading of Lemma~\ref{lem:u-lambdaone}'s proof), and by the Kubota--Leopoldt group ring at level~$2$; the two routes agree at every even
class of $\ell\in\{11,\dots,47,59,67,103,131,691\}$ ($344/344$ at $691$), $\mu=0$ everywhere, and $\lambda=1$ exactly at the seven pairs listed. Proposition~\ref{prop:u-eistrunc}'s four classes are among them;
nothing there uses this remark.

\begin{lemma}[The $\lambda$-criterion]\label{lem:u-lambdaone}
\PROVED{} (elementary). Let $a\not\equiv0$ be an irregular class ($v(B_k/k)\ge1$ for one, equivalently every, $k\equiv a$; Corollary~\ref{cor:u-irregular}), and
$k,k'\equiv a$ with $\ell\nmid k-k'$. Then $\mu(G_a)=0$ and $\lambda(G_a)=1$ iff $v\bigl(G_a(x_k)-G_a(x_{k'})\bigr)=1$. In that case $G_a=\varepsilon\cdot(X-\alpha_a)$,
$\varepsilon\in\Iw^{\times}$, $\alpha_a\in\ell\Zl$, $v(G_a(x_{k''}))=v(x_{k''}-\alpha_a)$ for every integer $k''$ --- which is $L_{k''}$ for $k''\ge2$ in the class, by
(U1-$k$) --- and, for the centred truncation, $k''\ge2$ in the class:
\[
  \Cg_{\Eis;K,m}(a)\cong\Iw_{K,m}/(X-\alpha_a)\cong\Z/\ell^{\min(K,\,m\,v(\alpha_a))},\qquad
  \Cg_{\Eis;K,m}(a;x_{k''})\cong\Z/\ell^{\min(K,\,m\,L_{k''})}.
\]
\end{lemma}

\begin{proof}
Write $G_a=\sum g_iX^{i}$. $g_0\equiv G_a(x_k)\equiv0\pmod\ell$ ($v(x_k)\ge1$). $G_a(x_k)-G_a(x_{k'})=(x_k-x_{k'})\sum_{i\ge1}g_ih_{i-1}(x_k,x_{k'})$ with
$v(x_k-x_{k'})=1$ and $v(h_{i-1})\ge i-1$, so the valuation is $1$ iff $g_1\in\Zl^{\times}$, which with $g_0\in\ell\Zl$ is $\mu=0$, $\lambda=1$. Weierstrass
preparation gives $\varepsilon$, $\alpha_a$. Then $\Iw_{K,m}/(X-\alpha_a)=(\Z/\ell^{K})/(\alpha_a^{m})$ and $\Iw_{K,m}(x)/(X-\alpha_a)=(\Z/\ell^{K})/((\alpha_a-x)^{m})$, with
$v(\alpha_a-x_{k''})=v(G_a(x_{k''}))$, which is $L_{k''}$ for $k''\ge2$ in the class (Prediction~\ref{pr:u-U1}, (U1-$k$)); Proposition~\ref{prop:u-truncbasic}(1).
\end{proof}

\begin{proposition}[The Eisenstein truncations at the slot-19 targets]\label{prop:u-eistrunc}
\PROVED{} from Lemma~\ref{lem:u-lambdaone}, Prediction~\ref{pr:u-U1} and the \COMP{} values below (exact Bernoulli numbers, in-house,
\file{docs/implementation/slot19/wp2/lambda\_check.py}, output \file{lambda\_check.out} beside it, $<1$\,s). At $(\ell,a)\in\{(37,32),(103,24),(691,12),(691,200)\}$,
with $k=a$, $k'=a+\ell-1$: $v(G_a(x_k))=v(G_a(x_{k'}))=1$ and $v(G_a(x_k)-G_a(x_{k'}))=1$. Hence at all four classes $\mu(G_a)=0$, $\lambda(G_a)=1$,
$v(\alpha_a)=1$, and for all $K,m\ge1$
\[
  \Cg_{\Eis;K,m}(a)\cong\Z/\ell^{\min(K,m)}.
\]
The root $\alpha_a$ modulo $\ell^{2}$ predicts the first weight of the class with $L_k\ge2$ ($x_k\equiv\alpha_a\pmod{\ell^{2}}$): $k=284$ at $37$ and $k=228$ at
$103$ --- the two depth-$2$ rungs of \S\ref{sec:u-test} and Remark~\ref{rmk:u-zero}, found there independently --- and $k=342252$ (class $12$) and $k=342440$
(class $200$) at $691$; the first is \S\ref{sec:u-test}'s ``$n=342250$''. \COMP{} (slot-19 referee F1, an independent route: $B_k\bmod\ell^{3}$ from the power sum
$S_k(\ell^{2})$, \file{docs/implementation/slot19/wp3/f1\_lambda.py}, output \file{f1\_lambda.out} beside it, $\approx25$\,s): the four readings reproduced; a third
weight $k''=a+2(\ell-1)$ linear modulo $\ell^{2}$, as $\lambda=1$ demands; the two rungs at $691$ \emph{computed}, not only predicted:
$v(B_k/k)=2$ at $k=342252$ and $k=342440$, and $1$ at $k\mp690$ (where \file{lambda\_check.out} reads ``n/a''); and every other weight scanned ($k<1500$ at $37$ and $103$,
$k<12000$ at $691$) reads $L_k=1$ (v0.3.0-pre.2).
\emph{Depth $3$} (v0.3.0-pre.3; slot-19 R19-2 and its referee F4, major M1). The root $\alpha_a$ modulo $\ell^{3}$ gives the first weight of the class with $L_k\ge3$:
\[
  k=37580\ (\ell=37),\qquad k=914250\ (\ell=103),\qquad k=49928412\ \text{(class $12$)},\quad k=159113510\ \text{(class $200$)}\ (\ell=691).
\]
\PROVED{} from Lemma~\ref{lem:u-lambdaone}, Theorem~\ref{thm:u-interp}(1), Lemma~\ref{lem:u-twistvalues}, Prediction~\ref{pr:u-U1} and \COMP{} values: five exact
Bernoulli numbers per class, $B_k$ at $k=a+j(\ell-1)$, $j=0,\dots,4$ (all $k\le3650$; slot-19 R19-2, \file{docs/implementation/slot19/r192/lambda\_trunc.py}; reproduced by
F4 with its own tangent-number Bernoulli numbers and Newton interpolation). At $37$ and $103$ Voronoi's congruence (\COMP, \file{r192/rung\_check.py}, and F4) reads
$v(B_k/k)=3$ exactly at the rung and $1$ at $k\mp(\ell-1)$; at $691$ the Voronoi check ($691^{4}$ terms) is an optional cross-check, not what the statement rests on,
which there gives $L_k\ge3$.
\end{proposition}

\begin{proof}
$v(\alpha_a)=v(g_0)$, and modulo $\ell^{2}$, $G_a(x)\equiv g_0+g_1x$ for $x\in\ell\Zl$; the two values determine $g_1\bmod\ell$ and $g_0\bmod\ell^{2}$, and the script
finds $g_0\not\equiv0\pmod{\ell^{2}}$ at all four. Lemma~\ref{lem:u-lambdaone}. That $\lambda=1$ at these four is also the literature's (\LIT{\cite{hho}: the
cyclotomic $\lambda$-invariants of all irregular pairs below $2^{31}$, \VERIFYloc}); nothing here uses it.
\emph{Depth $3$.} $G_a=\varepsilon(X-\alpha_a)$ with $\varepsilon\in\Iw^{\times}$, so $v(G_a(x_k))=v(x_k-\alpha_a)$ exactly ($v(\varepsilon(x_k))=0$), and $L_k=v(G_a(x_k))$ for
$k\ge2$ in the class ((U1-$k$), Lemma~\ref{lem:u-twistvalues}: $1-\ell^{k-1}$ is a unit); hence $L_k\ge3$ iff $x_k\equiv\alpha_a\pmod{\ell^{3}}$. As $\alpha\in\ell\Zl$ gives
$\alpha^{3}\in\ell^{3}\Zl$, $\alpha_a\bmod\ell^{3}$ is the root in $\ell\Z/\ell^{3}$ of $G_a\bmod(\ell^{3},X^{3})$, unique by Hensel ($g_1$ a unit); and
$G_a\bmod(\ell^{3},X^{3})$ is the interpolant through the $N=5=K+m-1$ values $G_a(x_k)=-(1-\ell^{k-1})B_k/k$ (Theorem~\ref{thm:u-interp}(1), $K=m=3$; the values exact).
As $u$ has order $\ell^{2}$ modulo $\ell^{3}$, $k\mapsto x_k\bmod\ell^{3}$ is a bijection from the class modulo $(\ell-1)\ell^{2}$ onto $\ell\Z/\ell^{3}$, so exactly one
$k\in[a,a+(\ell-1)\ell^{2})$ in the class has $L_k\ge3$; the script finds the four values displayed.
\end{proof}

\noindent\emph{(U12)'s checkable form, read} (\COMP, v0.3.0-pre.3; slot-19 R19-2, \file{docs/implementation/slot19/R19-2\_report.md}~\S\S3--5, \file{r192/}, predictions
committed first; refereed by F4). At the four classes, $K=3$, $m=1,2,3$, centred at $x_{k_0}$ and at the depth-$2$ weight ($24$ truncations), from the forms side
($q$-expansions, ordinary part): the $\mathfrak m_\Eis$-factor is $\Iw_{3,m}[Y]/(Y(Y+F_1))$ with $Y=T_2-\ee_2$ (Definition~\ref{def:u-eischar}) and
$F_1=\ee_2-a_2(f_X)$ (the sign fixed by F4, m10); $(F_1)=(\tilde G_a)=(X-\alpha_a)$ and the Hecke-side Eisenstein ideal $J=(\tilde G_a)$; the cokernel of (U12)'s exact
sequence tensored with $\Iw_{3,m}$ has order $\ell^{\min(3,m\,L_c)}$, $L_c$ the depth at the centre, at $24/24$; the truncation is $\Iw$-free of rank~$2$, generated by
$Y$. The $K+m-1$ digits of Theorem~\ref{thm:u-interp}(2) are sufficient and sharp (with $K+m-2$ digits the truncation changes in $23$ of the $24$ cases); (U11)'s
checkable form holds at every further weight; the forms-side characteristic polynomials reproduce all $753$ comparable rows of the tower catalogue
(\file{docs/matrix-engine/results/hida\_towers/}). A one-line extra gate (F4, m10): $F_1(x_k)/G_{32}(x_k)\equiv2\pmod{37}$ is constant along the eight tower weights
$k=32,68,\dots,284$ of $(37,32)$ (v0.3.0-pre.4, slot-19 WP-HL, which gates it: F4's script formed $a_2(f_X)-\ee_2=-F_1$, so its constant $35\equiv-2$ is the ratio in the
opposite sign; pre.3 paired $35$ with this $F_1$ and counted six weights).

\noindent\emph{The rungs from the cyclotomic layer, Bernoulli-free} (\COMP, v0.3.0-pre.4; slot-19 R19-7, \file{docs/implementation/slot19/R19-7\_report.md}~\S\S1--4, predictions
committed first). The Coates--Wiles measure of the $\omega^{a}$-part of the cyclotomic units of $\Q_\ell(\zeta_{\ell^{m}})$, built from one Iwasawa logarithm
$\log_\ell(1-\zeta_{\ell^{m}})$, has moments $\equiv(1-\ell^{k-1})B_k/k=-G_a(x_k)\pmod{\ell^{m}}$ (the congruence proved in R19-7~\S\S1.1--1.2; the equality by Lemma~\ref{lem:u-twistvalues}) and so reads
$G_a\bmod((1+X)^{\ell^{m-1}}-1,\ell^{m})$ without any Bernoulli number. Its unique zero modulo $\ell^{m}$ in the class is the depth-$m$ rung: $284$, $228$, $342252$, $342440$ at $m=2$
(the depth-$2$ rungs of Proposition~\ref{prop:u-eistrunc}), $37580$ and $914250$ at $m=3$ (two of the depth-$3$ rungs proved above), and at $m=4$, $(37,32)$, the depth-$4$ rung
$k=1072544$; $\alpha_{32}$ read back is $925$, $28305$, $281570$ modulo $37^{2},37^{3},37^{4}$ (those of R19-2's $\alpha$; the depth-$4$ zero and $\alpha\bmod37^{4}$ were predicted
from it before the run), $\alpha_{24}\equiv2266\pmod{103^{2}}$, and at $691$ $143037$, $272945$ modulo $691^{2}$. The depth-$3$ rungs at $691$ ($49928412$, $159113510$) need the
level-$3$ layer at $691$ (a field of degree $\approx3.3\cdot10^{8}$; R19-7~\S8, a plan for slot~20). \emph{Computed} (v0.4.0-pre.1; \COMP, the slot-19 run NF691 of \code{NumberFieldValidation}
with \code{NUMBERFIELD\_LAYER\_CELLS="691:3:12/200:3:11"}, commit \texttt{adc71e8d}, \file{docs/implementation/slot19/nf691\_run.md}, predictions committed first; $29.7$\,GiB, $1$\,h\,$15$\,min at
$8$ threads): the unique zero modulo $691^{3}$ is $k=49928412$ (class $12$) and $k=159113510$ (class $200$), with $\alpha_{12}\equiv139090008$, $\alpha_{200}\equiv279599330\pmod{691^{3}}$ ---
the depth-$3$ rungs proved above, now read on a second, Bernoulli-free route; every row modulo $691^{2}$ equals the level-$2$ one. Nothing above rests on these readings.

\begin{prediction}[(U14): the deeper rungs and $\alpha_a$ modulo $\ell^{m}$ at the slot-20 layer cells]\label{pr:u-rungs}
(New at v0.4.0-pre.1, hecke-extensions slot~20, WP-2; stated before slot-20 WP-RC20's layer windows run (rule~20).) Let $(\ell,a,m)\in\{(37,32,5),\,(103,24,4),\,(59,44,4)\}$.
\begin{enumerate}
\item \PROVED{} from Lemma~\ref{lem:u-lambdaone}, Lemma~\ref{lem:u-twistvalues}, Prediction~\ref{pr:u-U1} and the \COMP{} values below. At $(59,44)$, $\mu(G_{44})=0$ and $\lambda(G_{44})=1$:
  $v(G_{44}(x_{44}))=v(G_{44}(x_{102}))=v(G_{44}(x_{44})-G_{44}(x_{102}))=1$ (\COMP, exact Bernoulli numbers, \file{docs/implementation/slot20/wp2/lambda\_criterion.out}; at $(37,32)$, $(103,24)$ this
  is Proposition~\ref{prop:u-eistrunc}). Hence at all three classes $G_a=\varepsilon(X-\alpha_a)$, and exactly one weight $k\equiv a\pmod{\ell-1}$ in $[a,\,a+(\ell-1)\ell^{m-1})$ has $L_k\ge m$:
  the one with $x_k\equiv\alpha_a\pmod{\ell^{m}}$ (the depth-$3$ argument of Proposition~\ref{prop:u-eistrunc}, with $m$ in place of $3$).
\item \COMP{} (\file{docs/implementation/slot20/wp2/rung\_predict.py}, output \file{rung\_predict.out} and \file{rung\_predict\_stability.out} beside it; Python \code{Fraction}, exact Bernoulli
  numbers by the tangent numbers, in-house, no CAS, $\approx2$\,min). The values:
  \begin{center}\small
  \begin{tabular}{@{}llll@{}}
  \toprule
  $(\ell,a)$ & $m$ & $\alpha_a\bmod\ell^{m}$ & the depth-$m$ rung $k$ ($n=k-2$)\\
  \midrule
  $(37,32)$ & $5$ & $13400697$ & $55777784$\\
  $(103,24)$ & $4$ & $3736634$ & $60430740$\\
  $(59,44)$ & $2$, $3$, $4$ & $1711$, $147913$, $11649137$ & $914$, $86464$, $8162384$\\
  \bottomrule
  \end{tabular}
  \end{center}
  \emph{How.} $P$ is the Newton interpolant of $G_a$ through $N$ points $x_{k_j}$, $k_j=a+j(\ell-1)$, $j<N$, from the exact values $G_a(x_{k_j})=-(1-\ell^{k_j-1})B_{k_j}/k_j$. For $x\in\ell\Z_\ell$,
  $G_a(x)-P(x)=\prod_j(x-x_{k_j})\cdot G_a[x_{k_0},\dots,x_{k_{N-1}},x]$, and the divided differences of $G_a\in\Iw$ (a power series with $\Zl$-coefficients, part~1; v0.4.0-pre.2, F1 m-U1) are $\Zl$-integral (complete
  homogeneous symmetric polynomials in the nodes), so $v(G_a(x)-P(x))\ge N$ (\PROVED, elementary). With $N=m+2$ the root of $P$ in $\ell\Z/\ell^{m}$, found digit by digit, is
  $\alpha_a\bmod\ell^{m}$; $N=m+4$ gives the same values. The script reproduces every value already in this document or in [E]: $925$, $28305$, $281570$ and the rungs $284$, $37580$, $1072544$ at $37$;
  $2266$, $458453$ and $228$, $914250$ at $103$; $143037$, $272945$, $139090008$, $279599330$ and $342252$, $342440$, $49928412$, $159113510$ at $691$; and at $59$ the depth-$2$ rung $914$ (slot-19
  R19-3) and the depth-$3$ rung $86464$ with $\alpha_{44}\equiv147913\pmod{59^{3}}$ (slot-19 WP-NF2's cell arm, Bernoulli-free). The new values are consistent with them:
  $13400697\equiv281570\pmod{37^{4}}$, $3736634\equiv458453\pmod{103^{3}}$, $11649137\equiv147913\pmod{59^{3}}$.
\item \emph{Checkable form} (slot-20 WP-RC20, the layer windows: \code{NumberFieldValidation} with \code{NUMBERFIELD\_LAYER\_CELLS}, its \code{EFFECTIVE} line read first; cost from
  \code{nf::cw\_predict} before launch, at $37$, $m=5$: $7.1$\,GB, $\approx0.45$\,h at $8$ threads; at $103$, $m=4$: $11.5$\,GB, $\approx0.47$\,h; at $59$, $m=4$: predicted by WP-RC20 before launch).
  The level-$m$ Coates--Wiles measure of the $\omega^{a}$-part (the paragraph after Proposition~\ref{prop:u-eistrunc}) has, modulo $\ell^{m}$, a \emph{unique} zero in the class, equal
  to the rung of~(2), and its $\alpha$ read back equals $\alpha_a\bmod\ell^{m}$ of~(2); every row modulo $\ell^{m-1}$ equals the level-$(m-1)$ row. A mismatch is a finding against (1) or
  against the layer code, reported before any re-run. At $691$, $m\ge4$ is out of reach (the layer's $n_m\approx2.3\cdot10^{11}$; slot-19 WP-NF2) and is recorded with its cost only.
\item \emph{Outcome} (v0.4.0-pre.3; \COMP, slot-20 WP-RC20, \file{docs/implementation/slot20/RC20\_levelm\_report.md}; the shared slot-20 binary, \code{ENGINE\_THREADS=8}; each window
  \code{EFFECTIVE NumberFieldValidation arm=LAYER\_CELLS cells=1}, \code{NUMBERFIELD\_GATE checks=637 failing=0}). \textbf{Every prediction of~(2) and~(3) is met}: the level-$m$ measure has
  a unique zero modulo $\ell^{m}$ in the class, and it and $\alpha$ read back equal~(2) --- $55777784$, $\alpha_{32}\equiv13400697\pmod{37^{5}}$; $60430740$,
  $\alpha_{24}\equiv3736634\pmod{103^{4}}$; $8162384$, $\alpha_{44}\equiv11649137\pmod{59^{4}}$. The zero towers $[1874161,50653,1369,37,1]$, $[1092727,10609,103,1]$,
  $[205379,3481,59,1]$ are the depth-$m$ rungs of~(2) level by level, and every row agrees with the Bernoulli route and the control. Cost: $6.87$, $11.20$, $1.07$\,GB, $27$, $34$, $2$\,min
  (\code{nf::cw\_predict}'s figures within $5\,\%$ of the \code{time -v} peaks ($+4.4\,\%$, $+4.6\,\%$ at $37$, $103$), within $1\,\%$ of the high-water marks; F1$'$ m-U1). The prediction stays \PROVED; the readings are a second, Bernoulli-free route to it, as at $691$ above. No finding.
\end{enumerate}
\end{prediction}

\begin{remark}[What the truncations rest on]\label{rmk:u-lambdabasis}
\textbf{\LIT{} (Part~1's, unchanged).} The $\Iw$-freeness and control of $\hcusp(a)$: Hida, Theorem~\ref{thm:u-control}. The $\Iw$-freeness of $\Tu(a)$ and
$\Cg_\Eis(a)\cong\Iw/(G_a)$: Emerton's Theorem~5 through Lemma~\ref{lem:u-emerton} (Prediction~\ref{pr:u-U1}). The values of $G_a$: Washington, Lemma~\ref{lem:u-twistvalues}.
\textbf{Not used.} Wiles's $\Iw$-adic forms and the $\Iw$-adic Galois representations \cite{wiles_ordinary} (they would give the same algebra by duality); any
Gorenstein property of $\Tu(a)$ or $\hcusp(a)$; any main conjecture. \textbf{\PROVED{} here.} Proposition~\ref{prop:u-truncbasic}, Lemma~\ref{lem:u-wtprod},
Theorem~\ref{thm:u-interp}, Proposition~\ref{prop:u-lambdaprod}, Lemma~\ref{lem:u-lambdaone} (elementary, from the above). \textbf{\COMP.} The four $\lambda$-readings
of Proposition~\ref{prop:u-eistrunc} and the five exact Bernoulli numbers per class behind its depth-$3$ rungs; the one characteristic polynomial per class behind
(U13) at $37$ and $103$ (v0.3.0-pre.3); the tower data of Part~1 (\file{docs/matrix-engine/results/hida\_towers/}), which give $\chi^{\mathrm{ord}}_{n,p}$ modulo
$\ell^{m}$ and hence $c_{k,p}$ modulo $\ell^{m}$ (Lemma~\ref{lem:u-levelone}(3); the monic square root is unique for $\ell$ odd). \textbf{What a tower row is not.} A
row of level $m$ is one point of $\Iw/(\ell^{m},\Pk{k})$ (Lemma~\ref{lem:u-tower}): it never sees $X^{2}$; the $X$-direction comes only from \emph{several}
weights, through Theorem~\ref{thm:u-interp}, at the cost of $m-1$ digits.
\end{remark}

\begin{prediction}[(U11)--(U13): the truncated algebra at the slot-19 targets]\label{pr:u-lambda}
For slot-19 R19-2 (research) and WP-HL (code), stated before either runs (rule~20).
\begin{enumerate}
\item[(U11)] \PROVED{} (Proposition~\ref{prop:u-lambdaprod}(4), Theorem~\ref{thm:u-etale}(1)) \textbf{Étale classes.} In every class of Theorems~\ref{thm:u-ell13},
  \ref{thm:u-691classes} --- there $\bar c_{k,p}$ is squarefree for some $p\ne\ell$ --- $\hcusp_{K,m}(a)\cong\prod_iW_K(\F_{\ell^{\delta_i}})[X]/(X^{m})$ and
  $\hcusp_{K,m}(a)=\Iw_{K,m}[Y]/(F_p)$, $F_p$ the $\Iw$-adic characteristic polynomial of $T(p)$. \emph{Checkable form} (a failure is a defect of the code or of a tower datum): $F_p\bmod(\ell^{K},X^{m})$
  interpolated from $N=K+m-1$ consecutive weights modulo $\ell^{K+m-1}$ predicts $c_{k',p}\bmod\ell^{K+m-1}$ at every further weight $k'$ of the class
  (Theorem~\ref{thm:u-interp}(3)), the $\mathfrak m$-factors being matched across weights by their residual system (Lemma~\ref{lem:u-tower} at $m=1$). As
  Theorem~\ref{thm:u-interp}(3) holds for \emph{every} $g\in\Iw$, the check tests the existence of the $\Iw$-adic $F_p$ (Hida's control) and the implementation, not the
  étale structure (v0.3.0-pre.2, F1 m-U6).
\item[(U12)] \PROVED{} (Proposition~\ref{prop:u-eistrunc}) \textbf{The Eisenstein truncations.} At $(37,32)$, $(103,24)$, $(691,12)$, $(691,200)$:
  $\Cg_{\Eis;K,m}(a)\cong\Z/\ell^{\min(K,m)}$ and, centred at a weight, $\Z/\ell^{\min(K,mL_k)}$. \emph{Checkable form}: on the Hecke side,
  $0\to\Tu(a)/\textstyle\prod_j\Pk{k_j}\to(\hcusp(a)\times\Iw)/\prod_j\Pk{k_j}\to\Cg_\Eis(a)/\prod_j\Pk{k_j}\to0$ is exact ($\prod_j\Pk{k_j}$ is a non-zero-divisor on
  $\Iw/(G_a)$, as $G_a(x_{k_j})\ne0$), both terms on the left embed into $\prod_j(\hcusp_{k_j}\times\Zl)$ as in Proposition~\ref{prop:u-lambdaprod}(1), and the cokernel
  tensored with $\Iw_{K,m}$ must have order $\ell^{\min(K,m)}$; on the Bernoulli side, $\Iw_{K,m}/(\tilde G)$ with $\tilde G$ interpolated (Theorem~\ref{thm:u-interp}) from
  the values of Lemma~\ref{lem:u-twistvalues} must agree. The cells $(282,37)$ ($k=284$, $L_k=2$) and $(226,103)$ ($k=228$, $L_k=2$) are the first weights of
  their classes at which the truncation centred at the weight reads $\Z/\ell^{\min(K,2m)}$ rather than $\Z/\ell^{\min(K,m)}$.
\item[(U13)] \PROVED{} (at all four classes since v0.3.0-pre.3; at pre.1, pre.2 \CONJ{} at $37$ and $103$) \textbf{The Eisenstein-local algebra.} At the four classes of (U12), the cuspidal component through $\mathfrak m_\Eis$ is $\Iw$ (as \PROVED{} at $691$,
  Theorem~\ref{thm:u-691classes}), so $\Tu(a)_{\mathfrak m_\Eis}=\Iw\times_{\Iw/(G_a)}\Iw\cong\Iw[Y]/(Y(Y-G_a))$ and
  \[
    \Tu_{K,m}(a)_{\mathfrak m_\Eis}\cong\Iw_{K,m}[Y]/\bigl(Y(Y-X+\alpha_a)\bigr),
  \]
  free of rank $2$ over $\Iw_{K,m}$, local, with $\Tu_{K,m}(a)_{\mathfrak m_\Eis}/\ell\cong\Fl[X,Y]/(X^{m},Y(Y-X))$. \PROVED{} at $\ell=691$, classes $12$ and $200$ (the fibre product
  is Theorem~\ref{thm:u-691classes} with $J=(G_a)$; the presentation: $(1,1)$ and $Y:=(0,G_a)$ are a $\Iw$-basis with $Y^{2}=G_aY$; then
  $G_a=\varepsilon(X-\alpha_a)$, $Y\mapsto\varepsilon Y$). \PROVED{} at $(37,32)$ and $(103,24)$ by the same argument (v0.3.0-pre.3; slot-19 R19-2~\S3, the slot-19
  referee F4, m8), from Theorem~\ref{thm:u-etale}(1),(3), Lemma~\ref{lem:u-fibre}, (U1-$\Iw$) (\LIT: Hida's control, Emerton's Theorem~5) and a \COMP{} input at one
  weight: $\mathrm{charpoly}(T_2\mid S_{32})=X^{2}-39960X-2235350016$ has the roots $14$, $23$ modulo $37$, with $23\equiv1+2^{31}$ simple, and
  $\mathrm{charpoly}(T_2\mid S_{24})=X^{2}-1080X-20468736$ has the roots $70$, $83$ modulo $103$, with $83\equiv1+2^{23}$ simple. So $\bar c_{k,2}$ (a divisor of these) is
  squarefree with the Eisenstein value a simple root, and $v(B_k/k)=1$ at $k=32$, resp.\ $24$ (Proposition~\ref{prop:u-eistrunc}): the cuspidal component through
  $\mathfrak m_\Eis$ is étale of residue degree~$1$, i.e.\ $\Iw$ ($d_{\mathrm{Eis}}=1$ at every weight of the class, as R19-2 reads at all $13$, resp.\ $6$, tower weights),
  and the fibre product follows as at $691$. R19-2 also reads the presentation itself (the paragraph after Proposition~\ref{prop:u-eistrunc}).
\end{enumerate}
\emph{Code} (v0.3.0-pre.4, \COMP): slot-19 WP-HL's \code{HeckeLambdaValidation} gates the checkable forms of (U11)--(U13) at all $12$ truncations of $(37,32)$ and $(103,24)$, and with
\code{HECKE\_LAMBDA\_ALL=1} also at $(691,12)$, $(691,200)$; all held (the code map below).
\end{prediction}

\paragraph{Code map of part 3.} \textbf{Built} (v0.3.0-pre.4; slot-19 WP-HL, commit \texttt{e3a4be35}, after F1's ACCEPT of this part). \file{HeckeExt/HeckeHidaLambdaForms}:
the input arm~(Q) on $q$-expansions --- by the owner's ruling of 2026-10-03 the primary input, the deviation from~(i) below --- the Miller basis modulo $\ell^{W}$, exact products by a
multimodular NTT over $62$-bit primes with Garner into $\Z/\ell^{W}$, $T_p$ on $q$-expansions, the ordinary projector $T_\ell^{N}$; \file{HeckeExt/HeckeHidaLambda}: the $\Iw$ stage
--- a certified Fitting split, the monogenic route through Theorem~\ref{thm:u-interp} with its precision bound \emph{enforced} (refused below $K+m-1$ digits or weights, and on
non-consecutive weights), the general route of Proposition~\ref{prop:u-lambdaprod}(1) at $K+N-1$ (refused below), Bernoulli values by tangent numbers ($=$ WP-NF's exact ones), $\alpha$,
the rungs, the Kubota--Leopoldt value; record \code{hecke\_lambda/v1} (an FNV digest per line, not pinned; no golden line). The shipped tower route
(\file{HeckeExt/HeckeHidaTower}) is a \emph{cross-check} only, at $37$ and $103$ (arm~(T): $(30,37)$, $(66,37)$, $(22,103)$, $m\le2$). \code{HeckeLambdaValidation} (\code{standard},
$\approx7$\,s, $\approx70$\,MB, $138$ checks; the planned name \code{HeckeHidaLambdaValidation} is superseded): R19-2's readings at $(37,32)$, $(103,24)$ reproduced field by field
(the oracle is R19-2's data, never this code's output), (U11)--(U13) at all $12$ truncations, the Kubota--Leopoldt value $47910634$, the rungs $284$\,/\,$37580$, $228$\,/\,$914250$, the
F4 m10 ratio, the tower catalogue at $\ell\le13$ ($744/744$ rows), negative controls (the precision refusals, perturbed inputs caught); env \code{HECKE\_LAMBDA\_ALL=1} (read-only)
adds $(691,12)$, $(691,200)$ ($224$ checks, $\approx400$\,s, $0.15$\,GB). \COMP{} where a value is read; no statement of this part rests on the code. \emph{The plan, as written at
v0.3.0-pre.1}: slot-19 WP-HL (upper core; Opus; only on F1's ACCEPT of this part), an
arm of \file{HeckeExt/HeckeHidaTower} that (i) takes $N\ge K+m-1$ consecutive weights $k_0+j(\ell-1)$ of one class, computed modulo $\ell^{K+m-1}$ on the shipped
tower route (\HA{rmk:hidacheckable2}); (ii) interpolates the block characteristic polynomials (Proposition~\ref{prop:u-lambdaprod}(4)) and, where a block is not
monogenic, the subalgebra of Proposition~\ref{prop:u-lambdaprod}(1) at precision $K+\delta$; (iii) gates (U11)'s checkable form at the étale classes and (U12) against
in-house Bernoulli numbers (WP-NF's \code{BigInt}, or modulo $\ell^{K'}$); \code{HeckeHidaLambdaValidation} (\code{standard}, $\le8$\,GB; heavy cells by an env
interlock only; $(700,691)$ at $m=2$, $33.5$\,GiB, excluded). Research prototype: slot-19 R19-2 (from \file{results/hida\_towers/} and new cells only by prediction). \emph{Since v0.3.0-pre.3} (R19-2~\S6, refereed by F4):
the prototype ran on a $q$-expansion route; the shipped tower route cannot feed the $\Iw$-adic arm at $\ell=691$ ($(700,691)$ at $m=2$ already needs $33.5$\,GiB), so
R19-2 recommends a $q$-expansion input arm --- a deviation from (i) above that is the WP-HL planner's and the owner's decision. On either route the generator choice
and the component split are made \emph{on the ordinary part}: on the full $S_k$ at $k=140$, $176$, $284$ the Eisenstein eigenvalue $1+2^{k-1}$ of $T_2$ has
multiplicity~$2$ modulo $37$, with a non-ordinary companion, so on $H^{1}$ itself $T_2$ does not separate the $\mathfrak m_\Eis$-block there (F4, m9).
No golden line, digest or record key is touched by this part.

%----------------------------------------------------------------------
\section{The \texorpdfstring{$\Iw$}{Lambda}-adic modular symbols, truncated (part~4)}\label{sec:u-msym}
%----------------------------------------------------------------------
% NEW at v0.5.0-pre.1 (hecke-extensions slot 21, P8, text first; spec: the owner's padic-Hida note (private),
% roadmap target (ii), P8; plan docs/implementation/slot21/00-slot21-plan.md; implementation spec docs/implementation/slot21/P8_spec.md).
% Placed after part 3 so that no section number of v0.4.0 moves.  Contract only; the code (P8) is gated on the slot-21 Fable referee of
% this part (rule 22) and on P7 (commit 104223f5, HeckeExt/HeckeLambdaQExp, HeckeLambdaQExpValidation).

\paragraph{Status and scope of this part.} Written as a pre-release (v0.5.0-pre.1, \textbf{not signed}) for the slot-21 Fable referee, before any code of
slot-21 P8 (rules~20, 21); v0.5.0-pre.2 applies that referee's repairs (F4, \file{docs/implementation/slot21/P8\_F4\_referee\_report.md}: ACCEPT WITH REPAIRS; M1 the range of (BF),
M2 the normalisation of Hida's $P_n$ read on the page, M3 (U17)(a) as Harder's index on the ordinary block, and the minors m1--m7; no label added, renamed or removed; the
tags of 0.4.0 unchanged). Parts~1 and~3 read the $\Iw$-adic \emph{algebra} from its specialisations; this part does the same for the $\Iw$-adic
\emph{ordinary parabolic cohomology} of tame level~$1$ --- Hida's $\Iw$-adic modular symbols --- and fixes what slot-21 P9 (the two-variable measure) needs
from it: (i) the finite-weight objects and their structure per sign (Definition~\ref{def:u-msymk}, Lemma~\ref{lem:u-msymsigns}); (ii) that the ordinary
part does not see the difference between the contract's lattice $\Sym^{n}$ and Hida's lattice $L_n$ (Lemma~\ref{lem:u-latticeord}, \PROVED, new);
(iii) multiplicity one --- freeness of rank one over the component --- at every \emph{étale} component, at finite weight (Proposition~\ref{prop:u-msymfree},
\PROVED, elementary) and $\Iw$-adically and in every truncation (Theorem~\ref{thm:u-lmsymfree}, \PROVED{} from \LIT{} Hida); (iv) a normalisation of the
$\Iw$-adic generator by one linear functional, its unit check modulo $(\ell,X)$, and the period-free invariants it produces
(Proposition~\ref{prop:u-msymnorm}, \PROVED; the functional itself: Lemma~\ref{lem:u-msymmass}, \PROVED{} from [E]); (v) the predictions (U15)--(U17)
(Prediction~\ref{pr:u-msym}): (U15) \PROVED{} --- the match with P7's $\Iw$-adic $q$-expansions at eigenvector level; (U16) \CONJ{} --- the
$\Iw$-adic coherence of the normalised measures; (U17) \CONJ{} --- the Eisenstein gluing on $H^{1}$. Tame level~$1$, $\ell\ge5$, as in part~1. \emph{Not in this
part}: $\Ga_1(\ell^{r})$, wild characters and the $\omega$-type truncations (slot-21/22 P10), non-étale ordinary components (the ramified blocks of
Proposition~\ref{prop:u-loewy}(4)), whose modular symbols are module-valued and are refused by name (Remark~\ref{rmk:u-msymbasis}). (The ${\LL}=4$ cells of [L37] are not among them: they are \emph{non-ordinary}, $U_{37}$ nilpotent by [L37] \texttt{prop:nonord} (indeed $U_{37}\in\mathfrak m_A^{2}$, [L37] \texttt{prop:Uposition}), hence not objects of [U] at all.)

\paragraph{Notation of this part.} $\ell\ge5$, $a$ an even class, $k\ge3$ in the class, $n=k-2$. $\Vcoef_\ell:=\Sym^{n}\Z^{2}\otimes\Zl$ with $\rhon$ of
\HA{def:rhon} on all of $\mathrm{Mat}_2(\Z)$ (multiplicative, \HA{prop:symconv}(2)), $u_t=x^{n-t}y^{t}$, and $V'_n:=\bigoplus_t\binom nt\Zl u_t\subseteq\Vcoef_\ell$,
the image of Hida's lattice $L_n(\Zl)$ under $f\mapsto Df$ (\HA{rmk:stabconvention}(3)), a $\De$-stable lattice (\Espec{prop:lvlattice}); $m_n:=\max_tv\binom nt$,
so $\ell^{m_n}\Vcoef_\ell\subseteq V'_n$ and $m_n=0$ iff $V'_n=\Vcoef_\ell$. $\partial$ is the boundary of [C]'s pair (\HA{def:pair}), $H^{1}_c$ the cohomology of the
pair, and $H^{1}_{!}:=\im(H^{1}_c\to H^{1})=\ker(H^{1}\to H^{1}(\partial))$ (exactness of the sequence of the pair; for $\ell\ge5$ this is Hida's parabolic
cohomology $H^{1}_{P}$). $e_1:=\eord(T_\ell)$ at level~$1$ and $e_\ell:=\eord(U^{(\ell)}_\ell)$ at level $\ell$, the $\ell$-adic projectors of
\HA{rmk:ordtorsionfree} and \HA{lem:levelelltorsion}(3). $\inv^{\ast}$ is [C]'s involution (\HA{def:iota}), normalised so that the Eisenstein class lies in
$H^{1}_{+}$ (\HA{rmk:iotasign}); for $2\in\cring^{\times}$ every module below splits Hecke-stably into its $\pm$-parts (\HA{prop:iotasplit}, \Espec{rmk:lviotak}).
For a maximal ideal $\mathfrak m$ and a module $M$, $M_{\mathfrak m}$ is the $\mathfrak m$-part (Fitting / generalised eigenspace).

\begin{definition}[The finite-weight objects]\label{def:u-msymk}
For $\cring\in\{\Zl,\Z/\ell^{K}\}$ put
\[
  \mathsf S^{\pm}_k(\Zl):=e_1H^{1}_{!}(\Ga;\Vcoef_\ell)^{\pm},\qquad \mathsf H^{\pm}_k(\Zl):=e_1H^{1}(\Ga;\Vcoef_\ell)^{\pm},\qquad
  \mathsf S^{\pm}_k(\Z/\ell^{K}):=\mathsf S^{\pm}_k(\Zl)/\ell^{K},
\]
with the operators $T_p$ ($p\ne\ell$) and $U:=u_{\mathrm{st}}$, the unit root of $Y^{2}-T_\ell Y+\ell^{k-1}$ in the closure of $\Zl[T_\ell]$ (\HA{prop:stabilisation}(2),(3)),
so that $T_\ell=U+\ell^{k-1}U^{-1}$ on $e_1H$. The \emph{level-$\ell$ incarnation} is $e_\ell H^{1}(\Ga_0(\ell);\Vcoef_\ell)^{\pm}$, identified with $\mathsf H^{\pm}_k(\Zl)$
by $e_\ell\circ\mathrm{dg}_\ell$, with $U^{(\ell)}_\ell\leftrightarrow U$ (\HA{prop:stabilisation}(3),(4)); here $\inv^{\ast}$ at level $\ell$ is defined by the recipe of \HA{def:iota} ($\inv$
normalises $\Ga_0(\ell)$ and commutes with $g=\mathrm{diag}(1,\ell)$, so $\mathrm{dg}_\ell=[\Ga\,g\,\Ga_0(\ell)]$ commutes with it and the identification respects the signs). $\mathsf S^{\pm}_k(\Zl)$ is the \textbf{ordinary modular-symbol module of weight $k$ and sign $\pm$}; for a maximal ideal $\mathfrak m$
of $\hcusp_k$, an \textbf{ordinary eigensymbol} of $(\mathfrak m,\pm)$ is a generator of $\mathsf S^{\pm}_k(\Zl)_{\mathfrak m}$ over $\hcusp_{k,\mathfrak m}$ when that module is free of
rank one.
\end{definition}

\begin{lemma}[Structure over a field, per sign]\label{lem:u-msymsigns}
\PROVED{} from Lemma~\ref{lem:u-levelone}, \HA{rmk:iotasign}, \HA{Hi5} and \LIT{\cite[Lemma~6.4, pp.~587--588]{hida_invent}, read on the page for this text (and by the
slot-21 referee F4, pp.~587--588 on the extracted text: ``For any subfield $K$ of $\Cx$ or $\Omega$, $H^{1}_{P}(\Ga_1(M),L_n(K))$ is free of rank $2$ over the Hecke algebra
$h_{n+2}(\Ga_1(M);K)$ for each positive integer $M$'', quoted from Shimura [25, Th.~3.51 and Th.~8.4]; in the proof $V^{-}(R)\cong V^{+}(R)\cong h_{n+2}(\Ga_1(M);R)$): for a
field $K$ and $M\ge1$ --- $M=1$ included --- $H^{1}_{P}(\Ga_1(M),L_n(K))$ is free of rank~$2$ over the Hecke algebra $h_{n+2}(\Ga_1(M);K)$, and each of Hida's eigenspaces $V^{\pm}(K)$ of the
involution $[\varepsilon]$ is free of rank~$1$ over it; the matrix $\varepsilon$ is garbled in the extracted text and is not transcribed (\VERIFYloc): it normalises $\Ga_1(M)$ and has
determinant $-1$, and for the two candidates $\smat{-1}{0}{0}{1}$, $\smat{1}{0}{0}{-1}$ the operators agree for $n$ even ($-I$ acts by $(-1)^{n}$)}. At $M=1$ the cited statement is
classical (Eichler--Shimura and multiplicity one at level one, \HA{obl:oracle}), so nothing here depends on the hypotheses of Hida's \S4 (v0.5.0-pre.2, F4 m7). Let $k\ge3$ lie in the class $a$. Then $\hcusp_k$ acts on $\mathsf S^{\pm}_k(\Zl)$
(through $T_p\mapsto T_p$, $U\mapsto U$), $\Tu_k(a)$ acts on $\mathsf H^{+}_k(\Zl)$ and $\hcusp_k$ on $\mathsf H^{-}_k(\Zl)$, and over $\Qlb$:
\[
  \mathsf S^{\pm}_k\otimes\Qlb\cong\bigoplus_{f\in\Ek_k}\Qlb\,\varphi^{\pm}_f,\qquad \mathsf H^{+}_k\otimes\Qlb=(\mathsf S^{+}_k\otimes\Qlb)\oplus\Qlb\,\mathcal E_k,\qquad
  \mathsf H^{-}_k\otimes\Qlb=\mathsf S^{-}_k\otimes\Qlb,
\]
$\varphi^{\pm}_f$ an eigenvector with the eigensystem of $f$ (each $f\in\Ek_k$ exactly once in each sign), $\mathcal E_k$ the Eisenstein class ($T_p\mapsto1+p^{k-1}$, $U\mapsto1$).
So $\rank\mathsf S^{\pm}_k=d_k$, $\rank\mathsf H^{+}_k=d_k+1$, $\rank\mathsf H^{-}_k=d_k$, $\sum=h^{1}_{\mathrm{ord}}(n)$ (Lemma~\ref{lem:u-levelone}(3)).
\end{lemma}

\begin{proof}
Over $\Ql$, $L_n\otimes\Q=\Vcoef\otimes\Q$ through $D$ (\HA{rmk:stabconvention}(3)), Hecke- and $\inv$-equivariantly. Hida's Lemma~6.4 at $M=1$, $K=\Qlb$: each sign of
$H^{1}_{!}$ is free of rank one over $h_{k}(\SL_2(\Z);\Qlb)=\Qlb^{\{\text{eigenforms}\}}$ (multiplicity one at level~$1$), so every level-one eigenform occurs once in each
sign; $e_1$ keeps the $f$ with $v(a_\ell(f))=0$, i.e.\ $\Ek_k$, with $U\mapsto\alpha_f$ (Lemma~\ref{lem:u-levelone}(1)). For either candidate $\varepsilon$, $[\varepsilon]=\pm\inv^{\ast}$ ([C]'s $\inv=\mathrm{diag}(-1,1)$
and the overall sign of \HA{rmk:iotasign}); renaming the two eigenspaces does not change ``once in each sign''. $H^{1}/H^{1}_{!}\subseteq H^{1}(\partial)$ is,
after $e_1$ and $\otimes\Q$, the Eisenstein line (\HA{Hi5}; $E_k$ is ordinary), which lies in $H^{1}_{+}$ (\HA{rmk:iotasign}). The algebra statements: $\mathsf S^{\pm}_k$
and $\mathsf H^{\pm}_k$ are $\Zl$-free (\HA{rmk:ordtorsionfree}), so the image of $\Zl[T_p,U]$ in their endomorphisms embeds into the product of the eigencharacters
that occur, which is $\hcusp_k$ resp.\ $\Tu_k(a)$ by their definition (Definition~\ref{def:u-T}, Lemma~\ref{lem:u-levelone}(1)).
\end{proof}

\begin{lemma}[The ordinary part does not see the lattice]\label{lem:u-latticeord}
\PROVED{} from \HA{def:hecke} (\HA{eq:heckechain}), \HA{prop:symconv}(2), \HA{prop:levelcorrespondence}, \HA{def:bdryop}, \HA{prop:ordinary}(3) and
\Espec{prop:lvlattice} (the $\De$-stability of $V'_n$). Put $Q_n:=\Vcoef_\ell/V'_n\cong\bigoplus_t\bigl(\Zl/\tbinom nt\bigr)u_t$, a finite $\De$-module.
\begin{enumerate}
\item For every $A\in\mathrm{Mat}_2(\Z)$ with $\abs{\det A}=\ell^{N}$, $N\ge m_n$: $\rhon(A)\Vcoef_\ell\subseteq V'_n$, i.e.\ $\rho_{Q_n}(A)=0$.
\item $T_\ell^{\,m_n}=0$ on $H^{q}(\Ga;Q_n)$ and on $H^{q}(\partial;Q_n)$ for every $q$, and $(U^{(\ell)}_\ell)^{m_n}=0$ on $H^{q}(\Ga_0(\ell);Q_n)$. Hence $e_1$ and $e_\ell$
  kill all of them.
\item The inclusion $V'_n\subseteq\Vcoef_\ell$ induces isomorphisms $e_1H^{q}(\Ga;V'_n)\cong e_1H^{q}(\Ga;\Vcoef_\ell)$, $e_1H^{q}(\partial;V'_n)\cong e_1H^{q}(\partial;\Vcoef_\ell)$,
  hence $e_1H^{1}_{!}(\Ga;V'_n)\cong\mathsf S_k(\Zl)$, and the same at level $\ell$ with $e_\ell$; all of them equivariant for every $T_p$, for $U$ and for $\inv^{\ast}$. With
  $D$: $e_1H^{1}_{!}(\Ga;L_n(\Zl))\cong\mathsf S_k(\Zl)$, Hida's lattice against the contract's.
\item Consequently, for every $f\in\Ek_k$ the change-of-lattice index of \Espec{prop:lvlattice}(2) vanishes: $d^{\pm}_f=0$, and $c^{\pm}_f=e\,v_\ell(n!)$.
\end{enumerate}
\end{lemma}

\begin{proof}
(1) Write $A=\gamma\,\mathrm{diag}(d_1,d_2)\,\gamma'$ with $\gamma,\gamma'\in\GL_2(\Z)$ and $d_1\mid d_2$, $\abs{d_1d_2}=\ell^{N}$ (Smith). $V'_n$ is stable under $\rhon(B)$ for
\emph{every} $B\in\mathrm{Mat}_2(\Z)$, determinant $<0$ included: it is the image of the divided powers $\Gamma^{n}\Z^{2}\to\Sym^{n}\Z^{2}$ (\HA{rmk:stabconvention}(3)), a natural
transformation of functors of $\Z^{2}$, so every endomorphism of $\Z^{2}$ preserves the image (this is more than the $\De$-stability of \Espec{prop:lvlattice}, which covers
$\det>0$ only; v0.5.0-pre.2, F4 m1). So $\rhon(\gamma)$, $\rhon(\gamma')$ are automorphisms of $\Vcoef_\ell$ preserving $V'_n$, $\rhon(A)\Vcoef_\ell=\rhon(\gamma)\rhon(\mathrm{diag}(d_1,d_2))\Vcoef_\ell$,
and it suffices to treat $\mathrm{diag}(d_1,d_2)$. If $\ell\mid d_1$, $\rhon(\mathrm{diag}(d_1,d_2))$ is
divisible by $\ell^{n}$ and $n\ge m_n$ ($\binom nt<\ell^{n}$). Otherwise $d_1=\pm1$, $d_2=\pm\ell^{N}$, and $u_t\mapsto\pm\ell^{Nt}u_t$; for $t=0$ the component of $Q_n$ is $0$
($\binom n0=1$), for $t\ge1$, $Nt\ge m_n\ge v\binom nt$. (2) By \HA{eq:heckechain} the $N$-fold composite of the cochain map $\widetilde T$ of $\Tdc{g}$, $g=\mathrm{diag}(1,\ell)$,
sends $\phi$ to a $\Z$-linear combination of terms $\rho(\delta_{j_1})\cdots\rho(\delta_{j_N})\,\phi(c)=\rho(\delta_{j_1}\cdots\delta_{j_N})\,\phi(c)$ with $c$ a chain and
$\det(\delta_{j_1}\cdots\delta_{j_N})=\ell^{N}$; with coefficients in $Q_n$ and $N=m_n$ every term vanishes by~(1). At level $\ell$ the blocks of \HA{prop:levelcorrespondence} are
$\rhon$ of matrices of determinant $\ell$, on the boundary those of \HA{def:bdryop}; the same argument applies. $\eord=\lim U^{j!}$ is then $0$. (3) The long exact sequences of
$0\to V'_n\to\Vcoef_\ell\to Q_n\to0$ on $\Ga$ and on $\partial$ are sequences of finitely generated $\Zl$-modules, $T_\ell$-equivariant ($\widetilde T$ is natural in the coefficient
$\De$-module, so it commutes with the maps of the short exact sequence of cochain complexes and with the connecting maps). On a finite exact complex of finitely generated
$\Zl$-modules with a compatible endomorphism $T$, the image $B$ of $\Zl[X]$ ($X\mapsto T$) in the product of the endomorphism rings is a finite commutative $\Zl$-algebra, a
product of local rings, and $e:=\lim T^{j!}\in B$ is the idempotent of the factors in which $T$ is a unit; multiplication by an idempotent of $B$ is exact on complexes of
$B$-modules, and $eM=\eord M$ in the sense of \HA{rmk:ordtorsionfree} (the finite levels, \HA{prop:ordinary}(3), and the limit; v0.5.0-pre.2, F4 m2). So $e_1$ is exact on both long
exact sequences; the $Q_n$-terms vanish by~(2). The two isomorphisms commute with restriction to $\partial$, so they identify the kernels $H^{1}_{!}$. Every operator here is a
formula in the $\rhon(\delta)$ and integral chain data, so it commutes with the inclusion of $\De$-modules; $\inv$ acts diagonally on the $u_t$. For the last sentence: $D$ is an
isomorphism of $\mathrm{Mat}_2(\Z)$-modules $L_n(\Zl)\cong V'_n$, so it identifies $H^{1}_{!}(\Ga;L_n(\Zl))$ with $H^{1}_{!}(\Ga;V'_n)$ compatibly with [C]'s operators; Hida's
$T(p)$ and $T(\ell)$ at level~$1$ (the action through the main involution, \HA{rmk:stabconvention}(2)) agree with [C]'s $T_p$, $T_\ell$ on $H^{1}(\Ga;L_n(\Zl))\otimes\Q$ (both
act on the $f$-line by $a_p(f)$: Eichler--Shimura on either side), hence on the torsion-free $e_1H^{1}_{!}(\Ga;L_n(\Zl))$ and on every Hecke-module isomorphism class read on
it (v0.5.0-pre.2, F4 m3). (4) $d^{\pm}_f$ is the index of $H^{1}(\Ga;V'_n)[f]^{\pm}$ in $H^{1}(\Ga;\Vcoef_\ell)[f]^{\pm}$;
for $f$ ordinary both $f$-lines lie in the $e_1$-parts (an $f$-line of the free quotient lies in the direct summand $e_1(H/\mathrm{tors})$, the complement being torsion-free),
which coincide by~(3); and $e_1H^{1}(\Ga;V'_n)$ is torsion-free by~(3), so it is the ordinary part of its free quotient as in \HA{rmk:ordtorsionfree}(2). Then \Espec{prop:lvlattice}(2).
\end{proof}

\noindent\emph{Remarks on Lemma~\ref{lem:u-latticeord}.} (a) It upgrades slot-13 WP-5's finding (Hi8)(c) --- $D_\ast\colon M^{\mathrm H}_m(n)\to M_m(n)$ an isomorphism at $45/45$ cells
including the $31$ non-Lucas ones (slot-13 WP-5 report, (Hi8)(c) row) --- to a theorem, and it resolves the caution of \HA{rmk:stabconvention}(3) (``lattice-dependent: the
isomorphism class of a finite-level module'') and of Remark~\ref{rmk:u-gorenstein} (``nothing is transferred'') \emph{for ordinary parts}; neither text is edited here, and whether [C]
records it is the owner's call (a [C] erratum candidate, for the referee). (b) Checkable consequence: the five eigenlines at $(40,37)$ with $d_f=1$ (\Espec{rmk:lvlatticevalues}) are
non-ordinary. (c) The lemma is about ordinary parts only; non-ordinary blocks do see the lattice (\Espec{rmk:lvlatticevalues}: $d_f=1$ on $\omega^{b_j}$-twists). (d) (v0.5.0-pre.2,
F4.) Two consistency checks already in the record: $d_f=0$ on both signs at the $\mathfrak m_\Eis$-blocks of $(282,37)$, $(226,103)$, $(258,37)$ (\Espec{rmk:lvlatticevalues}; the cusp forms
there are ordinary), and the slot-13 (Hi8)(c) count of~(a). Not a contradiction of a statement of [C]: \HA{rmk:stabconvention}(3) classifies the finite-level datum as
\emph{lattice-dependent} in the sense ``not shown independent''; item~(3) shows it independent on ordinary parts. The literature transfer of Remark~\ref{rmk:u-gorenstein} now
applies to the contract's lattice on ordinary parts (the sentence added there).

\begin{proposition}[Multiplicity one at étale components, finite weight]\label{prop:u-msymfree}
\PROVED{} from Lemmas~\ref{lem:u-msymsigns}, \ref{lem:u-latticeord}, \ref{lem:u-fibre}, Theorem~\ref{thm:u-etale}, \HA{prop:bdrychar}, \HA{cl:chi}, \HA{cor:ordprecision}(1) and
\HA{prop:stabilisation}. Let $k\ge3$ lie in the class $a$ and let $\mathfrak m$ be a maximal ideal of $\hcusp_k$ with $\hcusp_{k,\mathfrak m}$ étale over $\Zl$, so
$\hcusp_{k,\mathfrak m}\cong W(\F_q)$, $q=\ell^{r}$, and $r$ eigenforms of $\Ek_k$ lie in $\mathfrak m$ (Theorem~\ref{thm:u-etale}(2)).
\begin{enumerate}
\item $\mathsf S^{\pm}_k(\Zl)_{\mathfrak m}$ is free of rank one over $\hcusp_{k,\mathfrak m}$, for both signs. Hence $\mathsf S^{\pm}_k(\Z/\ell^{K})_{\mathfrak m}$ is free of rank one over
  $W_K(\F_q)$, and for every $t\in\hcusp_{k,\mathfrak m}$ with $\hcusp_{k,\mathfrak m}=\Zl[t]$ (e.g.\ P7's $T_{p_0}$) it is cyclic over $(\Z/\ell^{K})[t]$, with annihilator the characteristic
  polynomial of $t$ on $\hcusp_{k,\mathfrak m}$. The same holds on Hida's lattice and at level $\ell$ (Lemma~\ref{lem:u-latticeord}(3), \HA{prop:stabilisation}(3),(4)).
\item \textbf{(The boundary away from the Eisenstein ideals.)} Let $\mathfrak m_b:=(\ell,\ T_p-p^{b}-p^{\,n+1-b}:p\ne\ell)$ for $b\in\Z/(\ell-1)\Z$, the ideals of the boundary
  characters $\chi_b$ of \HA{cl:chi} ($\mathfrak m_0$ is the Eisenstein ideal; $\mathfrak m_b=\mathfrak m_{b'}$ iff $b'\equiv b$ or $b'\equiv n+1-b$, \HA{cl:chi}(3)). If $\mathfrak m\ne\mathfrak m_b$
  for \emph{every} $b\in\Z/(\ell-1)\Z$ --- condition \textbf{(BF)} --- then
  $H^{q}(\partial;\Vcoef_\ell\otimes\cring)_{\mathfrak m}=0$, and $\mathsf S^{\pm}_k(\cring)_{\mathfrak m}=\mathsf H^{\pm}_k(\cring)_{\mathfrak m}=e_1H^{1}_c(\Ga;\Vcoef_\ell\otimes\cring)^{\pm}_{\mathfrak m}=
  e_1H^{1}(\Ga;\Vcoef_\ell\otimes\Z/\ell^{K})^{\pm}_{\mathfrak m}$; on $H^{1}_c$ it is read on the modular symbols $W(\cring)$ of \Espec{rmk:lvshapes}~(B).
\item \textbf{(The cuspidal component through $\mathfrak m_\Eis$.)} Let $a\not\equiv0$ and suppose the cuspidal component through $\mathfrak m_\Eis$ is étale of residue degree one (as at the
  four classes of (U13), Prediction~\ref{pr:u-lambda}), so $\Tu_{k,\mathfrak m_\Eis}=\Zl\times_{\Z/\ell^{L_k}}\Zl$. Then (1) holds for $\mathsf S^{\pm}_k(\Zl)_{\mathfrak m}$, both signs. On
  $H^{1}$, $\mathsf H^{+}_k(\Zl)_{\mathfrak m_\Eis}\cong\Zl\times_{\Z/\ell^{j}}\Zl$ for a unique $j=j_k\in[0,L_k]$, and it is free of rank one over $\Tu_{k,\mathfrak m_\Eis}$ iff $j_k=L_k$. On
  $H^{1}_c$, the boundary symbol $p_0=u_0-u_n$ spans the kernel of $e_1H^{1}_c(\Zl)\to e_1H^{1}(\Zl)$ (it lies in the minus part of [C], geometric-even, \Espec{rmk:lvshapes}~(B)); if the
  $H^{1}_c$-primitive cuspidal symbol $r$ of that sign satisfies $r\equiv u\,p_0\pmod{\ell^{d}}$ exactly, $u$ a unit (\Espec{rmk:lvlegc} (L1)), then the class of
  $w:=(r-u\,p_0)/\ell^{d}$ generates $\mathsf S^{-}_k(\Zl)_{\mathfrak m}$ and the image of $r$ is $\ell^{d}$ times a generator.
\end{enumerate}
\end{proposition}

\begin{proof}
(1) $M:=\mathsf S^{\pm}_k(\Zl)_{\mathfrak m}$ is $\Zl$-free and $\hcusp_k$ acts on it through $\hcusp_{k,\mathfrak m}$ (Lemma~\ref{lem:u-msymsigns}); $\hcusp_{k,\mathfrak m}=W(\F_q)$ is a discrete
valuation ring and $M$ is finitely generated and torsion-free over it, hence free, of rank $\rank_{\Zl}M/r=r/r=1$ (the $r$ eigenforms of $\mathfrak m$, once each in each sign). Modulo
$\ell^{K}$ by base change; cyclicity is the case $\hcusp_{k,\mathfrak m}=\Zl[t]\cong\Zl[Y]/(\mathrm{charpoly})$. (2) $H^{q}(\partial;\Vcoef_\ell\otimes\cring)_{\mathfrak m}=0$ for both $q$ iff
$H^{q}(\partial;\Vcoef_\ell\otimes\Fl)_{\mathfrak m}=0$ for both $q$ (dévissage along the sequences $0\to\Vcoef/\ell^{i}\to\Vcoef/\ell^{i+1}\to\Vcoef/\ell\to0$, all $q$ at once, and Nakayama for
$\cring=\Zl$). On $H^{q}(\partial;\Vcoef_\ell\otimes\Fl)$, $q=0,1$, the $T^{\partial}_p$, $p\ne\ell$, are simultaneously diagonalisable and every joint eigencharacter is some $\chi_b$,
$b\in\Z/(\ell-1)\Z$ --- \HA{lem:bdrytorus}(2),(3): the eigenvalue on a torus weight vector of weight $a^{\,n-j}d^{\,j}$ is $\chi_j(p)$ in degree~$0$ and $\chi_{n-j}(p)$ in degree~$1$ ---
with \HA{prop:bdrychar}(1) for the convention. Which $b$ occur is \HA{cl:chi}(1) (degree~$0$: exactly $0\le b\le\lfloor n/\ell\rfloor$, \PROVED) and \HA{cl:chi}(2) (degree~$1$: the same
multiset, \CONJ), which is why (BF) quantifies over all of $\Z/(\ell-1)\Z$ and not over $0\le b\le\lfloor n/\ell\rfloor$ (v0.5.0-pre.2, F4 M1: through pre.1 the range was
$0\le b\le\lfloor n/\ell\rfloor$, which rests on the \CONJ{} half of \HA{cl:chi}(2)); so under (BF) the $\mathfrak m$-parts of the boundary terms of the sequence of the pair vanish, and $H^{1}_c$, $H^{1}_{!}$,
$H^{1}$ agree after localisation. $e_1H^{1}(\Ga;\Vcoef_\ell\otimes\Z/\ell^{K})=e_1H^{1}(\Ga;\Vcoef_\ell)/\ell^{K}$ is \HA{cor:ordprecision}(1). (3) $\mathsf S^{\pm}_k$ carries only cuspidal
eigencharacters (Lemma~\ref{lem:u-msymsigns}), so the argument of (1) applies to its $\mathfrak m_\Eis$-part. On $\mathsf H^{+}_{k,\mathfrak m}$ the two eigenlines are the cuspidal and the
Eisenstein one; with $A_1,A_2$ the intersections of the lattice with them, the lattice modulo $A_1\oplus A_2$ is cyclic (a subgroup of $(\Ql/\Zl)$ through either projection),
$\cong\Z/\ell^{j}$, and after rescaling one coordinate by a unit the lattice is $\{(x,y):x\equiv y\ (\ell^{j})\}$; it is stable under $(\ell^{L},0)\in\Tu_{k,\mathfrak m}$ iff $j\le L$, and
$j$ is an invariant of the $\Tu_{k,\mathfrak m}$-module (the eigenlines are intrinsic); $j=L$ is $\Tu_{k,\mathfrak m}$ itself (Lemma~\ref{lem:u-fibre}). For $H^{1}_c$: over $\Zl$,
$H^{0}(\Ga;\Vcoef_\ell)=0$ ($n\ge1$) and $H^{0}(\partial;\Vcoef_\ell)=\Zl u_0$ with $T_p$-eigenvalue $1+p^{k-1}$ and $U$-eigenvalue $1$ (\HA{prop:bdrychar}(3)), whose image is
$\Zl p_0$ (\Espec{rmk:lvlegc}); if $r-u\,p_0=\ell^{d}w$, then $r\mapsto\ell^{d}[w]$, and $[w]$ is not divisible by $\ell$ since $r$ is not $\equiv$ a multiple of $p_0$ modulo $\ell^{d+1}$.
\end{proof}

\begin{definition}[Hida's $\Iw$-adic modular symbols]\label{def:u-lmsym}
Let $V^{0}(a)$ be the $a$-part of Hida's $\Iw$-module $V^{0}$ of \HA{thm:hidacontrol}(1) (the Pontryagin dual of the ordinary part of $\varinjlim_rH^{1}(X_r,\Ql/\Zl)$, i.e.\ of the
parabolic cohomology; \cite[Thm.~3.1, p.~562]{hida_invent}, $\Iw$-free of rank $2d(a)$ by \cite[p.~589, Step~II]{hida_invent}), with the Hecke operators of Hida's universal algebra and his
involution $[\varepsilon]$. Here ``the $a$-part'' is the component of $V^{0}$ on which $\hcusp(a)$ acts --- in Hida's natural indexing by the tame diamond character this is $V^{0}(a-2)$
(``$h^{0}_2(a)$ acts on $\mathcal V^{0}(a-2)$'', \cite[p.~593]{hida_invent}, read on the page for part~5): the weight-$2$ forms of the class $a$ have nebentypus $\omega^{a-2}$.
The $\Iw$-structure on it is Hida's twisted one ($1+X\mapsto u^{2}\langle u\rangle$ on the tower, Definition~\ref{def:u-towerlambda}), so that $\Pk{k}$ is the weight-$k$ point.
Labelling only: every statement of this part uses $\mathbf M(a)$ as the module $\hcusp(a)$ acts on (text v0.2 / [U] 0.6.0-pre.1, F-U5). Put
\[
  \mathbf M(a):=\Hom_\Iw\bigl(V^{0}(a),\Iw\bigr),\qquad \mathbf M(a)=\mathbf M(a)^{+}\oplus\mathbf M(a)^{-}\ (2\in\Iw^{\times}),
\]
with the transposed operators; for a maximal ideal $\mathfrak m$ of $\hcusp(a)$, $\mathbf M(a)^{\pm}_{\mathfrak m}$; and $\mathbf M_{K,m}(a)^{\pm}:=\mathbf M(a)^{\pm}\otimes_\Iw\Iw_{K,m}$,
$\mathbf M_{K,m}(a;x)^{\pm}$ likewise (Definition~\ref{def:u-trunc}).
\end{definition}

\begin{proposition}[Control for $\mathbf M(a)$]\label{prop:u-lmsymcontrol}
\PROVED{} from \LIT{\cite[Thm.~3.1, p.~562; Lemma~4.6, p.~571; Prop.~4.7, p.~572; Thm.~6.6, p.~590]{hida_invent}, transcribed with verified locators in \HA{thm:hidacontrol},
\HA{thm:hidaweight}; Prop.~4.7 (including its sentence ``the isomorphism in each direction is explicitly given by $e\circ\mathrm{res}$ and $e_0\circ\mathrm{Tr}$, which preserve parabolic
cohomology classes'') and Thm.~6.6 read on the page for this text; \textbf{and the normalisation of Hida's weight prime, read on the page by the slot-21 referee F4
(v0.5.0-pre.2; \VERIFYloc{} lifted)}: \cite[p.~545, Introduction]{hida_invent} --- ``the evaluation of power series in $\Lambda$ at the point $\varepsilon(u)u^{k}-1$ ($u=1+p\in\Gamma$)
gives an algebra homomorphism of $\Lambda$ into $\Omega$ for each finite order character $\varepsilon$ of $\Gamma$ and each integer $k$. Let $P_{k-2}$ be the prime ideal of $\Lambda$ which is
the kernel of this morphism'' --- so Hida's $P_n$ with $\varepsilon=1$ is the kernel of $X\mapsto u^{n+2}-1$, i.e.\ $P_n=\Pk{n+2}=\Pk{k}$ in the notation of this document, the index shift
being built into Hida's notation; \HA{thm:hidacontrol}(2) does not restate it} and Lemma~\ref{lem:u-latticeord}. For $k\ge3$ in the class there are isomorphisms, compatible with $T(p)\mapsto T_p$ ($p\ne\ell$) and $U\mapsto U$, carrying the two
$[\varepsilon]$-eigenspaces onto the two $\inv^{\ast}$-eigenspaces (in some order, fixed for the class):
\[
  \mathbf M(a)/\Pk{k}\mathbf M(a)\;\cong\;eH^{1}_{P}(\Phi_1;L_n(\Zl))\;\cong\;e_0H^{1}_{P}(\SL_2(\Z);L_n(\Zl))\;\cong\;\mathsf S_k(\Zl).
\]
\end{proposition}

\begin{proof}
Thm.~6.6 gives $\mathcal V^{0}(a)[P_n]\cong eH^{1}_{P}(\Phi_1;L_n(\Zl))\otimes\Ql/\Zl$ as Hecke modules; Pontryagin duality turns it into $V^{0}(a)/P_nV^{0}(a)\cong
\Hom_{\Zl}(eH^{1}_{P}(\Phi_1;L_n(\Zl)),\Zl)$ (the latter module is $\Zl$-free, Lemma~4.6). As $V^{0}(a)$ is $\Iw$-free, $\mathbf M(a)/P_n=\Hom_{\Zl}(V^{0}(a)/P_n,\Zl)$, the double dual,
i.e.\ $eH^{1}_{P}(\Phi_1;L_n(\Zl))$ with its own operators. Prop.~4.7 ($r=1$, $N=1$, parabolic version) gives the second isomorphism; at level~$1$ Hida's $T(p)$ is this document's $T_\ell$
(\HA{def:hidalattice}, variant (H): $\Ga\smat p001\Ga=\Ga\smat100p\Ga$). The third is $D$ and Lemma~\ref{lem:u-latticeord}(3). The $U$-compatibility: on the image of level one, Hida's $U$ and
[C]'s $U^{(\ell)}_\ell$ both act as the unit root (\HA{rmk:stabconvention}(4), \HA{prop:stabilisation}(3)). Restriction, trace and $D$ commute with $\varepsilon=\inv$.
\end{proof}

\begin{theorem}[$\Iw$-adic multiplicity one at étale components]\label{thm:u-lmsymfree}
\PROVED{} from Proposition~\ref{prop:u-lmsymcontrol}, Proposition~\ref{prop:u-msymfree}(1), Theorem~\ref{thm:u-etale}(1) and Nakayama. Let $\mathfrak m$ be a maximal ideal of $\hcusp(a)$ with
$\hcusp(a)_{\mathfrak m}$ étale over $\Iw$, so $\hcusp(a)_{\mathfrak m}\cong W(\F_q)[[X]]$, $q=\ell^{r}$.
\begin{enumerate}
\item $\mathbf M(a)^{\pm}_{\mathfrak m}$ is free of rank one over $\hcusp(a)_{\mathfrak m}$, for both signs.
\item For all $K,m\ge1$ and $x\in\ell\Zl$: $\mathbf M_{K,m}(a)^{\pm}_{\mathfrak m}\cong\hcusp_{K,m}(a)_{\mathfrak m}\cong W_K(\F_q)[X]/(X^{m})$ and $\mathbf M_{K,m}(a;x)^{\pm}_{\mathfrak m}\cong\hcusp_{K,m}(a;x)_{\mathfrak m}$.
\item A generator $\Phi^{\pm}$ of (1) --- a \textbf{$\Iw$-adic ordinary eigensymbol} of $(\mathfrak m,\pm)$ --- specialises at every $\Pk{k}$, $k\ge3$ in the class, to an ordinary eigensymbol
  $\varphi^{\pm}_k$ of $(\mathfrak m_k,\pm)$ in the sense of Definition~\ref{def:u-msymk}, $\mathfrak m_k$ the image of $\mathfrak m$; two generators differ by $\hcusp(a)_{\mathfrak m}^{\times}$.
\end{enumerate}
\end{theorem}

\begin{proof}
First, $\hcusp(a)$ acts on $\mathbf M(a)$: the $\Iw$-subalgebra of $\End_\Iw(\mathbf M(a))$ generated by the $T(p)$, $p\ne\ell$, and $U$ embeds into
$\prod_{k}\End(\mathbf M(a)/\Pk{k})=\prod_k\End(\mathsf S_k(\Zl))$ over the $k\ge3$ of the class ($\mathbf M(a)$ is $\Iw$-free and $\bigcap_k\Pk{k}\Iw=0$), hence into $\prod_k\hcusp_k$
(Lemma~\ref{lem:u-msymsigns}), as the $\Iw$-subalgebra generated by the tuples $(T_p)_k$, $(U)_k$ with $X$ acting as $(x_k)_k$; so does $\hcusp(a)$, with the same image
(Theorem~\ref{thm:u-control}, the argument of Proposition~\ref{prop:u-lambdaprod}(1)); so the two are identified, $T(p)\mapsto T(p)$, $U\mapsto U$, $\Iw$-linearly (v0.5.0-pre.2, F4 m4).
$\hcusp(a)$ is finite over the complete local $\Iw$, so $\mathbf M(a)^{\pm}=\bigoplus_{\mathfrak m}\mathbf M(a)^{\pm}_{\mathfrak m}$ and each summand is $\Iw$-free (a direct summand of the
$\Iw$-free $\mathbf M(a)$, over a local ring). Its $\Iw$-rank is the $\Zl$-rank of its specialisation at one $\Pk{k}$, which is $\mathsf S^{\pm}_k(\Zl)_{\mathfrak m_k}$ (Proposition~\ref{prop:u-lmsymcontrol})
of rank $r$ (Proposition~\ref{prop:u-msymfree}(1)). $\mathbf M^{\pm}_{\mathfrak m}/\mathfrak m_\Iw$ is a module over $\hcusp(a)_{\mathfrak m}/\mathfrak m_\Iw=\F_q$ of $\Fl$-dimension $r$, hence
one-dimensional; by Nakayama over the local ring $\hcusp(a)_{\mathfrak m}$ it is cyclic, $\hcusp(a)_{\mathfrak m}\twoheadrightarrow\mathbf M^{\pm}_{\mathfrak m}$, a surjection of free $\Iw$-modules of the
same rank $r$, hence an isomorphism. (2) is base change; (3) Proposition~\ref{prop:u-lmsymcontrol} and~(1).
\end{proof}

\begin{lemma}[The ball masses on $\mathsf S_k$, and their $\hcusp$-valued form]\label{lem:u-msymmass}
\PROVED{} from \Espec{def:lvmtt} and \Espec{prop:lvmtt}(1),(2),(6) (their proofs are linear in the symbol). For $x\in e_1H^{1}_c(\Ga;\Vcoef_\ell)$ with symbol $\Phi_x$
(\Espec{rmk:lvshapes}~(B)), $c_x(a,m)$ as in \Espec{def:lvmtt}, put
\[
  \mu_x\bigl(a+\ell^{m}\Zl\bigr):=c_{U^{-m}x}(a,m)\;-\;\ell^{k-2}\,c_{U^{-m-1}x}(a,m-1)\qquad(m\ge0),
\]
the formula of \Espec{def:lvmtt} with the unit root $\alpha$ replaced by the operator $U$ of Definition~\ref{def:u-msymk} --- on $e_1H^{1}_c$ itself $U$ is the unique unit root of
$Y^{2}-T_\ell Y+\ell^{k-1}$ in the closure of $\Zl[T_\ell]\subseteq\End(e_1H^{1}_c)$, which exists by Hensel's lemma exactly as in Lemma~\ref{lem:u-levelone}(1) ($T_\ell$ is a unit on the
ordinary part, the root $Y=T_\ell$ of $Y(Y-T_\ell)$ is simple modulo the radical of this finite $\Zl$-algebra), whether or not $e_1H^{1}_c$ is torsion-free; it induces the $U$ of
Definition~\ref{def:u-msymk} on $e_1H^{1}_{!}$ and acts as $1$ on $p_0$ (v0.5.0-pre.2, F4 m5).
\begin{enumerate}
\item $\mu_x$ is a $\Zl$-valued distribution on $\Zl$, $\Zl$-linear in $x$, with $\mu_{tx}$ for $t\in\hcusp_k$ given by $t$ acting on $x$; on a $T_\ell$-eigenclass it is \Espec{def:lvmtt}'s $\mu_x$.
\item On compact open sets inside $\Zl^{\times}$ it vanishes on $\Zl p_0$, so it descends to $\mathsf S^{\pm}_k(\Zl)$ there.
\item For $\mathfrak m$ as in Proposition~\ref{prop:u-msymfree} and a compact open $B\subseteq\Zl^{\times}$ there is a unique $\hcusp_{k,\mathfrak m}$-linear $\mu^{\hcusp}(B)\colon\mathsf S^{\pm}_k(\Zl)_{\mathfrak m}\to
  \hcusp_{k,\mathfrak m}$ with $\mathrm{Tr}_{\hcusp_{k,\mathfrak m}/\Zl}\bigl(c\,\mu^{\hcusp}_x(B)\bigr)=\mu_{cx}(B)$ for all $c$; it is the \textbf{$\hcusp$-valued mass} ($=\mu_x(B)$ when $r=1$).
\end{enumerate}
\end{lemma}

\begin{proof}
(1) Linearity of \Espec{prop:lvmtt}(1) in the symbol gives $\sum_bc_x(a+b\ell^{m},m+1)=c_{T_\ell x}(a,m)-\ell^{k-2}c_x(a,m-1)$ for every $x$, and $c_y(a,m)$ depends on $a\bmod\ell^{m}$; so
$\sum_b\mu_x(a+b\ell^{m},m+1)=c_{(T_\ell U^{-1}-\ell^{k-1}U^{-2})U^{-m}x}(a,m)-\ell^{k-2}c_{U^{-m-1}x}(a,m-1)=\mu_x(a+\ell^{m}\Zl)$, as $T_\ell=U+\ell^{k-1}U^{-1}$. $U$ commutes with
$\hcusp_k$; on an eigenclass $U=\alpha$. (2) $U$ acts as $1$ on $p_0$ (its $T_\ell$-eigenvalue is $1+\ell^{k-1}$), and \Espec{prop:lvmtt}(6) with $\alpha=1$ gives $D^{(1)}_{p_0}(a,m)=0$
for $m\ge1$; the kernel of $e_1H^{1}_c\to e_1H^{1}_{!}$ is $\Zl p_0$ (proof of Proposition~\ref{prop:u-msymfree}(3)). (3) $\hcusp_{k,\mathfrak m}=W(\F_q)$ is étale over $\Zl$, so its trace form is perfect
and $M\cong\hcusp_{k,\mathfrak m}$ (Proposition~\ref{prop:u-msymfree}(1)); $\mu^{\hcusp}$ is the image of the $\Zl$-linear $c\mapsto\mu_{cx}(B)$ under $\Hom_{\Zl}(W,\Zl)\cong W$.
\end{proof}

\begin{proposition}[Normalisation of the $\Iw$-adic eigensymbol]\label{prop:u-msymnorm}
\PROVED{} (elementary). In the situation of Theorem~\ref{thm:u-lmsymfree}, let $\lambda\colon\mathbf M(a)^{\pm}_{\mathfrak m}\to\Iw$ be $\Iw$-linear and $\lambda^{\hcusp}$ its
$\hcusp(a)_{\mathfrak m}$-linear form ($\mathrm{Tr}_{\hcusp/\Iw}(c\lambda^{\hcusp}(x))=\lambda(cx)$; $\hcusp(a)_{\mathfrak m}$ is finite étale over $\Iw$).
\begin{enumerate}
\item $\lambda$ is \textbf{admissible} if $\lambda^{\hcusp}(\Phi)\in\hcusp(a)_{\mathfrak m}^{\times}$ for one, equivalently every, generator $\Phi$; equivalently $\lambda^{\hcusp}\not\equiv0$ on the one-dimensional
  $\F_q$-space $\mathbf M^{\pm}_{\mathfrak m}/\mathfrak m\mathbf M^{\pm}_{\mathfrak m}$ --- \textbf{the unit check modulo $(\ell,X)$}. Then $\Phi^{\lambda}:=\lambda^{\hcusp}(\Phi)^{-1}\Phi$ is the unique generator with
  $\lambda^{\hcusp}(\Phi^{\lambda})=1$.
\item If $\lambda$ specialises at $\Pk{k}$ to $\lambda_k\colon\mathsf S^{\pm}_k(\Zl)_{\mathfrak m_k}\to\Zl$ (through Proposition~\ref{prop:u-lmsymcontrol}), then $\Phi^{\lambda}$ specialises to
  $\Phi^{\lambda}_k:=\lambda^{\hcusp}_k(\varphi_k)^{-1}\varphi_k$ for any eigensymbol $\varphi_k$ --- the normalisation ``$\lambda_0(\Pk{k})=1$'' --- and admissibility is decided at one weight:
  $\lambda^{\hcusp}_k(\varphi_k)\in W(\F_q)^{\times}$ at one $k$ iff at every $k$.
\item For any second such $\lambda'$, $\lambda'^{\hcusp}(\Phi^{\lambda})\in\hcusp(a)_{\mathfrak m}$ does not depend on $\Phi$, and its specialisation at $\Pk{k}$ is the \emph{period-free} ratio
  $\lambda'^{\hcusp}_k(\varphi_k)/\lambda^{\hcusp}_k(\varphi_k)\in W(\F_q)$, for any $\varphi_k$.
\end{enumerate}
\end{proposition}

\begin{proof}
$\mathbf M^{\pm}_{\mathfrak m}=\hcusp_{\mathfrak m}\Phi$ (Theorem~\ref{thm:u-lmsymfree}), $\lambda^{\hcusp}(t\Phi)=t\lambda^{\hcusp}(\Phi)$, and an element of the local ring $\hcusp_{\mathfrak m}$ is a unit iff it is
non-zero modulo $\mathfrak m$; $\mathfrak m$ contains $\mathfrak m_\Iw$ and the specialisation $\hcusp_{\mathfrak m}\to\hcusp_{k,\mathfrak m_k}$ is local, so ``unit'' is read at any one weight. The trace
form commutes with base change along $\Iw\to\Iw/\Pk{k}$. (3): the ratio is $\hcusp_{\mathfrak m}^{\times}$-invariant.
\end{proof}

\noindent\emph{The normaliser of record} (the convention of P8; Prediction~\ref{pr:u-msym}). $\lambda_0:=\mu^{\hcusp}(B_0)$ of Lemma~\ref{lem:u-msymmass}(3) for the first ball
$B_0=b_0+\ell^{M_0}\Zl\subseteq\Zl^{\times}$ in the order $M_0=1,2,\dots$, $b_0=1,\dots,\ell^{M_0}-1$ with $\ell\nmid b_0$, at which $\mu^{\hcusp}_{\varphi_c}(B_0)\not\equiv0\pmod{\mathfrak m}$ at the
centre weight $c$; refused by name (\texttt{no\_unit\_ball}) if none exists up to $M_0=M_{\max}$ (default $2$). That $\lambda_0$ is the specialisation of a $\Iw$-adic functional --- what
Proposition~\ref{prop:u-msymnorm}(2) needs --- is (U16), not proved here. On the sign of $\mathsf S_k$ that carries $p_0$ on $H^{1}_c$ at $\mathfrak m_\Eis$ the symbol to use is $w$
(Proposition~\ref{prop:u-msymfree}(3)); the $H^{1}_c$-primitive $r$ has all masses $\equiv0\pmod{\ell^{d}}$ on $\Zl^{\times}$ (\Espec{rmk:lvmeasplus}(ii)) and would fail the unit check.

\begin{prediction}[(U15)--(U17): $\Iw$-adic modular symbols at the slot-21 cells]\label{pr:u-msym}
(New at v0.5.0-pre.1; stated before slot-21 P8 runs (rule~20).) Cells: the étale cuspidal components of P7 --- $(\ell,a)=(11,2)$ (the $\Delta$ family), $(13,0)$, $(23,2)$ (residue degree
$2$), and the components through $\mathfrak m_\Eis$ at $(37,32)$, $(103,24)$ --- with $K=3$, $m=1,2,3$, $N=K+m-1$ consecutive weights $k_0+j(\ell-1)$ of the class computed modulo $\ell^{K+m-1}$
(Theorem~\ref{thm:u-interp}(2)), and at least three further weights; at $\ell\in\{5,7\}$ every class has $\mathsf S_k=0$ (Theorem~\ref{thm:u-ell13}).
\begin{enumerate}
\item[(U15)] \PROVED{} (Lemma~\ref{lem:u-msymsigns}, Proposition~\ref{prop:u-msymfree}(1),(2), Proposition~\ref{prop:u-lambdaprod}(4)) \textbf{The match with P7 at eigenvector level.} For P7's component of
  residual factor $g$ of $T_{p_0}$ ($\mathfrak m=(\ell,g(T_{p_0}))$, (BF) checked), at every computed weight $k$ and both signs: $\mathsf S^{\pm}_k(\Z/\ell^{K'})_{\mathfrak m}$ is free of rank $r$ over
  $\Z/\ell^{K'}$ and cyclic over $T_{p_0}$; its characteristic polynomial is P7's $F(x_k,Y)$; and $T_p=P_p(x_k,T_{p_0})$, $U=\alpha(x_k,T_{p_0})$ as operators on it, for every prime $p\le B$,
  $p\ne\ell$ --- for $r=1$: $T_p\varphi_k=a_p(x_k)\varphi_k$, $U\varphi_k=\alpha(x_k)\varphi_k$ --- modulo $\ell^{K'}$ against P7's forms data at the weight itself, and modulo
  $\ell^{\min(K,\,m\,v(x_k-x_c))}$ against P7's truncation centred at $x_c$. At level $\ell$ (cells with $\ell\le13$): $U^{(\ell)}_\ell=\alpha(x_k,T_{p_0})$ on $e_\ell H^{1}(\Ga_0(\ell);\Vcoef_\ell\otimes\Z/\ell^{K'})_{\mathfrak m}$,
  which has rank $2r$. A failure is a defect of the code or of P7.
\item[(U16)] \CONJ{} \textbf{$\Iw$-adic coherence of the normalised measures.} For every compact open $B\subseteq\Zl^{\times}$ there is $\nu^{\pm}_B\in\hcusp(a)_{\mathfrak m}$ with
  $\nu^{\pm}_B(x_k)=\mu^{\hcusp}_{\Phi^{\pm}_k}(B)$ for every $k\ge3$ in the class, $\Phi^{\pm}_k$ the eigensymbol normalised by $\lambda_0$ --- i.e.\ the normalised measures of the weights are the
  specialisations of one $\hcusp(a)_{\mathfrak m}$-valued measure, and $\lambda_0$ is admissible in the sense of Proposition~\ref{prop:u-msymnorm}. \emph{Checkable form}: (a) the unit check holds at
  every computed weight (Proposition~\ref{prop:u-msymnorm}(2)); (b) for each ball of level $M\le2$, the coordinates of $\mu^{\hcusp}_{\Phi_k}(B)$ in the basis $1,T_{p_0},\dots,T_{p_0}^{r-1}$,
  interpolated from the $N$ weights (Theorem~\ref{thm:u-interp}), predict every further weight modulo $\ell^{\min(K',N)}$ (Theorem~\ref{thm:u-interp}(3)); (c) the distribution relation holds for
  the normalised family at every weight (Lemma~\ref{lem:u-msymmass}(1), a code check). \emph{Expected literature}: the $\Iw$-adic modular-symbol measure of Greenberg--Stevens
  \cite{greenberg_stevens} and Kitagawa's construction (not held; \file{docs/implementation/slot21/missing\_literature.md} LIT-08). v0.5.0-pre.3: Greenberg--Stevens \textbf{read}
  (rendered from the image-only scan, \file{docs/implementation/slot21/gs\_reading\_notes.md}): Prop.~6.1 ($W^{0}$ $\Lambda$-free, $\Phi_\kappa=0\iff\Phi\in P_\kappa W^{0}$), Thm.~5.13 (control after
  localising at an arithmetic point, in characteristic~$0$) and (5.5), (5.10), Thm.~5.15 give a \emph{weaker} form of (U16): a local, characteristic-$0$, meromorphic coherence of the
  period-free ratios $\mu(B)/\mu(B_0)$; the integral $\hcusp(a)_{\mathfrak m}$-form with a unit normaliser is \emph{not} in Greenberg--Stevens. The tag stays \CONJ.
  \emph{Partial reason, not a proof}: for balls of level~$1$ and modulo $\ell$, $c_x(a,1)$ is the value at $\{a/\ell,\infty\}$ of the image of $x$ under $P\mapsto P(0,1)$, which is
  $\Ga_0(\ell)$-equivariant modulo $\ell$ into $\Fl$ with the character $d\mapsto d^{\,n}$, and \LIT{\cite[Cor.~4.5, p.~570]{hida_invent}, read on the page: $l_n\circ\mathrm{res}\colon
  eH^{1}(\Phi_1,L_n(\Z/p^{r}\Z))\cong eH^{1}(\Phi_r,\Z/p^{r}\Z(n))$, also for parabolic cohomology} makes the ordinary part at $r=1$ depend on $n$ only modulo $\ell-1$; transporting Hida's $T(p)=U^{\ast}$
  (\HA{rmk:stabconvention}(2)) to this document's $U$ is not carried out, so the tag stays \CONJ. At level $M$ the same mechanism needs $\Phi_M=\Ga_0(\ell^{M})$ --- where the
  $\Iw$-action lives (P10).
\item[(U17)] \CONJ{} (for (a): \PROVED{} from \STD{} at the weights with $n<\ell-1$, see below) \textbf{The Eisenstein gluing on $H^{1}$, and the reduced symbol.} At $(37,32)$ and $(103,24)$, at every computed weight: (a) $j_k=L_k$ in
  Proposition~\ref{prop:u-msymfree}(3), i.e.\ $\mathsf H^{+}_k(\Zl)_{\mathfrak m_\Eis}$ is free of rank one over $\Tu_{k,\mathfrak m_\Eis}$ (evidence: the cuspidal symbol of the other geometric sign is
  $\equiv s\,P^{\Eis}\pmod{\ell^{d}}$ to depth exactly $d$, \Espec{cl:lvmeascong}(2), and $\delta=0$ at $(10,691)$, Proposition~\ref{prop:u-eisblock}; expected from Harder's denominator of the
  Eisenstein class, \cite{harder_eisenstein}, not read: to be confirmed); (b) on $\mathsf S^{-}_k$ at $\mathfrak m_\Eis$ the reduced symbol $w$ passes the unit check (evidence: the reduced plus
  measure has $\mu=0$ at $(30,37)$, $(282,37)$, $(226,103)$, \Espec{rmk:lvmeasplus}(iii), an observation there).
  \emph{What (a) is} (v0.5.0-pre.2, F4 M3). With $M:=\mathsf H^{+}_k(\Zl)_{\mathfrak m_\Eis}$, $A_1=M\cap(\text{cusp line})=\mathsf S^{+}_k(\Zl)_{\mathfrak m_\Eis}$ ($H^{1}_{!}=\ker\mathrm{res}$ is
  saturated in $H^{1}$) and $A_2=\Zl E$, $E$ the primitive Eisenstein class, the proof of Proposition~\ref{prop:u-msymfree}(3) gives $\ell^{j_k}=[M:A_1\oplus A_2]$. Since $e_1$ is a direct-summand
  projector of the free module $H^{+}:=H^{1}(\Ga;\Vcoef_\ell)^{+}/\mathrm{tors}$ with $E\in e_1H^{+}$ and $(1-e_1)H^{+}\subseteq H^{+}_{!}$ ($H^{+}/H^{+}_{!}\otimes\Q$ is the Eisenstein line), and the
  other $\mathfrak m$-blocks of $e_1H^{+}$ lie in $H^{+}_{!}$ whenever $H^{1}(\partial;\Vcoef_\ell)$ is torsion-free --- true for $n<\ell$ (\HA{prop:bdrychar}(1): $\VR^{T}$ has rank $1$) ---
  one has $\ell^{j_k}=[H^{+}:H^{+}_{!}\oplus\Zl E]$ for $n<\ell$: \textbf{(a) is the index statement of \Espec{cl:lveis}(2) (Harder's denominator of the Eisenstein class) restricted to the
  ordinary $\mathfrak m_\Eis$-block}, with $v(B_k/k)=L_k$ by (U1). So (a) is \PROVED{} from \STD{} (\Espec{cl:lveis}(2), \cite{harder3}, \VERIFYloc{} there) for $n<\ell-1$ --- at the cells
  $(30,37)$ ($k=32$), $(22,103)$ ($k=24$) and $(10,691)$ --- and stays \CONJ{} for $n\ge\ell-1$ exactly as \Espec{cl:lveis}(2) does (the P8 weights $k\ge68$ at $37$, $k\ge126$ at $103$); and
  $\delta(\mathfrak m_\Eis)=0$ of Proposition~\ref{prop:u-eisblock} is \emph{equivalent} to $j_k=L_k$ (on $M$ the element $t=(\ell^{L},0)$ maps into $\ell M$ iff $j_k<L$, and on the rank-one
  minus block $t$ acts as $\ell^{L}$), so (a) is \COMP{} at $(10,691)$ as well. For the code, G11 of \file{P8\_spec.md} is a \emph{defect} at $k=32$, $k=24$ and a \emph{finding} at the
  other weights.
\end{enumerate}
\emph{Evidence} (v0.5.0-pre.3; \COMP, no tag moves). \textbf{P8} (\code{HeckeLambdaModSymValidation}, commit \texttt{96cdaaac}): at the four default classes $1126$ checks and $179$ finding gates,
$0$ failing, $0$ findings fired --- (U15) at every computed weight; (U16)(a) the unit check and (U16)(b) the prediction of every further weight; (U17)(a) $j_k=L_k$ at all $8$ weights of $37$
($j_k=2$ at $k=284$); env arm \code{HECKE\_LAMBDA\_MSYM\_ALL=1} (adds $(103,24)$, weights to $k=738$; $j_k=L_k$ there too, $L=2$ at $k=228$, as read in [RC] 0.1.0-pre.5's gluing table) $1830$ checks and $290$ finding gates, $0$ failing
(v0.5.0-pre.4). The referee F4's recommended $(691,12)$ (\file{P8\_spec.md} \S2, S5) is \emph{not} in the committed env arm. \textbf{P9} (\code{HeckeLambdaTwoVarValidation}, \Espec{rmk:lvtwovar}): the
interpolated masses predict three further weights modulo $\ell^{5}$ at all four classes; at $(11,2)$ the masses predicted at $x_2$ equal 11a's own normalised MTT masses on every unit
ball to level~$3$ modulo $11^{5}$ ((U16) extended to $k=2$, with P8's normaliser of record $2+11\Z_{11}$, resp.\ $3+11\Z_{11}$); at $(37,32)$ the two-variable Eisenstein cross products vanish
at $X=\alpha_{32}$ modulo $37^{3}$ (the full Kubota--Leopoldt precision), and at $(103,24)$ modulo $103^{2}$ at $X=\alpha_{24}$ (env arm). \textbf{P3} (\code{HeckeLValueCoefficientRingValidation},
\Espec{rmk:lvcoeffring}): at $\ell=29$ and the residually non-split component $\bar g=X^{2}+22X+15$ the level-$2$ unit-ball measures of $k=52,80,108$ are proportional modulo $\mathfrak m$ to
that of $k=24$ by a scalar in $\F_{29}^{\times}$ ($0/812$ balls mismatching) --- (U16) modulo $\mathfrak m$ at $r=2$ outside P8's cells. None of this proves (U16) or (U17)(b).
\textbf{P2, P4} (v0.5.0-pre.5; \Espec{prop:lveismeas}, \Espec{rmk:lvwild}): the exact Eisenstein measure and wild characters are computed at \emph{single} weights only --- no evidence for (U16) or
(U17); their bearing on part~3 is the two remarks after Lemma~\ref{lem:u-twistvalues}. Wild characters along a family (Remark~\ref{rmk:u-msymbasis}, P10) remain not built.
\end{prediction}

\begin{remark}[What part 4 rests on, and what it does not do]\label{rmk:u-msymbasis}
\textbf{\LIT.} Hida's $\Iw$-freeness and control of the ordinary parabolic cohomology (Invent.\ Thm.~3.1, Thm.~6.6), the level comparison (Prop.~4.7), multiplicity one over fields (Lemma~6.4):
Proposition~\ref{prop:u-lmsymcontrol}, Theorem~\ref{thm:u-lmsymfree}, Lemma~\ref{lem:u-msymsigns}. \textbf{\VERIFYloc.} Only the matrix $\varepsilon$ of Lemma~6.4 (garbled; both candidates give
the same operator); the normalisation of Hida's $P_n$ is read on the page since v0.5.0-pre.2 (p.~545: $P_{k-2}=\ker(X\mapsto u^{k}-1)=\Pk{k}$).
\textbf{\PROVED{} here, elementary.} Lemma~\ref{lem:u-latticeord}, Proposition~\ref{prop:u-msymfree}, Lemma~\ref{lem:u-msymmass}, Proposition~\ref{prop:u-msymnorm}, (U15). \textbf{\CONJ.} (U16); (U17)(b);
(U17)(a) for $n\ge\ell-1$ ((U17)(a) is \PROVED{} from \STD{} for $n<\ell-1$: \Espec{cl:lveis}(2) on the ordinary block).
\textbf{Not used.} Any Gorenstein property beyond étaleness; any Galois representation; any main conjecture; Greenberg--Stevens and Kitagawa (cited only as the expected source of (U16)).
\textbf{What a weight is not.} $\mathsf S^{\pm}_k(\Z/\ell^{K})$ is one point of $\mathbf M(a)/(\ell^{K},\Pk{k})$: at level $\Ga_0(\ell)$ the cohomology carries no $\Iw$-action, so the $X$-direction comes only
from several weights, through Theorem~\ref{thm:u-interp} and a normalisation --- the generators $\varphi_k$ are each defined up to $W(\F_q)^{\times}$, and the $\Iw$-adic datum is exactly which
tuple of them is the image of one $\Phi$; Proposition~\ref{prop:u-msymnorm} fixes it by $\lambda_0$, provided (U16). \textbf{Refused by name.} A non-étale component (residually non-squarefree
$T_{p_0}$ on $e\,S_k$ at every $p_0\le B$: P7's refusal), where $\mathbf M_{\mathfrak m}$ need not be free and a ``two-variable function'' is a Fitting ideal, not a function; a component failing (BF) on
$H^{1}_c$ (other than $\mathfrak m_\Eis$, handled by Proposition~\ref{prop:u-msymfree}(3)); $\Ga_1(\ell^{r})$ and wild characters (P10) --- see part~5, \S\ref{sec:u-tower}.
\end{remark}

\paragraph{Code map of part 4.} \textbf{Built} (v0.5.0-pre.3; slot-21 P8, commit \texttt{96cdaaac}, after F4's ACCEPT WITH REPAIRS of this part:
\code{HeckeLambdaModSymValidation}, \code{standard}, $\approx24$\,s, $0.29$\,GB, $1126$ checks and $179$ finding gates, record \code{hecke\_lambda\_msym/v1}; and slot-21 P9, \file{HeckeExt/HeckeLambdaTwoVar},
\code{HeckeLambdaTwoVarValidation}, \code{standard}, $\approx7.5$\,s, $71$ checks and $41$ finding gates, record \code{hecke\_lambda\_twovar/v1}, \Espec{rmk:lvtwovar}). \emph{The plan, as written at v0.5.0-pre.1}:
\textbf{Planned} (slot-21 P8; Opus; only after the Fable referee's ACCEPT of this part, rule~22; implementation spec
\file{docs/implementation/slot21/P8\_spec.md}). A new \file{HeckeExt/HeckeLambdaModSym} on top of P7's \file{HeckeExt/HeckeLambdaQExp} (\code{BuildLambdaQExpansion}, \code{SpecialiseAt},
\code{FactorSquarefreeModEll}) and WP-HL's \code{ComputeFormsWeight}, with the symbols, $T_p$, $\inv$ and the measure of slot-20 WP-MEAS / WP-LIN (\file{HeckeExt/LValue/BirchModularSymbol},
\file{ModularSymbolMeasure}); the level-$\ell$ arm on \file{Hecke/HeckeLevelEll}; record \code{hecke\_lambda\_msym/v1} (FNV per line, not pinned, no golden line); validation
\code{HeckeLambdaModSymValidation} (\code{standard}) gating (U15) and the checkable forms of (U16), (U17), with the lattice check of Lemma~\ref{lem:u-latticeord} at small cells and its consequence
at $(40,37)$. No statement of this part rests on the code.

%======================================================================
% v0.6.0-pre.1 (slot 22, P10-C): part 5, the text of docs/implementation/slot22/P10C_contract_text.tex.txt v0.2 part (B) verbatim, as
% refereed by F-U5 (P10C_F4_referee_report.md, ACCEPT WITH REPAIRS); its in-text "v0.6.0-pre.2" annotations name that refereed revision.
%======================================================================
\section{The tower \texorpdfstring{$\Ga_1(\ell^{r})$}{Gamma1(l^r)}, wild characters, and the \texorpdfstring{$\Omega$}{Omega}-type truncations (part~5)}\label{sec:u-tower}
%----------------------------------------------------------------------
% NEW at v0.6.0-pre.1 (hecke-extensions slot 22, P10-C, text first; plan docs/implementation/slot22/00-slot22-plan.md v0.2 "P10-C", owner
% ruling O3; implementation spec docs/implementation/slot22/P10_spec.md).  Placed after part 4 so that no section number of v0.5.x moves.

\paragraph{Status and scope of this part.} Written as a pre-release (v0.6.0-pre.1, \textbf{not signed}) for the slot-22 Fable referee F-U5, before any code of slot-22 P10
(rules~20--22). Parts~3 and~4 obtain the $\Iw$-adic algebra and the $\Iw$-adic modular symbols of a class from finitely many weights, by interpolation
(Theorem~\ref{thm:u-interp}): ``at level $\Ga_0(\ell)$ the cohomology carries no $\Iw$-action, so the $X$-direction comes only from several weights'' (Remark~\ref{rmk:u-msymbasis}).
This part reads the $\Iw$-action where Hida defines it --- on the ordinary cohomology of the tower $\Ga_1(\ell^{r})$, through the diamond operators --- at \emph{one} level and weight~$2$:
(i) the $\Iw$-structure on the tower, Hida's $\omega_{2,r}$ and the truncations $\Iw/(\ell^{K},\Om r)\cong(\Z/\ell^{K})[G_r]$ (Definition~\ref{def:u-towerlambda},
Lemma~\ref{lem:u-omegar}); (ii) the computed objects, on the tame-character groups $\GaT r$ (Definition~\ref{def:u-towerobj}); (iii) the dictionary between the diamonds and the classes
(Proposition~\ref{prop:u-towerorient}, \COMP{} at $r=1$ and gated); (iv) Hida's control read on these objects: freeness over the group ring, the tower, the Hecke algebra modulo
$(\ell^{K},\Om r)$, the weight fibres, the wild points (Theorem~\ref{thm:u-towercontrol}, \PROVED{} from \LIT{} Hida); (v) the Eisenstein family in closed form
(Proposition~\ref{prop:u-towereis}); (vi) multiplicity one and the computed structure constants $a_p^{(r)}$ at étale rank-one components, and the $X$-adic truncations they contain
\emph{without interpolation} (Proposition~\ref{prop:u-towerfree}, Corollary~\ref{cor:u-towertrunc}); (vii) the predictions (U18)--(U22) (Prediction~\ref{pr:u-tower}). Tame level~$1$,
$\ell\ge5$, as in part~1; $n=0$ (weight~$2$) throughout --- Hida's theorem makes weight~$2$ sufficient. \emph{Not in this part}: tame level $N>1$ (slot-23 P11); components glued to the
Eisenstein component ($L_k>0$, e.g.\ $(37,32)$) and non-étale components (refused by name, Remark~\ref{rmk:u-towerbasis}); Kitagawa's measure is a stretch with its own \CONJ{}
(Remark~\ref{rmk:u-towerkitagawa}).

\paragraph{Notation of this part.} $u=1+\ell$, $x_k=u^{k}-1$, $\Pk k=X-x_k$ (part~3). $\mathsf H_r:=\mu_{\ell-1}(\Z/\ell^{r})\le G_{\ell^{r}}=(\Z/\ell^{r})^{\times}$, the Teichm\"uller subgroup,
and $G_r:=G_{\ell^{r}}/\mathsf H_r\cong(1+\ell\Z)/(1+\ell^{r}\Z)$, cyclic of order $\ell^{r-1}$, generated by the class $g$ of $u$. $\GaT r:=\Ga_{\mathsf H_r}(\ell^{r})$ of \HA{def:nebgammaH}
($\GaT1=\Ga_0(\ell)$, as $\mathsf H_1=G_\ell$). For an even class $a$, $\chi_a:=\omega^{a-2}$, viewed as a character of $G_{\ell^{r}}$ through $(\Z/\ell^{r})^{\times}\to\F_\ell^{\times}$ and the
Teichm\"uller lift. For $p\ne\ell$, $s(p)\in\Zl$ with $\langle p\rangle_\ell=u^{s(p)}$ (as in Definition~\ref{def:u-eischar}). $O_j$, $\mathcal O_{j,e}$, $\pi$, $\psi_t$, $\chi_{i,t}$ are those of
\HA{def:nebwild}.

\begin{definition}[The $\Iw$-structure on the tower]\label{def:u-towerlambda}
For $r\ge1$ put $\Om r(X):=(1+X)^{\ell^{r-1}}-u^{2\ell^{r-1}}\in\Iw$ and $\IwT Kr:=\Iw/(\ell^{K},\Om r)$. On a $\Zl[G_r]$-module $M$ (here: $G_r$ acting by the diamonds of
\HA{prop:nebgammaH}(4)), the \textbf{tower $\Iw$-structure} is
\[
  \iota_r\colon\Iw\to\End(M),\qquad 1+X\mapsto u^{2}\langle u\rangle .
\]
\end{definition}

\begin{lemma}[$\Om r$ and the group ring]\label{lem:u-omegar}
\PROVED{} (elementary; Hida's notation \LIT{\cite[p.~552--553]{hida_invent}, read on the page for this text: $P_{\chi}=\iota(u)-\chi(u)$, $P_{k,\varepsilon}$ for $\chi(\gamma)=\gamma^{k}\varepsilon(\gamma)$,
$\omega_{k,r}=\prod_\varepsilon P_{k,\varepsilon}$ over the characters $\varepsilon$ of $\Ga/\Ga_r$})
\begin{enumerate}
\item $\Om r=\prod_{\varepsilon}\bigl((1+X)-\varepsilon(u)u^{2}\bigr)$ over the $\ell^{r-1}$ characters $\varepsilon$ of $G_r$ --- Hida's $\omega_{2,r}$ --- and $\Om r\equiv X^{\ell^{r-1}}\pmod\ell$: a distinguished
  polynomial of degree $\ell^{r-1}$.
\item $\iota_r$ induces $\Iw/\Om r\cong\Zl[G_r]$ and $\IwT Kr\cong(\Z/\ell^{K})[G_r]$, local with residue field $\F_\ell$; $1,X,\dots,X^{\ell^{r-1}-1}$ is a $\Z/\ell^{K}$-basis
  (Weierstrass division by $\Om r$). The \textbf{Weierstrass representative} of an element is its expression in this basis.
\item For $r'\le r$, $\Om{r'}\mid\Om r$, and $\Iw/\Om{r'}=\Zl[G_r]/(g^{\ell^{r'-1}}-1)$.
\item For an integer $k$: $\Om r(x_2)=0$ and $v(\Om r(x_k))=r+v(k-2)$ for $k\ne2$. Hence $\Iw/(\ell^{K},\Om r,\Pk k)=\Z/\ell^{K_r(k)}$, $K_r(k):=\min(K,\,r+v(k-2))$ ($K_r(2)=K$).
\item For $c\in\Z$ with $v(c-2)=0$ and $m\ge1$: $\Om r\in(\ell^{K},(X-x_c)^{m})$ iff $K\le r$ when $m=1$, and iff $m\le\ell^{r-1}$ and $K\le r-1-\lfloor\log_\ell(m-1)\rfloor$ when $m\ge2$.
  The same holds for $(\ell^{K},X^{m})$ ($c=0$). In that case $\Iw_{K,m}(x_c)$ (resp.\ $\Iw_{K,m}$) of Definition~\ref{def:u-trunc} is a quotient of $\IwT Kr$.
\end{enumerate}
\end{lemma}

\begin{proof}
(1) $\prod_{\zeta^{\ell^{r-1}}=1}(Y-\zeta u^{2})=Y^{\ell^{r-1}}-u^{2\ell^{r-1}}$ at $Y=1+X$; modulo $\ell$, $(1+X)^{\ell^{r-1}}\equiv1+X^{\ell^{r-1}}$ and $u\equiv1$. (2) $g\mapsto u^{-2}(1+X)$ sends
$g^{\ell^{r-1}}-1$ to $u^{-2\ell^{r-1}}\Om r$; Weierstrass preparation for the distinguished $\Om r$. (3) The roots of $\Om{r'}$ are among those of $\Om r$; $g^{\ell^{r'-1}}-1\mapsto$ a unit times
$\Om{r'}$. (4) $\Om r(x_k)=u^{2\ell^{r-1}}(u^{(k-2)\ell^{r-1}}-1)$ and $v(u^{m}-1)=1+v(m)$. (5) The Taylor coefficients of $\Om r$ at $x_c$ are $\binom{\ell^{r-1}}{i}(1+x_c)^{\ell^{r-1}-i}$ ($1\le
i\le\ell^{r-1}$, of valuation $r-1-v(i)$ for $i<\ell^{r-1}$ and $0$ at $i=\ell^{r-1}$) and $\Om r(x_c)$, of valuation $r+v(c-2)=r$; $\Om r$ lies in $(\ell^{K},(X-x_c)^{m})$ iff the
coefficients of order $<m$ have valuation $\ge K$. \COMP{} check: \file{docs/implementation/slot22/p10c/p10c\_predict.py} \S4, the criterion against the definition at
$\ell\in\{5,7,11,13\}$, $r\le3$, $K\le4$, $m\le\ell^{r-1}+2$: $0$ mismatches.
\end{proof}

\begin{definition}[The computed objects]\label{def:u-towerobj}
Let $a$ be an even class, $r,K\ge1$, $\cring\in\{\Zl,\Z/\ell^{K}\}$. On the $\GaT r$-model of \HA{def:nebgammaH} with character $\chi_a$, $n=0$, let $T_p$ ($p\ne\ell$), $U:=U_\ell$ and
$\langle d\rangle$ ($d\in G_r$) be the operators of \HA{prop:nebgammaH}(4),(5) (on $H^{1}$ by \HA{prop:nebgammaHheilbronn}), $e:=\lim_jU^{j!}$, and
\[
  \mathsf M_r(a;\cring):=e\,H^{1}(\GaT r;\cring(\chi_a)),
\]
a $\Iw$-module through $\iota_r$ and a module over the Hecke operators. For $b\in\Z/(\ell-1)$ let $\mathfrak m_b$ be the residual \textbf{boundary systems}
$T_p\mapsto\omega^{b}(p)+\omega^{a-2-b}(p)\,p$ ($p\ne\ell$), $U\mapsto1$; $\mathsf M^{\Eis}_r$ is the sum of the generalised eigenspaces of the $\mathfrak m_b$ in $\mathsf M_r$ and $\mathsf M^{\cusp}_r$
the sum of the others (condition (BF$'$) of a component $\mathfrak m$: $\mathfrak m\ne\mathfrak m_b$ for every $b$). $\hcusp_r(a;K)$ is the $\Iw$-subalgebra of $\End(\mathsf M^{\cusp}_r(a;\Z/\ell^{K}))$
generated by the $T_p$ and $U$ (with $\Iw$ acting through $\iota_r$), and $a^{(r)}_p$, $\alpha^{(r)}$ denote the elements of $\IwT Kr$ by which $T_p$, $U$ act when they act by scalars
(Proposition~\ref{prop:u-towerfree}).
\end{definition}

\begin{proposition}[The dictionary: diamonds, classes, and the orientation of $U$]\label{prop:u-towerorient}
Let $a$ be an even class, $r\ge1$, $\psi\colon G_r\to\mathcal O^{\times}$ a character. On $e H^{1}(\Ga_0(\ell^{r});\mathcal O(\chi_a\psi))$ --- by \HA{lem:nebshapiro}(a) the $\psi$-specialisation of
$\mathsf M_r(a;\cdot)$ --- the Hecke systems ($T_p$, $p\ne\ell$; $U$) are those of the $U$-ordinary weight-$2$ eigenforms of level $\ell^{r}$ and character $\chi_a\psi$ in the convention
$f|_2\sigma_d=\chi_a\psi(d)f$, $U\mapsto a_\ell$ --- the normalisation in which $E_2(1,\chi)$ has $a_p=1+\chi(p)p$, $a_\ell=1$, and $E_2(\chi,1)$ has $a_p=\chi(p)+p$, $a_\ell=\ell$, so that
``the ordinary Eisenstein line carries $1+\chi(p)p$'' is a statement about [C]'s model alone (v0.6.0-pre.2, F-U5 m3); their residual systems are those of the class-$a$ ordinary forms of level one (for $k\ge3$ in the class: $T_p\equiv a_p(f_k)$), together with
the Eisenstein system $1+\chi_a\psi(p)\,p$, $U=1$.
\textbf{Status.} \PROVED{} for the $T_p$-systems at $r=1$ from \HA{NB4}(c) (the untwisted constituent $\Sym^{a-2}$ lies in the $\omega^{a-2}$-part); \COMP{} for the $U$-orientation at $r=1$
(\file{docs/implementation/slot22/p10c/p10c\_orient.py}, a transcription of \HA{def:cosetrepsN}, \HA{eq:MHchi}, \HA{def:kappaN} and \HA{rmk:nebheilbronn}, no project code: at
$(\ell,\omega^{i})\in\{(5,\omega^2),(7,\omega^2),(7,\omega^4),(11,1),(11,\omega^2),(13,\omega^{10}),(13,\omega^2),(13,\omega^4)\}$ the $U$-ordinary part of the $\omega^{i}$-part has rank $1$, or $3$ at
$(11,1)$ and $(13,\omega^{10})$, its Eisenstein line carries $T_p=1+\omega^{i}(p)p$ and, at the seven cells with $\omega^{i}\ne1$ where the two differ, never $\omega^{i}(p)+p$, and at $(13,\omega^{10})$ the cusp system $T_2\equiv\tau(2)$, $T_3\equiv\tau(3)\pmod{13}$
is ordinary while at $(13,\omega^{2})$ the cusp form is not); \textbf{gated} on the project code at every default cell and every $r$ by (U18)(o) --- a failure is a defect of \emph{this text}
and stops P10 before any $\Iw$-adic reading.
\end{proposition}

\noindent\emph{Remark on the orientation} (why it has to be checked, not assumed; reworded at v0.6.0-pre.2, referee F-U5 m1--m3). \HA{rmk:stabconvention}(2): Hida's $T(p)$ at level $p^{r}$ is
this document's $U^{\ast}$, not $U_\ell$; and Hida lets $\sigma\equiv\smat\ast\ast0z$ act on cohomology through the main involution. On the $\Ga_1(\ell^{r})$ cohomology three conventions move
\emph{together} under either the passage $\alpha\mapsto\adjg\alpha$ or conjugation by the Atkin--Lehner element: $U\leftrightarrow U^{\ast}$, $\langle d\rangle\leftrightarrow\langle d\rangle^{-1}$,
and $T_p\leftrightarrow T_p\langle p\rangle^{\mp1}$ for $p\ne\ell$ ($\Ga_1\smat100p\Ga_1$ and $\Ga_1\smat p001\Ga_1$ are distinct double cosets, differing by a diamond; on forms, $a_p(f|W)=\bar\chi(p)a_p(f)$).
So no formal cancellation argument fixes the dictionary; under the \emph{swapped} dictionary the ordinary cusp system would sit on the $\chi_a^{-1}$-part twisted by $\chi_a^{-1}(p)$ (at $\ell=13$: on $\omega^{2}$
instead of $\omega^{10}$, with $T_p\equiv p^{2}\tau(p)$), and the ordinary Eisenstein line would read $\chi(p)+p$ or $1+\bar\chi(p)p$. The computation decides: in [C]'s $\omega^{i}$-character model the
$U_\ell$-ordinary Eisenstein line carries $1+\omega^{i}(p)p$ --- the system of $E_2(1,\omega^{i})$, whose $a_\ell=1$ is a unit, while $E_2(\omega^{i},1)$ has $a_\ell=\ell$ --- and the untwisted $\tau(p)$ is
ordinary on $\omega^{10}$ at $13$ (on $\omega^{0}$ at $11$); i.e.\ \emph{[C]'s $\chi$-model is the cohomology of the forms of nebentypus $\chi$ in the classical normalisation $f|_2\sigma_d=\chi(d)f$, and
[C]'s $U_\ell$ is the classical $U$}: the class $a$ sits on the $\chi_a=\omega^{a-2}$-part of [C]'s $U$ (Serre's untwisted congruence $f_k\leftrightarrow$ weight $2$, nebentypus $\omega^{k-2}$),
matching Hida's ``$h^{0}_2(a)$ acts on $\mathcal V^{0}(a-2)$'' \LIT{\cite[p.~593]{hida_invent}, read on the page}. Only the Eisenstein half of this is provable from the boundary; the cusp half
(which constituent of \HA{NB4}(c) is $U_\ell$-ordinary) is outside (NB4)(c)'s monoid (determinant prime to $\ell$) and stays \COMP{} + gated.

\begin{theorem}[Hida's control, read on the tower]\label{thm:u-towercontrol}
Let $a$ be an even class, $d:=d(a)=\rank_\Iw\hcusp(a)$ (Theorem~\ref{thm:u-control}), $r\ge1$, $K\ge1$.
\begin{enumerate}
\item \textbf{(Freeness.)} $\mathsf M_r(a;\Zl)$ is free over $\Iw/\Om r\cong\Zl[G_r]$ of rank $2d+1$, and $\mathsf M_r(a;\Z/\ell^{K})=\mathsf M_r(a;\Zl)/\ell^{K}$ is free over $\IwT Kr$ of rank $2d+1$.
  If the cuspidal and Eisenstein components of $\Tu(a)$ are disjoint ($L_k=0$; every class at $\ell\le13$, Theorem~\ref{thm:u-ell13}) then $\mathsf M_r=\mathsf M^{\cusp}_r\oplus\mathsf M^{\Eis}_r$
  with $\mathsf M^{\cusp}_r$ free of rank $2d$ and $\mathsf M^{\Eis}_r$ free of rank $1$.
\item \textbf{(The tower.)} For $1\le r'\le r$: $\mathsf M_r(a;\cring)/(g^{\ell^{r'-1}}-1)\cong\mathsf M_{r'}(a;\cring)$ \emph{as modules over the $T_p$, $U$ and the diamonds} (ranks, characteristic
  polynomials, and the structure constants $a^{(r)}_p\mapsto a^{(r')}_p$ of Proposition~\ref{prop:u-towerfree}: what (U18)(b) tests). The left side is $eH^{1}(\Ga_{\mathsf H'}(\ell^{r});\cring(\chi_a))$ by
  \HA{lem:nebshapiro}(b); the cohomological map realising the isomorphism in [C]'s $U$-convention is $e\circ\mathrm{dg}_{\ell^{r-r'}}$, \emph{not} $e\circ\mathrm{res}$ (Hida's Cor.~4.5 is stated for
  $e(T(p))=e(U^{\ast})$ and transports under $W_{\ell^{r}}$, which normalises $\Ga_1(\ell^{r'})\cap\Ga_0(\ell^{r})$ but not $\Ga_1(\ell^{r'})$, to the other degeneracy, exactly as in
  \HA{rmk:stabconvention}(2) at $r'=0$); that map is not written in [C] for $\Ga_{\mathsf H}$ with a character and is \emph{not} a statement of this text (v0.6.0-pre.2, F-U5 M2).
\item \textbf{(The Hecke algebra.)} $\hcusp(a)\to\End(\mathsf M^{\cusp}_r(a;\Z/\ell^{K}))$ ($T(p)\mapsto T_p$, $U\mapsto U$, $\Iw$ through $\iota_r$) factors through
  $\hcusp(a)/(\ell^{K},\Om r)$ and is onto $\hcusp_r(a;K)$; at every component $\mathfrak m$ with $\hcusp(a)_{\mathfrak m}$ étale, $\hcusp_r(a;K)_{\mathfrak m}\cong\hcusp(a)_{\mathfrak m}/(\ell^{K},\Om r)$.
\item \textbf{(Weight fibres.)} At such a component, for every integer $k$ in the class: the characteristic polynomial of $T_p$ on $\hcusp_r(a;K)_{\mathfrak m}\otimes_\Iw\Iw/\Pk k$ is $c_{k,p}$ restricted
  to $\mathfrak m$ (Lemma~\ref{lem:u-levelone}(2); for $k=2$ the level-$\ell$ weight-$2$ ordinary forms) modulo $\ell^{K_r(k)}$ (Lemma~\ref{lem:u-omegar}(4)); likewise $U$ against the unit roots.
\item \textbf{(Wild points.)} For $\psi=\psi_t\colon G_r\to\mu_{\ell^{j}}$ ($j\le r-1$) and $e\le K\varphi_j$: $\mathsf M_r(a;\Z/\ell^{K})\otimes_{\IwT Kr}\mathcal O_{j,e}$, $1+X\mapsto u^{2}\zeta^{t}$, is
  $eH^{1}(\Ga_0(\ell^{r});\mathcal O_{j,e}(\chi_a\psi_t))$, Hecke-equivariantly, free of rank $2d+1$ over $\mathcal O_{j,e}$.
\end{enumerate}
\textbf{Status.} (1)--(3) \PROVED{} from \LIT{\cite[Thm.~3.1(i)--(iii), p.~562; Cor.~4.5, p.~570; Thm.~1.2 and the display after it, p.~553; Lemma~4.6, p.~571; the twist, p.~593]{hida_invent}, all
read on the page for this text and re-read on the page images by the referee F-U5}, \STD{the ordinary weight-$2$ forms of level $\ell^{r}$ whose character has conductor $\le\ell^{r'}$ are old from level
$\ell^{r'}$: a newform of level $\ell^{r}$, $r\ge2$, with character of conductor dividing $\ell^{r-1}$ has $a_\ell=0$ \cite[Thm.~4.6.17]{miyake} \VERIFYloc}, \HA{lem:nebshapiro}, \HA{lem:nebwildtors} and
Theorem~\ref{thm:u-control} --- \textbf{given Proposition~\ref{prop:u-towerorient}} (the forms-side dictionary ``class $a$ $\leftrightarrow$ nebentypus $\omega^{a-2}$, untwisted, $U$-ordinary'' on
[C]'s model: \COMP{} at $r=1$, gated by (U18)(o); every Hida convention is the $W_{\ell^{r}}$-conjugate of [C]'s, so the rank $2d+1$ of the $\chi_a$-part rests on it; v0.6.0-pre.2, F-U5 M3); (4) \PROVED{}
from (3) and Lemma~\ref{lem:u-omegar}(4); (5) \PROVED{} from (1) and \HA{lem:nebshapiro}(a).
\end{theorem}

\begin{proof}
(1) Let $Y_r$ be the open modular curve of $\Ga_1(\ell^{r})$ and $\mathcal W^{0}_r:=eH^{1}(Y_r;\Ql/\Zl)$, $W^{0}$ the Pontryagin dual of $\varinjlim_r\mathcal W^{0}_r$. By Hida's Thm.~3.1(ii) $W^{0}$ is
$\Iw$-free, the group $Z\supseteq\Ga=1+\ell\Zl$ acting through $\Ga_0(\ell^{r})\to(\Z/\ell^{r})^{\times}$, $\smat abcd\mapsto d$ (p.~562); by Thm.~3.1(i) $\mathcal W^{0}_r=(\varinjlim\mathcal W^{0})^{\Ga_r}$, so
$W^{0}/(\gamma_r-1)W^{0}$, $\gamma_r=u^{\ell^{r-1}}$, is the dual of $\mathcal W^{0}_r$ and is free over $\Zl[\Ga/\Ga_r]=\Zl[G_r]$. The dual of $\mathcal W^{0}_r=eH^{1}(Y_r;\Zl)\otimes\Ql/\Zl$ ($eH^{1}(Y_r;\Zl)$ is
$\Zl$-free: Lemma~4.6, p.~571, and $eH^{0}=0$, $U$ acting on $H^{0}$ by $\ell$) is $\Hom_{\Zl}(eH^{1}(Y_r;\Zl),\Zl)$; a $\Zl$-dual of a free $\Zl[G]$-module is free ($\Hom_{\Zl}(\Zl[G],\Zl)\cong\Zl[G]$).
So $eH^{1}(\Ga_1(\ell^{r});\Zl)$ is $\Zl[G_r]$-free. Its $\chi_a$-part for the tame diamonds ($\abs{\mathsf H_r}=\ell-1$ is a unit; \HA{prop:nebgammaH}(6)) is $\mathsf M_r(a;\Zl)$, free over $\Zl[G_r]$. Its rank is
the $\Iw$-rank of the matching component of $W^{0}$: by the twist of p.~593 the class $a$ of the Hecke algebra corresponds to the natural tame character $a-2$, which on [C]'s side is the
$\chi_a$-part (Proposition~\ref{prop:u-towerorient}); there $\rank_\Iw V^{0}=2\rank_\Iw h^{0}=2d$ (Thm.~3.1(iii), component by component through Thm.~6.6 and Lemma~6.4, as in
Proposition~\ref{prop:u-lmsymcontrol}), and $\rank W^{0}-\rank V^{0}=\tfrac12\varphi(\ell)=\tfrac{\ell-1}2$ in total at $N=1$ (Thm.~3.1(iii)), one in each of the $\tfrac{\ell-1}2$ even classes, which each
carry the $\Iw$-adic Eisenstein series (Remark~\ref{rmk:u-hidaM}, \LIT{\cite[Thm.~7.2, p.~272]{hida_ens}}). Modulo $\ell^{K}$: $H^{2}=0$ and $eH^{0}=0$, so $eH^{1}(\Z/\ell^{K})=eH^{1}(\Zl)/\ell^{K}$. The
splitting: under disjointness the Hecke algebra of $\mathsf M_r$ is a product, and the residual systems of the cuspidal factor are not boundary systems (the boundary of $\GaT r$ carries the
Eisenstein systems $\varphi(p)+\psi(p)p$ of $E_2(\varphi,\psi)$, $\varphi\psi=\chi_a\cdot(\text{wild})$, residually $\omega^{b}(p)+\omega^{a-2-b}(p)p$ --- \STD{\cite[Thm.~4.5.2, 4.6.2]{diamond_shurman}
\VERIFYloc}; \HA{lem:bdrytorus} is level~$1$ with trivial character and does \emph{not} cover this, F-U5 m8; (BF$'$) is checked, (U18)(c)); the ranks are those of $V^{0}$ and $W^{0}/V^{0}$.
(2) Hida's Cor.~4.5 (p.~570, first display, $A=\Zl$, $n=0$): $\mathrm{res}\colon eH^{1}(\Ga_1(\ell^{r'});\Zl)\cong eH^{1}(\Ga_1(\ell^{r'})\cap\Ga_0(\ell^{r});\Zl)$ for $r>r'>0$; take $\chi_a$-parts; and
\HA{lem:nebshapiro}(b) for $\GaT r\trianglelefteq\Ga_{\mathsf H'}(\ell^{r})$, $\mathsf H'$ the preimage of $\mathsf H_{r'}$, quotient $\langle g^{\ell^{r'-1}}\rangle$; then reduce modulo $\ell^{K}$ as in (1).
(3) Hida p.~553: $h^{0}(N;\mathcal O)/\omega_{k,r}h^{0}\cong h^{0}_k(\Ga_1(Np^{r});\mathcal O)$ for $k\ge2$, with the $\Iw$-structure of p.~552 ($l\equiv1\bmod Np$ acts by $l^{2}T(l,l)$). At $k=2$ and on
cohomology $T(l,l)=\langle l\rangle$ (\HA{prop:nebgammaH}(5) at $n=0$ with $\chi_a(l)=1$ for $l\equiv1\bmod\ell$), so $l^{2}T(l,l)=\iota_r((1+X)^{s(l)})$: the $\Iw$-structures agree, and $\omega_{2,r}=\Om r$
(Lemma~\ref{lem:u-omegar}(1)). The weight-$2$ ordinary Hecke algebra of $\Ga_1(\ell^{r})$ acts faithfully on $eH^{1}_{!}(\Ga_1(\ell^{r});\Zl)$ (Eichler--Shimura and multiplicity one over a field,
\LIT{\cite[Lemma~6.4, pp.~587--588]{hida_invent}}, as in Lemma~\ref{lem:u-msymsigns}), hence $\hcusp(a)/\Om r\to\End(\mathsf M^{\cusp}_r(a;\Zl))$ is injective with image the $\Zl$-algebra of the operators;
reduce modulo $\ell^{K}$ (onto). At an étale component $\mathsf M^{\cusp}_{r,\mathfrak m}$ is free over $\hcusp(a)_{\mathfrak m}/\Om r$ (Proposition~\ref{prop:u-towerfree}), so the reduction stays faithful.
(4) Specialisation commutes with characteristic polynomials; $\hcusp(a)/\Pk k=\hcusp_k$ (Theorem~\ref{thm:u-control}) and $(\ell^{K},\Om r,\Pk k)=(\ell^{K_r(k)},\Pk k)$. (5) \HA{lem:nebshapiro}(a) with
$\Ga''=\GaT r$, $\Ga'=\Ga_0(\ell^{r})$, $P=\mathcal O_{j,e}$ through $g\mapsto\zeta^{t}$, and (1); $e$ commutes (\HA{lem:nebshapiro}); the torsion of \HA{lem:nebwildtors}(1) is not ordinary.
\end{proof}

\noindent\emph{The route of record for (1), without Thm.~3.1 and without any convention transport} (promoted from ``a second route'' by the referee F-U5, v0.6.0-pre.2; the Thm.~3.1 argument above
is kept as the second route): $\rank_{\Zl}\mathsf M_r(a;\Zl)=\sum_\psi\rank_{\mathcal O}eH^{1}(\Ga_0(\ell^{r});\mathcal O(\chi_a\psi))=\ell^{r-1}(2d+1)$ (Shapiro over $\Qlb$) by Hida's
Thm.~1.2 ($h^{0}/P_{2,\varepsilon}\cong h^{0}_2(\Phi_r,\varepsilon)$, rank $d$ on the class-$a$ part), Lemma~6.4 (rank two over the Hecke algebra over a field) and one ordinary Eisenstein line per
character (for $\psi$ of conductor $<\ell^{r}$ the old Eisenstein space $E(\ell^{j}z)$ has $U$-ordinary part of rank one, as for cusp forms); $\mathsf M_r(a;\Zl)$ is $\Zl$-free (the applicability
paragraph of \HA{lem:nebchainrref}); $\mathsf M_r/(g-1)=eH^{1}(\Ga_0(\ell^{r});\Zl(\chi_a))$ by \HA{lem:nebshapiro}(b), of rank $2d+1$ by the \STD{} input of the Status line (ordinary $=$ old from level
$\ell$, one $U$-ordinary line per form, \HA{prop:stabilisation}'s argument) and torsion-free (\HA{lem:nebwildtors}'s argument), so $\dim_{\F_\ell}\mathsf M_r/(\ell,g-1)=2d+1$; Nakayama over the local
ring $\Zl[G_r]$ gives a surjection $\Zl[G_r]^{2d+1}\to\mathsf M_r$ of $\Zl$-free modules of equal rank, hence an isomorphism. The orientation enters only through ``class $a$ $\leftrightarrow$
nebentypus $\omega^{a-2}$'' on the forms side (Hida p.~589: $r(a)=\rank h^{0}_2(\Phi_1,\omega^{a})$), which is what (U18)(o) gates.

\begin{proposition}[The Eisenstein family on the tower]\label{prop:u-towereis}
\PROVED{} from Theorem~\ref{thm:u-towercontrol}(1),(5), Definition~\ref{def:u-eischar} and \STD{the ordinary weight-$2$ Eisenstein series $E_2(1,\chi)$ of character $\chi$: $a_p=1+\chi(p)p$
($p\nmid\ell$), $a_\ell=1$; \cite[Thm.~4.6.2]{diamond_shurman} \VERIFYloc}. Under disjointness, $\mathsf M^{\Eis}_r(a;\Z/\ell^{K})\cong\IwT Kr$ with
\[
  T_p=\ee_p(X)=1+\omega(p)^{a-1}\langle p\rangle_\ell^{-1}(1+X)^{s(p)},\qquad U=1,\qquad \langle p\rangle=\omega^{a-2}(p)\,g^{s(p)}\quad\text{in }\IwT Kr,
\]
$\ee_p$ the $\Iw$-adic Eisenstein eigenvalue of part~1 with the same $a$. At $X=\zeta^{t}u^{2}-1$ it is $1+p\,\chi_a\psi_t(p)$; at $X=x_k$ it is $1+p^{k-1}$ (Lemma~\ref{lem:u-eischar}).
\end{proposition}

\begin{proof}
$\mathsf M^{\Eis}_r$ is free of rank one over the local ring $\IwT Kr$ (Theorem~\ref{thm:u-towercontrol}(1)), so $T_p$ acts by an element $\epsilon\in\IwT Kr$; at the $\Zl$-level ($\mathsf M^{\Eis}_r(a;\Zl)$,
rank one over $\Zl[G_r]\subseteq\prod_\psi\mathcal O_\psi$, injective) its value at every $\psi$ is the $T_p$-eigenvalue of the ordinary Eisenstein line of character $\chi_a\psi$
(Theorem~\ref{thm:u-towercontrol}(5), Proposition~\ref{prop:u-towerorient}), i.e.\ $1+p\chi_a\psi(p)$, which is $\ee_p(\psi(u)u^{2}-1)$: $\omega(p)^{a-1}\langle p\rangle_\ell^{-1}(\psi(u)u^{2})^{s(p)}=
\omega(p)^{a-1}\langle p\rangle_\ell^{-1}\psi(p)\langle p\rangle_\ell^{2}=p\,\omega^{a-2}(p)\psi(p)$. So $\epsilon-\ee_p$ vanishes at every point of $\Zl[G_r]$, hence is $0$; reduce. The diamond: $\langle p\rangle=
\langle\omega(p)\rangle\langle\langle p\rangle_\ell\rangle$ and $\langle\zeta\rangle=\chi_a(\zeta)$ on the $\chi_a$-model for $\zeta\in\mathsf H_r$ (\HA{prop:nebgammaH}(4) with the extension $\chi_a$).
\COMP{} check of the identity: \file{p10c\_predict.py} \S6, $144/144$ cells $(\ell,a,p)$, $\ell\le13$, $p\le29$.
\end{proof}

\begin{proposition}[Multiplicity one along the tower, and the computed structure constants]\label{prop:u-towerfree}
\PROVED{} from Theorem~\ref{thm:u-towercontrol}(1),(3), Theorem~\ref{thm:u-lmsymfree}(1) and \LIT{Hida, Lemma~6.4} (as in Lemma~\ref{lem:u-msymsigns}). Let $\mathfrak m$ be a component of $\hcusp(a)$ with
$\hcusp(a)_{\mathfrak m}$ étale, $\cong W(\F_q)[[X]]$, satisfying (BF$'$).
\begin{enumerate}
\item $\mathsf M^{\cusp}_r(a;\Z/\ell^{K})_{\mathfrak m}$ is free of rank $2$ over $\hcusp_r(a;K)_{\mathfrak m}\cong W_K(\F_q)[X]/(\Om r)$; with $\inv^{\ast}$ (\HA{def:iota}; $\inv$ normalises $\GaT r$ and fixes
  $a_\gamma$) each sign is free of rank one.
\item If $q=\ell$ ($\hcusp(a)_{\mathfrak m}=\Iw$; e.g.\ $(11,2)$, $(13,0)$, Theorem~\ref{thm:u-ell13}): every $T_p$ acts on $\mathsf M^{\cusp}_{r,\mathfrak m}$ by the scalar $a^{(r)}_p\in\IwT Kr$, the image of
  $a_p(X):=T(p)\in\hcusp(a)_{\mathfrak m}=\Iw$, and $U$ by $\alpha^{(r)}$, the image of $\alpha(X):=U$; in particular $T_p$ lies in the image of the diamond algebra, $T_p=a^{(r)}_p(u^{2}\langle u\rangle-1)$.
  The $a^{(r)}_p$ are the \textbf{structure constants of the tower}: their Weierstrass representatives are polynomials in $X$ of degree $<\ell^{r-1}$ over $\Z/\ell^{K}$; under
  $\IwT Kr\to\IwT K{r'}$ they map to $a^{(r')}_p$.
\item The $\Iw$-adic $q$-expansion along the tower: $a_{p^{s+1}}=a_pa_{p^{s}}-p\,\omega^{a-2}(p)g^{s(p)}a_{p^{s-1}}$ and multiplicativity, in $\IwT Kr$ (the operator identity $T(p,p)=\langle p\rangle$ of
  \HA{prop:nebgammaH}(5) at $n=0$, and $p\langle p\rangle=\omega(p)^{a-1}\langle p\rangle_\ell^{-1}(1+X)^{s(p)}$ under $\iota_r$).
\end{enumerate}
\end{proposition}

\begin{proof}
(1) Directly (F-U5 m10, v0.6.0-pre.2; the base-change route through Theorem~\ref{thm:u-lmsymfree}(1) and the identification $\mathsf M^{\cusp}_{r,\mathfrak m}(\Zl)=\mathsf M_{\mathfrak m}\otimes\Iw/\Om r$,
$\mathsf M_{\mathfrak m}$ the $\mathfrak m$-part of the $\Iw$-dual of $V^{0}(a-2)$, needs Hida's duality transported by $W_{\ell^{r}}$ and is kept only as a second route): by Theorem~\ref{thm:u-towercontrol}(1),(2)
and \HA{lem:nebshapiro}(b), $\mathsf M^{\cusp}_{r,\mathfrak m}/(\ell,g-1)$ is the $\mathfrak m$-part of the cuspidal weight-$2$ cohomology of $\GaT1=\Ga_0(\ell)$ modulo $\ell$, free of rank $2$ over
$\hcusp_{2,\mathfrak m}/\ell=\F_q$ (Lemma~6.4 at level $\ell$ over $\F_q$, as in Lemma~\ref{lem:u-msymsigns}); so $\mathsf M^{\cusp}_{r,\mathfrak m}/\mathfrak m_{h}$ has $\F_q$-dimension $2$, Nakayama over the local
$\hcusp_r(a)_{\mathfrak m}\cong W(\F_q)[X]/(\Om r)$ gives a surjection $\hcusp_{r,\mathfrak m}^{2}\to\mathsf M^{\cusp}_{r,\mathfrak m}$ of $\Zl$-free modules of equal rank $2[\F_q:\F_\ell]\ell^{r-1}$, hence an isomorphism; the
signs by $\inv^{\ast}$ ($2\in\Zl^{\times}$), each of rank one by the same count at $r=1$. (2) A scalar endomorphism of a free module of
rank two over a commutative ring is determined by the ring element; $\hcusp(a)_{\mathfrak m}=\Iw$ acts through $\Iw\to\IwT Kr$. (3) The Hecke relations of \HA{thm:nebrelations}(2) on each specialisation,
transported (the $\Zl[G_r]$-level module is torsion-free, so a relation holding at every $\psi$ holds), and $p\langle p\rangle$ computed as in Proposition~\ref{prop:u-towereis}.
\end{proof}

\begin{corollary}[Truncations from one level, without interpolation]\label{cor:u-towertrunc}
\PROVED{} from Proposition~\ref{prop:u-towerfree}(2) and Lemma~\ref{lem:u-omegar}(5). In the situation of Proposition~\ref{prop:u-towerfree}(2):
\begin{enumerate}
\item For $(K,m,c)$ with $\Om r\in(\ell^{K},(X-x_c)^{m})$ the reduction of $a^{(r)}_p$ modulo $(X-x_c)^{m}$ is the truncation $a_p\bmod(\ell^{K},(X-x_c)^{m})$ of
  Proposition~\ref{prop:u-lambdaprod}(4) --- read from \emph{one} level $\ell^{r}$ at weight~$2$; no interpolation, no consecutive weights, no loss of precision. For $X$-adic truncations
  ($c=0$): $\Iw_{1,m}$ for $m\le\ell^{r-1}$, $\Iw_{2,m}$ for $m\le\ell^{r-2}$, \dots; at $r=3$: $\Iw_{1,\ell^{2}}$, $\Iw_{2,\ell}$, $\Iw_{3,1}$, of which $\Iw_{1,\ell^{2}}$ and $\Iw_{2,\ell}$ exceed the reach
  $N=K+m-1\le\ell$ of Theorem~\ref{thm:u-interp}(2).
\item In general $a^{(r)}_p$ and the interpolated truncation of part~3 have equal images in $Q:=\Iw/(\ell^{K},\Om r,(X-x_c)^{m})$ (both are images of $a_p(X)$). The lengths of $Q$ at the P7 centres
  are in \file{p10c\_predict.py} \S5.
\item For $k$ in the class, $a^{(r)}_p(x_k)\equiv a_p(f_k)\pmod{\ell^{K_r(k)}}$; and if $\Om r\in(\ell,(X-x_{k})^{2})$ ($r\ge2$, $v(k-2)=0$), the first-order coefficient of $a^{(r)}_p$ at $x_k$ modulo
  $\ell$ is that of $a_p(X)$.
\end{enumerate}
\end{corollary}

\begin{prediction}[(U18)--(U22): the tower at $\ell\le13$, $r\le3$]\label{pr:u-tower}
(New at v0.6.0-pre.1; stated before slot-22 P10 runs (rule~20).) Cells: every even class at $\ell\in\{5,7\}$ ($d=0$: the $S_k=0$ controls); at $\ell=11$ the class $a=2$ (the $\Delta$ family,
$11a$ at weight~$2$) and the control $a=4$; at $\ell=13$ the class $a=0$ (the $\Delta$ family) and the control $a=4$; $r\in\{1,2,3\}$ ($r=3$ at $\ell\in\{11,13\}$ in the env arm of
\file{P10\_spec.md}), $K=3$. Numbers: \file{P10\_spec.md} \S7, from \file{p10c\_predict.py} (output of record \file{out\_predict.txt}).
\begin{enumerate}
\item[(U18)] \PROVED{} (Theorem~\ref{thm:u-towercontrol}(1),(2), Proposition~\ref{prop:u-towereis}; (a)--(c) \emph{given} (o), which is \COMP{} at $r=1$: a failure of (o) is a defect of this text
  and stops P10, F-U5 M3) \textbf{Ranks, freeness, the tower.} (o) \emph{Orientation} (Proposition~\ref{prop:u-towerorient}): the
  ordinary cusp system lies in the $\chi_a$-part and is untwisted; the ordinary Eisenstein line carries $1+\chi_a\psi(p)p$. (a) $\rank_{\Z/\ell^{K}}\mathsf M_r(a)=\ell^{r-1}(2d+1)$: at $\ell\in\{5,7\}$ and the
  controls $\ell^{r-1}$; at $(11,2)$ $3,33,363$; at $(13,0)$ $3,39,507$ for $r=1,2,3$; free over $\IwT Kr$ (the $F_\ell$-dimension of $\mathsf M_r/(\ell,g-1)$ is $2d+1$). (b) The tower: $\mathsf M_r/(g^{\ell^{r'-1}}-1)\cong
  \mathsf M_{r'}$ and $a^{(r)}_p\mapsto a^{(r')}_p$. (c) (BF$'$) at $(11,2)$, $(13,0)$; the Eisenstein closed form of Proposition~\ref{prop:u-towereis} on $\mathsf M^{\Eis}_r$, every cell.
\item[(U19)] \PROVED{} (Theorem~\ref{thm:u-towercontrol}(3),(4), Proposition~\ref{prop:u-towerfree}, Corollary~\ref{cor:u-towertrunc}) \textbf{The structure constants against the forms side.} At $(11,2)$, $(13,0)$:
  $T_p$ acts by scalars $a^{(r)}_p$ (off-diagonal coordinates $0$ in a $\IwT Kr$-basis); $a^{(r)}_p(x_k)\equiv a_p(f_k)\pmod{\ell^{K_r(k)}}$ at $k=2$ ($11a$; the $\omega^{10}$-form of level $13$), $12$ ($\Delta$),
  $22$ ($\Delta E_{10}$), $24$ (the ordinary line of $S_{24}$) and $112$ resp.\ $158$; the first-order coefficients at $x_{12}$ modulo $\ell$ at $r\ge2$; and the images of $a^{(r)}_p$ and of P7's truncations
  $(K,m)=(3,1),(3,2),(3,3)$ centred at $x_{12}$ agree in $Q$.
\item[(U20)] \PROVED{} (Theorem~\ref{thm:u-towercontrol}(5), \HA{NB8}, \HA{NB9}) \textbf{Wild points.} At $\psi_1$ of order $\ell$ ($r=2$, $e=3(\ell-1)$) and of order $\ell^{2}$ ($r=3$, at $\ell=5$ with $e=40$): the
  character model over $\mathcal O_{j,e}$ has the lengths of \HA{NB8}, its ordinary part is free of rank $2d+1$, and $T_p$, $U$ on it are $a^{(r)}_p$, $\alpha^{(r)}$ evaluated at $g\mapsto\zeta$.
\item[(U21)] \PROVED{} from \COMP{} (\file{p10c\_predict.py} \S3: $a'_2(x_{12})\equiv2\pmod{11}$, $a'_3(x_{12})\equiv12\pmod{13}$, both non-zero) \textbf{The Hecke algebra without diamonds.} At $(11,2)$, $(13,0)$ and every
  $r,K$, the $\Z/\ell^{K}$-algebra generated by the $T_p$ ($p\le30$, $p\ne\ell$) on $\mathsf M^{\cusp}_r$ is all of $\IwT Kr$: $t:=a_{p_0}-a_{p_0}(x_{12})$ lies in $\mathfrak m_\Iw\setminus(\mathfrak m_\Iw^{2}+\ell\Iw)$, so
  $\Iw=\Zl[[t]]$ and $t$ is nilpotent in the Artinian $\IwT Kr$ --- the ordinary Hecke algebra of the $11^{r-1}$ resp.\ $13^{r-1}$ weight-$2$ forms of the family is the group ring $(\Z/\ell^{K})[C_{\ell^{r-1}}]$,
  generated by one $T_p$.
\item[(U22)] \CONJ{} (stretch, P10(c); Remark~\ref{rmk:u-towerkitagawa}) \textbf{Kitagawa's measure on the tower} equals P9's interpolated two-variable measure in the common quotient.
\end{enumerate}
\end{prediction}

\begin{remark}[What part 5 rests on, and what it does not do]\label{rmk:u-towerbasis}
\textbf{\LIT.} Hida, Invent.\ 85: Thm.~3.1 (freeness, invariants, ranks), Cor.~4.5 (raising the $\Ga_0$-level does not change the ordinary part), Thm.~1.2 with the $\omega_{k,r}$ display (control
at $k=2$ for $\Ga_1(\ell^{r})$), the $\Iw$-structure of p.~552 and its twist against the natural action, p.~593 --- all read on the page for this text --- and Lemma~6.4 (read for part~4).
\textbf{\PROVED{} here.} Lemma~\ref{lem:u-omegar}, Corollary~\ref{cor:u-towertrunc}, Proposition~\ref{prop:u-towereis} (given \STD{} for $E_2(1,\chi)$), Proposition~\ref{prop:u-towerfree}; and [C]'s new
lemmas. \textbf{\COMP{} and gated.} The orientation (Proposition~\ref{prop:u-towerorient}): \COMP{} at $r=1$ (re-derived independently by the referee F-U5 at all eight cells), gated at every cell and
every $r$ by (U18)(o); Theorem~\ref{thm:u-towercontrol}(1),(5), Proposition~\ref{prop:u-towereis}, \HA{NB8}'s rank clause and (U18)(a)--(c), (U20) are \PROVED{} \emph{given} it (F-U5 M3). \textbf{\CONJ.} (U22). \textbf{Hypotheses, each named where used}: ordinarity (only $e$-parts are
claimed); étaleness of the component (Proposition~\ref{prop:u-towerfree}; non-étale components are \emph{refused by name}, \code{tower\_nonetale}: there $\mathsf M_{\mathfrak m}$ need not be free and the
structure constants are a Fitting ideal); disjointness of the cuspidal and Eisenstein components (Theorem~\ref{thm:u-towercontrol}(1); a glued component as at $(37,32)$ is \emph{refused by name},
\code{tower\_eisenstein\_glued} --- it needs the gluing of (U17) along the tower, not stated here); multiplicity one over fields (Lemma~\ref{lem:u-msymsigns}); the freeness of the $\Iw$-adic symbols
(Theorem~\ref{thm:u-lmsymfree}). \textbf{Not used.} Any Gorenstein property beyond étaleness; Galois representations; Kitagawa (only in (U22)). \textbf{What a level is.} $\mathsf M_r(a)$ is
$\mathbf M(a)/(\Om r)$ together with the Eisenstein part: one level of the tower determines the $\Iw$-adic object modulo $\Om r$, which is \emph{not} a truncation of the form $(\ell^{K},X^{m})$ but
contains them as Corollary~\ref{cor:u-towertrunc} says.
\end{remark}

\begin{remark}[Kitagawa's measure on the tower: stretch, recalled]\label{rmk:u-towerkitagawa}
\CONJ{} (the construction is \emph{recalled}: Kitagawa (LIT-08) is not held; Greenberg--Stevens \cite{greenberg_stevens} is held image-only, read for (U16) in a weaker, local,
characteristic-$0$ form). For $\Phi\in\mathsf M^{\cusp}_r(a;\Z/\ell^{K})^{\pm}_{\mathfrak m}$ and a ball $b+\ell^{M}\Zl\subseteq\Zl^{\times}$, $M\le r$, put $\mu_\Phi(b+\ell^{M}\Zl):=$ the value of
$U^{-M}\Phi$ on the modular symbol $\{\infty\}-\{b/\ell^{M}\}$ of $\GaT r$ (the $n=0$ case of \Espec{def:lvmtt}'s functional, read on the $\GaT r$-model through the
boundary-symbol identification of \Espec{rmk:lvshapes}~(B) at level $\ell^{r}$, not yet written). \emph{Expected} (Kitagawa, recalled): a $\IwT Kr$-valued distribution, whose
specialisation at $x_k$ is the weight-$k$ measure of part~4 modulo $\ell^{K_r(k)}$ up to the normalisation of Proposition~\ref{prop:u-msymnorm}. (U22) compares it with P9's interpolated
two-variable measure (\Espec{rmk:lvtwovar}) in $Q$. Nothing else in this part depends on it.
\end{remark}

\paragraph{Code map of part 5.} \textbf{Planned} (slot-22 P10; Opus; only after F-U5's ACCEPT, rule~22; implementation spec \file{docs/implementation/slot22/P10\_spec.md}): the $\GaT r$-model on
WAKE-1's \code{GammaHLevelN} with the tame scalar (\HA{rmk:nebcosetlevel}), the chain-ring cokernel (\HA{lem:nebchainrref}) and the Heilbronn $U_\ell$, $T_p$
(\HA{prop:nebgammaHheilbronn}), the ordinary part by a certified block-Krylov iteration of $U$; \file{HeckeExt/HeckeLambdaTower} (the reading over $\IwT Kr$: the free basis, the
$a^{(r)}_p$, $\alpha^{(r)}$, their Weierstrass representatives, the comparisons with P7 (\file{HeckeExt/HeckeLambdaQExp}) and with WP-HL's \code{ComputeFormsWeight}); record
\code{hecke\_lambda\_tower/v1} (FNV per line, not pinned, no golden line); validation \code{HeckeLambdaTowerValidation} (\code{standard}; env arm \code{HECKE\_LAMBDA\_TOWER\_R3=1} for
$r=3$ at $\ell\in\{11,13\}$) gating (U18)--(U21). No statement of this part rests on the code.

%----------------------------------------------------------------------
\section{The Eisenstein component over \texorpdfstring{$\Iw$}{Lambda} at irregular pairs (part~6)}\label{sec:u-eislambda}
%----------------------------------------------------------------------
% NEW at v0.7.0-pre.1 (hecke-extensions slot 25, package U6; card docs/implementation/slot24/heckez/06-work-packages.md "U6", mathematics
% docs/implementation/slot24/heckez/05-hida-targets.md v0.3.1 Sec. 1.2, 2.2-2.4; placement docs/implementation/slot25/00-slot25-plan.md v0.2
% wave 3).  Placed after part 5 so that no section or statement number of 0.6.0 moves.  Text first, before the code of H1 (rule 22).

\paragraph{Status and scope of this part.} Written as a pre-release (v0.7.0-pre.1, \textbf{not signed}; refereed ACCEPT WITH MINORS, repaired at v0.7.0-pre.2) for an adversarial referee, before the code of the
slot-25/26 package H1 (rules~20--22). It answers, for the Eisenstein component of a class $a\not\equiv0$, the questions ``is $\Tu(a)_{\mathfrak m}$ Gorenstein, a complete
intersection, regular; is the Eisenstein ideal principal, and by which operator'' over $\Iw$, from finitely many finite-weight data: (i) \textbf{one weight decides the ring-theoretic
type} (Lemma~\ref{lem:u-H0}); (ii) at cuspidal rank $g_a=1$ the structure is forced (Proposition~\ref{prop:u-H11}); (iii) at the one known pair of cuspidal rank $2$,
$(\ell,a)=(547,486)$, the cuspidal family is a regular local ring $\Zl[[Y]]$ of degree $2$ over $\Iw$ (Proposition~\ref{prop:u-H12}), with multiplicity one at weight $486$ at a
non-étale component, outside Theorem~\ref{thm:u-lmsymfree}; (iv) the predictions (U23)--(U24). Tame level~$1$, $\ell\ge5$, as in part~1.
\textbf{Literature and novelty.} No novelty is claimed for the finite-level facts at $(547,486)$: the pair, its cuspidal rank $2$ and its uniqueness for $\ell<10^{5}$ are
\cite[p.~68]{calegari_eisdef} and \cite[\S1.1]{bellaiche_pollack}; the finite-level Gorenstein and DVR properties are known consequences of \cite{kurihara1993}
(Remark~\ref{rmk:u-H12lit}). The $\Iw$-adic statements of Proposition~\ref{prop:u-H12} were not found in the sources searched (slot-24 LIT, \file{docs/implementation/slot24/heckez/LIT-report.md}
\S3, \S3a: Calegari 2006, Bella\"\i che--Pollack 2019, Kurihara 1993, Wake 2015 and arXiv:1401.3764, Wake--Wang-Erickson 2018 and arXiv:1707.01894, Pollack--Wake arXiv:2401.14518);
that is a statement about a search, not a novelty claim. \emph{Not in this part}: the class $a\equiv0$ (Remark~\ref{rmk:u-classzero}: no Eisenstein congruence there); the
cuspidal congruence modules of the ramified families of Proposition~\ref{prop:u-loewy}(4) (slot-26 H2); tame level $N>1$.

\paragraph{Notation of this part.} $a\not\equiv0$ is an even class, $\mathfrak m$ the Eisenstein maximal ideal of $\Tu(a)$ (the one over the Eisenstein component; it always exists, as the quotient of $\Tu(a)$ by the Eisenstein component is $\Iw$, which is local;
its image in $\hcusp(a)$ is a proper ideal, i.e.\ $\mathfrak m$ lies on a cuspidal component, iff the pair is irregular, Corollary~\ref{cor:u-irregular}), and also its images in $\hcusp(a)$ and in $\Tu_k(a)$, $\hcusp_k$. $x_k:=u^{k}-1$, so $\Pk k=X-x_k$.
$I^{0}(a)\subseteq\hcusp(a)$ is the Eisenstein ideal of Lemma~\ref{lem:u-emerton}(3) and $G_a$ the series of Prediction~\ref{pr:u-U1}; $\ee_q$ is Definition~\ref{def:u-eischar}.
The \emph{cuspidal rank} is $g_a:=\rank_\Iw\hcusp(a)_{\mathfrak m}$; by Theorem~\ref{thm:u-control} it is $\rank_{\Zl}\hcusp_{k,\mathfrak m}$ at every $k\ge3$ in the class, the number
of ordinary normalised eigenforms of weight $k$ congruent to $E_k$, counted over $\Qlb$. By Prediction~\ref{pr:u-U1}, $\rank_\Iw\Tu(a)_{\mathfrak m}=1+g_a$.

\begin{lemma}[One weight decides the ring-theoretic type]\label{lem:u-H0}
\PROVED{} from Theorem~\ref{thm:u-control}, Proposition~\ref{prop:u-truncbasic}(2) and \STD{a local ring that is finite and free over the regular local ring $\Iw$ is
Gorenstein, resp.\ a complete intersection, iff its quotient by a regular sequence is: \cite[Prop.~3.1.19(b), Thm.~2.3.4(a)]{bruns_herzog}\VERIFYloc}.
Let $a\not\equiv0$, $k\ge3$ in the class, and $\mathfrak m$ a maximal ideal of $\Tu(a)$ (not necessarily the Eisenstein one). Write $R$ for $\Tu(a)_{\mathfrak m}$ and $R_k$ for
$\Tu_k(a)_{\mathfrak m}$, or $R=\hcusp(a)_{\mathfrak m}$ and $R_k=\hcusp_{k,\mathfrak m}$ when $\mathfrak m$ comes from $\hcusp(a)$.
\begin{enumerate}
\item $R/\mathfrak m_\Iw R\cong R_k/\ell R_k$.
\item $R$ is Gorenstein, resp.\ a complete intersection, iff the Artinian ring $R_k/\ell R_k$ is (Gorenstein: a socle of dimension $1$ over the residue field).
\item If $R_k/\ell R_k$ is generated, as an algebra over its residue field, by the image of one $t\in R$, and the residue field is $\Fl$, then $R=\Iw[t]\cong\Iw[Y]/(F)$ with $F$
  the characteristic polynomial of $t$ over $\Iw$, monic of degree $\rank_\Iw R$; $R$ is a complete intersection.
\end{enumerate}
\end{lemma}

\begin{proof}
(1) $\mathfrak m_\Iw=(\ell,X)=(\ell,\Pk k)$ as $x_k\in\ell\Zl$. $\Tu(a)/\Pk k\Tu(a)\cong\Tu_k(a)$ (Proposition~\ref{prop:u-truncbasic}(2)) and $\hcusp(a)/\Pk k\cong\hcusp_k$
(Theorem~\ref{thm:u-control}); a finite $\Iw$-algebra is the product of its localisations at its maximal ideals, each of which contains $\mathfrak m_\Iw\supseteq(\Pk k)$, so the
quotients commute with localisation. (2) $R$ is a direct factor of a free $\Iw$-module, hence $\Iw$-free, and $\ell$, $\Pk k$ is a regular sequence on $\Iw$, hence on $R$ (all in the
maximal ideal of $R$). The cited facts apply to $R\to R/(\ell,\Pk k)R=R_k/\ell R_k$. (3) By (1), $\Iw[t]+\mathfrak m_\Iw R=R$, so $\Iw[t]=R$ by Nakayama ($R$ is finite over $\Iw$).
Cayley--Hamilton for $t$ acting on the free module $R$ gives $F(t)=0$, so $\Iw[Y]/(F)\to R$ is a surjection of free $\Iw$-modules of the same rank $r=\deg F$, hence an isomorphism.
For the complete-intersection property, replace $t$ by $t-c$, $c\in\Zl$ a lift of the residue of $t$ (the residue field is $\Fl$; this changes neither $\Iw[t]$ nor
$\Iw[Y]/(F)$ up to $Y\mapsto Y+c$): then $t\in\mathfrak m_R$, $F$ is distinguished, and, $R$ being complete and local, $R\cong\Zl[[X,Y]]/(F)$, a regular local ring modulo one
non-zero-divisor.
\end{proof}

\begin{remark}[The census reading]\label{rmk:u-H0census}
\PROVED{} from Lemma~\ref{lem:u-H0}, given the certificate named. The census column \code{gorenstein} at an Eisenstein row $(n,\ell)$, $k=n+2$, is a statement about the
$\Iw$-adic $\Tu(a)_{\mathfrak m}$ whenever the row's algebra $A$ is certified to be $\Tu_k(a)_{\mathfrak m}/\ell$. The certificate is generation by saturation (the slot-24 HeckeZ
plan's Prop.~2.2, [C] 1.25.0-pre.3 \HA{prop:heckeZgen}(1), a pre-release): $\dim_{\Fl}A=1+g_a$, with
$1+g_a$ known a priori, \textbf{provided $k\le\ell+1$ or $T_\ell$ is among the operators of the row, and the row's block has dimension $1+2g_a$} (the row's \code{module\_D}), so
that it is the $\mathfrak m$-part. A census block is cut out by finitely many operators, so it may also contain a form that is Eisenstein only at those; then the lattice is larger than
$H_{\mathfrak m}$, the a-priori rank of \HA{prop:heckeZgen}(1) is not $1+g_a$, and $\dim A=1+g_a$ certifies nothing (at $(484,547)$, (D4) has \code{module\_D=5}). Without $T_\ell$ the block cut out by the $T_p$, $p\ne\ell$, at the residual
Eisenstein system contains every eigenform $f$ with $\bar\rho_f^{\mathrm{ss}}\cong\mathbf1\oplus\omega^{k-1}$, ordinary or not, while $g_a$ counts the ordinary ones. For
$2\le k\le\ell+1$ a non-ordinary level-one eigenform has $\bar\rho_f|_{I_\ell}\cong\omega_2^{k-1}\oplus\omega_2^{\ell(k-1)}$, irreducible
(\STD{Fontaine's theorem, as stated in \cite[Thm.~2.6]{edixhoven}}\VERIFYloc), so the two blocks coincide; this covers $(484,547)$ ($k=486<548$) and every irregular-pair row at
$k=a<\ell$. For $k>\ell+1$ a non-ordinary congruent form is possible and the row certifies nothing about $\Tu(a)_{\mathfrak m}$ unless $T_\ell$ is read ($T_\ell\equiv1+\ell^{k-1}\equiv1$
is then a unit on the block, which makes it ordinary).
\end{remark}

\begin{proposition}[Cuspidal rank one: the structure is forced]\label{prop:u-H11}
\PROVED{} from Prediction~\ref{pr:u-U1} (\PROVED{} from \LIT{Emerton}), Lemma~\ref{lem:u-fibre} and Theorem~\ref{thm:u-control}. Let $(\ell,a)$ be irregular with $g_a=1$. Then
$\hcusp(a)_{\mathfrak m}=\Iw$, $I^{0}(a)_{\mathfrak m}=G_a\Iw$, and
\[
  \Tu(a)_{\mathfrak m}=\Iw\times_{\Iw/(G_a)}\Iw\;\cong\;\Iw[t]/(t^{2}-G_at),\qquad t=(G_a,0),
\]
a complete intersection, hence Gorenstein. For a prime $q\ne\ell$: $I^{0}(a)_{\mathfrak m}=(T(q)-\ee_q)$ iff $v\bigl(a_q(f_k)-(1+q^{k-1})\bigr)=v(B_k/k)$ for \emph{one} $k\ge3$ in
the class, $f_k$ the ordinary eigenform of weight $k$ congruent to $E_k$ (and then for every such $k$).
\end{proposition}

\begin{proof}
A local $\Iw$-algebra, free of rank $1$, is $\Iw$. Localised at $\mathfrak m$, (U1-$\Iw$) gives $\hcusp(a)_{\mathfrak m}/I^{0}(a)_{\mathfrak m}\cong\Iw/(G_a)$, so $I^{0}_{\mathfrak m}=(G_a)$,
and Lemma~\ref{lem:u-fibre} with $J=G_a\Iw$ gives the fibre product ($\Tu(a)\subseteq\hcusp(a)\times\Iw$, localised at $\mathfrak m$; the factor $\Iw$ is local and lies over
$\mathfrak m$). $\{(x,y):x\equiv y\ (G_a)\}=\Iw(1,1)\oplus\Iw t$ and $t^{2}=G_at$; $\Iw[t]/(t^{2}-G_at)$ is a regular ring modulo a monic polynomial. For the generator,
$T(q)-\ee_q\in I^{0}_{\mathfrak m}$, so $T(q)-\ee_q=wG_a$ in $\hcusp(a)_{\mathfrak m}=\Iw$; it generates iff $w\in\Iw^\times$, iff $w(x_k)\in\Zl^\times$ for one (every) $k$, as
$w(x_k)\equiv w(0)\pmod\ell$. At $X=x_k$, $T(q)-\ee_q\mapsto a_q(f_k)-(1+q^{k-1})$ (Theorem~\ref{thm:u-control}, Lemma~\ref{lem:u-eischar}) and
$v(G_a(x_k))=v(B_k/k)$ ($1-\ell^{k-1}$ is a unit).
\end{proof}

\noindent
\emph{Where $g_a=1$} (\COMP, uncertified forms-side scripts: slot-17 R17-1 \file{gw\_scan.py} and its independent re-derivation in the slot-24 HeckeZ audit P2a). At the weight
$k=a<\ell+1$ every residually Eisenstein form is ordinary (Remark~\ref{rmk:u-H0census}), so $g_a$ is at most the multiplicity of the root $1+2^{a-1}$ of $\chi(T_2\mid S_a)\bmod\ell$,
and $g_a\ge1$ at an irregular pair (Corollary~\ref{cor:u-irregular}); a multiplicity $1$ therefore proves $g_a=1$. Of the $155$ irregular pairs with $\ell<2000$, every one has
multiplicity $1$ except $(547,486)$ (multiplicity $2$). \LIT{\cite[\S1.1, arXiv v1, p.~2]{bellaiche_pollack} (page image verified, slot-24 LIT \S3)}: cuspidal rank one holds at every
irregular pair with $p<10^{5}$ except $(547,486)$, where the cuspidal algebra has $\Zl$-rank $2$; the scan is a check of the project's code against that published result, and
Proposition~\ref{prop:u-H11} applies at every other pair below $10^{5}$ (their $\mathbb T^{c}_{\mathfrak m^{c}}$ read as $\hcusp(a)_{\mathfrak m}$ from their notation, not
re-derived). Example: at $(n,\ell)=(10,691)$, $k=12$, $\tau(2)-(1+2^{11})=-2073=-3\cdot691$ and $v_{691}(B_{12}/12)=1$, so $I^{0}(12)_{\mathfrak m}=(T(2)-\ee_2)$, which re-derives the
Eisenstein block of Theorem~\ref{thm:u-691classes}. The generator test at the other pairs is one integer per pair (slot-25/26 H1, on the C2 chains of the HeckeZ plan).

\begin{proposition}[Cuspidal rank two at $(547,486)$]\label{prop:u-H12}
\PROVED{} given the data (D1)--(D4) below (\COMP), from Lemma~\ref{lem:u-H0}, Prediction~\ref{pr:u-U1}, Lemma~\ref{lem:u-fibre}, Theorem~\ref{thm:u-control},
Corollary~\ref{cor:u-irregular}, Lemma~\ref{lem:u-eischar}, Lemma~\ref{lem:u-emerton}(3), Lemma~\ref{lem:u-msymsigns}, Remark~\ref{rmk:u-H0census} with [M1] Prop.~1(2),
\HA{rmk:ordtorsionfree}, and \STD{Weierstrass preparation in $\Zl[[X,Y]]$; Krull--Schmidt}. Let $\ell=547$, $a=486$, $\mathfrak m$ the Eisenstein maximal ideal, and $Y:=T(2)-\ee_2\in\hcusp(a)_{\mathfrak m}$.
\begin{enumerate}
\item $\hcusp(a)_{\mathfrak m}=\Iw[Y]\cong\Iw[Y]/(F)$, $F=Y^{2}+bY+c$, with $b\in\mathfrak m_\Iw$, $c=w\,G_a$, $w\in\Iw^\times$, and $G_a=w_1(X-\xi)$ with $w_1\in\Iw^\times$,
  $\xi\in\ell\Zl$ ($\lambda(G_a)=1$, $\mu(G_a)=0$). Hence $\hcusp(a)_{\mathfrak m}\cong\Zl[[Y]]$ is a \textbf{regular local domain of dimension $2$}: a single cuspidal Hida family
  of $\Iw$-degree $g_a=2$, on which $X-\xi$ has Weierstrass degree $2$ in $Y$.
\item $I^{0}(a)_{\mathfrak m}=(Y)$: the Eisenstein ideal is principal, generated by $T(2)-\ee_2$, and $\hcusp(a)_{\mathfrak m}/(Y)\cong\Iw/(G_a)\cong\Zl$.
\item $\Tu(a)_{\mathfrak m}=\Iw[Y']\cong\Iw[Y]/\bigl(Y\,F(Y)\bigr)$, $Y'=(Y,0)=(T(2),\ee_2)-\ee_2$: a complete intersection, hence Gorenstein, of $\Iw$-rank $3$, with special fibre
  $\Fl[y]/(y^{3})$.
\item At every $k\ge3$ in the class, $\hcusp_{k,\mathfrak m}=\Zl[T_2]\cong\Zl[x]/(f_k)$ with $f_k(x+1+2^{k-1})=x^{2}+b(x_k)x+c(x_k)$, and
  $v\bigl(c(x_k)\bigr)=v(B_k/k)$. Where $v(B_k/k)=1$, $v(\operatorname{disc}f_k)=1$: $\hcusp_{k,\mathfrak m}$ is the ring of integers of a \textbf{ramified} quadratic extension of
  $\Ql$, and the two ordinary eigenforms congruent to $E_k$ are Galois-conjugate.
\item \textbf{(The rung.)} Where $v(B_k/k)\ge2$ --- in particular at $k^{*}=55086$, where $v_{547}(B_{k^{*}}/k^{*})=2$ (\COMP, below) --- the order $\hcusp_{k,\mathfrak m}$ is
  \textbf{not maximal}: its index in the maximal order of $\hcusp_{k,\mathfrak m}\otimes\Q$ is $\ge547$. Whether the two forms there are conjugate over a ramified or an unramified
  extension, or rational, depends on $b(x_k)$, which (D1)--(D4) do not determine.
\item \textbf{(Multiplicity one at weight $486$.)} Let $H:=\eord H^{1}(\Ga;\Sym^{484}\Z^{2})\otimes\Zl$ localised at $\mathfrak m$, $H=H^{\varepsilon}\oplus H^{-\varepsilon}$ by the
  involution, $\varepsilon$ the sign that carries the Eisenstein class ($\varepsilon=+$ in [C]'s convention: Lemma~\ref{lem:u-msymsigns}, \HA{rmk:iotasign}). Then $H^{\varepsilon}\cong\Tu_{486}(a)_{\mathfrak m}$ and $H^{-\varepsilon}\cong\hcusp_{486,\mathfrak m}$, free of
  rank one, with $H^{\varepsilon}/\ell\cong\Fl[y]/(y^{3})$ and $H^{-\varepsilon}/\ell\cong\Fl[y]/(y^{2})$; and $\delta(\mathfrak m)=0$ on this lattice. This is multiplicity one at a
  \emph{non-étale} component, outside the hypothesis of Theorem~\ref{thm:u-lmsymfree}, at the one weight $486$; the $\Iw$-adic statement is (U24) below.
\end{enumerate}
\emph{The data} (slot-24 HeckeZ stage 05 \S2.4; scratch data outside the tree):
\textbf{(D1)} $T_2$ on $S_k(\SL_2(\Z))$ modulo $547^{4}$ at $k=486,1032,1578,2124$: the root $1+2^{k-1}$ has multiplicity $2$ modulo $547$, and its Hensel factor $f_k$ has
$v(\operatorname{disc}f_k)=1$ and $v(f_k(1+2^{k-1}))=1$; at $k=486$, $f=x^{2}+18353242915\,x+40544750686\bmod547^{4}$ (\file{blk547.py}, pure Python, uncertified; controls
$g=1$ at six pairs and $\tau(2)$ at $(691,12)$).
\textbf{(D2)} On cohomology, \code{HeckeHidaComponentsValidation} with \code{HECKE\_HIDA\_COMPONENTS\_CELLS="484:547:3"}: the Eisenstein block has dimension $5=1+2g_a$, $L_{\mathrm{coh}}=1$.
\textbf{(D3)} $v_{547}(B_k/k)=1$ at $k=486,1032$, and $v\bigl(G_a(x_{486})-G_a(x_{1032})\bigr)=1$ (\file{bern547.py}, power sums modulo $\ell^{4}$; controls: the rungs $284$ at $37$,
Remark~\ref{rmk:u-zero}, and $342252$ at $691$, Proposition~\ref{prop:u-691}); and, evaluated directly at precision $547^{3}$ by the slot-24 audit P2a,
$v_{547}(B_{55086}/55086)=2$, $=1$ at $55086\pm546$.
\textbf{(D4)} The census-4.0 row $(484,547)$, \code{eis(0,485)}: \code{length\_A=3}, \code{hilbert=[1,1,1]}, \code{embdim=1}, \code{gorenstein=true}, \code{module\_D=5},
\code{module\_loewy=[2,2,1]}, \code{min\_generators=2}.
(D1) and (D3) are uncertified scripts; (D2) is a shipped validation and (D4) a signed census row. The tag becomes ``\PROVED{} given certified data'' when H1 reproduces
(D1), (D3) in C++ and the Lean checker K9 of the HeckeZ plan checks them.
\end{proposition}

\begin{proof}
\emph{Step 1 (weight $486$).} By (D1) the root $\bar e$, $e:=1+2^{485}$, of $\chi(T_2\mid S_{486})$ modulo $\ell$ has multiplicity $2$; let $E$ be the $T_2$-generalised eigenspace
of $\bar e$ in $S_{486}\otimes\Zl$, of rank $2$, and $f:=f_{486}$ the characteristic polynomial of $T_2$ on it (its Hensel factor). $f(x+e)=x^{2}+\beta x+\gamma$ with $\beta\equiv\gamma\equiv0$
($e$ is a double root modulo $\ell$) and $v(\gamma)=1$: $f$ is Eisenstein at $e$, so irreducible over $\Ql$, and $\Zl[x]/(f)$ is a DVR, the maximal order of $\Ql[x]/(f)$.
Every eigenform congruent to $E_{486}$ has $a_2\equiv e$, so the block $\hcusp_{486,\mathfrak m}$ lies in $E$. One such eigenform $f_1$ exists ($g_a\ge1$, Corollary~\ref{cor:u-irregular}); the
other eigenform in $E$ is its $\Ql$-conjugate (the two roots of the irreducible $f$ are conjugate), which is ordinary and congruent to $E_{486}$ at every prime too (the conjugation fixes
the unique prime above $\ell$ of the ramified quadratic field). Hence $g_a=2$, $E$ is the block, and $\Zl[T_2]\subseteq\hcusp_{486,\mathfrak m}$ forces equality. (Neither (D2) nor
Remark~\ref{rmk:u-H0census} is used here; (D2), the cohomological block of dimension $5=1+2g_a$, is a check.) So $\hcusp_{486,\mathfrak m}/\ell=\Fl[\bar y]/(\bar y^{2})$, $\bar y$ the image of $T_2-e$.
\emph{Step 2 ((1), (2)).} $\ee_2(x_{486})=e$ (Lemma~\ref{lem:u-eischar}), so $Y$ maps to $T_2-e$, and Lemma~\ref{lem:u-H0}(1), (3) give $\hcusp(a)_{\mathfrak m}=\Iw[Y]/(F)$,
$F=Y^{2}+bY+c$; $F\equiv Y^{2}$ modulo $\mathfrak m_\Iw$ gives $b,c\in\mathfrak m_\Iw$. $Y\in I^{0}(a)$ (Lemma~\ref{lem:u-emerton}(3)), so
$\Iw/(c)=\hcusp(a)_{\mathfrak m}/(Y)\twoheadrightarrow\hcusp(a)_{\mathfrak m}/I^{0}(a)_{\mathfrak m}\cong\Iw/(G_a)$ ((U1-$\Iw$), localised), a map of $\Iw$-algebras; hence $c\in(G_a)$,
$c=wG_a$. At $X=x_{486}$: $c(x_{486})=f(e)$ has valuation $1$ (D1) and $v(G_a(x_{486}))=v(B_{486}/486)=1$ (D3), so $w(x_{486})$, hence $w$, is a unit; then $\Iw/(c)\to\Iw/(G_a)$ is a
surjection between isomorphic Noetherian rings, so injective, and $I^{0}_{\mathfrak m}=(Y)$. By (D3), $G_a(0)\equiv G_a(x_{486})\equiv0\pmod\ell$ and the difference quotient
$\bigl(G_a(x_{486})-G_a(x_{1032})\bigr)/(x_{486}-x_{1032})$ is a unit ($v(x_{486}-x_{1032})=v(u^{546}-1)=1$), and it is $\equiv$ the $X$-coefficient of $G_a$ modulo $\ell$; so
$G_a=w_1(X-\xi)$ by Weierstrass preparation, $v(\xi)\ge1$. Then $c=w_2(X-\xi)$, $w_2:=ww_1\in\Iw^\times$. $F$ is distinguished in $Y$, so $\Iw[Y]/(F)=\Zl[[X,Y]]/(F)$; modulo $(\ell,Y)$,
$F\equiv\bar w_2(X)\,X$ (as $\xi\equiv0$), so $F$ is distinguished of degree $1$ in $X$, $F=(\text{unit})(X-\varphi(Y))$, and $\Zl[[X,Y]]/(F)\cong\Zl[[Y]]$. In it
$X-\xi=-w_2^{-1}(Y^{2}+bY)$ with $b\equiv0$ modulo $(\ell,Y)$, so $X-\xi$ has Weierstrass degree $2$ in $Y$.
\emph{Step 3 ((3)).} $\Tu(a)_{\mathfrak m}=\hcusp(a)_{\mathfrak m}\times_{\Iw/(G_a)}\Iw$ (Lemma~\ref{lem:u-fibre} with (U1-$\Iw$), localised as in Proposition~\ref{prop:u-H11}), the gluing map
sending $Y\mapsto0$. $Y'\in\Tu(a)$, and for $h=h_0+Yh_1(Y)\in\hcusp(a)_{\mathfrak m}=\Iw\oplus\Iw Y$ and $(h,\lambda)\in\Tu(a)_{\mathfrak m}$: $\lambda-h_0\in G_a\Iw$, and
$(h,\lambda)=h_0+Y'h_1(Y')+(\lambda-h_0)\,G_a^{-1}w^{-1}F(Y')$, because $F(Y')=(F(Y),F(0))=(0,c)=(0,wG_a)$. So $\Tu(a)_{\mathfrak m}=\Iw[Y']$. $Y'F(Y')=(YF(Y),0)=0$, and
$\Iw[Y]/(YF)\to\Tu(a)_{\mathfrak m}$ is a surjection of free $\Iw$-modules of rank $3$ (Prediction~\ref{pr:u-U1}: rank $1+g_a$), hence an isomorphism; $YF\equiv Y^{3}$ modulo
$\mathfrak m_\Iw$. (The fibre $\Fl[y]/(y^{3})$ is the census row (D4), \code{hilbert=[1,1,1]}: a consistency check, not an input.)
\emph{Step 4 ((4), (5)).} Theorem~\ref{thm:u-control} and Lemma~\ref{lem:u-H0}(1): $\hcusp_{k,\mathfrak m}=\hcusp(a)_{\mathfrak m}/\Pk k=\Zl[Y]/(F(x_k;Y))$, $Y\mapsto T_2-(1+2^{k-1})$.
$v(c(x_k))=v(G_a(x_k))=v(B_k/k)$. $\operatorname{disc}=b(x_k)^{2}-4c(x_k)$ with $v(b(x_k)^{2})\ge2$: at $v(B_k/k)=1$ it has valuation $1$, $F(x_k;\cdot)$ is Eisenstein, and the ring is a
ramified DVR, of rank $2$ with the two roots conjugate. At $v(B_k/k)\ge2$ both roots $y_1,y_2$ of $F(x_k;\cdot)$ have $v\ge1$ (Newton polygon: $v(b(x_k))\ge1$, $v(c(x_k))\ge2$), and
$F(x_k;\cdot)$ is squarefree ($\hcusp_k$ is reduced). If it is irreducible, $y_1/\ell$ is integral and $\Zl[y_1]=\Zl\oplus\Zl y_1\not\ni y_1/\ell$; if it splits,
$\Zl[Y]/(F)=\{(p(y_1),p(y_2))\}$ has index $\ell^{v(y_1-y_2)}\ge\ell$ in $\Zl^{2}$. Either way the index is $\ge\ell$.
\emph{Step 5 ((6)).} $H$ is $\Zl$-free (\HA{rmk:ordtorsionfree}), and $\Tu_{486}(a)_{\mathfrak m}$ acts on it faithfully. By (D4), $\dim_{\Fl}A=3=1+g_a=\rank\Tu_{486}(a)_{\mathfrak m}$
for the image $A$ of $\Tu_{486}(a)_{\mathfrak m}$ in $\End(H/\ell)$, so $A=\Tu_{486}(a)_{\mathfrak m}/\ell\cong\Fl[y]/(y^{3})$ and $\delta(\mathfrak m)=0$ (generation by saturation,
Remark~\ref{rmk:u-H0census}; [M1] Prop.~1(2)). $H/\ell$ is an $\Fl[y]/(y^{3})$-module of dimension $5$, with $2$ generators and Loewy layers $(2,2,1)$, so $H/\ell\cong\Fl[y]/(y^{3})\oplus
\Fl[y]/(y^{2})$. The involution commutes with the Hecke operators ($\ell$ odd), and $H^{\pm}\otimes\Q$ has ranks $1+g_a=3$ and $g_a=2$ (each cuspidal system occurs once in each sign,
the Eisenstein system once, in the sign $\varepsilon$; Lemma~\ref{lem:u-msymsigns}). By Krull--Schmidt, $H^{\varepsilon}/\ell\cong\Fl[y]/(y^{3})$ and $H^{-\varepsilon}/\ell\cong\Fl[y]/(y^{2})$. Both
are cyclic, so $H^{\pm}$ are cyclic by Nakayama; $H^{\varepsilon}$ is faithful over $\Tu_{486}(a)_{\mathfrak m}$ (every system of the block occurs in it) and $H^{-\varepsilon}$ over
$\hcusp_{486,\mathfrak m}$, through which $\Tu_{486}(a)_{\mathfrak m}$ acts on it. A cyclic faithful module is free of rank one.
\end{proof}

\begin{remark}[The literature at $(547,486)$: known consequences, not results of {[U]}]\label{rmk:u-H12lit}
(The sentence in the \LIT{} tag is the owner-signed P-K2 of 2026-11-05; the rest of this remark is new at 0.7.0-pre.1.) Let $T^{E}$ be Kurihara's algebra: the finite-level cuspidal Hecke algebra of weight $k$ (all $T_l$, $T_p$ included) completed at the
Eisenstein maximal ideal, \cite[pp.~282--283]{kurihara1993}; at $(547,486)$ it is $\hcusp_{486,\mathfrak m}$ (Lemma~\ref{lem:u-levelone}). The pair is in Kurihara's range ($0<486<546$),
and \LIT{$T^{E}$ is Gorenstein by \cite[Rmk.~4.5(3)]{kurihara1993}, as $g=2$; a DVR by \cite[Thm.~0.4]{kurihara1993} given Vandiver at $547$ (page images verified, slot-24 LIT \S3a,
[GT] \texttt{lit:hz-kurihara-gor})}. In detail: Rmk.~4.5(3), ``at most two liftings of $G_k$ mod $p$'', makes $T^{E}$ monogenic over $\Zl$, hence Gorenstein; $n=v(B_{486}/486)=1$ gives
$\mathfrak IT^{E}=\mathfrak m$ (his Lemma~3.1); and under $A^{\omega^{62}}=0$ (Vandiver at $547$, an input from outside that paper, not read here) his Thm.~0.4(ii) makes $\mathfrak m$
principal. This is the case $k=486$ of Proposition~\ref{prop:u-H12}(4), which Step~1 of its proof obtains from (D1) without Vandiver. Kurihara does not name $547$; the pair is named first by
\cite[p.~68]{calegari_eisdef} (``$\mathbf T_{\mathfrak I}/p\simeq\mathbf F_p[x]/x^{2}$'' at one finite level), and its cuspidal rank $2$ and uniqueness for $p<10^{5}$ are
\cite[\S1.1, \S3.9]{bellaiche_pollack}. Kurihara proves unconditionally only the implication ``$T^{E}$ Gorenstein $\Rightarrow$ the $\omega^{1-k}$-part of the class group is cyclic'' (Thm.~0.3) and gives
\emph{no} non-Gorenstein example. None of these sources treats the $\Iw$-adic ring $\hcusp(a)_{\mathfrak m}$, the principality of $I^{0}(a)$ over $\Iw$, or the rung.
\end{remark}

\begin{prediction}[(U23), (U24) The $547$ family]\label{pr:u-H1}
\begin{enumerate}
\item[(U23)] \textbf{(The rung order.)} Stated here as a test of the implementation, before the computation (slot-25/26 H1-R, the forms side on all of $S_{k^{*}}$,
  $\dim S_{k^{*}}=4590$, modulo $547^{3}$, with $T_2$, $T_3$ only). At $k^{*}=55086$:
  \PROVED{} given (D1)--(D4), from Proposition~\ref{prop:u-H12}(4), (5): the characteristic polynomial $f_{k^{*}}$ of $T_2$ on the ordinary block $\hcusp_{k^{*},\mathfrak m}$ has degree
  $2$ and $v(\operatorname{disc}f_{k^{*}})\ge2$, $\Zl[T_2|_{\mathrm{block}}]$ has index $\ge547$ in its normalisation, and $f_{k^{*}}$ divides the Hensel factor of
  $\chi(T_2\mid S_{k^{*}})$ at $1+2^{k^{*}-1}$.
  \COMP{} (\file{docs/implementation/slot25/U6\_referee.md} \S1, \file{rung.py}, an uncertified referee script, modulo $547$; H1-R reproduces it): that Hensel factor has degree
  exactly $2$, so it \emph{is} $f_{k^{*}}$ (at $k^{*}>\ell+1$ a non-ordinary form with $a_2\equiv1+2^{k^{*}-1}$ is not excluded a priori, Remark~\ref{rmk:u-H0census}; none occurs).
  H1 reads $b$, $c$ modulo $(\ell^{3},X^{3})$, so it predicts $f_{k^{*}}$ modulo $547^{3}$; the sharp form of the test is that the predicted polynomial equals the Hensel factor.
  \textbf{The anchor} (v0.7.0-pre.3, written before H1-R runs, rule~22; \COMP: slot-25 H1, \code{HECKE\_LAMBDA\_CELLS="547:486:3:3"}, \file{docs/implementation/slot25/H1\_report.md}
  \S2, reproduced independently by Newton interpolation in the audit \file{H1\_audit.md} \S1):
  \[
    f_{k^{*}}\equiv x^{2}+56702273\,x+83587893\pmod{547^{3}},\qquad v_{547}(\operatorname{disc}f_{k^{*}})=2,\quad v_{547}\bigl(f_{k^{*}}(1+2^{k^{*}-1})\bigr)=2 .
  \]
  H1-R's Hensel factor of $\chi(T_2\mid S_{k^{*}})$ at $1+2^{k^{*}-1}$, reduced modulo $547^{3}$, must be this polynomial.
  A Hensel factor of degree $\ne2$, or $v(\operatorname{disc})\le1$, at $k^{*}$ is a defect of the implementation, of (D1)--(D4) or of that computation.
\item[(U24)] \textbf{($\Iw$-adic multiplicity one at a non-étale component.)} \CONJ{} At $(547,486)$, $\mathbf M(a)^{\pm}_{\mathfrak m}$ (Definition~\ref{def:u-lmsym}, the
  ordinary parabolic $\Iw$-adic modular symbols of part~4) is free of rank one over $\hcusp(a)_{\mathfrak m}$, for both signs. Expected from
  Proposition~\ref{prop:u-H12}(6) by Nakayama and control (Proposition~\ref{prop:u-lmsymcontrol}, \HA{thm:hidacontrol}(1), (2), Lemma~\ref{lem:u-latticeord}, Lemma~\ref{lem:u-levelone});
  the chain of identifications is not written out, so it is not tagged \PROVED. ($\Tu(a)_{\mathfrak m}$ acts on $\mathbf M(a)^{\pm}_{\mathfrak m}$ through $\hcusp(a)_{\mathfrak m}$,
  as the specialisations are cuspidal; a $\Iw$-adic analogue of $H^{\varepsilon}\cong\Tu_{486}(a)_{\mathfrak m}$ would need a $\Iw$-adic full $H^{1}$, which [U] does not define.)
\end{enumerate}
\end{prediction}

\begin{remark}[Comparison: Mazur's Eisenstein ideal at prime level (slot-25 H3)]\label{rmk:u-H3prime}
\COMP{} (\file{docs/implementation/slot25/H3\_report.md}, \file{H3\_cells.tsv}; refereed ACCEPT\_WITH\_MINORS with independent spot checks, \file{H3\_referee.md}). Outside the scope of
[U] (prime level $N$, weight $2$; [U] cites no statement of it), recorded because it is the finite-level model of Proposition~\ref{prop:u-H12}(4)--(5). On the $180$ cells $(N,\ell)$,
$\ell\mid\nu=\mathrm{num}((N-1)/12)$, of the H3 list, the cuspidal Eisenstein-local ring $\mathbb T^{0}_{\mathfrak m}=\Zl[\eta_q]\cong\Zl[y]/(F)$ ($q$ a good prime that passes the
valuation check $v_\ell F(0)=v_\ell(\nu)$) has rank $e=1,2,3$ at $148$, $29$, $3$ cells, $e\ge2$ exactly at the Merel cells, and $\deg F=(\dim W-1)/2$ on every cell. Where
$\ell\Vert\nu$ the ring is a totally ramified DVR (the analogue of (4)); where $\ell^{2}\mid\nu$ and $e=2$ it is \textbf{not maximal}, of two of the three kinds that (5) leaves open (ramified, unramified, rational; no ramified non-maximal cell occurs):
$(401,5)$ and $(701,5)$ give the index-$5$ order $\Z_5+5\Z_{25}$ (unramified splitting field), $(751,5)$ the split fibre product $\Z_5\times_{\F_5}\Z_5$. So the trichotomy of (5) is not
vacuous in two of its branches. The agreement of $\deg F$ with $\dim W$ is evidence for [C]'s prediction \HA{LN4} (``equality is part of the prediction'') on these cells, not a proof: \HA{LN4} stays a
prediction, and [C]'s \HA{LN6} reading is unchanged.
\end{remark}

\begin{remark}[Steps 1--2 of Proposition~\ref{prop:u-H12} at any cuspidal rank $g_a\ge2$]\label{rmk:u-H12general}
\PROVED{} from Lemma~\ref{lem:u-H0}(1), (3), Prediction~\ref{pr:u-U1}, Lemma~\ref{lem:u-emerton}(3), Theorem~\ref{thm:u-control}, Lemma~\ref{lem:u-eischar} and
\STD{Weierstrass preparation in $\Zl[[X,Y]]$; $q$-expansion duality $S_k(\Zl)_{\mathfrak m}\cong\Hom_{\Zl}(\hcusp_{k,\mathfrak m},\Zl)$}. Let $(\ell,a)$ be irregular, $\mathfrak m$ the
Eisenstein maximal ideal, $k\ge3$ in the class, $g:=g_a\ge2$, $q\ne\ell$ a prime, $Y:=T(q)-\ee_q\in\hcusp(a)_{\mathfrak m}$, and $B$ the $\mathfrak m$-part of $\eord S_k(\Zl)$ (rank $g$),
cut out by the full residual system, $T_\ell$ included. Suppose (a) the image $\bar y$ of $T_q-(1+q^{k-1})$ in $\End(B/\ell)$ satisfies $\dim_{\Fl}\Fl[\bar y]=g$; (b)
$v\bigl(\det(T_q-(1+q^{k-1})\mid B)\bigr)=v(B_k/k)$; (c) $v\bigl(G_a(x_k)-G_a(x_{k+\ell-1})\bigr)=1$. Then the conclusions (1), (2) of Proposition~\ref{prop:u-H12} hold with
$T(2)$ replaced by $T(q)$ and $2$ by $g$: $\hcusp(a)_{\mathfrak m}=\Iw[Y]\cong\Iw[Y]/(F)\cong\Zl[[Y]]$, $F$ monic of degree $g$, a regular local domain of dimension $2$,
$I^{0}(a)_{\mathfrak m}=(Y)$, and $X-\xi$ has Weierstrass degree $g$ in $Y$.
\emph{Proof.} By duality $B$ is a faithful $\hcusp_{k,\mathfrak m}$-module of rank $g=\rank_{\Zl}\hcusp_{k,\mathfrak m}$, and $B/\ell\cong\Hom(\hcusp_{k,\mathfrak m}/\ell,\Fl)$ is faithful over
$\hcusp_{k,\mathfrak m}/\ell$, so (a) says that $\bar y$ generates $\hcusp_{k,\mathfrak m}/\ell$ (dimension $g$); the residue field is $\Fl$. This replaces Step~1. $Y$ maps to
$T_q-(1+q^{k-1})$ at $x_k$ (Lemma~\ref{lem:u-eischar}), so Lemma~\ref{lem:u-H0}(1), (3) give $\hcusp(a)_{\mathfrak m}=\Iw[Y]\cong\Iw[Y]/(F)$, $F=Y^{g}+b_{g-1}Y^{g-1}+\dots+b_1Y+c$,
and $F\equiv Y^{g}$ modulo $\mathfrak m_\Iw$ ($\bar y$ is nilpotent in the local Artinian ring), so $b_i,c\in\mathfrak m_\Iw$. Step~2 then runs verbatim: $Y\in I^{0}(a)$
(Lemma~\ref{lem:u-emerton}(3)) gives $\Iw/(c)\twoheadrightarrow\Iw/(G_a)$, so $c=wG_a$; $F(x_k;\cdot)$ is the characteristic polynomial of $T_q-(1+q^{k-1})$ on $\hcusp_{k,\mathfrak m}$
(Theorem~\ref{thm:u-control}), equal to that on $B$ (duality), so $v(c(x_k))=v(\det)=v(B_k/k)=v(G_a(x_k))$ by (b), $w$ is a unit and $I^{0}_{\mathfrak m}=(Y)$; (c) and
$v(x_k-x_{k+\ell-1})=1$ make the $X$-coefficient of $G_a$ a unit, $G_a=w_1(X-\xi)$, $c=w_2(X-\xi)$, and modulo $(\ell,Y)$, $F\equiv\bar w_2(X)X$, so
$\Zl[[X,Y]]/(F)\cong\Zl[[Y]]$, in which $X-\xi=-w_2^{-1}(Y^{g}+\dots+b_1Y)$ has Weierstrass degree $g$ in $Y$. \qed

Steps~3--5 (the full algebra, the specialisations, multiplicity one) use data of $(547,486)$ and are not claimed here. At $g=1$ the hypotheses (a) and (c) are not needed and the
statement is Proposition~\ref{prop:u-H11}'s generator test. Where (a), (b) or (c) fails no conclusion is drawn. This is the rule of the record's G8 word \code{regular\_deg\_g}
(\HA{sec:recordschema} item~8, the G8 bullet: (a) and (b) for the first $q$ that passes, (c) by the $\lambda$ test; otherwise \code{open}). Below $\ell<10^{5}$ the only pair with
$g_a\ge2$ is $(547,486)$ \cite[\S1.1]{bellaiche_pollack}, where the remark is Proposition~\ref{prop:u-H12}(1), (2).
\end{remark}

\paragraph{Code map of part 6.} \textbf{Built} (v0.7.0-pre.3; slot-25 H1, commit \texttt{006f60a9}, audited PASS\_WITH\_MINORS in \file{docs/implementation/slot25/H1\_audit.md}):
\file{HeckeExt/HeckeHidaEisLambda.\{h,cpp\}} (the ordinary Eisenstein block of $\eord S_k$ at one weight, the generator test, K9's integers, the G8 source),
\code{HeckeHidaEisLambdaValidation} (Proposition~\ref{prop:u-H11} on the $11$ C2 chain pairs, (U1) on the forms side and on cohomology, the record's G8), and the record's G8 under the
non-default settings value \code{heckez.G8 = hida} (\HA{sec:recordschema} item~8; never golden). As planned (slot-25/26 H1, after this part's referee ACCEPT, rule~22): the read-only env arm \code{HECKE\_LAMBDA\_CELLS="ell:a:K:m,\dots"} on
\code{HeckeLambdaValidation} (\code{EFFECTIVE} line) for the truncation of $\Tu(486)_{\mathfrak m}$ at $\ell=547$, $K=m=3$ (Theorem~\ref{thm:u-interp}: $5$ weights at precision $5$),
reproducing (D1) and (D3) and reading $b$ modulo $(\ell^{3},X^{3})$; Proposition~\ref{prop:u-H11}'s generator test on the C2 chains (one integer per pair); the Lean checker K9 (a few integers
per pair, \code{decide +kernel}). H1-R (a second workstation, owner-launched) is the test of (U23). Slot-25 HZ-LOC (\code{HeckeZLocalValidation}) supplies the per-maximal-ideal certificate of
Remark~\ref{rmk:u-H0census} for census rows. No statement of this part rests on code.

%----------------------------------------------------------------------
\section{The ramified ordinary families at tame level one: the discriminant series (part~7)}\label{sec:u-H2}
%----------------------------------------------------------------------
% NEW at v0.7.4-pre.1 (hecke-extensions slot 26, package Z6 = H2; design docs/implementation/slot26/Z6_design.md v0.2, audit Z6_design_audit.md;
% mathematics docs/implementation/slot24/heckez/05-hida-targets.md v0.3.2 Sec. 4).  Placed after part 6 so that no statement number of 0.7.3 moves.
% Text first, before the code (rule 22).

\paragraph{Status and scope of this part.} Written as a pre-release (v0.7.4-pre.1) for a referee before the code of the slot-26 package Z6; v0.7.4-pre.2
applied the referee's changes and added the computed data (Remark~\ref{rmk:u-H2data}); \textbf{since v0.8.0} (signed by the owner 2026-11-11) this part is released with the rest.
Proposition~\ref{prop:u-loewy}(4) lists eleven ordinary non-Eisenstein blocks of Loewy length~$2$ at tame level~$1$, at ten cells (two blocks at $(68,157)$). This
part shows that each is a single regular family, ramified over $\Iw$ (Lemma~\ref{lem:u-DF}). It defines the family's discriminant series $D_F$ and states what a
truncation of $D_F$ modulo $(\ell^{K},(X-x_c)^{m})$ decides (Lemma~\ref{lem:u-DFread}): either a sub-tower of weights on which the specialisation is a non-maximal order, or
ramification of constant depth at every weight. Only in the first case is the family isomorphic to $\Zl[[Z]]$. The computed truncations are the prediction (U25).
Tame level~$1$, $\ell\ge5$, $x_k=u^{k}-1$ as in part~3.
\textbf{Literature.} Hida's congruence module of a family \LIT{\cite[Thm.~0.1, p.~329]{hida_ajm88}; the locator checked against the text of Hida's paper (read-only copy)} is, for a monogenic quadratic family after base change to its normalisation, the different, of norm $(D_F)$;
it interpolates the algebraic part of $L(1,\mathrm{Ad}\,f_k)$. No statement of this part uses it.

\begin{definition}[The discriminant series]\label{def:u-DF}
Let $\mathfrak m$ be a maximal ideal of $\hcusp(a)$ with $\rank_\Iw\hcusp(a)_{\mathfrak m}=2$. For $t\in\hcusp(a)_{\mathfrak m}$ with characteristic polynomial $Y^{2}-\tau Y+\nu$ on the free
$\Iw$-module $\hcusp(a)_{\mathfrak m}$, put $D_t:=\tau^{2}-4\nu\in\Iw$. If $t$ generates $\hcusp(a)_{\mathfrak m}$ over $\Iw$, write $D_F:=D_t$, where $F$ is the family
$\operatorname{Spec}\hcusp(a)_{\mathfrak m}$. $D_F$ is determined up to $(\Iw^{\times})^{2}$: if $t,t'$ both generate,
$t'=\alpha+\beta t$ with $\beta\in\Iw^{\times}$ and $D_{t'}=\beta^{2}D_t$ (the proof of Lemma~\ref{lem:u-DF}(3), which uses only that
$\{1,t\}$ is a $\Iw$-basis). A generator always exists: the fibre $\hcusp(a)_{\mathfrak m}/\mathfrak m_\Iw$ is a local $\Fl$-algebra of dimension~$2$, hence $\Fl[y]/(y^{2})$ or $\mathbb F_{\ell^{2}}$,
generated by one element, and by Nakayama a lift of that element generates $\hcusp(a)_{\mathfrak m}$ over $\Iw$.
\end{definition}

\begin{lemma}[The eleven families]\label{lem:u-DF}
\PROVED{} given Proposition~\ref{prop:u-loewy}(4) (\COMP), from Lemma~\ref{lem:u-H0}, Lemma~\ref{lem:u-levelone}, Theorem~\ref{thm:u-control}, Proposition~\ref{prop:u-lambdaprod}(4), Lemma~\ref{lem:u-DFread} and
\STD{$q$-expansion duality $S_k(\Zl)_{\mathfrak m}\cong\Hom_{\Zl}(\hcusp_{k,\mathfrak m},\Zl)$}. Let $\mathfrak m$ be one of the eleven blocks of Proposition~\ref{prop:u-loewy}(4), at the cell
$(n,\ell)$, $k_0=n+2$, $a\equiv k_0$ (in all ten cells $k_0<\ell-1$, so $a\not\equiv0$).
\begin{enumerate}
\item $\hcusp(a)_{\mathfrak m}$ is $\Iw$-free of rank $2$ with $\hcusp(a)_{\mathfrak m}/\mathfrak m_\Iw\cong\Fl[y]/(y^{2})$; the residue $\bar\lambda_p$ of $T(p)$ at $\mathfrak m$ is the
  eigenvalue of $T_p$ on the $\mathfrak m$-part of $\eord S_k$ modulo $\ell$ at every $k\ge3$ in the class.
\item For a prime $q\ne\ell$, $T(q)$ generates $\hcusp(a)_{\mathfrak m}$ over $\Iw$ iff $T_q-\bar\lambda_q\ne0$ on the $\mathfrak m$-part of $\eord S_k/\ell$ for one (equivalently every)
  $k\ge3$ in the class; $T(2)$ does. Then $\hcusp(a)_{\mathfrak m}=\Iw[Y]/(F_q)$, and $F_q(x_k;Y)$ is the characteristic polynomial of $T_q$ on the $\mathfrak m$-part of $\eord S_k(\Zl)$.
\item If $t,t'$ both generate, then $D_{t'}=\beta^{2}D_t$ with $\beta\in\Iw^{\times}$. So the ideal $(D_F)$ is the discriminant ideal of $\hcusp(a)_{\mathfrak m}$ over $\Iw$, and its
  $\mu$- and $\lambda$-invariants and its zeros do not depend on the generator.
\item $D_F\in\mathfrak m_\Iw\setminus\mathfrak m_\Iw^{2}$, and $\hcusp(a)_{\mathfrak m}\cong\Zl[[X,Z]]/(Z^{2}-D_F/4)$ is a regular local domain of dimension~$2$: a single family, normal and
  ramified over $\Iw$ (kind~(ii) of Proposition~\ref{prop:u-loewy}(2), now over $\Iw$). It is isomorphic to $\Zl[[Z]]$ iff case~(A) of Lemma~\ref{lem:u-DFread} holds for $D_F$;
  in case~(B), $\ell$ lies in the square of its maximal ideal, and it is not (for instance $\Zl[\sqrt\ell][[X]]$ when $D_F=4\ell$).
\item At every $k\ge3$ in the class, $\hcusp_{k,\mathfrak m}\cong\Zl[x]/(F_q(x_k;x))$ has discriminant $D_F(x_k)\ne0$ up to a unit square. It is the ring of integers of a ramified
  quadratic extension of $\Ql$ iff $v(D_F(x_k))=1$; if $v(D_F(x_k))=v\ge2$ it is not maximal, of index $\ell^{\lfloor v/2\rfloor}$ in its normalisation.
\end{enumerate}
\end{lemma}

\begin{proof}
(1) Theorem~\ref{thm:u-control} and Lemma~\ref{lem:u-H0}(1) with Proposition~\ref{prop:u-loewy}(4) at $k_0$; the maximal ideals of $\hcusp(a)$ and of $\hcusp_k$ correspond with
$T(p)\mapsto T_p$, and the residue field is $\Fl$. (2) By Lemma~\ref{lem:u-H0}(3) and (1), $t$ generates iff its image in $\hcusp_{k,\mathfrak m}/\ell\cong\Fl[y]/(y^{2})$ is not
a scalar. By duality (and Lemma~\ref{lem:u-levelone}(1)) the $\mathfrak m$-part of $\eord S_k(\Zl)$ is a faithful $\hcusp_{k,\mathfrak m}$-module,
free of rank $2$ over $\Zl$, and its reduction is faithful over $\hcusp_{k,\mathfrak m}/\ell$; so the image is a scalar iff
$T_q-\bar\lambda_q=0$ there. At $k_0$, $\Zl[T_2]$ has discriminant of
valuation~$1$ (Proposition~\ref{prop:u-loewy}(4)), so it is the maximal order of a ramified quadratic extension and $T_2-\tilde\lambda_2$ is a uniformiser, not $0$ modulo $\ell$. The characteristic polynomial
specialises (Proposition~\ref{prop:u-lambdaprod}(4) for $\hcusp$), and the transpose of multiplication by $t$ is the action on the dual. (3) $\{1,t\}$ is a basis;
$t'=\alpha+\beta t$, and $\bar\beta\ne0$ since $t'$ generates modulo $\mathfrak m_\Iw$; the discriminant of $Y^{2}-\tau Y+\nu$ changes by $\beta^{2}$. (4) $v(D_F(x_{k_0}))=1$ by (2) and
Proposition~\ref{prop:u-loewy}(4), while every element of $\mathfrak m_\Iw^{2}=(\ell^{2},\ell X,X^{2})$ takes values in $\ell^{2}\Zl$ on $\ell\Zl$. $D_F\in\mathfrak m_\Iw$ since the fibre is not
reduced. Completing the square ($\ell$ odd), $\hcusp(a)_{\mathfrak m}\cong\Zl[[X]][Z]/(Z^{2}-D_F/4)\cong\Zl[[X,Z]]/(Z^{2}-D_F/4)$, which is regular as $Z^{2}-D_F/4\notin(\ell,X,Z)^{2}$
and a domain as $D_F$ is not a square in $\Iw$ (it lies in $\mathfrak m_\Iw\setminus\mathfrak m_\Iw^{2}$). Its maximal ideal is $(\ell,X,Z)$ and its cotangent space is spanned by the images of $\ell,X,Z$
modulo the relation $Z^{2}-D_F/4\equiv-(d_0+d_1X)/4$ (Lemma~\ref{lem:u-DFread}'s $d_i$). In case~(A) $d_1$ is a unit, the relation eliminates $X$, so $\ell$ and $Z$ form a regular system of
parameters with $\ell\notin\mathfrak m^{2}$, and $\Zl[[Z]]\to\hcusp(a)_{\mathfrak m}$ is an isomorphism (Weierstrass: $X-\xi\in(Z^{2})\cdot$unit). In case~(B) $v(d_0)=1$, the relation gives $\ell\in\mathfrak m^{2}$;
in $\Zl[[Z]]$, $\ell\notin\mathfrak m^{2}$. (5) Theorem~\ref{thm:u-control} and (2). $\hcusp_k\otimes\Q$ is étale at level one (Lemma~\ref{lem:u-levelone}(2)), so $D_F(x_k)\ne0$. For $v(\delta)=v$,
$\sqrt\delta/\ell^{\lfloor v/2\rfloor}$ is integral over $\Zl$, and $\delta/\ell^{2\lfloor v/2\rfloor}$ has
valuation $0$ or $1$, so $\Zl[\sqrt\delta/\ell^{\lfloor v/2\rfloor}]$ is the normalisation ($\ell$ odd; $\Zl\times\Zl$ when that quotient is a unit
square).
\end{proof}

\begin{lemma}[What a truncation of $D_F$ decides]\label{lem:u-DFread}
\PROVED{} (elementary; \STD{Weierstrass preparation in $\Iw$}). Let $D=\sum_i d_iX^{i}\in\mathfrak m_\Iw\setminus\mathfrak m_\Iw^{2}$.
\begin{enumerate}
\item[(A)] If $d_1\in\Zl^{\times}$, then $D=w\,(X-\xi)$ with $w\in\Iw^{\times}$, $\xi=u^{s}-1$ for a unique $s=s_F\in\Zl$, and $v(D(x_k))=1+v_\ell(k-s_F)$ for every integer $k$. In
  every class the weights with $v(D(x_k))\ge j+1$ are those with $k\equiv s_F\pmod{\ell^{j}}$, a non-empty sub-tower for every $j\ge1$.
\item[(B)] If $d_1\in\ell\Zl$, then $v(d_0)=1$ and $v(D(x_k))=1$ for every integer $k$. Either $D=\ell\cdot w$ ($\mu=1$, $\lambda=0$), or $\mu=0$, $\lambda\ge2$ and the zeros of $D$
  have valuation $1/\lambda$.
\item[(R)] \textbf{(Reading.)} Let $K,m\ge2$, $x_c\in\ell\Zl$, and $\tilde D=\sum_{i<m}e_iW^{i}$, $W:=X-x_c$, with $\tilde D\equiv D\pmod{(\ell^{K},W^{m})}$. Then (A) holds iff $e_1$ is a unit.
  In (A), $\tilde D$ has a root $\tilde\zeta\in\ell\Zl$ modulo $\ell^{\min(K,m)}$, every such root satisfies $x_c+\tilde\zeta\equiv\xi\pmod{\ell^{\min(K,m)}}$, and so $\tilde D$ determines $s_F$
  modulo $\ell^{\min(K,m)-1}$. In (B), if some $e_i$, $2\le i<m$, is a unit, the least such $i$ is $\lambda$ and $\mu=0$; otherwise $\tilde D$ is consistent both with $\mu=1$ and
  with $\mu=0$, $\lambda\ge m$, and leaves the invariants undecided.
\end{enumerate}
\end{lemma}

\begin{proof}
$x_k\equiv k\ell\pmod{\ell^{2}}$ and $v(x_k-x_{k'})=1+v(k-k')$ (the proof of Lemma~\ref{lem:u-tower}); $k\mapsto u^{k}$ is an isomorphism $\Zl\to1+\ell\Zl$.
$D\notin\mathfrak m_\Iw^{2}$ means $v(d_0)=1$ or $d_1\in\Zl^{\times}$. (A) Weierstrass: $\lambda=1$, the distinguished factor is $X-\xi$, $\xi\in\ell\Zl$ as $d_0\in\ell\Zl$; $v(D(x_k))=v(x_k-\xi)=v(u^{k}-u^{s})=1+v(k-s)$.
Non-emptiness: the Chinese remainder theorem for $\ell^{j}$ and $\ell-1$. (B) $v(d_0)=1$, and $v(d_1x_k+d_2x_k^{2}+\dots)\ge2$; Weierstrass as stated, the distinguished factor Eisenstein. (R) Recentring by $x_c\in\mathfrak m_\Iw$
preserves $\mu$, $\lambda$, and $e_1\equiv D'(x_c)\equiv d_1\pmod\ell$. In (A), $e_0\in\ell\Zl$ and $e_1$ is a unit, so Hensel gives $\tilde\zeta\in\ell\Zl$; then
$D(x_c+\tilde\zeta)=\tilde D(\tilde\zeta)+O(\ell^{K})+O(\tilde\zeta^{m})\equiv0\pmod{\ell^{\min(K,m)}}$ and $v(D(x_c+\tilde\zeta))=v(x_c+\tilde\zeta-\xi)$. In (B), $e_i\equiv
d_i\pmod\ell$ for $i<m$. If every $e_i$, $1\le i<m$, lies in $\ell\Zl$, then $\tilde D$ itself ($\mu=1$, as $v(e_0)=1$) and $\tilde D+W^{m}$
($\mu=0$, $\lambda=m$) both lie in $\mathfrak m_\Iw\setminus\mathfrak m_\Iw^{2}$ and have the truncation $\tilde D$.
\end{proof}

\begin{prediction}[(U25) The eleven discriminant series]\label{pr:u-H2}
Stated before the computation (rule~22), at the eleven blocks of Proposition~\ref{prop:u-loewy}(4) and, in~(c), at one known answer.
\begin{enumerate}
\item[(a)] \PROVED{} given Proposition~\ref{prop:u-loewy}(4) (\COMP), from Lemmas~\ref{lem:u-DF}, \ref{lem:u-DFread} and Theorem~\ref{thm:u-interp}, applied to
  $g(W):=D_{T(2)}(x_{k_0}+W)\in\Zl[[W]]$ at the nodes $w_j=x_{k_j}-x_{k_0}\in\ell\Zl$ (its proof and that of Lemma~\ref{lem:u-wtprod} use only
  $v(w_j)\ge1$ and $v(w_a-w_b)=1$ for $a\ne b$): at each of the eleven blocks, $D_{T(2)}$ modulo
  $(\ell^{3},(X-x_{k_0})^{3})$ is determined by the characteristic polynomials of $T_2$ on the $\mathfrak m$-parts of $\eord S_{k_j}$ modulo $\ell^{5}$, $k_j=k_0+j(\ell-1)$, $0\le j\le4$;
  $e_0,e_1,e_2$ are the coefficients of $D_{T(2)}$ in $W=X-x_{k_0}$, and $v(e_0)=1$; $D_F$ falls in exactly one of the cases (A), (B) of Lemma~\ref{lem:u-DFread}; in (A) the truncation determines $s_F$ modulo $\ell^{2}$, hence the least
  weights $k^{(2)}_F$, $k^{(3)}_F$ ($\ge3$) of the class with $k\equiv s_F$ modulo $\ell$, resp.\ $\ell^{2}$ ($k_0$ is never one: $v(D_F(x_{k_0}))=1$); in (B) it determines $\lambda=2$ or leaves the invariants undecided. The case,
  $s_F$, $\mu$ and $\lambda$ are invariants of $D_F$ by Lemma~\ref{lem:u-DF}(3). The truncations read through $T(2)$ and through a second generator --- for a prime $q\ne2,\ell$, one of $T(q)$, $T(2)+T(q)$ generates (Lemma~\ref{lem:u-DF}(2): if
  $T(q)$ is a scalar modulo $\mathfrak m_\Iw$, then $T(2)+T(q)$ is not) --- differ by the square of a unit of $\Iw$; and the interpolation polynomial of
  the five values predicts $D_{T(2)}(x_{k_5})$ modulo $\ell^{5}$ (Theorem~\ref{thm:u-interp}(3)). The truncation from $k_2,k_3,k_4$ agrees with that from $k_0,\dots,k_4$ modulo
  $(\ell^{2},(X-x_{k_0})^{2})$; this is a consequence of the $\ell$-integrality of the five-node interpolant (it holds for any five values whose interpolant is $\ell$-integral) and is
  \emph{not} evidence: it detects no error of the values, and the implementation uses it as a self-check of the code only.
\item[(b)] \PROVED{} given Proposition~\ref{prop:u-loewy}(4) (\COMP), from Lemma~\ref{lem:u-DF}(5) and Lemma~\ref{lem:u-DFread}(A): in case (A), $v(\operatorname{disc})=1+v(k^{(2)}_F-s_F)\ge2$ at
  $k^{(2)}_F$ and $\hcusp_{k^{(2)}_F,\mathfrak m}$ is not maximal; at every weight with $k\not\equiv s_F\pmod\ell$ it is a ramified discrete valuation ring. In case (B) every $\hcusp_{k,\mathfrak m}$ is a
  ramified discrete valuation ring.
\item[(c)] \textbf{(A known answer.)} \PROVED{} given the data (D1)--(D4) of Proposition~\ref{prop:u-H12}, from its
parts~(1), (4), (5) and Lemma~\ref{lem:u-DFread}(A): at $(\ell,a)=(547,486)$ the Eisenstein-local
  $\hcusp(a)_{\mathfrak m}$ has rank~$2$, is generated by $T(2)$, and $D_F=b^{2}-4wG_a$ with $b\in\mathfrak m_\Iw$, $G_a=w_1(X-\xi)$; so $D_F$ is in case (A), its zero is $\equiv\xi$
  modulo $\ell^{2}$, $s_F\equiv s_a\pmod\ell$, and $k^{(2)}_F=55086$, the rung of Proposition~\ref{prop:u-H12}(5).
\end{enumerate}
\emph{As a test of the implementation} (slot-26 Z6, \code{HeckeHidaDiscLambdaValidation}, the forms side on $\eord S_k$ modulo $\ell^{5}$, $K=m=3$):
the computation reproduces (a) at the eleven blocks and (c) at $(547,486)$, and, under \code{HECKE\_HIDA\_H2\_ALL=1}, checks (b) directly at $k^{(2)}_F$ for $\ell\in\{67,71,73\}$ in case (A).
(U25) makes no guess at which case occurs at which block (nothing in [U] predicts it); so (a) and (b) can fail only through a defect of the implementation or of
Proposition~\ref{prop:u-loewy}(4), as for (U23).
The computed values ($e_0,e_1,e_2$, the case, $s_F$ modulo $\ell^{2}$, $k^{(2)}_F$, $k^{(3)}_F$ per block) are Remark~\ref{rmk:u-H2data}; they are pinned
in the validation as oracle constants.
\end{prediction}

\begin{remark}[The eleven discriminant series, computed]\label{rmk:u-H2data}
\COMP{} (slot-26 Z6: \code{HeckeHidaDiscLambdaValidation}, implementation commit \code{11507416}, binary sha256 \code{82c5b9b2\dots abacf}; the rows are pinned as oracle
constants from commit \code{4b35b0ba} on, \code{pinned=11/11}; report \file{docs/implementation/slot26/Z6\_report.md} \S3). At each block, $D=[e_0,e_1,e_2]$ is $D_{T(2)}$ modulo
$(\ell^{3},W^{3})$, $W=X-x_{k_0}$, read from $k_j=k_0+j(\ell-1)$, $0\le j\le4$, modulo $\ell^{5}$; $\bar r_2$ is the residue of $T(2)$ at $\mathfrak m$.
\begin{center}\small
\begin{tabular}{rrrlcrrrc}
cell $(n,\ell)$ & $k_0$ & $\bar r_2$ & $[e_0,e_1,e_2]$ & case & $s_F\bmod\ell^{2}$ & $k^{(2)}_F$ & $k^{(3)}_F$ & $\lambda$\\\hline
$(30,67)$ & 32 & 14 & $[67670,81595,32906]$ & A & 2079 & 2012 & 293864 & 1\\
$(26,131)$ & 28 & 52 & $[2028797,1769733,956238]$ & A & 8181 & 13028 & 643138 & 1\\
$(26,139)$ & 28 & 30 & $[903500,1799153,2635989]$ & A & 18627 & 3754 & 617578 & 1\\
$(32,479)$ & 34 & 472 & $[61589820,24027696,6939087]$ & A & 163322 & 24412 & 43298230 & 1\\
$(36,181)$ & 38 & 178 & $[4850076,877126,5105143]$ & \textbf{B} & -- & -- & -- & \textbf{2}\\
$(36,349)$ & 38 & 171 & $[21863454,40091683,4242504]$ & A & 553 & 63722 & 22046534 & 1\\
$(38,73)$ & 40 & 61 & $[248054,218975,21318]$ & A & 27 & 976 & 69304 & 1\\
$(44,227)$ & 46 & 5 & $[5546291,7722249,7001584]$ & A & 17357 & 38014 & 4706496 & 1\\
$(52,71)$ & 54 & 1 & $[97341,43135,167397]$ & A & 3084 & 1664 & 255134 & 1\\
$(68,157)$ & 70 & 96 & $[2470552,253158,98053]$ & A & 22529 & 23314 & 145774 & 1\\
$(68,157)$ & 70 & 111 & $[880770,2234752,1764920]$ & A & 5411 & 24094 & 2938642 & 1\\
\end{tabular}
\end{center}
At every block: rank~$2$ at every weight, $T(2)$ generates (Lemma~\ref{lem:u-DF}(2)), $v(e_0)=1$; the truncation read through $t'=T(3)$ (which generates at all eleven) gives
the same case and $s_F$ and a unit-square ratio; the held-out weight $k_5$ agrees modulo $\ell^{5}$ wherever computed (all eleven under \code{HECKE\_HIDA\_H2\_ALL=1}).
Ten blocks are in case (A), $\lambda=1$, $\hcusp(a)_{\mathfrak m}\cong\Zl[[Z]]$ (Lemma~\ref{lem:u-DF}(4)), with the rung sub-tower $k\equiv s_F\pmod\ell$; at $\ell\in\{67,71,73\}$ the
rung check of (U25)(b) holds: $v(\operatorname{disc})=2=1+v(k^{(2)}_F-s_F)$ at $k^{(2)}_F=2012$, $1664$, $976$ respectively, a non-maximal order.
\textbf{At $(36,181)$ the family is in case (B)}: $e_1=877126=181\cdot4846$ is not a unit and $e_2\equiv38\pmod{181}$ is, so $\mu=0$, $\lambda=2$ by Lemma~\ref{lem:u-DFread}(R) (not
undecided); $v(D_F(x_k))=1$ at every weight of the class, so every specialisation $\hcusp_{k,\mathfrak m}$ is a ramified discrete valuation ring and there is no rung; and by
Lemma~\ref{lem:u-DF}(4) $\hcusp(a)_{\mathfrak m}$ is a regular local ring with $\ell\in\mathfrak m^{2}$, \emph{not} isomorphic to $\Zl[[Z]]$. It is read independently through
$T(3)$ (case (B), $\lambda=2$, unit-square ratio), the held-out weight $k_5=938$ agrees modulo $181^{5}$, and an independent implementation sharing no code with HACS (the full
$S_k$, the Hensel factor of $(Y-\bar r_2)^{2}$, exact rational interpolation; slot-26 Z6 code audit \S2) gives the same $[e_0,e_1,e_2]$ and $v(D(x_{k_j})-D(x_{k_0}))=2$ at
$j=1,\dots,6$. The known answer of (c) agrees: at $(547,486)$, $\chi(T_2\mid\text{block})$ equals Proposition~\ref{prop:u-H12}'s $f_k$ modulo $547^{5}$ at the five weights, case (A), $k^{(2)}_F=55086$.
\end{remark}

\paragraph{Code map of part 7.} \emph{Built} (slot-26 Z6, after this part's referee; commits \code{11507416}, \code{4b35b0ba}): \file{HeckeExt/HeckeHidaDiscLambda.\{h,cpp\}} (the per-block Fitting split at a
non-Eisenstein residual system, the truncation of $D_{T(2)}$, the reading of Lemma~\ref{lem:u-DFread}), \code{HeckeHidaDiscLambdaValidation} (standard; the read-only arms
\code{HECKE\_HIDA\_H2\_ALL=1} and \code{HECKE\_HIDA\_H2\_PROBE=1}; the eleven rows of Remark~\ref{rmk:u-H2data} pinned). No statement of this part rests on code; Remark~\ref{rmk:u-H2data} is \COMP.

\bibliographystyle{alpha}
\bibliography{gl2_refs,hida_hecke_algebra_refs}

%======================================================================
% Revision history (v0.7.0-pre.4): the date field of 0.7.0-pre.3 and the changelogs of parts 1--6, moved here verbatim.
%======================================================================
\appendix
\sloppy
\section{Revision history}\label{app:u-history}
[U] keeps [C]'s practice: a published version's text never changes under its own number. This appendix is the record of the
versions. Its first subsection is the release record, one paragraph per version, newest first (up to 0.7.0-pre.3 it was the date
field of the title page); the second holds the changelogs of parts~1--7 (those of parts~1--6 up to 0.7.0-pre.3 stood at the end of each part); the
third quotes the earlier title page.

\subsection{Release record}
Version 0.8.1 --- 2026-11-11 (hecke-extensions slot~26, \textbf{signed by the owner 2026-11-11} (``I sign the commit update. Fix the two in [U] please.''): status sentences only; no statement, label or value moves. Two sentences of 0.8.0 still called part~7 a pre-release, \textbf{not signed}. In \S\ref{sec:u-status}: ``Part~7 \dots{} is a pre-release, not signed, written for a referee before the code of the slot-26 package Z6 \dots; parts~1--6 are as released in 0.7.3 \dots'' now reads ``Part~7 \dots{} was a pre-release through v0.7.4-pre.2, written for a referee \dots; since v0.8.0 (signed by the owner 2026-11-11) it is released with the rest; parts~1--6 are as released in 0.7.3 \dots''. In the ``Status and scope'' paragraph of part~7: ``A pre-release (v0.7.4-pre.1, not signed), written for a referee before the code of the slot-26 package Z6; v0.7.4-pre.2 (not signed) applies the referee's changes and adds the computed data'' now reads ``Written as a pre-release (v0.7.4-pre.1) for a referee before the code of the slot-26 package Z6; v0.7.4-pre.2 applied the referee's changes and added the computed data; since v0.8.0 (signed by the owner 2026-11-11) this part is released with the rest.'' Process rule (owner, 2026-11-11): the commit that records an owner signature may also change the document's own ``pre-release, not signed'' status sentences for the part being signed to ``signed'', and nothing else in the body. Companion pointers: [C] 1.27.0, [E] 1.25.3.)

Version 0.8.0 --- 2026-11-11 (hecke-extensions slot~26, \textbf{signed by the owner 2026-11-11} (``I sign [U] 0.7.4-pre.2. I agree with the version number.''): the owner-signed release of v0.7.4-pre.2; text unchanged --- header, date, this record and the changelog only; no statement, label or value moves; numbered 0.8.0 rather than 0.7.4 per the referee's m11 (every earlier new part took a minor version), accepted by the owner. Touches of signed 0.7.3 text, both signed with this release: the Part~6 release sentence of \S\ref{sec:u-status} (m10, quoted in the record of 0.7.4-pre.2 below), and the printed key of \cite{hida_aif2}, which becomes [Hid88b] (was [Hid88]) at its eight sites in parts~3--4 and this appendix, the source text unchanged, because the new bibliography entry for Hida's 1988 paper is [Hid88a]. Notes of the confirm-read: (U25) ``stated before the computation (rule~22)'' is true of the statement as issued at v0.7.4-pre.1; the \COMP{} tag of Remark~\ref{rmk:u-H2data} names the binary of commit \code{11507416}, the pinned gate being of the \code{4b35b0ba} build, the same eleven rows. Companion pointers: [C] 1.27.0, [E] 1.25.3.)

Version 0.7.4-pre.2 --- 2026-11-11 (hecke-extensions slot 26, Z6 fix round, \textbf{a pre-release, not signed}: the referee's m1--m10 (report \file{U074pre1\_referee.md}, ACCEPT WITH CHANGES) and the notes n1, n3, n4, n6 in part~7 and the front matter; the code audit's F1 (the subset agreement of (U25)(a) is a consequence, not evidence); new Remark~\ref{rmk:u-H2data} (\COMP: the eleven computed rows; $(36,181)$ is in case (B), $\lambda=2$, not $\cong\Zl[[Z]]$); the code map \emph{Built}; a bibliography entry for Hida's 1988 paper (its \VERIFYloc{} lifted). No statement, label or number of 0.7.4-pre.1 moves; one tag gains a citation (Lemma~\ref{lem:u-DF}: Lemma~\ref{lem:u-levelone}), one gains a source ((U25)(c)). \textbf{Touches signed 0.7.3 text} (referee m10, a status sentence of \S\ref{sec:u-status}): ``Part~6 \dots{} is a pre-release, not signed, \dots; parts~1--5 are as released in 0.6.0 \dots. It cites [C] v1.24.0 by label, and [C] 1.25.0-pre.3's \HA{sec:heckeZ} (a pre-release; the 0.7.0 release must cite a released [C] 1.25.0) by label.'' now reads ``Part~6 \dots{} was a pre-release through v0.7.0-pre.4 \dots; since v0.7.0 (signed by the owner 2026-11-07) it is released with the rest; it cites [C] v1.24.0 by label and [C] v1.25.0 (\HA{sec:heckeZ}).'' --- for the owner's sign-off. \textbf{Version:} the referee (m11) recommends that the signed release of part~7 be numbered 0.8.0 (every earlier new part took a minor version); the number stays 0.7.4-pre.2 here, and the owner chooses 0.7.4 or 0.8.0 at signing. Companion pointers: [C] 1.27.0, [E] 1.25.3.)

Version 0.7.4-pre.1 --- 2026-11-11 (hecke-extensions slot 26, Z6, \textbf{a pre-release, not signed}: new part~7, \S\ref{sec:u-H2} --- the eleven ramified ordinary families of Proposition~\ref{prop:u-loewy}(4): regular over $\Iw$, the discriminant series $D_F$, what its truncation decides, the prediction (U25); the abstract's count (U1)--(U25) and one clause; a Part~7 sentence in \S\ref{sec:u-status}; no statement, label, number or tag of 0.7.3 moves; see the changelog of part~7. Companion pointers: [C] 1.27.0, [E] 1.25.3.)

Version 0.7.3 --- 2026-11-10 (hecke-extensions slot 26, \textbf{signed by the owner 2026-11-10}: the owner-signed release of 0.7.3-pre.1; no change of mathematical content --- header, date and this record only; no statement, label, number or tag moves. Companion pointers: [C] v1.26.4, [E] 1.25.3.)

Version 0.7.3-pre.1 --- 2026-11-10 (slot 26, a \textbf{pre-release, not signed}: title page only, no change of mathematical content: the date line of the title page names only the version, the date and this appendix; the words ``signed by the owner 2026-11-09'' leave the title page, and the sign-off record of every version stays in this record; no statement, label, number or tag moves. Companion pointers: [C] v1.26.4, [E] 1.25.3 (the same title-page change; signed with 0.7.3).)

Version 0.7.2 --- 2026-11-09 (hecke-extensions slot 26, \textbf{signed by the owner 2026-11-09}: the owner-signed release of 0.7.2-pre.1; no change of mathematical content --- header, date and this record only; no statement, label, number or tag moves. Companion pointers: [C] v1.26.3, [E] 1.25.2.)

Version 0.7.2-pre.1 --- 2026-11-09 (slot 26, a \textbf{pre-release, not signed}: text-only: private path removed, no change of mathematical content (release audit of 2026-11-09, finding S1): the seven paths into the author's scratch directory were reduced to the scripts' file names, or to ``outside the tree''; the two citations of the September cluster run's report name it ``the WP-C10 report'' instead of a file name that carried the cluster's name; no statement, label, number or tag moves. Companion pointers: [C] v1.26.3 (the same scrub; signed with 0.7.2), [E] 1.25.2.)

Version 0.7.1 --- 2026-11-09 (hecke-extensions slot 26, \textbf{signed by the owner 2026-11-09}: the owner-signed release of 0.7.1-pre.1; no change of content --- header, date and this record only; no statement, label, number or tag moves. Companion pointers: [C] v1.26.1, [E] 1.25.1.)

Version 0.7.1-pre.1 --- 2026-11-09 (slot 26, a \textbf{pre-release, not signed}: no change of mathematical content: machine names and private paths removed for the public release (owner ruling 2026-11-09).)

Version 0.7.0 --- 2026-11-07 (hecke-extensions slot 25, \textbf{signed by the owner 2026-11-07}: the owner-signed release of 0.7.0-pre.4; no change of content --- header, date and this record only; no statement, label, number or tag moves. Companion pointers: [C] v1.25.0, [E] 1.25.0.)

Version 0.7.0-pre.4 --- 2026-11-07 (\textbf{a pre-release, not signed, not yet seen by the owner}: front matter and layout only. The title page has a one-line date, the uniform author block and a generic abstract that describes all six parts; the date field of 0.7.0-pre.3 (the release record below) and the changelogs of parts~1--6, which stood at the end of each part (in part~5 with the paragraphs ``Wording repair at 0.6.0-pre.2 (part~4)'' and ``Release candidate 0.6.0''), are moved here verbatim; the earlier author line and abstract are quoted in the last subsection. The sections of parts~2--6 are retitled by subject; their earlier titles were ``Part 2: tame level \texorpdfstring{$N$}{N} --- CM components, CM \texorpdfstring{$\cap$}{cap} GL(2), CM \texorpdfstring{$\cap$}{cap} Eisenstein'', ``Part 3: the \texorpdfstring{$\Iw$}{Lambda}-adic algebra, truncated'', ``Part 4: \texorpdfstring{$\Iw$}{Lambda}-adic modular symbols, truncated'', ``Part 5: the tower \texorpdfstring{$\Ga_1(\ell^{r})$}{Gamma1(l^r)}, wild characters, and the \texorpdfstring{$\Omega$}{Omega}-type truncations'' and ``Part 6: the Eisenstein component over \texorpdfstring{$\Iw$}{Lambda} at irregular pairs''. No statement, label, number or status tag changes; the status paragraph of \S\ref{sec:u-status} and the code maps are unchanged).

Version 0.7.0-pre.3 --- 2026-11-06 (hecke-extensions slot 25, H1 merge, \textbf{a pre-release, not signed, not yet seen by the owner}: in part~6 the explicit (U23) anchor polynomial, Remark~\ref{rmk:u-H12general} (Steps 1--2 of Proposition~\ref{prop:u-H12} at any $g_a\ge2$) and the code map \emph{Built}; no result changes, no statement number moves; see the changelog of part~6);

Version 0.7.0-pre.2 --- 2026-11-06 (hecke-extensions slot 25, U6 fix round, \textbf{a pre-release, not signed}: the referee's F1--F13 in part~6 and one status sentence of \S\ref{sec:u-status}; no result changes; see the changelog of part~6);

Version 0.7.0-pre.1 --- 2026-11-06 (hecke-extensions slot 25, U6, \textbf{a pre-release, not signed}: new part~6, \S\ref{sec:u-eislambda} --- the Eisenstein component over $\Iw$ at irregular pairs: one weight decides the ring-theoretic type, cuspidal rank one, the regular family at $(547,486)$, the rung, multiplicity one at weight $486$, the predictions (U23)--(U24); Remark~\ref{rmk:u-classzero} credits Ohta's adaptation of Harder--Pink and Kurihara (owner-signed P-K1); a Bella\"\i che--Pollack citation in Proposition~\ref{prop:u-loewy}(3); no statement, label, number or tag of 0.6.0 moves; see the changelog of part~6);

Version 0.6.0 --- 2026-10-16 signed by the owner (text of 2026-10-12, slot 22 CLOSE-T);

Version 0.6.0 (release candidate) --- 2026-10-12 (hecke-extensions slot 22, CLOSE-T: the \textbf{release candidate} from 0.6.0-pre.3, \textbf{owner sign-off pending}, not yet signed; header, date and changelog only --- no statement, label, number or tag moves);

Version 0.6.0-pre.3 --- 2026-10-12 (hecke-extensions slot 22, \textbf{a pre-release, not signed}: CLOSE-T T1 repairs --- two wording points of the F-T1 release referee (a paragraph title, a version label in Definition~\ref{def:u-lmsym}); no statement, label, number or tag moves);

Version 0.6.0-pre.2 --- 2026-10-08 (hecke-extensions slot 22, \textbf{a pre-release, not signed}: one wording repair in part~4's ``Not in this part'' sentence --- the Loewy-$4$ cells of [L37] are non-ordinary, not non-étale ordinary components; no statement, label, number or tag moves);

Version 0.6.0-pre.1 --- 2026-10-06 (hecke-extensions slot 22, \textbf{a pre-release, not signed}: new part~5, \S\ref{sec:u-tower} --- the tower $\Ga_1(\ell^{r})$, wild characters
and the $\Om{}$-type truncations, as refereed by F-U5; an addendum in Definition~\ref{def:u-lmsym} (labelling only); no statement, label, number or status tag of 0.5.0 moved;
see the changelog of part~5);

Version 0.5.0 --- 2026-10-06 (hecke-extensions slot 21, the close: \textbf{release} from v0.5.0-pre.8, released with the slot-21 close under the owner's decision D16
(accepted as recommended, 2026-10-06); part~4 as refereed by F4, F4m and the F4 of P5; pointers to [C] v1.21.0 and [E] 1.21.0; no statement, label, number or status tag moved; see the changelog of part~4);

Version 0.5.0-pre.8 --- 2026-10-06 (hecke-extensions slot 21, \textbf{a pre-release, not signed}: the Fable referee F4's repairs to the third remark after the proof of
Lemma~\ref{lem:u-twistvalues} --- ``$\lambda(G_b)=0$ iff $\ell\nmid B_b/b$'' is \PROVED{} given $\mu(G_b)=0$, and the step $\lambda_p=\sum_b\lambda(G_b)$ behind \cite{hho} is tagged \STD; no statement, label, tag or number changed; see the changelog of part~4);

Version 0.5.0-pre.7 --- 2026-10-06 (hecke-extensions slot 21, \textbf{a pre-release, not signed}: the slot-21 P5 result as a third unnumbered remark after the proof of
Lemma~\ref{lem:u-twistvalues} --- the Iwasawa invariants of $G_b$ at the irregular pairs used, re-derived from exact Bernoulli numbers; no statement, label, tag or number changed; see the changelog of part~4);

Version 0.5.0-pre.6 --- 2026-10-06 (hecke-extensions slot 21, \textbf{a pre-release, not signed}: the measure-side referee F4m's repairs to the two remarks after the proof of
Lemma~\ref{lem:u-twistvalues} --- the sign of the two factors, the residue's constant and the wild extension proved; no statement, label or number changed; see the changelog of part~4);

Version 0.5.0-pre.5 --- 2026-10-06 (hecke-extensions slot 21, \textbf{a pre-release, not signed}: the slot-21 P2 and P4 results as two remarks after the proof of
Lemma~\ref{lem:u-twistvalues} (no number moves) and one evidence sentence after Prediction~\ref{pr:u-msym}; no statement, label or tag changed; see the changelog of part~4);

Version 0.5.0-pre.4 --- 2026-10-05 (hecke-extensions slot 21, \textbf{a pre-release, not signed}: one evidence sentence after Prediction~\ref{pr:u-msym}, the $(103,24)$ env arm and the
absent $(691,12)$; see the changelog of part~4);

Version 0.5.0-pre.3 --- 2026-10-05 (hecke-extensions slot 21, \textbf{a pre-release, not signed}: status text only --- Proposition~\ref{prop:u-691}'s Bernoulli
numbers reproduced in-house (the validations named), Greenberg--Stevens read for (U16) (a weaker form; the tag stays), the slot-21 evidence of P8, P9 and P3 after
Prediction~\ref{pr:u-msym}, the code map of part~4 built; no statement, label or tag changed; see the changelog of part~4);

Version 0.5.0-pre.2 --- 2026-10-05 (hecke-extensions slot 21, the F4 repairs to part~4, \textbf{a pre-release, not signed}: the slot-21 Fable referee of part~4
(\file{docs/implementation/slot21/P8\_F4\_referee\_report.md}, ACCEPT WITH REPAIRS) applied --- (BF) over all $b\in\Z/(\ell-1)\Z$, Hida's $P_n=\Pk{n+2}$ read on p.~545, (U17)(a) as Harder's
index on the ordinary block (\PROVED{} from \STD{} for $n<\ell-1$), seven minors; no label added, renamed or removed; see the changelog of part~4);

Version 0.5.0-pre.1 --- 2026-10-05 (hecke-extensions slot 21, P8, \textbf{a pre-release, not signed}, for the slot-21 Fable referee: \textbf{part~4}, the
$\Iw$-adic modular symbols truncated, \S\ref{sec:u-msym} --- the ordinary modular-symbol modules per sign, that their ordinary parts do not see the lattice, multiplicity one at étale
components at finite weight and $\Iw$-adically, the normalisation by one functional with its unit check modulo $(\ell,X)$, the predictions (U15)--(U17); no statement of 0.4.0 changed;
see the changelog of part~4);

Version 0.4.0 --- 2026-10-05 (hecke-extensions slot 20, the close: \textbf{release} from v0.4.0-pre.3 under the owner's ruling O6 of 2026-10-04, on the
slot-20 referees F1 (ACCEPT) and F1$'$ (ACCEPT; its optional m-U1 applied); pointers to [C] v1.20.0 and [E] 1.20.0; no statement, label or status tag moved; see the changelog of part~3);

Version 0.4.0-pre.3 --- 2026-10-04 (hecke-extensions slot 20, the close delta for the Fable referee F1$'$, \textbf{a pre-release, not signed}: Prediction~\ref{pr:u-rungs} gains item~(4),
the outcome --- (U14) \textbf{met} at all three layer cells by slot-20 WP-RC20 (\COMP); no statement, label or tag moved; pointers [C] v1.20.0-pre.2, [E] 1.20.0-pre.6);

Version 0.4.0-pre.2 --- 2026-10-04 (hecke-extensions slot 20, WP-3 repairs after the Fable referee F1 (\file{docs/implementation/slot20/WP-3\_F1\_review.md}: [U] ACCEPT, the
optional m-U1 applied), \textbf{a pre-release, not signed}: Prediction~\ref{pr:u-rungs}(2) names $G_a\in\Iw$ as the series whose divided differences are integral; no statement, label or tag moved;
pointers [C] v1.20.0-pre.1, [E] 1.20.0-pre.2);

Version 0.4.0-pre.1 --- 2026-10-04 (hecke-extensions slot 20, WP-2, \textbf{a pre-release, not signed}, for the slot-20 Fable referee F1 (WP-3): part~3 gains
Prediction~\ref{pr:u-rungs} ((U14): the depth-$5$ rung at $37$, the depth-$4$ rungs at $103$ and $59$, and $\alpha_a$ modulo $\ell^{m}$, \PROVED{} from exact Bernoulli numbers, stated before
slot-20 WP-RC20's layer windows), and the $691$ depth-$3$ rungs read on the Bernoulli-free layer route (slot-19 NF691, \COMP); one label added (\texttt{pr:u-rungs}), none renamed or removed,
no tag moved; pointers [C] v1.20.0-pre.1, [E] 1.20.0-pre.1; see the changelog of part~3);

Version 0.3.0 --- 2026-10-04 (hecke-extensions slot 19, the close: \textbf{release} from v0.3.0-pre.4 under the owner's ruling O3, on the
slot-19 referees F1 (part~3, ACCEPT WITH REPAIRS), F4 (ACCEPT WITH REPAIRS) and F1-delta (ACCEPT); part~3 binding with the rest; pointers to [C] v1.19.0 and
[E] 1.19.0; no statement, label or status tag moved; see the changelog of part~3);

Version 0.3.0-pre.4 --- 2026-10-04 (hecke-extensions slot 19, WP-2 continued, \textbf{a pre-release, not signed}: the code map of part~3 \textbf{built} (slot-19
WP-HL, \code{HeckeLambdaValidation}); the F4 m10 ratio gate corrected in sign and count (WP-HL); the depth-$3$ and depth-$4$ rungs read from the Coates--Wiles measure
of the layer $\Q_\ell(\zeta_{\ell^{m}})$ (slot-19 R19-7, \COMP); no label added, renamed or removed, no tag moved; see the changelog of part~3);

Version 0.3.0-pre.3 --- 2026-10-04 (hecke-extensions slot 19, WP-2 continued, \textbf{a pre-release, not signed}: the slot-19 research results R19-2,
R19-3 folded into part~3 with the tags of the Fable referee F4 (ACCEPT WITH REPAIRS) --- Proposition~\ref{prop:u-eistrunc}'s depth-$3$ rungs \PROVED{} from five exact
Bernoulli numbers per class (F4's major M1) and R19-2's \COMP{} reading of (U12); (U13) \PROVED{} at $(37,32)$ and $(103,24)$ (F4 m8); Proposition~\ref{prop:u-691}'s
rungs gain a cross-reference to the $\Z/691^{2}$ Kummer fields; no label added, renamed or removed; see the changelog of part~3);

Version 0.3.0-pre.2 --- 2026-10-03 (hecke-extensions slot 19, WP-3 repairs, \textbf{a pre-release, not signed}: the slot-19 Fable referee F1's
six minors to part~3 (ACCEPT WITH REPAIRS, textual only) --- Lemma~\ref{lem:u-twistvalues}'s Washington locator verified on the page (the $u^{s}$ convention) and its
\VERIFYloc{} lifted; the semicolon notation of Definition~\ref{def:u-trunc}; Proposition~\ref{prop:u-lambdaprod}(3)'s generating set includes the weight tuple;
Lemma~\ref{lem:u-lambdaone}'s reading as $L_{k''}$ for $k''\ge2$ in the class; Proposition~\ref{prop:u-eistrunc}'s $691$ rungs computed (F1); (U11)'s checkable form;
no statement changed in substance, no label added, renamed or removed, no tag moved; see the changelog of part~3);

Version 0.3.0-pre.1 --- 2026-10-03 (hecke-extensions slot 19, WP-2, \textbf{a pre-release, not signed}, for the slot-19 Fable referee F1 (WP-3);
owner ruling O5 of 2026-10-03: \textbf{part~3}, the $\Iw$-adic algebra truncated modulo $(\ell^{K},X^{m})$, \S\ref{sec:u-lambda} --- what the truncation is, that
finitely many weights determine it and at which precision, the Eisenstein truncations at $37$, $103$, $691$, the predictions (U11)--(U13); the scope sentence of
\S\ref{sec:u-status} and \S\ref{sec:u-scope} amended; no statement of 0.2.1 changed; see the changelog of part~3);

Version 0.2.1 --- 2026-10-02 (hecke-extensions slot 17, the close: \textbf{release} from v0.2.1-pre.2 under the owner's ruling O3 of
2026-10-01; the referee F1 (slot-17 WP-3) was ACCEPT WITH REPAIRS, $0$ blocking, its three minors applied at pre.2; over pre.2 the [C] / [E] pointers only, now
[C] v1.17.0 and [E] 1.17.0; no statement, label or status tag moved; see the changelog of part~2);

Version 0.2.1-pre.2 --- 2026-10-02 (hecke-extensions slot 17, WP-2 resumed, \textbf{a pre-release, not signed}: the slot-17 referee F1's three minors
--- the $\ell^{2}$ argument of Proposition~\ref{prop:u-691} moved out of its \LIT{} citation, the placement sentence of \S\ref{sec:u-scope}'s \COMP{} reworded, the [C]
pointer note; no statement, label or status tag moved; see the changelog of part~2);

Version 0.2.1-pre.1 --- 2026-10-01 (hecke-extensions slot 17, WP-2, drafted during the slot-16 close, \textbf{a pre-release, not signed}: locators
and one \COMP{} sentence --- the last two Washington \VERIFYloc{} sites re-pointed to verified first-edition locators, and the slot-16 WP-11 sentence on the
Frobenius-layer hidden blocks beside \S\ref{sec:u-scope}'s (U-1); no statement, label or status tag of 0.2.0 moved; referee slot-17 WP-3; see the changelog of
part~2);

Version 0.2.0 --- 2026-10-01 (hecke-extensions slot 16, the close: \textbf{owner-signed release} from v0.2.0-pre.2, owner ruling O2 of
2026-10-01; the referee (slot-16 WP-14) was ACCEPT WITH REPAIRS, every repair applied at pre.2 and confirmed; over pre.2 locators only, verified by slot-16
WP-8; no statement, label or status tag moved; part~2 is binding with the rest; see the changelog of part~2);

v0.1.1 --- 2026-10-01 (\textbf{owner-signed}, ruling O2 of 2026-10-01: the text of v0.1.1-pre.6 signed as 0.1.1, no change of content; recorded at the
slot-16 close with 0.2.0; see the changelog);

v0.2.0-pre.2 --- 2026-10-01 (hecke-extensions slot 16, WP-12 resumed, \textbf{a pre-release, not signed}: the WP-14 referee's repairs
(ACCEPT WITH REPAIRS; U-1 the hypothesis $h_K=1$, U-2--U-7); see the changelog of part~2);

v0.2.0-pre.1 --- 2026-10-01 (hecke-extensions slot 16, WP-12, \textbf{a pre-release, not signed}: \textbf{part~2},
tame level $N$ --- where CM components can live, CM $\cap$ Eisenstein, CM $\cap$ $\GL(2)$, the predictions (U8)--(U10), \S\ref{sec:u-tameN}; proposal (U-1)
in \S\ref{sec:u-scope} for the referee; no statement of v0.1.x changed; referee slot-16 WP-14);

v0.1.1-pre.6 --- 2026-09-30 (slot-14 inventory patch list, \textbf{a pre-release, not signed}: $k=914$ in
Remark~\ref{rmk:u-zero}, and the referee's rewritten \COMP{} sentence after Prediction~\ref{pr:u-U2}; no statement, label or status tag moved;
see the changelog);

v0.1.1-pre.5 --- 2026-09-29 (slot-14, owner rulings of 2026-09-29 on the contract-currency audit, \textbf{a pre-release, not
signed}: a dated \textbf{erratum} to Proposition~\ref{prop:u-loewy}(4) --- nine non-reduced cells, not twelve; the (U4) ladder at $n=4150$ and
$(4830,2411)$; (U4) exactly sharp and (U7) on all $287$ weight-shift rungs; new Proposition~\ref{prop:u-HP} (Harder--Pink's sub-tower, a second route)
and Remark~\ref{rmk:u-dihedral}; no existing tag moved; see the changelog);

v0.1.1-pre.4 --- 2026-09-29 (slot-14 WP-C14, \textbf{a pre-release, not signed}: two \COMP{} additions as the Fable referee of
WP-C14 accepted them --- the depth-$2$ cells in Remark~\ref{rmk:u-zero}, and (U4)/(U7) at every $\ell$ with weight-shift data after \S\ref{sec:u-pred}'s list;
no statement, label or status tag moved; see the changelog);

v0.1.1-pre.3 --- 2026-09-28 (slot-14 WP-U1 audit repairs, \textbf{a pre-release, not signed}: the Fable audit's repairs R1--R6 ---
locators and wording only, every new locator verified on the page; no statement, label or status tag moved; see the changelog);

v0.1.1-pre.2 --- 2026-09-28 (slot-14 WP-U1, \textbf{a pre-release, not signed}: (U1) proved from Emerton's Theorem~5 --- both halves re-tagged
\PROVED{} from \LIT{}; the class $a\equiv0$ from Ohta's Theorem~(1.5.5); a Fable audit follows; see the changelog);

v0.1.1-pre.1 --- 2026-09-28 (slot-14 references edit, \textbf{a pre-release, not signed}: citations only ---
Emerton, Ohta and Citro--Ghitza verified against the owner's library; no statement, label or status tag moved; see the changelog);

v0.1.0 --- 2026-09-28 (hecke-extensions slot 14, package WP-13/Q5, \textbf{owner-signed
release}: from v0.1.0-pre.2, owner decision 2026-09-28 ``Go for publication and release'' ([S14] Q5 at the
recommended option). [U] becomes \textbf{binding}, like [C] and [E]: owner-released, cited by label, code
gated on it. Scope tame level $1$ only; tame $N\ge3$ is slot~15. No statement changed from pre.2; every
VERIFY locator stays marked VERIFY, including Theorem~\ref{thm:u-nocm}(b)'s dependence on Wiles. Release
note: \file{docs/implementation/slot14/WP-13\_release\_note.md} \S6);

v0.1.0-pre.2: 2026-09-28 (hecke-extensions slot 14, packages WP-8 and WP-9, \textbf{a pre-release,
non-binding} (owner question Q5 at its default): the first version with the referee's repairs. Reports:
\file{docs/implementation/slot14/WP-8\_report.md}, \file{docs/implementation/slot14/WP-9\_referee.md})

\subsection{Changelogs of the parts}
Part~1 is Sections~\ref{sec:u-status}--\ref{sec:u-scope}; parts~2--7 are Sections~\ref{sec:u-tameN}--\ref{sec:u-H2}. The changelogs are quoted verbatim in their original wording.

\paragraph{Changelog.} v0.1.0-pre.1 (2026-09-28, slot-14 WP-8): first version.
v0.1.0-pre.2 (2026-09-28, slot-14 WP-9, repairs half): the referee's findings applied. \textbf{B1}: Theorem~\ref{thm:u-etale}(3)'s second half needs
$\Cg_\Eis(a)\ne0$ (supplied at $k=12,200$ by $v(B_k/k)=1$). \textbf{B2}: Theorem~\ref{thm:u-691classes} no longer claims $J$ principal; the fibre product is over
$\Iw/J$, Remark~\ref{rmk:u-Jprincipal} (principal iff $\Tu(a)_{\mathfrak m_\Eis}$ is $\Iw$-free). \textbf{S1} Lemma~\ref{lem:u-deligne} (Deligne's bound), cited in
Propositions~\ref{prop:u-components}, \ref{prop:u-eisspec} and Theorem~\ref{thm:u-691classes}; \textbf{S2} $I_{P,k}$ defined (Definition~\ref{def:u-congmod}); \textbf{S3}
(U6) names $\Tu_{k,\mathfrak m}/\ell$ and is restricted to what is proved; \textbf{S4} (U7)'s blocks are cut by $U_\ell$ too; \textbf{S5} the weight congruence over
$\F_\lambda$; \textbf{S6} Proposition~\ref{prop:u-loewy}(2)'s trichotomy and (4), the ramified ordinary blocks of [M1] \S5.5, and (U2)'s wording. \textbf{m1--m7}: the $d_k$
warning and the $k\ge3$ convention, (U1)'s direction, the three-point Bernoulli route, the forms-side cross-check, the cluster-report pointers in (U4), Ribet's
construction. New \S\ref{sec:u-test}.
\textbf{v0.1.0} (2026-09-28, slot-14 WP-13/Q5, \textbf{owner-signed release}): [U] becomes \textbf{binding}, like [C] and [E] --- owner-released, cited by label,
code gated on it (owner decision 2026-09-28 ``Go for publication and release'', [S14] Q5 at the recommended option); \textbf{scope} recorded in
\S\ref{sec:u-status} and \S\ref{sec:u-scope}: tame level $1$ only, tame $N\ge3$ deferred to slot~15. Released \textbf{unbuilt} (\S\ref{sec:u-scope}'s code map,
unchanged). No statement changed from pre.2: every finding of the WP-9 referee stands as applied, and every \VERIFYloc{} locator (rule A2, [S14] Q6 at its
default) stays marked \VERIFYloc{}, \textbf{including Theorem~\ref{thm:u-nocm}'s no-CM proof, whose case (b) depends on Wiles \cite[Thm.~2.1.4]{wiles_ordinary}}
(case (a) and the rest of the theorem do not need it; the crystalline alternative of Remark~\ref{rmk:u-cmroute} remains a remark, itself \VERIFYloc{}). [C]
and [E] citations repointed to v1.14.0 and v1.15.0 (\S\ref{sec:u-status}), the versions released alongside this one. Report:
\file{docs/implementation/slot14/WP-13\_release\_note.md} \S6.
\textbf{v0.1.1-pre.1} (2026-09-28, slot-14 references edit, \textbf{not signed}): citations only, after the owner added Emerton's and Ohta's papers to
the library. Prediction~\ref{pr:u-U1}'s source line: the locators \cite[Thm.~5]{emerton_eis} and \cite[Thm.~(1.5.5)]{ohta_cong} verified and quoted,
``none is in the project's library'' corrected (Mazur--Wiles is there too, locator still \VERIFYloc); Proposition~\ref{prop:u-691}: the verified
cross-check \cite{citro_ghitza} for the two irregular classes of $691$; bibliography entries completed (\file{hida\_hecke\_algebra\_refs.bib}). No statement,
label or status tag moved: (U1) stays \CONJ. \emph{For the owner}: Emerton's Theorem~5 is, up to the dictionary between his $\mathbf T^{\mathrm{ord}}_{\mathrm{par},(i)}$,
$L_\ell(\omega^{i},1-s)$ and our $\hcusp(a)$, $G_a$ (not written out here), the $\Iw$-adic half of (U1), with no Gorenstein hypothesis; whether (U1) is
re-tagged is an owner decision (compare proposal P24 of [D1]).
\textbf{v0.1.1-pre.2} (2026-09-28, slot-14 WP-U1, \textbf{not signed}; owner decision: write the matching out, re-tag if it holds; a Fable audit follows).
New Lemma~\ref{lem:u-emerton} (the dictionary: Emerton's component $i$ is our class $a$ with no shift, $\mathbf T^{\mathrm{ord}}_{\mathrm{par},(a)}=\hcusp(a)$,
$I^{0}_{\Eis,(a)}=I^{0}(a)$, $L_\ell(\omega^{a},1-s)\mapsto G_a$), a proof of Prediction~\ref{pr:u-U1} (Steps~1--4: $P_\Eis$ is generated by the Eisenstein
differences; the kernel of $\Iw\to\Cg_\Eis(a)$ is both $J$ and Emerton's pull-back; his Theorem~5; Proposition~\ref{prop:u-eisspec}), and
Remark~\ref{rmk:u-ohta} (Ohta's theorem at $N=1$, its weight-$2$ shift, its main-conjecture route for $\psi=\mathbf1$, why no Gorenstein hypothesis
enters). \textbf{Tags moved}: Prediction~\ref{pr:u-U1} \CONJ{} $\to$ \PROVED{} from \LIT{}, both halves, (U1-$\Iw$) and (U1-$k$) for every $k\ge2$ in a class
$a\not\equiv0$, $\ell\ge5$ ((U1) of \S\ref{sec:u-pred} likewise); Remark~\ref{rmk:u-classzero} \CONJ{} $\to$ \PROVED{} from \LIT{} (Ohta's vanishing clause).
Remark~\ref{rmk:u-Jprincipal} decided ($J=(G_a)$; $\Tu(a)$ is $\Iw$-free); Theorem~\ref{thm:u-691classes}'s parenthetical, Remark~\ref{rmk:u-zero} and
\S\ref{sec:u-test}'s (U1) item reworded accordingly. No other statement moved. Supersedes [D1]'s P24 split (its (U1-$\Iw$) ``stays \CONJ'') and WP-Q2 \S4's
``Gorensteinness or freeness gives the other'' (dated addenda there). Standalone matching: \file{docs/implementation/slot14/WP-U1\_emerton\_matching.md}.
\textbf{v0.1.1-pre.3} (2026-09-28, slot-14 WP-U1 audit repairs, \textbf{not signed}): the Fable audit
\file{docs/implementation/slot14/WP-U1\_fable\_audit.md} (verdict: the re-tag stands) applied; every new locator read on the page. \textbf{R1}
Lemma~\ref{lem:u-emerton}(2): the identification $\mathbf T^{\mathrm{ord}}_{\mathrm{par},(a)}=\hcusp(a)$ made checkable --- \cite[(1.11a), (1.15a)]{hida_ens},
\cite[(1.2), (1.3), Thm.~1.3, Thm.~2.2]{hida_aif2} (new bib entry). \textbf{R2} a paragraph after the lemma recording what the tag rests on: Emerton's
Corollary~2 filled in (finite-level duality in weights $k\not\equiv0\pmod{\ell-1}$, Katz's density and $q$-expansion principle, \VERIFYloc) and Lemma~3 by
\cite[Thm.~2.2]{hida_aif2}. \textbf{R3} Remark~\ref{rmk:u-classzero}: the ``well known'' vanishing is \cite[Prop.~6.16, Prop.~13.22]{washington1st} plus Nakayama
(first edition, new bib entry; its \VERIFYloc{} removed); Ohta's inputs named ((3.4.2), the freeness of $X^{+}$, [15]~(5.1.2), [16]~(2.3.6)); the main
conjecture enters his Appendix~A only for $(\theta,\psi)\ne(\omega^{-2},\mathbf1)$; the identification of $I^{*}(\omega^{-2},\mathbf1)$ with $I^{0}(0)\otimes r$ and
faithfully flat descent. \textbf{R4} Remark~\ref{rmk:u-ohta}: the conductor $u=1$ for $a\equiv2$; ``Corollary~2 at $R=\Iw_{(a)}$'' in place of ``Proposition~1
and Corollary~2''. \textbf{R5} Lemma~\ref{lem:u-emerton}(4): the existence of $G_a\in\Iw$ is \cite[Thm.~7.10, p.~123]{washington1st} (verified; \cite[Thm.~5.11, p.~57]{washington1st} for the values (verified, slot-16 WP-8; this locator updated at v0.2.0)). \textbf{R6} a footnote: Emerton's p.~2 description of $V$ is not used. No statement, label or status tag moved.
\textbf{v0.1.1-pre.4} (2026-09-29, slot-14 WP-C14, \textbf{not signed}): two \COMP{} additions, as the Fable referee
\file{docs/implementation/slot14/wpc14/REFEREE\_I3b\_I7b\_I8\_I9.md}~\S5 accepted them. Remark~\ref{rmk:u-zero}: one sentence pointing to the depth-$2$
cells $k=284$ ($\ell=37$) and $k=228$ ($\ell=103$), computed on cohomology ([EX] \texttt{s14:depthlaw}; the law itself stays where it was, no new
proposition; $k=914$ at $\ell=59$ is not added, its rung not having finished). \S\ref{sec:u-pred}: a paragraph on (U4) and (U7) at every $\ell$ with
weight-shift data (WP-C14 I-8), with the referee's changes (sharpness at the new $\ell$ at $m=1$ only; $g_E=1$ at the $11$ classes as data). I-8's
proposed addition to Prediction~\ref{pr:u-U2} was rejected as worded by the referee and is not applied. No statement, label or status tag moved.
\textbf{v0.1.1-pre.5} (2026-09-29, owner rulings of that day on the Fable audit \file{docs/implementation/slot14/AUDIT\_contract\_currency\_2026-09-29.md},
\textbf{not signed}). \textbf{Q-A}: [U] is binding, as recorded at v0.1.0 (\S\ref{sec:u-status}; the [C] pointer moves to v1.14.1).
\textbf{Erratum} (dated 2026-09-29; slot-14 WP-14c \S6, refereed in \file{docs/implementation/slot14/wpc16/\allowbreak REFEREE\_J8\_WP-C16.md} Part~II, whose
corrected text II.3 is applied): Proposition~\ref{prop:u-loewy}(4) listed $12$ non-reduced rows of [M1]'s screened list; three of them,
$(46,383)$, $(50,677)$, $(54,67)$, are reduced ($v_\ell(\operatorname{disc}c_{k,2})=2$, the double root lifts modulo $\ell^{2}$, two residual systems),
carried over from WP-9's filter without the Loewy-length condition. The item now names the nine non-reduced cells and the three reduced rows; the
sentence on $(68,157)$ is kept (it was true as written) and sharpened; the item's conclusion (kind (ii) only, no cuspidal $\cap$ cuspidal instance) and
Prediction~\ref{pr:u-U2}'s stand. \S\ref{sec:u-test}'s (U3) item: ``the $12$ cells'' $\to$ the nine cells and the two blocks of $(68,157)$; the (U2)
item and Prediction~\ref{pr:u-U2} name the nine cells and $(68,157)$. \textbf{Q-B} (the second cluster run signed off): L5 of
\file{WP-L1\_binding\_doc\_proposals.md} --- (U4)'s ladder confirmed at $n=4150$ ($\ell=691$) and $n=4830$ ($\ell=2411$). [D1] \textbf{P19}: (U7)'s
parenthetical gains $\theta^{12}\Delta$ at $\ell=13$, and a \COMP{} paragraph after \S\ref{sec:u-pred}'s list gives (U7) on all $287$ rungs
($281$ / $353$ / $745$, the counts of [EX] \texttt{s14:sept} in place of the $266$-rung $260$ / $321$ / $607$). \textbf{P21}: in the same paragraph,
(U4) exactly sharp on the $7385$ pairs at $\ell\in\{7,11,13\}$ ($m=1$, $2$, $3$ on $6575$, $778$, $32$). \emph{The count}: [D1]'s $6662$ was the $266$-rung
run; the audit asked to re-derive $7385$ (WP-L1's L6) against I-8's $7412$. Both are right, on different sets: $7385$ is the weight-shift run's
pairs at $7$, $11$, $13$ (\file{wpl1/w1\_287/hida\_pairs.tsv}, reproduced row for row by I-8), and $7412=7385+27$ adds I-8's pairs at $23$, $31$,
$37$, $103$, $131$ (sharp at $m=1$ only), as pre.4's paragraph states. \textbf{P23}: new Proposition~\ref{prop:u-HP} (\PROVED{} from \LIT{Harder--Pink},
with the repair of WP-Q2 \S3.2), after Theorem~\ref{thm:u-L691}, as a second route to a special case of Prediction~\ref{pr:u-U1}; bibliography
entry \texttt{harder\_pink}. \textbf{J-8's [U] remark}: new Remark~\ref{rmk:u-dihedral} after Remark~\ref{rmk:u-cmroute} --- (i) \COMP{} (the
ordinary dihedral part has $\Fl$-dimension $h(-\ell)-1$ on $8$ primes; the $\Iw$-rank not claimed); (ii) with the referee's change: the non-ordinary
dihedral blocks are non-reduced, $\F_q[\varepsilon]/(\varepsilon^{2})$, anticyclotomic (\COMP), and not CM by the \emph{level} (\STD{}), not by
Theorem~\ref{thm:u-nocm}; Proposition~R stays in [EX]. \textbf{Labels added}: \texttt{prop:u-HP}, \texttt{rmk:u-dihedral}; none renamed or removed;
no existing status tag moved.
\textbf{v0.1.1-pre.6} (2026-09-30, slot-14 inventory patch list, \file{docs/implementation/slot14/INVENTORY\_findings\_vs\_latex\_2026-09-29.md}
\S3 items 6--7, \textbf{not signed}): two referee-accepted items that had not been written down. Remark~\ref{rmk:u-zero}: $k=914$ ($\ell=59$,
the cell $(912,59)$) joins the depth-$2$ cells computed on cohomology, as the referee's proviso asked once its rung had finished (it has,
[EX] \texttt{s14:depthlaw}). After Prediction~\ref{pr:u-U2}: the \COMP{} sentence in the referee's rewritten wording
(\file{wpc14/REFEREE\_I3b\_I7b\_I8\_I9.md} \S5, ``reject as worded, rewrite accepted''): in the $320$ weight-shift rungs every ordinary cuspidal
block has $g=1$, and the ramified rank-$2$ blocks of Proposition~\ref{prop:u-loewy}(4) are the next targets. No statement, label or status tag moved.
\textbf{v0.1.1} (signed 2026-10-01, \textbf{owner-signed}; owner ruling O2 of 2026-10-01, ``sign [U] 0.1.1''): the text of v0.1.1-pre.6 is signed as 0.1.1 with
\textbf{no change of content} --- pre.1--pre.6 above, among them (U1) proved from Emerton's Theorem~5 (pre.2, re-tag audited, pre.3), the erratum to
Proposition~\ref{prop:u-loewy}(4) (pre.5), Proposition~\ref{prop:u-HP} and Remark~\ref{rmk:u-dihedral} (pre.5). The signature was approved on 2026-10-01 together
with the decision to build part~2 on it and is recorded at the slot-16 close, with the release of 0.2.0 (\S\ref{sec:u-tameN}'s changelog); 0.2.0 is built on the 0.1.1 text, and every change it makes to it is listed in the
changelog of part~2. Report: \file{docs/implementation/slot16/contract\_release\_report.md}.

\paragraph{Changelog of part 2.} \textbf{v0.2.0-pre.1} (2026-10-01, slot-16 WP-12, \textbf{a pre-release, not signed}; referee WP-14): new
\S\ref{sec:u-tameN} --- Definition~\ref{def:u-hN}, Theorem~\ref{thm:u-controlN} (\LIT, \VERIFYloc), Theorem~\ref{thm:u-cmN} (\PROVED{} from \STD/\LIT),
Proposition~\ref{prop:u-cmexist}, Proposition~\ref{prop:u-cmeis}, Remark~\ref{rmk:u-cmgl2}, Remark~\ref{rmk:u-pred2} ((U8)--(U10)); in \S\ref{sec:u-scope}
the out-of-scope sentence points here, and proposal (U-1) of \file{docs/implementation/slot15/WP-12\_U\_E\_proposals.md} is added as a \CONJ{}
sentence \textbf{for the referee} (owner ruling O4 of 2026-10-01). Bibliography: \texttt{hida\_tilouine}. No statement of v0.1.x changed, no label renamed or
removed, no tag moved.
\textbf{v0.2.0-pre.2} (2026-10-01, slot-16 WP-12 resumed, \textbf{a pre-release, not signed}): the repairs answering the Fable referee
\file{docs/implementation/slot16/WP-14\_referee.md} ([U] ACCEPT WITH REPAIRS; $0$ blocking, $1$ major, $6$ minor), every one as worded. \textbf{U-1} (major)
Proposition~\ref{prop:u-cmeis}(2) assumes $h_K=1$ (else the congruence is with $E_k(\mu,\mu\chi_K)$), and (U10) records $h_K=1$ at its three fields. \textbf{U-2}
Proposition~\ref{prop:u-cmeis}'s set-up stated as three clauses on $\bar\rho_\lambda$, and $\bar\varepsilon:=\bar\rho_\lambda$ after unramifiedness. \textbf{U-3}
Theorem~\ref{thm:u-cmN}'s proof: the lift of $c$ chosen in $D_\ell$. \textbf{U-4} Wiles Thm.~2.1.4, p.~561, verified (L16-3) at Theorems~\ref{thm:u-cmN},
\ref{thm:u-nocm} and in the bib note; Theorem~\ref{thm:u-controlN} cites Hida Cor.~3.2 (p.~250) verbatim with the direct-factor sentence. \textbf{U-5} (U-1) in
\S\ref{sec:u-scope} re-tagged \PROVED{} for $\theta^{j}$-images (\HA{prop:thetatwist}), citing [EX] \texttt{s15:hidden} (\textsf{CMPD}). \textbf{U-6} Remark~\ref{rmk:u-cmgl2}(2):
one divisibility Hida--Tilouine, the converse Mazur--Tilouine; the \texttt{mazur\_tilouine} bib note amended. \textbf{U-7} (U10): both Eisenstein systems, dimension
$\ge4$; Theorem~\ref{thm:u-cmN}'s proof writes $\psi=\tilde\chi$. No label added, renamed or removed.
\textbf{v0.2.0} (2026-10-01, the slot-16 close, \textbf{owner-signed release}; owner ruling O2 of 2026-10-01, ``release [C] v1.16.0 and [U] 0.2.0 since WP-14
accepted and the repairs are confirmed; sign [U] 0.1.1''): released from v0.2.0-pre.2. The referee's repairs U-1--U-7 are confirmed, $16$ of $16$ findings with
[C]'s, by \file{docs/implementation/slot16/WP-14\_repair\_confirmation.md}. Over pre.2, \textbf{locators only} (slot-16 WP-8,
\file{docs/implementation/slot16/WP-8\_citations.md}~\S\S1--2, on L16-3's readings): Corollary~\ref{cor:u-irregular} cites Kummer's congruence as
\cite[Cor.~5.14, p.~61]{washington1st} and the weight congruence as \cite[Thm.~2(iv), p.~22]{swinnerton_dyer}, \cite[n$^\circ$~1.2, Th.~2, p.~416-04]{serre_bourbaki416}, all
verified (\VERIFYloc{} removed there); the paragraph after Lemma~\ref{lem:u-emerton} cites the same weight congruence, verified; v0.1.1-pre.3's R5 entry
records Washington's Theorem~5.11 (the values) at p.~57 of the first edition, verified. The status paragraph of \S\ref{sec:u-status} and this part's record the
signatures; [C] and [E] pointers move to v1.16.0 and v1.16.2. No statement, label or status tag moved. \textbf{Not in this release} (new mathematics, for a referee):
slot-16 WP-11's proposed sentence on the F1 hidden blocks beside Proposition~\ref{prop:u-loewy}, and any statement beyond the readings of (U8)--(U10).
\textbf{v0.2.1-pre.1} (2026-10-01, slot-17 WP-2, drafted during the slot-16 close, \textbf{a pre-release, not signed}; referee slot-17 WP-3, release as 0.2.1 at
the slot-17 close under the owner's ruling O3): \textbf{locators and one \COMP{} sentence}. (1) The two Washington \VERIFYloc{} sites the 0.2.0 release left
(\file{docs/implementation/slot16/contract\_release\_report.md}~\S2): Lemma~\ref{lem:u-emerton}'s Kubota--Leopoldt interpolation now cites
\cite[Thm.~5.11, p.~57]{washington1st} (verified, slot-16 WP-8), and Proposition~\ref{prop:u-691}'s Kummer congruence modulo $\ell^{2}$ cites
\cite[Thm.~5.11, Thm.~5.12, Cor.~5.14]{washington1st} with the two-line argument from Thm.~5.12 written out (Cor.~5.14 verified as a congruence modulo~$\ell$ by
L16-3; the $\ell^{2}$ form is the argument, not a quoted locator); the chain sentence of Proposition~\ref{prop:u-691} follows. (2) In \S\ref{sec:u-scope}, beside
(U-1): slot-16 WP-11's sentence on the Frobenius-layer hidden blocks $(220,17)$, $(246,19)$ --- not free over the full Hecke algebra, $\mathsf A_1=\Fl$ against
$\Talg_{\mathfrak m}/\ell=\Fl[x]/x^{2}$ --- as \COMP{} (gated by \code{HeckeNebentypusValidation}), placed there rather than inside Proposition~\ref{prop:u-loewy} because
the blocks are non-ordinary. No statement, label or status tag moved; [C] and [E] pointers stay at v1.16.0 and v1.16.2 (their 1.17.0-pre.1 drafts, with
\HA{rmk:sattor}, are in the tree for the same referee).
\textbf{v0.2.1-pre.2} (2026-10-02, slot-17 WP-2 resumed, \textbf{a pre-release, not signed}): the repairs of the slot-17 Fable referee F1
(\file{docs/implementation/slot17/WP-3\_referee\_F1.md}; [U] ACCEPT WITH REPAIRS, $0$ blocking, $3$ minor), each as worded. m-U1: in Proposition~\ref{prop:u-691} the
two-line $\ell^{2}$ argument from Thm.~5.12 stands in the text, the \LIT{} citation names only the theorems and what was verified (whether Washington's Cor.~5.14
states the $\ell^{a+1}$ form remains unread). m-U2: the \COMP{} sentence of \S\ref{sec:u-scope} says it is recorded beside (U-1), not at
Proposition~\ref{prop:u-loewy}, and why. m-U3: the citation paragraph notes that at the release the pointers read [C] v1.17.0 and the then-current [E]
(\HA{rmk:sattor} exists only from [C] 1.17.0-pre.1 on; 1.17.0-pre.3 is in the tree). No statement, label or status tag moved.
\textbf{v0.2.1} (2026-10-02, the slot-17 close, WP-18): \textbf{the release} from v0.2.1-pre.2 under the owner's ruling O3 of 2026-10-01, as v0.2.0 was released
under O2. Over pre.2 only the citation paragraph's pointers, now [C] v1.17.0 and [E] 1.17.0 (m-U3). No statement, label or status tag moved.

\paragraph{Changelog of part 3.} \textbf{v0.3.0-pre.1} (2026-10-03, slot-19 WP-2, \textbf{a pre-release, not signed}; referee slot-19 WP-3 (Fable F1); owner ruling O5
of 2026-10-03): new \S\ref{sec:u-lambda} --- Definition~\ref{def:u-trunc}, Proposition~\ref{prop:u-truncbasic}, Lemma~\ref{lem:u-wtprod}, Theorem~\ref{thm:u-interp},
Proposition~\ref{prop:u-lambdaprod}, Lemma~\ref{lem:u-twistvalues}, Lemma~\ref{lem:u-lambdaone}, Proposition~\ref{prop:u-eistrunc} (\COMP: \file{docs/implementation/slot19/wp2/lambda\_check.py}),
Remark~\ref{rmk:u-lambdabasis}, Prediction~\ref{pr:u-lambda} ((U11), (U12) \PROVED, (U13) \PROVED{} at $691$ and \CONJ{} at $37$, $103$). In \S\ref{sec:u-status} and
\S\ref{sec:u-scope} the out-of-scope sentence on the $\Iw$-adic algebra is amended (its truncations are in scope; the whole algebra is not), and the code map
points here. Bibliography: \texttt{hho}. No statement of v0.2.1 changed, no label renamed or removed, no tag moved.
\textbf{v0.3.0-pre.2} (2026-10-03, slot-19 WP-3 repairs, \textbf{a pre-release, not signed}; report \file{docs/implementation/slot19/WP-3\_F1\_referee\_report.md},
repair table \file{docs/implementation/slot19/WP-3\_repairs.md}): m-U1 Lemma~\ref{lem:u-twistvalues}, Thm.~7.10 (1st ed., p.~123) quoted, even $\chi$, $1+q_0=u$, both
\VERIFYloc{} removed (Proposition~\ref{prop:u-eistrunc}'s \cite{hho} stays \VERIFYloc, unused); m-U2 Definition~\ref{def:u-trunc}, $\Tu_{K,m}(a;x)$ etc.; m-U3
Proposition~\ref{prop:u-lambdaprod}(3), the $\Iw$-subalgebra with the weight tuple; m-U4 Lemma~\ref{lem:u-lambdaone}, $k''\ge2$ in the class; m-U5
Proposition~\ref{prop:u-eistrunc}, F1's power-sum route and scans, \file{docs/implementation/slot19/wp3/f1\_lambda.py} + \file{.out} tracked; m-U6 (U11)'s checkable
form, and the $\tilde D_i$ in $\Q_\ell$ in Theorem~\ref{thm:u-interp}(2). Tags unchanged.
\textbf{v0.3.0-pre.3} (2026-10-04, slot-19 WP-2 continued, \textbf{a pre-release, not signed}; the slot-19 research results R19-2, R19-3, refereed by the Fable
referee F4, \file{docs/implementation/slot19/R19-6\_F4\_referee\_report.md}, ACCEPT WITH REPAIRS; change table \file{docs/implementation/slot19/WP-2\_pre3\_changes.md}):
Proposition~\ref{prop:u-eistrunc} gains the depth-$3$ rungs $37580$, $914250$, $49928412$, $159113510$, \PROVED{} from five exact Bernoulli numbers per class with
their proof (F4's major M1: R19-2 had the two $691$ rungs as ``predicted''), and the \COMP{} reading of (U12)'s checkable form by R19-2 (F4 m10's sign of $F_1$ and
its ratio gate); (U13) \CONJ{} $\to$ \PROVED{} at $(37,32)$, $(103,24)$ (F4 m8); Proposition~\ref{prop:u-691} a cross-reference to R19-3's $\Z/691^{2}$ Kummer fields;
Remark~\ref{rmk:u-lambdabasis}'s \COMP{} list; the code map of part~3: R19-2's route recommendation and F4 m9 (the ordinary part). No label added, renamed or removed;
one tag moved ((U13) at $37$, $103$).
\textbf{v0.3.0-pre.4} (2026-10-04, slot-19 WP-2 continued, \textbf{a pre-release, not signed}; change table \file{docs/implementation/slot19/WP-2\_pre5\_changes.md}): the code
map of part~3 \textbf{built} (slot-19 WP-HL; arm~(Q) primary by the owner's ruling of 2026-10-03; \code{HeckeLambdaValidation}, the planned name superseded) and a \emph{Code} line in
Prediction~\ref{pr:u-lambda}; the F4 m10 ratio corrected to $F_1/G_{32}\equiv2$ at eight weights (WP-HL); a \COMP{} paragraph after Proposition~\ref{prop:u-eistrunc}: the rungs and
$\alpha$ modulo $\ell^{m}$ read from the level-$m$ Coates--Wiles measure (slot-19 R19-7), the depth-$4$ rung $1072544$ at $(37,32)$. No label added, renamed or removed; no tag moved.
\textbf{v0.3.0} (2026-10-04, slot-19 close, WP-18; \textbf{the release} from v0.3.0-pre.4 under the owner's ruling O3): the slot-19 F1-delta read the pre.3 $\to$ pre.4 delta of [U]
as ACCEPT (\file{docs/implementation/slot19/F1delta\_referee\_report.md}; its optional m-U1 not applied, \file{docs/implementation/slot19/F1delta\_repairs.md}). Release lines
only: part~3 binding with the rest; the citation paragraph's pointers read [C] v1.19.0 and [E] 1.19.0. No statement, label or status tag moved.
\textbf{v0.4.0-pre.1} (2026-10-04, slot-20 WP-2, \textbf{a pre-release, not signed}; referee slot-20 WP-3 (Fable F1); list of new and changed labels
\file{docs/implementation/slot20/WP-2\_drafts.md}): a new Prediction~\ref{pr:u-rungs} ((U14)) --- at $(37,32)$, $m=5$, the rung $55777784$ and $\alpha_{32}\equiv13400697\pmod{37^{5}}$; at
$(103,24)$, $m=4$, the rung $60430740$ and $\alpha_{24}\equiv3736634\pmod{103^{4}}$; at $(59,44)$ $\lambda=1$ (\COMP) and, $m=2,3,4$, the rungs $914$, $86464$, $8162384$ with
$\alpha_{44}\equiv11649137\pmod{59^{4}}$ --- \PROVED{} from exact Bernoulli numbers (\file{docs/implementation/slot20/wp2/rung\_predict.py}), with the checkable form that slot-20 WP-RC20's layer
windows gate; the paragraph after Proposition~\ref{prop:u-eistrunc}: the $691$ depth-$3$ rungs and $\alpha$ modulo $691^{3}$ read on the level-$3$ layer (slot-19 NF691, \COMP). One label added;
none renamed or removed; no tag moved; no golden line, digest or record key touched.
\textbf{v0.4.0-pre.2} (2026-10-04, slot-20 WP-3 repairs, \textbf{a pre-release, not signed}): the slot-20 Fable referee F1 (\file{docs/implementation/slot20/WP-3\_F1\_review.md}) read [U]
as ACCEPT and reproduced (U14) by an independent route; its optional m-U1 is applied (Prediction~\ref{pr:u-rungs}(2): ``the divided differences of $G_a\in\Iw$''). No statement, label or
status tag moved (\file{docs/implementation/slot20/WP-3\_repairs.md}).
\textbf{v0.4.0-pre.3} (2026-10-04, slot-20 close delta, \textbf{a pre-release, not signed}; referee F1$'$; list \file{docs/implementation/slot20/close\_delta.md}): Prediction~\ref{pr:u-rungs}
gains item~(4), the outcome of slot-20 WP-RC20's layer windows --- (U14) \textbf{met} at $(37,32,5)$, $(103,24,4)$, $(59,44,4)$ (\COMP). No statement changed; no label added, renamed or
removed; no tag moved.
\textbf{v0.4.0} (2026-10-05, slot-20 close, WP-18; \textbf{the release} from v0.4.0-pre.3 under the owner's ruling O6 of 2026-10-04): the slot-20 F1$'$ read the pre.2 $\to$ pre.3 delta of [U]
as ACCEPT (\file{docs/implementation/slot20/F1prime\_review.md}); its optional m-U1 (the cost figures of Prediction~\ref{pr:u-rungs}(4) against the \code{time -v} peaks) was applied by slot-20
WP-17 (\file{docs/implementation/slot20/WP-17\_repairs.md}). Release lines only: the citation paragraph's pointers read [C] v1.20.0 and [E] 1.20.0. No statement, label or status tag moved.

\paragraph{Changelog of part 4.} \textbf{v0.5.0} (2026-10-06, slot-21 release, the close): released with the slot-21 close from v0.5.0-pre.8 under the owner's decision D16
(accepted as recommended); release lines only --- the status paragraph's sentence on part~4, the pointers [C] v1.21.0 and [E] 1.21.0, the date line. No statement, label, number or status tag moved.
\textbf{v0.5.0-pre.8} (2026-10-06, slot~21, \textbf{a pre-release, not signed}; the Fable referee F4 of P5, \file{docs/implementation/slot21/rc/F4\_referee\_p5\_report.md}):
the third remark after Lemma~\ref{lem:u-twistvalues} repaired --- ``$\lambda(G_b)=0$ iff $\ell\nmid B_b/b$'' \PROVED{} given $\mu(G_b)=0$; the step $\lambda_p=\sum_b\lambda(G_b)$ tagged \STD{} (Iwasawa, Vandiver
from \cite{hho} item~(2)); \cite{hho} \S7 read by the referee; the seven pairs re-derived with the referee's own code. No statement, label, tag or number changed.
\textbf{v0.5.0-pre.7} (2026-10-06, slot~21, \textbf{a pre-release, not signed}; slot-21 P5, \file{docs/implementation/slot21/P5\_spec.md} v0.2): a third
unnumbered remark after Lemma~\ref{lem:u-twistvalues} (part~3) --- $\lambda(G_b)=1$ at the irregular pairs used (\LIT{\cite{hho}}), $\mu(G_b)=0$ (\STD{Ferrero--Washington, recalled}), both
re-derived from exact Bernoulli numbers (\COMP). No statement, label, tag or number changed. \textbf{v0.5.0-pre.6} (2026-10-06, slot~21, \textbf{a pre-release, not signed}; the measure-side Fable referee F4m,
\file{docs/implementation/slot21/rc/F4\_referee\_measure\_report.md}): the two remarks after Lemma~\ref{lem:u-twistvalues} (part~3) repaired --- (i) each $G$-factor is \emph{minus} the branch
series ($E_1=-x\,\zeta_\ell$; the product unchanged; Washington Thm.~12.2 read on the page), the residue's constant \PROVED{} from the factorisation; (ii) the wild extension written out as a proof
from Thm.~7.10 / 5.11 (\COMP{} $\to$ \PROVED, with \COMP). No statement, label or number changed. \textbf{v0.5.0-pre.5} (2026-10-06, slot~21, \textbf{a pre-release, not signed}): the slot-21 P2 and P4 results --- two unnumbered remarks after the proof of
Lemma~\ref{lem:u-twistvalues} (part~3; the $G$-factors of the Eisenstein measure, and the wild extension $G_b(\zeta^{t}u^{r}-1)=-B_{r,\omega^{b-r}\psi_t}/r$, \COMP) and one sentence in the Evidence
paragraph (single weights: no evidence for (U16), (U17)). No statement, label or tag changed; no number moves. \textbf{v0.5.0-pre.4} (2026-10-05, slot~21, \textbf{a pre-release, not signed}): the Evidence paragraph gains the $(103,24)$ env arm ($j_k=L_k$ to $k=738$,
$L=2$ at $228$) and the note that F4's $(691,12)$ is not in the committed env arm. No statement, label or tag changed. \textbf{v0.5.0-pre.3} (2026-10-05, slot~21, \textbf{a pre-release, not signed}): the code map \textbf{built} (P8, P9); an \emph{Evidence} paragraph
after Prediction~\ref{pr:u-msym} (P8, P9, P3; \COMP; no tag moves); (U16)'s literature sentence --- Greenberg--Stevens read, a weaker (local, characteristic-$0$, meromorphic) form,
not the integral one. No statement, label or tag changed. \textbf{v0.5.0-pre.2} (2026-10-05, slot-21 F4 repairs after the Fable referee of part~4, \textbf{a pre-release, not signed};
\file{docs/implementation/slot21/P8\_F4\_referee\_report.md}, ACCEPT WITH REPAIRS, $0$ blocking, $3$ major, $7$ minor, all applied): \textbf{M1} Proposition~\ref{prop:u-msymfree}(2): (BF)
quantifies over all $b\in\Z/(\ell-1)\Z$ (the pre.1 range $0\le b\le\lfloor n/\ell\rfloor$ rested on the \CONJ{} half of \HA{cl:chi}(2); the proof now cites \HA{lem:bdrytorus}(2),(3)); \textbf{M2}
Proposition~\ref{prop:u-lmsymcontrol}: Hida's $P_n$ read on p.~545, $P_{k-2}=\ker(X\mapsto u^{k}-1)=\Pk{k}$, \VERIFYloc{} lifted (the only remaining \VERIFYloc{} of part~4 is the matrix $\varepsilon$ of
Lemma~6.4); \textbf{M3} (U17)(a) identified with \Espec{cl:lveis}(2) on the ordinary $\mathfrak m_\Eis$-block, hence \PROVED{} from \STD{} for $n<\ell-1$ ($k=32$ at $37$, $k=24$ at $103$,
$(10,691)$) and \CONJ{} beyond, with $\delta(\mathfrak m_\Eis)=0\iff j_k=L_k$; \textbf{m1} Lemma~\ref{lem:u-latticeord}(1): $V'_n$ is $\mathrm{Mat}_2(\Z)$-stable (divided powers), not only
$\De$-stable; \textbf{m2} its (3): the exactness of $e_1$ on finitely generated $\Zl$-modules written out; \textbf{m3} its (3): Hida's level-one operators agree with [C]'s on the
torsion-free ordinary part (Eichler--Shimura on both sides); \textbf{m4} Theorem~\ref{thm:u-lmsymfree}: the identification of the algebra acting on $\mathbf M(a)$ with $\hcusp(a)$ written out;
\textbf{m5} Lemma~\ref{lem:u-msymmass}: $U$ on $e_1H^{1}_c$ by Hensel; \textbf{m6} Remark~\ref{rmk:u-gorenstein}: a pointer sentence (the literature transfer now applies on ordinary parts);
\textbf{m7} Lemma~\ref{lem:u-msymsigns}: Lemma~6.4's statement quoted (``for each positive integer $M$''), classical at $M=1$. Remark~(d) after Lemma~\ref{lem:u-latticeord}; the status
paragraph; Remark~\ref{rmk:u-msymbasis}'s \VERIFYloc{} line. No statement of 0.4.0 changed except the pointer sentence of Remark~\ref{rmk:u-gorenstein}; no label added, renamed or removed; no tag of 0.4.0 moved.
\textbf{v0.5.0-pre.1} (2026-10-05, slot-21 P8, \textbf{a pre-release, not signed}; for the slot-21 Fable referee): new \S\ref{sec:u-msym} ---
Definition~\ref{def:u-msymk}, Lemma~\ref{lem:u-msymsigns}, Lemma~\ref{lem:u-latticeord} (new: ordinary parts are lattice-independent), Proposition~\ref{prop:u-msymfree}, Definition~\ref{def:u-lmsym},
Proposition~\ref{prop:u-lmsymcontrol}, Theorem~\ref{thm:u-lmsymfree}, Lemma~\ref{lem:u-msymmass}, Proposition~\ref{prop:u-msymnorm}, Prediction~\ref{pr:u-msym} ((U15) \PROVED, (U16), (U17) \CONJ),
Remark~\ref{rmk:u-msymbasis}, the code map. Bibliography: \texttt{greenberg\_stevens} (not read). Hida, Invent.\ 85: Lemma~6.4 (pp.~587--588), Cor.~4.5 (p.~570), Prop.~4.7 (p.~572), Thm.~6.6 (p.~590)
read on the page. No statement, label or status tag of 0.4.0 changed; the status paragraph of \S\ref{sec:u-status} gains one sentence.

\paragraph{Wording repair at 0.6.0-pre.2 (part~4)} (2026-10-08, slot~22 DOC22; not signed). Part~4's ``Not in this part'' sentence listed ``the ${\LL}=4$ cells of [L37]'' among the non-étale components. Those cells (type $\mathrm{eis}(a,a{+}5)$, $\ell=37$) are non-ordinary ($U_{37}$ nilpotent, [L37] \texttt{prop:nonord}; indeed $U_{37}\in\mathfrak m_A^{2}$, [L37] \texttt{prop:Uposition}), so they are not objects of [U]; the entry is removed and the point stated in a parenthesis. Source: \file{docs/implementation/slot22/nonreduced\_ordinary\_survey.md}. Nothing else changed.

\paragraph{Release candidate 0.6.0} (2026-10-12, slot~22 CLOSE-T; \textbf{owner sign-off pending}). The text of 0.6.0-pre.1 to pre.3 as stated in their entries (part~5 refereed by F-U5, \file{P10C\_F4\_referee\_report.md}; the pre.2 and pre.3 wording points read by \file{docs/implementation/slot22/F-T1\_release\_referee.md}); only the version is changed. Companions [C] v1.22.0, [E] 1.22.0.

\paragraph{Changelog of part 5.} \textbf{v0.6.0-pre.1 as inserted} (2026-10-06, slot~22, \textbf{a pre-release, not signed}): this part is the refereed text v0.2 (whose entries follow, labelled by the text's own revision
names); in this document it is v0.6.0-pre.1. Inserted verbatim except two cross-document references to [C] (\verb|\eqref{eq:MHchi}| in Proposition~\ref{prop:u-towerorient}'s status, \verb|[C] \S\ref{sec:nebwildsec}| in the
entry below) written as \verb|\HA{...}|, since a \verb|\ref| into [C] is undefined here; the addendum of \file{docs/implementation/slot22/P10C\_addenda.md}~(2) in Definition~\ref{def:u-lmsym} (labelling only, verified on Hida p.~593 by F-U5);
the pointer in Remark~\ref{rmk:u-msymbasis}; one sentence in the status paragraph of \S\ref{sec:u-status}; bibliography key \texttt{miyake} (Thm.~4.6.17, locator not verified). No statement, label, number or status tag of 0.5.0 changed. \textbf{v0.6.0-pre.2} (2026-10-06, slot~22, \textbf{a pre-release, not signed}; the Fable referee F-U5, \file{docs/implementation/slot22/P10C\_F4\_referee\_report.md}, ACCEPT WITH REPAIRS, $0$ blocking / $3$ major / $13$ minor, all applied): \textbf{M1} [C] \HA{lem:nebshapiro}'s $C$-action on $P$ through the inverse of the quotient map (the application unchanged); \textbf{M2} Theorem~\ref{thm:u-towercontrol}(2) restated as a module isomorphism, Cor.~4.5's $e\circ\mathrm{res}$ is $e\circ\mathrm{dg}$ in [C]'s convention; \textbf{M3} (1), (5), Proposition~\ref{prop:u-towereis}, (U18)(a)--(c), (U20), \HA{NB8} tagged \PROVED{} \emph{given} Proposition~\ref{prop:u-towerorient}; the route of record for (1) is the Thm.~1.2 + Nakayama argument with the named \STD{} input (ordinary $=$ old); \textbf{m1--m3} the orientation remark reworded (no formal cancellation; the computation decides); \textbf{m8} the boundary systems cite \STD{} not \HA{lem:bdrytorus}; \textbf{m10} Proposition~\ref{prop:u-towerfree}(1) proved directly. Hida pp.~549, 552, 553, 562, 570, 571, 589, 590, 593 re-read on the page images by the referee. No statement of 0.5.x changed. \textbf{v0.6.0-pre.1} (2026-10-06, slot~22 P10-C, \textbf{a pre-release, not signed}; for the Fable referee F-U5): new \S\ref{sec:u-tower} ---
Definitions~\ref{def:u-towerlambda}, \ref{def:u-towerobj}, Lemma~\ref{lem:u-omegar}, Proposition~\ref{prop:u-towerorient}, Theorem~\ref{thm:u-towercontrol}, Propositions~\ref{prop:u-towereis}, \ref{prop:u-towerfree},
Corollary~\ref{cor:u-towertrunc}, Prediction~\ref{pr:u-tower} ((U18)--(U21) \PROVED, (U22) \CONJ), Remarks~\ref{rmk:u-towerbasis}, \ref{rmk:u-towerkitagawa}, the code map; one pointer in
Remark~\ref{rmk:u-msymbasis}. Hida, Invent.\ 85, pp.~549, 552--553, 562, 570, 593 read on the page. Companion: \HA{sec:nebwildsec} (v1.22.0-pre.1). No statement, label or status tag of 0.5.x
changed.

\paragraph{Changelog of part 6.} \textbf{v0.7.0-pre.1} (2026-11-06, slot~25 U6, \textbf{a pre-release, not signed}; for an adversarial referee): new \S\ref{sec:u-eislambda} ---
Lemma~\ref{lem:u-H0}, Remark~\ref{rmk:u-H0census}, Proposition~\ref{prop:u-H11}, Proposition~\ref{prop:u-H12}, Remark~\ref{rmk:u-H12lit} (the owner-signed P-K2 of 2026-11-05),
Prediction~\ref{pr:u-H1} ((U23) \PROVED{} given data, (U24) \CONJ), Remark~\ref{rmk:u-H3prime} (\COMP), the code map; the mathematics of the slot-24 HeckeZ stage 05 v0.3.1 \S1.2,
\S2.2--\S2.4 (Lemma H0, Props.~H1.1, H1.2), re-proved here with two changes: Proposition~\ref{prop:u-H12}(3) is proved from the fibre product directly (the census row (D4) becomes a
check), and (6) adds the sign $\varepsilon$ ($H^{\varepsilon}\cong\Tu_{486}(a)_{\mathfrak m}$). Elsewhere: Remark~\ref{rmk:u-classzero} credits Ohta's adaptation of Harder--Pink and
Kurihara (P-K1, owner-signed 2026-11-05); Proposition~\ref{prop:u-loewy}(3) gains the Bella\"\i che--Pollack citation (phenomenology catalogue \S11 item~4, \textbf{not yet seen by the
owner}); bibliography keys \texttt{kurihara1993}, \texttt{calegari\_eisdef}, \texttt{bellaiche\_pollack}, \texttt{bruns\_herzog\_cm}. No statement, label, number or status tag of
0.6.0 moved. Report: \file{docs/implementation/slot25/U6\_report.md}.
\textbf{v0.7.0-pre.2} (2026-11-06, slot~25 U6 fix round, \textbf{a pre-release, not signed}): the referee \file{docs/implementation/slot25/U6\_referee.md} (ACCEPT WITH MINORS, F1--F14, no
MAJOR) applied; no result changes. (U23): the degree-$2$ clause on all of $S_{k^{*}}$ is \COMP{} (the referee's computation), the block statements stay \PROVED{} (F1); (U24) restated as
``$\mathbf M(a)^{\pm}_{\mathfrak m}$ free over $\hcusp(a)_{\mathfrak m}$, both signs'' (F2); Proposition~\ref{prop:u-H12}, Step~1 reordered, $g_a=2$ from Corollary~\ref{cor:u-irregular} and
conjugation, (D2) a check (F3); Remark~\ref{rmk:u-H0census} gains the block-dimension condition (F4) and the locator Thm.~2.6 (F12); Lemma~\ref{lem:u-H0}(3), proof: $t$ translated
first (F5); Proposition~\ref{prop:u-H12}'s tag completed, $\varepsilon=+$ via Lemma~\ref{lem:u-msymsigns} (F6); Remark~\ref{rmk:u-H12lit}: the owner-signed scope limited to the one sentence,
pp.~282--283, ``unconditionally'' (F7); the notation of $\mathfrak m$ (F9); Remark~\ref{rmk:u-H3prime}: ``two of the three kinds'' (F10); Proposition~\ref{prop:u-loewy}(3): the
tag of the inserted sentences and ``arXiv v1, p.~2'' (F11). \textbf{Released front matter, not yet seen by the owner}: the status paragraph of \S\ref{sec:u-status} no longer calls
part~5 a pre-release (``since v0.6.0 \dots\ released with the rest'', F8). With Z0 on \texttt{hecke-extensions}: \HA{prop:heckeZgen}, \HA{sec:heckeZ} cited by label; the
bibliography key \texttt{bruns\_herzog\_cm} replaced by [C]'s \texttt{bruns\_herzog}. No statement number or label of 0.7.0-pre.1 moves.
\textbf{v0.7.0-pre.3} (2026-11-06, slot~25, the H1 merge after its audit \file{docs/implementation/slot25/H1\_audit.md}, PASS\_WITH\_MINORS; \textbf{a pre-release, not signed,
not yet seen by the owner}): Prediction~\ref{pr:u-H1} (U23) gains the explicit anchor $x^{2}+56702273x+83587893\bmod547^{3}$, $v(\operatorname{disc})=2$, before H1-R runs (F9);
new Remark~\ref{rmk:u-H12general}, Steps~1--2 of Proposition~\ref{prop:u-H12} at any $g_a\ge2$ (F5), placed last so that no statement number moves; the code map \emph{Built}. No result,
label or status tag of pre.2 moves.

\paragraph{Changelog of part 7.} \textbf{v0.7.4-pre.1} (2026-11-11, slot~26 Z6, \textbf{a pre-release, not signed}; for a referee, before the code): new \S\ref{sec:u-H2} --- Definition~\ref{def:u-DF}, Lemmas~\ref{lem:u-DF}, \ref{lem:u-DFread}, Prediction~\ref{pr:u-H2} (U25), the code map; from the slot-24 HeckeZ stage 05 v0.3.2 \S4, whose two-way question (``$\ell\cdot$unit or a zero'') is replaced by Lemma~\ref{lem:u-DFread}'s cases (A), (B), and whose ``$\cong\Zl[[Y]]$'' holds only in case (A) (Lemma~\ref{lem:u-DF}(4)).
\textbf{v0.7.4-pre.2} (2026-11-11, slot~26 Z6 fix round, \textbf{a pre-release, not signed}): the referee's MINOR items --- m1 the generality of Lemma~\ref{lem:u-DF}(3) in
Definition~\ref{def:u-DF}, (U25)(c) from Proposition~\ref{prop:u-H12}(1), (4), (5) and Lemma~\ref{lem:u-DFread}(A); m2 the proof of Lemma~\ref{lem:u-DFread}(R): $\tilde D$ and $\tilde D+W^{m}$
instead of $D+X^{M}$ (the statement was right); m3 the centred use of Theorem~\ref{thm:u-interp} in (U25)(a); m4 $D_{T(2)}$ in (U25)(a); m5 which second generator exists; m6 the abstract
(``eleven'', not ``the eleven''); m7 $W:=X-x_c$ in Lemma~\ref{lem:u-DFread}(R) and (U25)(a) ($Z$ keeps the meaning of Lemma~\ref{lem:u-DF}(4)); m8 Lemma~\ref{lem:u-levelone} in the tag of
Lemma~\ref{lem:u-DF} and the proof of (2); m9 the normalisation in the proof of (5); m10 the release line of part~6 in \S\ref{sec:u-status} (signed 0.7.3 text, see the release record).
Notes: n1 the Hida locator (Thm.~0.1, p.~329) verified, \VERIFYloc{} removed, bibliography entry \code{hida\_ajm88}; n3 a generator always exists (Definition~\ref{def:u-DF}); n4 ``$\ge3$''
and $k_0$ in (U25)(a); n6 (U25) can fail only through a defect (after (U25)). Not adopted: n2 (swapping Lemmas~\ref{lem:u-DF} and \ref{lem:u-DFread} would move numbers); n5 is moot (no
block is undecided). Code audit F1: the subset agreement in (U25)(a) is stated as a consequence of integrality, not as evidence. New Remark~\ref{rmk:u-H2data} (\COMP) with the eleven rows,
the finding that $(36,181)$ is in case (B), $\lambda=2$, and the code map \emph{Built}.
\textbf{v0.8.0} (2026-11-11, slot~26 Z6, \textbf{signed by the owner 2026-11-11}): the owner-signed release of v0.7.4-pre.2; text unchanged; no statement, label or value moves; numbered 0.8.0 per the referee's m11, accepted by the owner; see the release record.
\textbf{v0.8.1} (2026-11-11, slot~26, \textbf{signed by the owner 2026-11-11}): two status sentences of part~7 (\S\ref{sec:u-status} and the ``Status and scope'' paragraph) say ``released'' instead of ``pre-release, not signed''; no statement, label or value moves; see the release record.

\subsection{The title page of 0.7.0-pre.3: author line and abstract, verbatim}
\begin{quote}
Orbifold Cohomology of $\GL(2)$\\ companion to \texttt{gl2\_contract.tex} [C] and \texttt{hecke\_extensions.tex} [E]
\end{quote}
\begin{quote}
Let $\ell\ge5$. At tame level~$1$, Hida's ordinary Hecke algebra $\Tu$ is a finite torsion-free algebra over
$\Iw=\Zl[[X]]$, split into the even weight classes $a$ modulo $\ell-1$; its irreducible components are one
Eisenstein component per class and the components of the cuspidal algebra $\hcusp(a)$. This document (i) fixes
these objects and reads them at finite weight (\S\ref{sec:u-setting}--\S\ref{sec:u-finite}): every quotient
$M/(P_k,\ell^m)M$ of a finitely generated $\Iw$-module depends only on the weight tower $\mathrm{Tw}_m$ of [C], so
the characteristic polynomials of the $T_p$ on the ordinary cusp forms modulo $\ell^m$, the mod-$\ell$ local
algebras and the truncated congruence lengths are tower invariants, and a single squarefree characteristic
polynomial modulo $\ell$ proves the whole class \emph{étale} over $\Iw$; (ii) proves that the Eisenstein component
meets the cuspidal locus in a class $a\not\equiv0$ exactly at the irregular pairs, with the congruence length
$L_k$ at weight $k$ squeezed between $v_\ell(B_k/k)$ and $v_\ell$ of the Eisenstein value of every cuspidal
characteristic polynomial, and --- through Emerton's theorem, matched to these objects --- equal to $v_\ell(B_k/k)$, the $\Iw$-adic congruence module being
$\Iw/(G_a)$ (\S\ref{sec:u-eis}); (iii) proves that \textbf{no component is CM} at tame level~$1$
(\S\ref{sec:u-cm}); (iv) links saturation, Gorensteinness and Loewy length (\S\ref{sec:u-sat}). On the project's
published tower data and one new cell it determines the component structure of every class of $\ell\in\{5,7,11,13\}$
and of six classes at $\ell=691$, and proves $L_k=1$ at $(k,\ell)=(12,691)$, $(200,691)$, $(702,691)$.
\S\ref{sec:u-pred} lists the checkable predictions (U1)--(U7) for the weight ladders $n+j(\ell-1)\ell^{m-1}$.
\end{quote}

\end{document}
